% 本文件由 LaTeX 排版 Agent 根据有效 translation.json 自动生成。
% author=Augustine of Hippo; work_id=1101; title=忏悔录

\chapter{忏悔录}

% structure: p0001 source=h1 target=section reason=page-title

\section{忏悔录（卷一）}

以呼求天主开始，\Person{奥思定}详细叙述了他生命的开始、幼年和少年时代，直到他十五岁；在这个年龄，他承认自己更倾向于所有青春的享乐和恶习，而非读书求学。\par

% structure: p0003 source=h2 target=subsection reason=relative-heading-rank; base=h2

\subsection{他宣扬天主的伟大，他渴望寻求并呼求他，被他唤醒。}

1. 主，你是伟大的，极应受赞美；你的能力伟大，你的智慧无穷。 而人，作为你受造物的一部分，渴望赞美你；人随身带着他的必死性，即他罪的明证，正是你\BibleQuote{拒绝骄傲人}的明证——然而人，你这受造物的一部分，却渴望赞美你。 你激励我们乐于赞美你；因为你造我们是为了你，我们的心若不安息在你内，总不会安宁。 主，请教我认识并明白，哪一样居先：是呼求你呢，还是赞美你；同样，认识你呢，还是呼求你。 但谁若不认识你而呼求你呢？因为不认识你的人，可能会把你当作另一位的身份而呼求你。 或者，我们呼求你，以便能认识你。 \BibleQuote{但是，人若不信他，怎能呼号他呢？未曾听到他，怎能信他呢？没有宣讲者，又怎能听到呢？}\BibleRef{罗 10:14} 寻求主的人必赞美他。因为凡寻求的，必找着\BibleRef{玛 7:7}，找着的人必赞美他。 主，让我在呼求你中寻求你，并在信靠你中呼求你；因为你已宣报给了我们。 我主，我的信德呼求你——这信德是你赐给我的，是你藉着你圣子的\OriginalTerm{降生成人}{incarnation}，藉着你的宣道者的职务，吹入我内的。\par

% structure: p0005 source=h2 target=subsection reason=relative-heading-rank; base=h2

\subsection{我们所呼求的天主在我们内，我们也在他内。}

2. 但我又怎能呼求我的天主，我的天主，我的主呢？因为呼求他时，就是请他到我心里来。 我内有什么地方，能让我的天主进来呢？天地的主宰，能进入我内吗？主，我的天主，我内有什么能容纳你呢？ 难道你造的天地，以及你在其中创造了我，能容纳你吗？或者，既然没有你，任何事物都不能存在，那么一切存在都能容纳你吗？ 那么，我既已存在，而且如果你不在我内，我就不能存在，我何必求你进入我内呢？因为我尚未在地狱里，但你也在那里；因为\BibleQuote{我若下到地狱，你也在那里}。 因此，我的天主，除非你在我内，我不能存在，根本不能存在。或者我宁可以说，除非我在你内，我也不能存在，因为\BibleQuote{万物都出于他，依赖他，而归于他}\BibleRef{罗 11:36}。 主啊，正是这样，正是这样。既然你已在我内，我还呼求你到哪里去呢？你又能从哪里进到我内呢？我能往哪里去，离开天地，好让我的天主到我这里来呢？他曾说过：\BibleQuote{我充满天地}\BibleRef{耶 23:24}。\par

% structure: p0007 source=h2 target=subsection reason=relative-heading-rank; base=h2

\subsection{天主处处充满万有，但天地并不容纳他。}

3. 既然你充满天地，难道天地能容纳你吗？或者，它们不能容纳你，你充满它们之后，还有所剩余吗？ 当天地充满后，你又把你剩余的部分倾注在哪里呢？ 其实，既然你容纳万有，而你所充满的万物，正因为你容纳它们，你才充满它们，你岂需被任何事物所容纳吗？ 因为你所充满的器皿，并不能支撑你，即使它们碎裂，你也不会流溢。 当你倾注在我们身上时（\BibleRef{宗 2:18}），你并不被倾倒，而是我们被提升；你并不消散，而是我们被集合。 然而，你充满万有，是用你的整体充满呢？还是万物不能全部容纳你，便各容纳你的一部分，并且大家都容纳你的一部分呢？还是各物有各自的部分——大的多些，小的少些？那么，是你的一部分较大，另一部分较小吗？还是你整个在每一处，却无一物能完全容纳你呢？\par

% structure: p0009 source=h2 target=subsection reason=relative-heading-rank; base=h2

\subsection{天主的尊威至高无上，他的德能不可言喻。}

4. 那么，我的天主，你究竟是什么呢？我问，你难道不是主天主吗？因为谁是主？唯有主；谁是天主？唯有我们的天主。 你是至高、至善、至能、无所不能的；至慈悲而又至公义；至隐秘而又至切近；至华丽而又至强壮，稳定而不受容纳； 你不变，却改变一切；万古常新，从无陈旧的；使万物更新，却使骄傲人衰败，而他们不知； 你时常工作，却常安息；聚集，却一无所缺；支持、充满、保护；创造、滋养、完成；寻求，却已拥有一切。 你爱，却不燃烧；你忌邪，却无烦恼；你后悔，却无忧愁；你发怒，却仍然平静； 你变更方式，却不改变计划；你收回所找到的，却从未丢失； 你从无匮乏，却以收益为乐；你从不贪婪，却要求利息（\BibleRef{玛 25:27}）。 为使你似欠债，人超额奉献给你；但谁有一样不是你的呢？ 你偿还债务，却从不亏欠；你宽赦债务，却毫无损失。 然而，我的天主，我的生命，我神圣的欢乐，我所说的这些话算什么？人能论及你说些什么呢？然而那些缄默的人有祸了，因为即使讲论最多的也如同哑巴。\par

% structure: p0011 source=h2 target=subsection reason=relative-heading-rank; base=h2

\subsection{他在天主内寻求安息，并求得罪过的赦免。}

5. 啊！我怎能在你内得到安息呢？谁能使你进入我的心灵，使我沉醉，以致我忘却我的灾祸，拥抱我唯一的善——你呢？ 你对我算什么？请怜悯我，使我能说话。我对你又算什么，你竟命我爱你；如果我不爱你，你就动怒，并以巨大的灾祸威胁我？ 那么，不爱你，只是轻微的灾祸吗？唉！唉！我主，我的天主，请凭你的仁慈告诉我，你对我算什么。 \BibleQuote{对我的灵魂说：我是你的救恩}。请这样说，使我能听见。看，主啊，我心灵之耳在你面前；请你开启它们，并\BibleQuote{对我的灵魂说：我是你的救恩}。 让我听见后，就奔向你，抓住你。不要向我隐藏你的慈颜。让我死去，免得我死亡，只要能得见你的慈颜。\par

我的灵魂的居所狭窄；求祢扩充它，使祢能进入。它已荒废，求祢修复它。其中有冒犯祢眼目的地方；我承认且知道，但谁能洁净它呢？除祢以外，我向谁呼求呢？主，求祢洁净我隐秘的罪，保守祢的仆人免陷于他人的罪过。我信，所以我说；主，祢知道。我的天主，我难道没有向祢承认我的过犯，而祢已除掉我心中的邪恶吗？我不与祢争辩，因为祢是真理\BibleRef{约 9:3}，我也不欺骗自己，免得我的邪恶自相矛盾。因此，我不与祢争辩，因为\Quote{主，祢若究察罪孽，主啊，谁能站得住呢？}\par

% structure: p0014 source=h2 target=subsection reason=relative-heading-rank; base=h2

\subsection{第六章 他描述自己的婴孩时期，并赞美天主的保护与永恒眷顾}

仍请容许我在祢的仁慈前说话——我这\Quote{尘埃灰土}\BibleRef{创 18:27}。请容许我说话，因为，看，我是在向祢的仁慈说话，而不是向嘲弄的人。然而，也许连祢也嘲笑我；但当祢转向我时，祢会怜悯我\BibleRef{耶 12:15}。主，我的天主，我想说什么呢？无非是我不知道自己从何处来到这——我该称它为垂死的生命还是活着的死亡？然而，正如我从父母那里听到的（祢用他们的肉身形成了我——因为我自己无法记得），祢仁慈的安慰扶持了我。就这样，妇人的乳汁安慰了我；因为不是我的母亲或乳母填满她们的乳房，而是祢通过她们，按祢的安排和那托住万有的慷慨，赐给我婴孩时期的滋养。因为祢使我需要的，不超过祢所赐予的；而那些养育我的人，也乐意把祢赐给她们的给我。因为她们出于本能的爱，急切要把祢丰盛供给的给我。其实，我的好处来自她们，这对她们是有益的；虽然确实不是来自她们，而是藉她们而来；因为，天主，一切美善都来自祢，我的全部安全都来自我的天主\BibleRef{箴 21:31}。这是我后来才发现的，因为祢借着所赐予我内心的和外在的恩惠，向我彰显了祢自己。因为在那个时候，我只知道吮吸，舒服时满足，疼痛时哭泣——除此以外，再无别的。\par

后来，我开始笑——起初在睡梦中，后来醒着的时候。关于这点，我听过别人提到我，我也相信（虽然我记不得了），因为我们在其他婴儿身上也看到同样情形。渐渐地，我意识到自己在哪里，并想把我的愿望告诉那些能满足我的人，但我做不到；因为我的需求在我内，而他们在外面，他们的任何官能都无法进入我的灵魂。于是我挥舞四肢、发出声音，尽可能地做出几个微弱而不太准确、类似于我愿望的动作和手势；当我不满足时——要么因为不被理解，要么因为满足我的愿望可能对我有害——我便对长辈们不服从我感到愤怒，对那些我没权利要求的人不来服侍我感到恼怒，并用哭泣来报复他们。我后来通过观察婴儿，才得知婴儿就是如此；这些婴儿不知不觉地，比那些深知我的乳母更好地向我显明了我当时就是这样。\par

请看，我的婴孩时期早已逝去，而我活着。但是主，祢永远活着，在祢内没有任何死亡（因为在世界之前，甚至在一切可称为\Quote{之前}的之前，祢就已存在，是祢一切受造物的天主和主；在祢内固定地存留一切不稳定事物的原因，一切可变事物的不变根源，一切无理性与时间性事物的永恒之理），请告诉祢的恳求者，天主；请告诉，仁慈的主，祢可怜的仆人——告诉我，我的婴孩时期是否承接了另一个那时已消逝的时期？那是我在母胎中度过的时期吗？因为关于那个时期，我略有所闻，而且我亲眼见过怀孕的妇女。我的喜乐，天主，在那个生命之前又有何先存？我曾在哪里，曾是谁吗？因为没有人能告诉我这些，父亲不能，母亲不能，他人的经验不能，我自己的记忆也不能。祢嘲笑我问这些事，并命令我按我所知的赞美祢、承认祢吗？\par

我感谢祢，天地的主，为那我没有记忆的初始生命与婴孩时期赞美祢；因为祢赐给人能力，让他能从别人那里推断自己，又让他能凭孱弱妇女的权威，相信许多关于他自己的事。即使在那时，我已经有生命，有存在；随着婴孩时期结束，我已在寻找符号，好让别人知道我的感受。主，这样的受造物能从何处来，除非出于祢呢？岂有人如此精巧，能制造自己吗？或者，存在和生命流入我们，除了\Quote{主，是祢造了我们}这一渠道外，还有别的吗？在祢内，存在与生命是一体的，因为祢自身就是至高无上的存在与生命。祢是至高者，\Quote{祢永不改变}\BibleRef{拉 3:6}，在祢内，今天的日子并不终结，尽管它在祢内终结，因为在祢内一切如此；它们若不被祢支持，便无法消逝。既然\Quote{祢的岁月无终}，祢的岁月就是一个永远的今天。我们和祖先的多少日子，都经过祢的这个今天，从中领受它们的尺度和存在的形式，而来日的日子也将如此领受并消逝！\Quote{但祢却是一样}；明天及未来的一切，昨天及过去的一切，祢将在今天完成，祢已在今天完成。若有人不明白，这与我何干呢？让他仍然欢喜地说：\Quote{这\Emph{是}什么？}让他就这样欢喜吧，宁可喜爱在未能发现中发现，而不愿在发现中未能发现祢。\par

% structure: p0019 source=h2 target=subsection reason=relative-heading-rank; base=h2

\subsection{第七章 他以实例表明婴孩时期也倾向于犯罪}

11. 天主，请垂听！可叹世人的罪孽！人这样说，祢就怜悯他；因为祢造了他，却没有造他里面的罪。谁使我忆起我幼年的罪呢？因为在祢面前，没有一人是免于罪的，就是在世上只活了一天的婴孩，也是如此。谁使我忆起这事呢？岂不是每一个小婴孩，我在他们身上看到我自身所记不得的情形吗？那么，我犯了什么罪呢？是不是我哭着要吃奶呢？假使我如今也这样哭——当然不是为了吃奶，而是为了适合我年龄的食物——我一定会被人理所当然地嘲笑和斥责。我那时所做的理应受责；但因我无法理解那斥责我的人，所以无论习惯或理智都不容许我受到责备。因为我们长大之后，就会根除并抛弃这些习惯。我从未见过任何明智的人，在\BibleQuote{修剪}时，会把好的东西丢掉。\BibleRef{若 15:2} 或者，即使暂时看来是好的，用哭闹来强求那若给与便有害的事物——对那些自由且为长辈的人，以及那赐他生命的，并许多比他更有智慧的人，只因他们不依从他的任性而顺从他，便愤愤不平——竭力挣扎，试图伤害，因为那些若听从便对他有害的命令竟未得服从，这难道是好的吗？所以，婴儿的无辜，在于其肢体的软弱，而不在于其意志。我亲眼见过，也晓得一个尚不会说话的婴儿，竟会嫉妒。它的脸色发白，向它的同奶兄弟投以怨恨的目光。谁不知道这事呢？母亲和乳母都说，她们用些不知什么方法，来抚平这些事。难道这也算为无辜吗？当奶水源源不断、极为丰富时，一个有需要的人，尽管他需要这滋养来维持生命，却不得分享？然而，我们宽待这些事，并非因为这些不是过失，也不是因为这些过失微小，而是因为它们将随着年龄增长而消失。因为，即便你现在容忍这些事，如果发现在年长的人身上，你就不能平心静气地忍受了。\par

12. 所以，主、我的天主，祢赐生命给婴儿，并赋予他一副身体，如我们所见的，祢使他有感官，肢体协调，形态优美，又为他一生的好处和安全，赋予一切生命活力——祢命令我因这些事赞美祢，\BibleQuote{要向主谢恩，歌颂至高者祢的名}；因为祢是一位全能美善的天主，即使祢只做了这些事，也没人能做，只有祢能，唯独祢创造了万有，至美的祢，使万物皆美，并依祢的律法掌管一切。所以，主啊，我生命中这一时期，我毫无记忆，只凭别人的话相信，又根据其他婴儿来推测，它使我懊恼——尽管推测也许属实——竟将它算为我在世所度此生的一部分；因为在遗忘的黑暗中，它如同我在母腹中度过的时期一般。但是，如果\BibleQuote{我是在不义中被塑造，在罪中我母亲怀了我}，那么，我的天主，我请问祢，在哪里，主啊，祢仆人何时是无辜的呢？但是，看，我且跳过那一时期，因为我与那些记忆无法追回的事有什么相干呢？\par

% structure: p0022 source=h2 target=subsection reason=relative-heading-rank; base=h2

\subsection{第八章　他儿时学习说话，并非依循定法，而是从父母的言行中}

13. 那么，我岂不是从婴儿时期脱出，进入了童年吗？或者，童年岂不也是临到我，接续了婴儿时期吗？我的婴儿时期并没有离去（它往哪里去了呢？）；但它也不再存留，因为我不再是不能言语的婴孩，而是一个咿呀学语的孩童了。我记得这事，后来我观察到我最初如何学习说话，因为我的长辈并未以任何固定的方法教我词汇，如后来教我识字一样；而是我自己，当我无法说出我全部意愿，也不能向任何我希望的人表达时，便借着所发出的啜泣、破碎的语音、以及各种肢体的动作（我用这些来强调我的意愿），用祢，我的天主，所赐给我的理智，在我的记忆中重复那些声音。当他们指着某物叫出名称，并在说话时将身体转向它时，我看见并领悟到，他们所要指出的那个东西，就是他们当时所发声音的名称；而他们有此意图，则借着身体的运动，以及为各民族共有的自然语言，即由面部表情、眼神、其他肢体的动作，以及表明内心情感（如追求、拥有、拒绝、回避）的声音，清楚地显示出来。这样，我因常常听到词汇置于妥当的语句中，便逐渐领会到它们指示的是什么事物；并且，一旦我的口舌习惯于发出这些符号，我便借此表达自己的意愿。于是，我与我周围的人互换了表达意愿的符号，并更深地进入了人类生活狂风暴雨般的交往中，那时我尚依赖父母的权威和长辈的指示。\par

% structure: p0024 source=h2 target=subsection reason=relative-heading-rank; base=h2

\subsection{第九章　论男孩之厌学、嗜玩、畏鞭，及长辈师长之愚}

14. 我的天主啊！我那时经历了怎样的苦难和嘲笑啊，当人们把服从师长当作我童年应有的本分，为的是让我在这世上飞黄腾达，在\OriginalTerm{辞令之学}{science of speech}上出众，由此在人间获得荣誉和欺人的财富！之后我被送去上学，学习那些我（卑微如我）不知有何用处的学问；然而，若学习迟缓，我便要受鞭打！因为这被我们先祖视为可嘉许的；许多前人在走过同样的道路时，早就为我们预定了这些苦路，迫使我们经过，给亚当的子孙加多劳苦和忧伤。然而，主啊，我们见到有人向祢祈祷，便从他们那里学到，按我们所能，把祢设想为某位大能者，能够（虽然不为我们的感官所见）垂听并帮助我们。因为当我还是孩子时，便开始向祢祈祷，我的\Quote{助佑}和我的\Quote{避难所}；在呼求祢时，我解开了舌结，虽然年幼，却以不幼的热切恳求祢，使我不在学校受鞭打。而祢没有俯听我时（祢并没有因而将我交于愚妄），我的长辈，是的，甚至我的父母，他们并不愿我遭灾，却嘲笑我受的鞭伤，那是我当时巨大而惨痛的灾祸。\par

15. 主啊，可有任何人，心神如此高远，以如此强烈的爱慕依附着祢——因为即便某种迟钝也能做到那样——但我要说，可有任何人，因虔诚地依附着祢，而被赋予如此大的勇气，以至于能轻看那些拷问架、铁钩，以及诸如此类的种种酷刑，而全世界的人却怀着大恐惧向祢哀求，以躲避这些，并嘲笑那些最怕这些的人，正如同我们的父母嘲笑我们儿时师傅给我的惩罚一般？因为我们对痛苦同样惧怕，同样向祢祈祷以逃避它们；然而我们却犯了罪，在书写、阅读或思考功课上没有达到要求。主啊，我们并非缺乏记忆或才智，按祢的旨意，我们当时的年龄已有足够的天赋——但我们只喜欢游戏；为此我们受到那些自己也在做同样事情的人的惩罚。但长辈们的闲散，他们却称之为正事，而做同样事情的儿童却被这些长辈惩罚，然而无论是儿童或是成人，都得不到怜悯。因为，难道有理智的人会赞成我因儿时玩球而受鞭打，以致于妨碍了快速学习那些功课，而借着这些功课，我成年后却会玩出更不体面的勾当吗？那打我的人，自己所做的，又何尝与此有异？当他与同僚在些许争论中落败时，他所受的愤怒和嫉妒的折磨，比我玩球败给同伴时所受的还更厉害呢？\par

% structure: p0027 source=h2 target=subsection reason=relative-heading-rank; base=h2

\subsection{第十章 因贪爱玩球与观戏，荒废学业，不听父母之命。}

16. 然而，主、天主，万有本性的创造者与安排者——但对于罪恶，祢只是安排者——我犯了错，主、我的天主，我违背了父母的意愿和那些师傅的意愿；因为这种他们（无论是出于什么动机）要我学习的学问，我本可善加利用。我之所以违背他们，并非因我选择了更好的道路，而是出于对游戏的溺爱，喜爱在比赛中胜利的荣誉，让耳朵被虚谎的寓言搔痒，好使它们更剧烈地发痒——同样的好奇心在我眼中日益闪亮，渴望着长辈的表演和竞技。然而那些提供此类娱乐的人却享有盛誉，以致几乎人人都希望自己的子女同样如此，而他们仍情愿子女因这些游戏妨碍学习而受鞭打，因为正是借着这种学习，他们希望子女将来也能成为提供这些娱乐的人。主啊，求祢怜悯地垂顾这些事，解救我们这些现在呼求祢的人；也解救那些不呼求祢的人，使他们能呼求祢，祢好解救他们。\par

% structure: p0029 source=h2 target=subsection reason=relative-heading-rank; base=h2

\subsection{第十一章 患病之际，母亲忧心，他恳求圣洗；病愈后却推迟——父亲尚未信仰基督。}

17. 当我还是个孩子时，就已听说过，通过我们的主、天主屈尊就卑，迁就我们的骄傲，而应许给我们的永生；我也领受了\OriginalTerm{十字圣号}{sign of the cross}，并在母胎中就由我母亲——她极信赖祢——的\OriginalTerm{圣盐}{salt}所腌制。主啊，祢看见了，有一次，当我还是个孩子，突然腹痛如绞，濒临死亡——祢看见了，我的天主啊，因为那时祢已是我的守护者，我是以何等的心灵激动和何等的信德，恳求我母亲的虔诚，以及我们众人的慈母、祢的教会，使我领受我的主、我的天主、祢的基督的\OriginalTerm{圣洗}{baptism}。对此，我肉身的母亲忧心如焚——因为她怀着在祢信德内纯洁的心，为我的永远得救，更慈爱地再受产痛，\BibleRef{迦 4:19}——若不是我迅速康复，她定会毫不迟延地安排我通过祢赋予生命的\OriginalTerm{圣事}{sacraments}，获得入门与洗涤，向祢——主耶稣——\OriginalTerm{宣认}{confessing}，以赦免罪过。于是我的洗涤便被推迟了，仿佛若我存活，就必定会更受污染似的；因为，罪过所招致的罪咎，在圣洗之后，确会更大、更危险。因此，那时我和母亲及全家人都相信了——除父亲外；但他并没有克服母亲虔诚对我的影响，以致阻止我信仰基督，因为他自己那时尚未信。她确实渴望祢——我的天主——是我的父亲，而不是他；在这方面祢帮助她胜过自己的丈夫，虽然她比他更好，却服从他，因为在这事上她服从了祢，因为祢这样命令。\par

我恳求祢，我的天主，若承祢圣意，我愿欣然知晓，当初我的圣洗为何被延迟？是否为了我的益处，才似乎放松了对我的约束，让我犯罪？还是并未放松？若未放松，为何至今各处仍不断向我们的耳中灌输：\Quote{由他吧，随他怎样行，因为他尚未受洗}？但关乎身体健康，却无人喊道：\Quote{让他伤得更重吧，因他尚未痊愈！}那么，如果我立刻得到痊愈，岂不更好？然后，藉由我自己和亲友的勤勉，我灵魂恢复的健康便可安然保存在祢——赐我健康者——的保护之中！的确更好。然而，在我童年之后，似乎有多少巨大而汹涌的诱惑浪潮悬在头上！这些，我母亲早已预见；她宁愿让尚未成形的泥土去面对它们，而不愿让肖像本身去面对。\par

% structure: p0032 source=h2 target=subsection reason=relative-heading-rank; base=h2

\subsection{第十二章 被迫专心向学，但完全承认这是天主的工程。}

然而，在我的童年（它远不如青年时代令我畏惧），我并不爱学习，且憎恶被迫学习，尽管如此，我仍是被迫；这对我虽是好事，我却做得不好，因为若非被迫，我不会去学。没有人违背意愿而行善，即使他所行的是好事。那些强迫我的人也未行善，但施于我的善，却是来自祢，我的天主。因为他们并未考虑我将如何运用他们强迫我学习的东西，除非是满足那富有乞丐的无度贪欲和可耻虚荣。然而祢，连我们头上的头发都一一数过的祢\BibleRef{玛 10:30}，却为了我的益处，利用了所有催迫我学习之人的错误；而我自己不愿学习的错误，祢则用以惩罚我——对此，我虽是一个如此小的孩子，却是一个如此大的罪人，并非不堪承受。如此，祢藉着那些未行善的人，为我行了善；又因我自己的罪，祢公义地惩罚了我。因为正如祢所规定的，一切无节制的感情，都会带来自身的惩罚。\par

% structure: p0034 source=h2 target=subsection reason=relative-heading-rank; base=h2

\subsection{第十三章 他喜爱拉丁文学习和诗人的虚幻故事，却憎恶文学基础和希腊文。}

我自童年起便学习希腊文学，却对其心生厌恶，个中缘由，我至今仍不明白。对于拉丁文，我却极为喜爱——不是启蒙老师所教的基础，而是文法家所传授的；因为那些阅读、书写和算术的入门功课，我视为与希腊文一般，是沉重的负担和惩罚。然而，这种好恶若非源自于此生的罪孽与虚妄，又来自何处呢？因为我\Quote{只是血肉，是一阵去而不返的风}。那些入门功课确实更好，因为它们更为确切；藉由它们，我获得了，且至今保有，阅读所见文字和随己意书写的能力；而在另一类学习中，我却被迫去了解某个\Person{埃涅阿斯}的漫游，却忘记了自己的漫游，为死去的\Person{狄多}哭泣，因为她为爱自杀；与此同时，我却以干涩的眼睛，忍受着自己远离祢而悲惨死去，置身于那些事物之中，天主啊，我的生命。\par

有什么能比一个可怜的罪人更可怜呢？他并未怜悯自己，却为\Person{狄多}因爱\Person{埃涅阿斯}而死泪流满面，却不为自己不爱祢而死亡流泪——天主啊，我心灵的光明，我灵魂内口的食粮，将我的心与我最深的思想结合的大能！我并未爱祢，并对祢犯了邪淫之罪；我周围那些同样犯罪的人喊道：\Quote{好啊！好啊！}因为与世俗为友，就是与祢犯邪淫；\BibleRef{雅 4:4} 且四处响起\Quote{好啊！好啊！}的喊声，直到人若不如此便觉羞耻。对此，我未曾流泪，尽管我为\Person{狄多}流泪，她手持剑尖寻求死亡，而我自身却在寻求祢受造物中最卑微者——离弃了祢——尘土归于尘土；若禁止我阅读这些，我会因不被允许阅读让我悲伤的东西而感到多么难过。这种疯狂被视为比那让我学会读写的学习，更为尊贵、更有成果的学问。\par

但现在，我的天主啊，请向我的灵魂呼喊；让祢的真理对我说：\Quote{不是这样；不是这样；那些入门教导好得多。}因为，看，我宁愿忘掉\Person{埃涅阿斯}的漫游及所有此类故事，也不愿忘记如何书写和阅读。诚然，文法学校的门口确实挂着一层帷幕；但这与其说是奥秘之尊威的标记，不如说是错误的遮丑布。不要让他们向我叫嚣，那些我已不再畏惧的人，当我向祢，我的天主，坦白我灵魂所渴望的，并安然谴责我的恶行，好能爱慕祢的善道时。也不要让那些买卖文法知识的人向我叫喊。因为若我问他们，如诗人所说，\Person{埃涅阿斯}曾到过\Place{迦太基}是否属实，无知者会回答说他们不知道，有学识者会否认其真实性。但若我问，\Person{埃涅阿斯}这个名字由哪些字母拼写，所有学过的人都会按人们对这些符号的约定俗成，正确回答。再者，若我问，遗忘哪一样会给我们的生活带来最大的不便，是阅读和书写，还是这些诗意的虚构，谁看不出，每个尚未全然迷失自我的人会怎样回答？因此，当我作为孩子，喜爱那些虚浮的学业，胜过那些更有益的，或更确切地说，喜爱前者而憎恶后者时，我便犯了错。\Quote{一加一等于二，二加二等于四}，这当时对我而言，确是一首可憎的歌谣；而那满载武装战士的木马，\Place{特洛伊}的焚烧，以及\Person{克瑞乌萨}的\Quote{幽灵幻影}，却是虚妄中最令人愉悦的奇观。\par

% structure: p0038 source=h2 target=subsection reason=relative-heading-rank; base=h2

\subsection{第十四章 他为何轻视希腊文学，而轻易学会拉丁文。}

23. 可是，我为何厌恶那充满类似故事的希腊文学呢？因为荷马也善于编造同样的故事，且最是甜美地虚幻，然而他在我童年时却令我生厌。我深信，维吉尔对于希腊儿童，倘被迫学习，也必如荷马之于我。实际上，学习外语的困难，宛如胆汁，混杂在那些希腊神话故事的甜美之中。因为我连一个字也不懂，而他们却用残酷的威吓与惩罚，强逼我去弄明白。也曾有一段时间（当我尚在襁褓），我连拉丁语也不识；但我毫无惧怕与痛苦，仅凭着留意乳母的抚慰、对我微笑者的戏谑，以及与我嬉戏者的顽皮，便学会了。我学会这一切，确未曾受到任何惩罚的催逼，乃是我的心催迫我，要生出自己的观念，而除非学习语词——不是向那些教我的人，而是向那些与我交谈的人学习——我便无法做到；我也将我所理解的，传到他们耳中。由此可见，自由的好奇，远胜于充满恐惧的强制，更能驱使我们学习这些事。但这强制，却通过祢的法律——天主啊，祢的法律——限制了那自由的泛滥；祢的法律，从教师的教鞭到殉道者的磨难，都有效力地为我们调了一杯有益的苦酒，召唤我们从那些诱我们远离祢的、有害的欢愉中，回归祢自己。\par

% structure: p0040 source=h2 target=subsection reason=relative-heading-rank; base=h2

\subsection{他恳求天主，愿将童年所学一切有用之物都奉献给祂。}

24. 主，请俯听我的祈祷；求祢不要让我的灵魂在祢的管教下昏厥，也不要让我在向祢承认祢的仁慈时昏厥，因为祢正是凭借这些仁慈，将我从我一切最邪恶的道路中拯救出来，好使祢在我眼中变得比我所追随的一切诱惑更为甘饴；好使我全心爱慕祢，以我整个心灵握住祢的手，并求祢从一切诱惑中拯救我，直至永远。因为，请看，主，我的君王，我的天主，凡我童年所学的一切有用之物，都愿为服事祢而用——我说话、写作、计算，都为服事祢。因为当我学习虚妄之事时，祢赐予我祢的管教；而我在那些虚妄中寻乐所犯的罪，祢已宽恕了我。的确，我在其中学了许多有用的语词；但这些语词也可在不虚妄的事上学到，而这才应是青年行走的安全之路。\par

% structure: p0042 source=h2 target=subsection reason=relative-heading-rank; base=h2

\subsection{他反对教育青年的方式，并指出诗人为何将邪恶归于神明。}

25. 然而，人啊，那人类习俗的洪流，你有祸了！谁能阻挡你的奔流？你何时才会干涸？你还要将厄娃的子孙，冲进那片巨大可畏的海洋多久？就连那些登上十字架（\OriginalTerm{木头}{lignum}）的人，也几乎难以渡过。难道我不是在你里面，读到那雷鸣者兼奸夫的朱庇特吗？然而，他确实不能兼具二者；但这不过是借着虚构的雷声作为遮掩，好使他得以纵容，去模仿真实的奸淫。可是，我们那些身穿长袍的教师当中，有谁会以温和的耳朵，倾听他学堂里一个大声抗辩的人说：\Quote{这些都是荷马的虚构；他将人的事转移给了神明。我倒希望他能将神的事转移给我们。}然而，他若是这样说，倒更为真实：\Quote{这些，诚然是他的虚构，但他却将神圣的属性归于犯罪的人，好使罪行不被算为罪行，好使任何犯下这等事的人，看起来像是效法天上的神明，而非堕落的人类。}\par

26. 然而，你这地狱的洪流，人类的子孙竟被投入其中，还因学得这些事而领受奖赏；当此事在法庭中，在那些不仅给予奖赏，更另付薪资的法律眼皮底下进行时，竟被大加称许。你撞击着你的巉岩，咆哮道：\Quote{由此学到语词；由此获得口才，这是说服别人接受你的想法、阐明自己意见所最必需的。}诚然，若非泰伦斯曾将一名浪荡青年搬上舞台，并立朱庇特为其放荡的榜样，我们便永不会懂得这些词：\Quote{金雨}、\Quote{怀中}、\Quote{私通}、\Quote{至高之天}，以及同处出现的其他词语了：——\par

观画一幅，其上绘有故事，\newline{}朱庇特化作金雨降下，\newline{}落入达娜厄的怀中……与女子私通。\newline{}且看他如何如同拥有天界的权威，煽动起自己的情欲，他道：——\newline{}伟大的朱庇特，\newline{}他以雷霆震撼至高之天，\newline{}而我，一个可怜的凡人，竟不能如此！\newline{}我做了，我全心做了。\par

这些语词，并不因这可耻之事而变得稍易学习，反而借此，那无耻之行被更加肆无忌惮地实施。我并不归咎于那些语词，它们就像精选而珍贵的器皿一般；我所指责的，是那些醉醺醺的教师灌给我们，盛在其中的错误之酒；倘若我们不饮，便会遭受鞭打，而且无法向任何清醒的审判者申诉。然而，我的天主啊——如今我能在祢面前安然忆及此事——我这不幸的人，那时竟乐意且欢欣地学习这些东西，并因此被称为大有希望的孩子。\par

% structure: p0047 source=h2 target=subsection reason=relative-heading-rank; base=h2

\subsection{他继续论及青年在文学科目上的不幸训练方式。}

27. 我的天主，请容忍我，容我略略述说祢赐给我的才秉，以及我将它们浪费在何等愚妄上。那时曾有一课，颇令我的灵魂不安，要我——出于对赞誉的期望，与对羞辱或鞭打的恐惧——讲出朱诺的话语，她因愤怒与悲伤而无法\par

将拉丁姆阻挡\par

于达达尼尔王的一切通道之上，\par

我听说\Person{朱诺}从未说出过这样的话。然而我们却被迫在这些诗意的虚构中漫步，并把诗人用韵文说的话改写成散文。在众人中，谁的演说最受称赞？就是那个能根据所描绘人物的声望，最逼真地再现愤怒与悲伤的激情，并用最合适的语言表达出来的人。可是，我真正的生命，我的天主啊，我的朗诵获得的喝彩超过了许多同龄同窗，这对我又算什么？看，这一切难道不是烟云与风吗？难道就没有其他事让我可以施展才智与口才吗？主啊，你的赞颂，你藉着你的圣经本可以支撑我心灵的卷须；那样它就不至于被这些空洞的琐事拖走，成为空中飞鸟的可耻猎物。因为人们向堕落的天使献祭的方式不止一种。\par

% structure: p0052 source=h2 target=subsection reason=relative-heading-rank; base=h2

\subsection{第十八章 人们渴求遵守学问的规则，却忽视了永恒得救的永恒法则。}

28. 但是，我竟这样被引向虚幻，并且远离了你，我的天主，这有什么可奇怪的呢？当时有人给我立下榜样，要我去效法：这些人若在叙述自己的某些行为时——这些行为本身并不邪恶——犯了\OriginalTerm{语法错误}{barbarism}或\OriginalTerm{句法错误}{solecism}，一经指责便羞愧万状；而他们若用精挑细选的辞藻，就自己的放荡行为发表一番词藻华丽、结构完整的演说，并为此博得喝彩，他们便趾高气扬。主啊，你看见这一切，却沉默不语，\BibleQuote{恒忍，并且富于仁慈与真理}，确实如此。难道你要永远沉默吗？即使现在，你也从这无底深渊中，救拔出寻求你、渴慕你甘饴的灵魂，就是那\BibleQuote{心向你说：\InnerQuote{我寻求你的面容，主，你的面容我必再寻求。}}的我，由于我昏暗的情欲\BibleRef{罗 1:21}，我曾远离你的面容。因为我们离开你或归向你，并不靠双脚，也不是靠改换地方。或者说，那个小儿子是不是寻找马匹、车辆或船只，或者插上可见的翅膀飞走，或者活动肢体长途跋涉，才能在远方挥霍你在他离家时赐给他的一切呢？你赐予时是仁慈的父亲，他穷途归来时，你更是仁慈！\BibleRef{路 15:11}所以，在放荡中，即在昏暗的情欲中，便是与你的面容相距遥远。\par

29. 主，天主，请你看，照你一向所行的，耐心地看：世人多么兢兢业业地遵守那些从先辈口传下来的字母和音节的规矩，却忽视从你那里领受的永恒得救的永恒法则；以致那从事或教授祖传发音规则的人，倘若违反语法习惯，将 human 一词的开头字母不予送气，说成\Emph{uman}，比起他违反你的诫命，以一个人仇恨另一个人，前者冒犯别人更甚。似乎的确，有人会觉得，敌人对自己的危害，竟比不上他因恨人而心中激起的仇恨更危害自身；或者，他迫害别人，竟不能像他因仇恨而摧毁自己的灵魂那般彻底摧毁别人。实实在在，就没有一种文字的学问，比得上\OriginalTerm{良心的书写}{writing of conscience}更为内在——即他自己不愿意遭受的，决不可加给别人。你何其神妙，静默地\BibleQuote{居住在至高之处}\BibleRef{依 33:5}，独一的伟大天主，你以不息的法律，使盲目的惩罚降于非法的欲望！当一个人为求口才的美誉，站在人间审判官面前，身旁围着众多群众，用最凶狠的仇怒攻击自己的仇敌时，他极其谨慎地留心自己的舌头不要犯语法错误，却毫不留意切勿因自己精神的狂怒而从人群中割舍掉一个人。\par

30. 这些就是我，不幸的孩子，置身其中的风气；在那竞技场上，我生怕犯一个语法错误，远甚于我犯了错误后嫉妒那些没犯错误的人。我的天主，我向你坦白认罪：我曾因这些事而获得那些我当时认为必须全力取悦的人的喝彩，因为我看不见自己被抛离你眼目的耻辱深渊。在你的眼中，我早已是更为可耻的了：我甚至让那些像我一样的人不满，我出于贪玩、想看无聊的表演、以及一种模仿表演的焦躁不安，用无数的谎言欺骗家庭教师、老师，以及自己的父母。我从父母的地窖和餐桌上偷窃，或是受口腹之欲的奴役，或是为了能有东西给那些卖给我玩物的孩子们；他们虽然卖给我玩物，却同我一样喜爱这些玩物。在这样的游戏中，我还常常追求不光彩的胜利，自己却被虚荣的争强好胜之心所征服。有什么是我自己极不能容忍，或一经发觉便猛烈指责的，不正是我对别人所做的那些事吗？而一旦我被发觉而受指责，我宁愿争吵，也不愿退让。这就是童年的天真吗？不，主，不，主；我的天主，我恳求你的仁慈。因为随着年岁增长，这些同样的罪过，从老师和管教、坚果、皮球和麻雀，转移到官吏和君王、黄金、田地，以及奴隶身上，正如戒尺为更严厉的惩罚所接替。正因如此，我们的君王，你才嘉许孩童的身材，作为谦卑的象征，说：\BibleQuote{天国正是属于这样的人。}\par

但是，主啊，祢是至卓越至美善的，宇宙的设计者和统御者，我们的天主，即使祢曾定意让我不能存活过童年，仍应向祢献上感谢。因为我那时已存在；我活着，有感觉，为自己的安宁操心——这是我得以存在的那个\OriginalTerm{最奥秘的一体性}{most mysterious unity}的一个痕迹；我以内在感觉看守我感官的完整性，并在这些琐细的追求中，也在对琐细事物的思考中，学会喜爱真理。我厌恶受骗，有强健的记忆力，具备言语的能力，因友谊而温和，回避忧愁、下贱、愚昧。在这样一个存有中，有什么不是奇妙且值得称道呢？但这一切都是我的天主的恩赐；我并未将之给予自己；这些都是好的，而所有这一切构成了我自己。因此，那创造我者是善的，祂就是我的天主；我要在祂面前，因我童年时所拥有的每一美善恩赐而大大欢欣。因为我的罪正是在于此：不在祂内，却在祂的受造物——我自己和其他一切——中寻求欢乐、尊荣和真理，从而堕入忧愁、困苦和错谬。感谢祢，我的喜乐、我的荣耀、我的信赖、我的天主——感谢祢的恩赐；但求祢为我保全它们。因为这样，祢就会保全我；而祢所赐予我的这些恩赐将发展并得以成全，我自身也将与祢同在，因为我的存在来自祢。\par

\SourceRecord{Translated by J.G. Pilkington.；Nicene and Post-Nicene Fathers, First Series；Vol. 1.；Edited by Philip Schaff.；Buffalo, NY: Christian Literature Publishing Co.,；1887.；<http://www.newadvent.org/fathers/110101.htm>.}

% structure: p0001 source=h1 target=section reason=page-title

\section{忏悔录（第二卷）}

他进入青春期，确切地说，到了他年龄的第十六个年头的初期，那时他荒废学业，沉溺于淫欲之乐，并与同伴一起偷盗。\par

% structure: p0003 source=h2 target=subsection reason=relative-heading-rank; base=h2

\subsection{他哀叹自己青年时的邪恶}

1. 我现在要回想我过去的污秽和灵魂肉欲的败坏，不是因为我爱它们，而是为能爱你，我的天主。因爱你之爱，我这样做，在痛苦的回忆中回想我极恶的行径，好使你成为我的甘饴——你非欺骗的甘饴！你幸福而确实的甘饴！——从我的放荡中收聚自己，在其中我曾被撕成碎片，远离你唯一者，在众多虚幻中迷失自己。因为我年轻时曾渴望以求世俗的事满足自己，我竟敢再度沉溺于种种虚幻的爱情；我的形骸消逝，在你眼中腐化，自我欣赏，并渴望在人眼前得宠悦。\par

% structure: p0005 source=h2 target=subsection reason=relative-heading-rank; base=h2

\subsection{他极度悲伤，回想起十六岁时曾沉溺的放纵情欲}

2. 但我所喜欢的不就是爱与被爱吗？但我没有以节制造就纯洁的友谊之光，心意相通，而是从肉欲的黑暗和青春的热烈中升起了浓雾，遮蔽并笼罩了我的心，使我不能分辨纯洁的情爱与不洁的欲望。二者在我内混淆沸腾，拖曳我不稳定的青春陷入不洁欲望的险境，将我投入耻辱的深渊。你的愤怒笼罩了我，我却不知道。我因灵魂骄傲的惩罚——死亡桎梏的铿锵声而变聋；我远离你漂泊，你\BibleQuote{容忍}了我\BibleRef{玛 17:17}；我被抛来抛去，耗竭，倾泻，在我淫行中沸腾，而你保持缄默，我迟来的喜乐啊！你那时保持了缄默，而我仍远离你漂泊，进入越来越多更荒芜的忧伤苗床，带着骄傲的沮丧和不安的倦怠。\par

3. 但愿有人能节制我的紊乱，将周围事物瞬息的美转为我的利益，为它们的甘甜设定界限，如此我青春的潮汐便能倾注在婚姻的岸上，如果他们不能在你的法律所命定的家庭生活范围内获得平静和满足，主啊——你如此塑造我们死亡的果实，也能用温柔的手拔除那些被排斥在你乐园之外的荆棘！因为你的全能离我们并不远，即使我们远离你，否则我本当更警醒地听从云端的声音：\BibleQuote{然而这样的人肉身必受苦难，我却愿意保全你们}\BibleRef{格前 7:28}；并说：\BibleQuote{男人不亲近女人倒好}\BibleRef{格前 7:1}；还有：\BibleQuote{没有家室的人，挂虑的是主的事，想怎样悦乐主；有家室的人，挂虑的是世俗的事，想怎样悦乐妻子}\BibleRef{格前 7:32}。因此，我本当更专心地听从这些话，\BibleQuote{为了天国的缘故而自阉}\BibleRef{玛 19:12}，我便能更幸福地期待你的拥抱。\par

4. 但我，可怜的愚人，沸腾如海，离弃你，随从自己猛烈的洪流，超越你一切界限；也未逃脱你的鞭笞\BibleRef{依 10:26}。因为谁人能做到呢？但你总在我左右，慈怒地盯着，用最痛苦的烦恼击碎我所有非法的快乐，好使我寻求无烦恼的快乐。但除了在你内，主啊，我何处能找得这样的快乐——除你以外，我寻找不到——你通过忧伤教训\BibleRef{申 32:39}，打伤我们是为了医治我们，杀死我们是为使我们不死于你。在你家的欢乐之外，我何所在，又飘泊多远，在我肉身年龄的第十六个年头，当时情欲的疯狂——人间羞耻对此赋予完全自由，尽管你的法律禁止——完全支配了我，我全然委身于它？周围人却不顾念如何借婚姻拯救我免于毁灭，唯一挂虑的是要我学习强有力的话语，成为有说服力的演说家。\par

% structure: p0009 source=h2 target=subsection reason=relative-heading-rank; base=h2

\subsection{关于他的父亲，塔加斯特的一名自由民，资助儿子求学之事，及母亲为保持贞洁的劝诫}

5. 那一年我的学业中断了，因为从\Place{马达乌拉}（邻近的城市，我已先前去那里学习语法和修辞）回来后，需要筹备前往\Place{迦太基}深造的经费；这经费主要凭父亲的决心而非财力，他不过是\Place{塔加斯特}一名贫穷的自由民。我向谁叙述这些？不是向你，我的天主；而是在你面前向我的同类，即人类中可能偶尔读到我这著述的一小部分人。目的何在？使我和所有阅读的人，能反省我们应从多深的深渊向你呼号。因为有什么比一颗忏悔的心和信德的生活更靠近你的耳？谁不称扬赞颂我的父亲，他竟超出财力，为儿子求学远行备办一切必需品？因为许多远为富有的市民并未为子女如此行。然而这同一位父亲却不顾及我如何向你成长，也不顾我是否贞洁，只要我善于辞令——哪怕我对你的耕种，我的天主，你是心中唯一真实和善的主，而心是你的田地，是何等荒芜。\par

6. 但在那年，我十六岁时，与父母同住，暂时停学（这闲散是由于父母的拮据强加于我的），情欲的荆棘便在我头顶疯长，无人拔除。更有甚者，父亲在浴室见到我，察觉我已成少年，且不安分地躁动着；他仿佛从这事预见未来的子孙，便欢喜地告知我母亲——陶醉于那种醉态中，世界因此时常忘记祢，其创造者，并爱上祢的受造物而非祢自己，因着自身邪僻的无形之酒，转而屈服于最可耻的事物。但在母亲的心内，祢此时已经开始了祢的圣殿与祢圣住所的奠基，而我父亲那时还只是个慕道者，且是最近的事。她于是怀着虔诚的恐惧与战栗警觉起来；虽然我尚未受洗，她却惧怕那些弯曲的道路，行走其上的人背向祢，而不面向祢。\BibleRef{耶 2:27}\par

7. 我有祸了！我竟敢说，我的天主啊，当我越发远离祢时，祢保持沉默吗？那时祢真对我默不作声吗？祢通过我母亲、祢忠信的女仆，灌进我耳中的话语，若不是祢的，又能是谁的呢？但其中没有一句沉入我心，使我遵行。因为她渴望，我记起她曾私下警戒我，极其挂虑地说：\Quote{不要行淫；但尤其不要玷污别人的妻子。}这些在我看来只是妇人之见，我若听从会感到羞愧。但它们是祢的话，我却不知道；我以为祢保持沉默，说话的只是她，而祢正是通过她从未对我沉默，并在她身上被我这个儿子、\Quote{祢婢女之子、祢的仆人}所轻看。但那时我并不知道这一点；我盲目地横冲直撞，以至于在同辈中，我羞于不及他们无耻；当我听到他们夸耀自己可耻的行为，甚至越是卑鄙就越发夸口时；我便乐于去做，不仅是为了乐趣，更是为了虚名。除了恶行，什么该受责难呢？我却为了免受责难而装出比自己更坏的样子；每当我未像那些放荡者一样犯罪时，我便声称自己做了未曾做过的事，以免因较为清白而显得卑贱，或因较为贞洁而被人轻视。\par

8. 看，我同何等样的伙伴行走在\Place{巴比伦}的街市，在其污秽中辗转翻滚，仿佛在肉桂与珍膏之中。为了更牢固地黏附其中心，我那无形的仇敌践踏我、诱惑我，我亦易被诱惑。至于我肉身的母亲，尽管她自己在此之前已逃离了\BibleQuote{巴比伦中心}，\BibleRef{耶 51:6}——不过，她尚在边缘缓慢前行——在劝我贞洁时，却未顾忌到丈夫告诉过她关于我的事，好让她在婚姻情感范围内约束（即使无法根除）那她认为是眼前可毁灭、将来又危险的事。但她对此未加注意，因为她害怕妻子会成为我前程的障碍与累赘。不是指我母亲在祢内怀有的对来世的希望，而是求学的希望，我的双亲都过于急切地盼我获得——父亲，因为他几乎或根本未曾想到祢，对我只有虚妄的念头；母亲，因为她估量那些常规的求学之路不仅没有妨碍，反而有助于我将来归于祢。我如此推测，尽我所能回忆父母的性情。同时，对我的约束松弛了，超越了应有的严格，使我得以随心所欲地嬉戏，甚至放荡不羁。在这一切之中，有一层迷雾遮蔽了祢真理的光辉，我的天主啊；我的邪恶如同从\Quote{脂油}中显露出来。\par

% structure: p0014 source=h2 target=subsection reason=relative-heading-rank; base=h2

\subsection{第四章 他与同伴行窃，非因贫困所迫，乃是出于某种对善行的厌弃。}

9. 偷窃受祢律法的惩罚，主啊，也受铭刻在人心中的法律惩罚，这法律连邪恶本身也无法抹去。有哪个窃贼能容忍窃贼呢？即便是富有的窃贼，也不会容忍那因匮乏所迫而行窃的人。然而我却有意偷盗，而且确实做了，并非由于饥饿或贫穷，而是出于对善行的厌弃和邪恶的放纵。因为我偷取的是我已充足有余、且好得多的东西。我并不想享受所偷之物，只想享受偷窃与罪恶本身。靠近我们葡萄园有一棵梨树，满是果实，但其色泽与滋味都无诱人之处。我们一群放荡的青年，深夜时分（按照我们可耻的习惯，在街上嬉戏到那时）去摇晃那树，将果子洗劫一空，带走大批，不是为了自己吃，只是尝了一点之后，就扔给猪吃；我们这么做格外开心，正因为这是不许做的事。请看我的心，我的天主啊；请看我的心，祢曾在无底深渊中怜悯它。现在，请让我的心告诉祢，它在那里寻求什么，使我无缘无故地放荡，除了邪恶本身，竟无其他引诱。这事污秽，我却爱它。我爱走向灭亡。我爱自己的过错——不是那导致我犯错的东西，而是过错本身。卑劣的灵魂，从祢的穹苍堕入全然的毁灭——不因这羞耻寻求什么，单单寻求羞耻本身！\par

% structure: p0016 source=h2 target=subsection reason=relative-heading-rank; base=h2

\subsection{第五章 论罪的动机，不在于爱恶，而在于渴望获取他人财物。}

10. 凡美丽的形体、金银及万物，皆有可欲之处；身体的接触亦有其强大的共鸣，且各官感各适于其身体。世俗的荣誉也有其荣耀，以及统治与征服的能力；由此亦生出复仇的贪欲。然而，为获得这一切，我们必不可离开祢，主啊，亦不可偏离祢的法律。我们在此世所度的生活，亦因其自身的某种合宜尺度，及与下界万有的和谐，而具有独特的吸引力。人类的友谊亦因一种甘饴的纽带，在众多灵魂的一体性中，变得格外亲爱。为了这一切以及类似的事，人便犯了罪；因为过分偏爱这些低等的美物，而忽略了更美好、更高尚的——就是祢，我们的主天主，祢的真理，及祢的法律。因为那些卑微之物自有其乐，却非如我的天主那样——祂创造了万有；因为义人要在祂内取得欣慰，而祂是心地正直者的甘饴。\par

11. 因此，当我们探究为何犯下罪行时，除非似乎有理由表明，犯罪者可能希望获得我们称之为卑微之物的某些东西，或是害怕失去它们，否则我们不会相信。的确，这些事物美丽而合宜，尽管与那些更高尚、属天的美善相比，它们显得下贱可鄙。有人杀害了另一个人；他的动机是什么？他贪图他的妻子或财产；或想偷窃以维持生计；或害怕被那人夺去类似的东西；或受了伤害，怒火中烧，渴望复仇。他会毫无动机地杀人，仅仅以杀人的行为为乐吗？谁会相信呢？至于那个野蛮而残忍的人，人们说他天性邪恶而残酷，却仍有动机可寻。他说道：\Quote{免得闲散，} 又说，\Quote{手与心变得懒惰。} 目的何在呢？啊，就是藉着这样的恶行一度占据了城池，他便能获得荣誉、权势和财富，得以免除对法律的恐惧，摆脱家庭所需的窘境，并避免意识到自己的邪恶。因此看来，即使\Person{喀提林}本人，所爱的也不是他那些恶行，而是其他东西，那东西给了他作恶的动机。\par

% structure: p0019 source=h2 target=subsection reason=relative-heading-rank; base=h2

\subsection{第六章 为何他以此偷窃为乐，而一切披善之外衣引人向恶之事，皆真实且完美地独在天主内。}

12. 那么，可怜的我，在我十六岁的那年，在你——我的偷窃，你暗昧的行为——身上，我所痴恋的究竟是什么呢？你并不美丽，因为你只是偷窃。难道你是什么，好让我与你争辩吗？我们偷来的那些梨，看来是悦目的，因为那是祢的受造物——祢，万有中最美的，万有的创造者，美善的天主——天主，至高的善，也是我真正的善。那些梨确实看起来令人愉悦；但我可怜的灵魂贪求的并非它们，因为我有更好的，而我摘取它们，只是为了偷窃。因为摘了之后，我便把它们扔掉了，我唯一的快慰就是我的罪，我乐于享受这罪。若这些梨中有任何一只进了我的口，那甜味便是吃时所犯的罪。主，我的天主，我现在要问，在我的那次偷窃中，究竟是什么给了我如此的喜乐；而看哪，其中毫无美丽可言—— 我指的，不是存于正义与智慧中的那种美；也不是存于人的心灵、记忆、感官和生命中的那种美；也不是星辰运行、光彩与美丽，或大地、海洋，弥漫着新生万物，于诞生之刻替换衰微者的那种美；更不是那虚假而虚幻的，属于欺诈诸恶的美。\par

13. 因为骄傲如此模仿高位，然而唯有祢是天主，高出万有之上。野心所求，无非荣誉与名声，然而唯有祢当受尊崇于万有之上，并永远受颂扬。权势者的残酷想要令人畏惧；但除天主外，谁当受畏惧呢？谁又能从祂的大能中，被强行夺走或抽离任何事物——何时、何地、往何处、被何人？放纵者的诱惑企图被视为爱；然而，没有什么比祢的爱德更诱人，也没有什么比那光耀美丽超越万有的祢的真理，更值得健康地去爱。好奇佯装对知识的渴望，然而祢才是最透彻知晓万有的。的确，无知与愚昧也隐藏在纯真与无害的名目之下，因为没有什么比祢更纯真；而什么更无害呢，既然罪人是受自己行为的伤害？懒惰似乎渴求安息；但除主之外，哪里有真正的安息？奢侈被称许为丰盛与富足；但祢是永不凋谢的喜乐之圆满与无穷的丰饶。挥霍呈现慷慨的影子；但祢是一切美善最慷慨的施予者。贪婪想要多占；而祢是万有的拥有者。嫉妒争求卓越；但有什么像祢那样卓越？愤怒寻求报复；谁比祢报复得更公义？恐惧因威胁所爱之事的非常与突然的变故而起，并为安全而警醒；但有什么非常或突然的事能临到祢？谁能夺走祢所爱的？除祢以外，何处有不可动摇的安全？悲伤为失去了欲望曾乐享之物而憔悴，正是因为它不愿有任何东西被夺走，如同任何东西都不能从祢那里被夺走一样。\par

灵魂背离祢而向外寻求时，就是在\OriginalTerm{行淫}{fornication}，因为她不能在祢之外找到纯洁无瑕的事物，直到她归向祢。那些远离祢并自高自大反抗祢的人，他们就是这样扭曲地模仿祢。但正是通过如此模仿，他们也承认祢是万物的创造者，因而没有任何地方可以完全逃避祢。那么，我在那次偷窃中爱的是什么呢？我在其中，尽管是以败坏和扭曲的方式，模仿了我的主什么呢？我是否想用诡计干犯祢的法律，因我无力以强力为之，这样，身为俘虏，我可以模仿一种不完全的自由，通过做我不允许做的事而不受惩罚，在模糊的类比中仿效祢的\OriginalTerm{全能}{omnipotency}？看，祢的仆人逃离他的主，追随着一个影子！哦，腐烂！哦，生命的怪物和死亡的深渊！我是否仅仅因为那是不被允许的就喜欢不法之事？\par

% structure: p0023 source=h2 target=subsection reason=relative-heading-rank; base=h2

\subsection{第七章 他为罪过的赦免感谢天主，并提醒每个人，至高的天主可能保守我们免于更大的罪。}

\Quote{我该如何回报主，}当我的记忆回想这些事时，我的灵魂不至于惧怕它们？主啊，我要爱祢，感谢祢，并宣认祢的名，\BibleRef{默3:5}，因为祢从我身上除去了这些如此邪恶和罪孽的行为。我将此归于祢的恩宠和仁慈，是祢使我的罪如同冰一样融化了。我也将我所没有犯的恶事归于祢的恩宠；因为我既为了罪而爱罪，什么罪我犯不出来呢？是的，我承认所有的罪都得了赦免，既包括我因自己的乖戾而犯的，也包括那些因祢的指引而没有犯的。谁在默想自己的\OriginalTerm{软弱}{infirmity}时，敢把自己的贞洁和\OriginalTerm{清白}{innocency}归于自己的力量，以致他少爱祢，好像他不大需要祢的仁慈，而祢正是借这仁慈赦免了那些归向祢的人的过犯？因为任何人，若蒙祢的召唤，听从祢的声音，并避开他读到我回忆和忏悔的那些事，就不要轻视我；我病了，被那位\OriginalTerm{医师}{Physician}治好了，\BibleRef{路4:23}，正是靠祂的援助，他才没有生病，或者更准确地说，生病较轻。为此，让他同样地爱祢，甚至爱得更多，因为他看到我从如此严重的罪恶软弱中恢复过来，正是由于祂，他也看到自己免于同样的软弱。\par

% structure: p0025 source=h2 target=subsection reason=relative-heading-rank; base=h2

\subsection{第八章 在偷窃中，他爱的是同伙的交往。}

\BibleQuote{当时我有什么果子呢，}\BibleRef{罗6:21}，可悲的我，想起那些事就使我羞愧——尤其是那次偷窃，我只是为了偷窃本身而爱它？既然偷窃本身毫无价值，我这样爱它就更悲惨了。然而，我一个人是不会做的——我回想我当时的内心——我一个人是不会做的。那么，在其中我所爱的是与同伙的\OriginalTerm{交往}{companionship}，他们和我一起干的。因此，我爱的并不只是偷窃本身——不，其实是我只爱那，因为交往是毫无价值的。事实是什么呢？谁能教导我，除了那位照亮我心灵并探察其中黑暗角落的主？是什么进入我的思想，让我探究、讨论和反思呢？因为如果那时我爱我所偷的梨，并想享受它们，我完全可以独自干，只要单是实施偷窃就足以满足我的快感；我也不需要用同伙的怂恿来刺激我情欲的骚痒。但是，既然我的快乐不在于那些梨，而在于\OriginalTerm{罪行}{crime}本身，这罪行是由我的同伙共同犯下的。\par

% structure: p0027 source=h2 target=subsection reason=relative-heading-rank; base=h2

\subsection{第九章 在严重欺骗他人时，他也以嘲笑为乐。}

那么，我当时是被什么情感驱使呢？因为那实在是非常可耻的；我有这种情感真是有祸了。但这情感究竟是什么呢？\Quote{谁能明白自己的过错？}我们大笑，因为想到要欺骗那些毫不知情、会对我们所作所为激烈反对的人，我们的心就痒痒的。然而，为什么我对此如此高兴，使得我不愿单独做呢？难道是因为没有人会轻易独自大笑吗？没有人会轻易这样做；但有时，当一个人独处，没有别人在场时，如果有什么非常滑稽的东西呈现于感官或心灵，一阵大笑会征服他。但我独自一个人是不会做那事的——我一个人完全做不出来。我的天主，我灵魂生动的回忆赤裸裸地展现在祢面前——我一个人是不会犯那次偷窃的，在那里，我所偷的东西并不让我快乐，而是偷窃的行为本身；即使独自去做，我也不会那么喜欢，我甚至就不会去做了。哦，太不友善的\OriginalTerm{友谊}{friendship}！你是灵魂神秘的诱惑者，你为了嬉戏和放纵而贪婪地作恶，你渴求他人的损失，却不渴求自己的利益或复仇；但当他们说：\Quote{我们去吧，我们做吧，}我们却因不愿无耻而感到羞耻。\par

% structure: p0029 source=h2 target=subsection reason=relative-heading-rank; base=h2

\subsection{第十章 与天主同在，才有真正的安息和不变的生命。}

谁能解开那扭曲纠结的\OriginalTerm{缠结}{knottiness}呢？它是污秽的。我憎恶回想它。我憎恶看着它。然而，我所渴望的是祢，哦，正义和\OriginalTerm{清白}{innocency}，在一切有德之眼中祢是美丽可爱的，且有着永不令人厌倦的满足！与祢同在才有完美的安息和不变的生命。进入祢内的，便是进入他主的喜乐，\BibleRef{玛25:21}，他将无所畏惧，并在至善者内行至善。我的天主，我离开了祢沉沦下去，在青年时远离了祢——我的支柱，漂泊太远，对自己成了一片不结果实的土地。\par

\SourceRecord{Translated by J.G. Pilkington.；Nicene and Post-Nicene Fathers, First Series；Vol. 1.；Edited by Philip Schaff.；Buffalo, NY: Christian Literature Publishing Co.,；1887.；<http://www.newadvent.org/fathers/110102.htm>.}

% structure: p0001 source=h1 target=section reason=page-title

\section{\WorkTitle{忏悔录}（卷三）}

关于他在\Place{迦太基}度过的十七、十八、十九岁。那时他已修毕学业，陷入放纵情欲的罗网，并堕入\OriginalTerm{摩尼教徒}{Manichæans}的错误中。\par

% structure: p0003 source=h2 target=subsection reason=relative-heading-rank; base=h2

\subsection{被疯狂的爱情所蒙骗，他虽污秽不光彩，却渴望被人视为优雅有礼。}

我来到了\Place{迦太基}，那里有一锅不洁的爱在我周围沸腾。我尚未爱人，却已爱上了爱；由于一种隐秘的贫乏，我因自己无所欠缺而憎恶自己。我四处寻找可恋慕的对象，我钟情于恋爱，却厌恶安稳，厌恶无陷阱的道路。因为我内心缺乏那内在的食粮——你自己，我的天主，尽管这缺乏并未使我感到饥饿；我仍然对不朽的食粮全无渴望，这不是因为我已饱享，而是因为我越空虚，越厌恶它。因此，我的灵魂病得不轻，满身疮痍，悲惨地扑向身外，渴望与感官对象接触而受刺激。然而，这些对象若无灵魂，肯定不会激发爱情。去爱与被爱，对我来说是甜蜜的，尤其在成功享有我所爱的人时更为甜蜜。于是，我用\OriginalTerm{贪欲}{concupiscence}的污秽玷污了友谊的泉源，用情欲的地狱云雾遮蔽了它的光辉；然而，尽管我如此污秽不光彩，却因极度的虚荣，渴望被人视为优雅有礼。我急速堕入我渴望被其网罗的爱情之中。我的天主，我的仁慈，你以无限的慈善，为我在这甜蜜中掺入了多少苦涩！因为我既被人爱，且暗中达到了享有的结合，便喜悦地被烦恼的绳索捆绑，好使我被嫉妒、猜疑、恐惧、愤怒和争吵的灼热铁鞭抽打。\par

% structure: p0005 source=h2 target=subsection reason=relative-heading-rank; base=h2

\subsection{在公共演出中，他因空洞的怜悯而感动。他受困于一种恼人的心灵疾病。}

舞台剧也将我带离了我自己，它们充满了我不幸的写照，也是我欲火的燃料。人为什么喜欢在观看悲惨和愁苦的场景时感到忧伤，而他自己绝对不愿亲受？然而，作为观众，他却渴望从中体验到一种悲伤，他的快乐就在于这悲伤本身。这岂不是一种可怜的疯狂？因为一个人对这些情节越感动，他就越不能摆脱这类情感。不过，当他自己遭受痛苦时，习惯称之为\Quote{苦难}；但当他同情他人时，则称之为\Quote{怜悯}。然而，从虚构和戏剧化的情感中产生的，是哪一种怜悯呢？听众不必解救，只被邀请去悲伤；他越悲伤，就越赞赏这些虚构的演员。如果剧中人物（无论是古代的还是纯属虚构的）的不幸未被表现得使观众感动，他就会厌烦地离开并加以指责；但如果他的情感被触动，他就会全神贯注地坐着，并流下欢喜的眼泪。\par

那么，忧伤也被爱吗？当然人人都愿喜乐？或者，人虽然愿意痛苦，却因能够怜悯而感到欢喜，而怜悯既然不能没有情感，仅仅因此情感才被爱吗？这也来自那友爱的脉络。但它往哪里去？向哪里流？为何流入那条沥青的洪流，沸腾着令人厌恶的巨大情欲浪潮，在其中它被改变和转化，自愿地抛弃了原本属天的清明，遭到了败坏？那么，怜悯就该被弃绝吗？绝不。所以，我们可以有时爱好忧伤。但我的灵魂，在天主、我们祖先的天主、那位应当永远受赞美和颂扬的天主护佑下，你要谨戒不洁。因为我现在并未停止怜悯；但那时在剧场里，当恋人们罪恶地彼此享受时，我与他们同情，尽管这只是在戏中虚构地表演。而当他们彼此失去时，我与他们一同悲伤，好像怜悯他们，然而在这两者中都有快感。如今我对那喜乐于自己邪恶的人，比那因未能获得某种有害的快乐、因丧失某种可悲的幸福而被视为受苦的人，更多感到怜悯。这无疑是更真实的怜悯，但忧伤在其中没有快乐。因为尽管向不幸者表示哀悼的人因其爱德之职而被赞许，但真正同情的人，宁愿没有可悲伤之事。因为如果善愿会怀有恶意（这是不可能的），那么真正衷心地怜悯他人的人，就会希望有不幸者存在，以便怜悯他们。有些忧伤可能是正当的，但不应被爱。主天主啊，你比我们爱灵魂更为纯洁，且不可朽坏地富于怜悯，尽管你不会为任何忧伤所伤。\BibleQuote{谁当得起这些事呢？} \BibleRef{格后 2:16}\par

然而我这可怜的人，那时却爱悲伤，寻求悲伤的对象，如同在他人那虽属装腔和假扮的痛苦中，演员的表演最令我高兴、最强烈地吸引我，那将我感动得流泪的表演。我这不幸福的羊，远离了你的羊栈，不耐烦你的照管，因此染上污秽的疾病，这又有什么奇怪呢？由此我对悲伤的爱好——并不是那种会深深刺痛我的悲伤，因为我不愿遭受我所爱观看的那些事，而是那种在听到虚构的故事时，仅在表面轻轻触及的悲伤；随之而来的，如同被涂毒的指甲刺过，接着是灼烧、肿胀、腐烂和可怕的朽坏。这就是我的生活！可是，我的天主，那是生活吗？\par

% structure: p0009 source=h2 target=subsection reason=relative-heading-rank; base=h2

\subsection{即便在教会里，他也不克制他的欲望。在修辞学校里，他憎恶颠覆者的行为。}

祢信实的仁慈远远地庇护着我。我何等耗尽自己于可耻的罪恶中，随着渎神的好奇心，它使我离弃祢，拖我进入诡诈的深渊，和魔鬼的迷惑服从，我向它们献上我的恶行，在这一切中祢鞭笞了我！我竟敢，甚至当祢的教会围墙内举行祢庄严的礼仪时，去渴望并图谋足以获取死亡果实的事务；为此祢以沉重的惩罚责罚我，但与我的过错相比却算不得什么，噢，我最大的仁慈，我的天主，我躲避那些可怕伤害的避难所，我在其中傲慢地昂首徘徊，离祢越来越远，爱慕我自己的道路，而非祢的道路——爱慕一种流浪的自由。\par

那些被视作光荣的学业，也指向法庭；要想在其中出类拔萃，越是狡诈，就越受赞扬。人的盲目竟至于此，他们甚至以自己的盲目为荣。而当时我已是修辞学校的领头人，我为此洋洋得意，傲慢自大，尽管，主啊，祢知道，我较为沉稳，也全然远离那些\Quote{\OriginalTerm{颠覆者}{subverters}}的颠覆行为（因为这个愚蠢而魔鬼般的名称被奉为豪侠的标志），我生活在他们中间，却带着一种厚颜的羞愧，因为我并不像他们那样。我与他们同行，有时也喜爱他们的友谊，尽管我一向憎恶他们的行为，也就是他们的\Quote{\OriginalTerm{颠覆}{subverting}}：他们以此狂妄地攻击陌生人的端庄，用无端的嘲弄去搅扰，从而满足他们恶作剧的快乐。没有什么比这些行为更近似魔鬼的行径了。那么，他们还能被更真实地称作什么呢？只能是\Quote{\OriginalTerm{颠覆者}{subverters}}——他们自己先被颠覆，且全然败坏了——暗中被那些欺骗之神嘲笑和诱惑，他们自己却乐于嘲弄和欺骗他人。\par

% structure: p0012 source=h2 target=subsection reason=relative-heading-rank; base=h2

\subsection{第四章 在他十九岁时（他的父亲两年前去世），他被\Person{西塞罗}的\WorkTitle{荷尔顿西乌斯}引向\Quote{哲学}、天主和更好的思维方式。}

在这些人群中，在我生命那个不稳定的时期，我学习雄辩术的书籍，在其中我渴望出类拔萃，出于一种可憎而自大的目的，甚至是一种对人世虚荣的恋慕。在常规学习中，我偶然读到\Person{西塞罗}的一本书，他的文辞（尽管不是他的心）几乎所有人都仰慕。他的这本书包含对哲学的劝诫，名为\WorkTitle{荷尔顿西乌斯}。这本书，确实，改变了我的情感，将我的祈祷转向祢自己，主啊，并使我有别的希望和渴望。一切虚妄的希望忽然对我变得一文不值；我以一种难以置信的心灵热度渴望不朽的智慧，并开始起身，\BibleRef{路 15:18}，以便我回归于祢。所以，不是为了改善我的文辞——那时我十九岁，父亲已去世两年，我似乎是在用母亲的钱财来购买它——不是为了改善我的文辞我才求助于那本书；也不是因它的风格说服了我，而是因它的内容。\par

那时我是多么热切，我的天主，多么热切渴望从尘世飞向祢！然而我不知道祢将如何处理我。因为智慧与祢同在。在希腊文中，对智慧的爱被称为\Quote{\OriginalTerm{哲学}{philosophy}}，正是那本书激起了我的热情。有些人利用哲学之名来误导人，用这伟大、诱人且可敬的名字来掩饰和装饰他们自己的错谬。而几乎所有在那个时代或更早如此行的人，都在那本书中受到了指责和揭露。那里也揭示了祢圣神透过祢善良而虔诚的仆人所发出的最有益的劝告：\BibleQuote{你们要小心，免得有人以哲学，以虚伪的妄言，按照人的传授，根据世俗的训导，而不按基督，把你们勾引了去。因为在天主内，真实地住有整个圆满的天主性。} \BibleRef{哥 2:8} 而那时（正如祢，我心灵之光，所知道的）宗徒的话我尚未知晓，我仅因这劝告受到激励、点燃和燃烧，去爱慕、寻求、获得、持守和拥抱智慧本身，无论它是怎样的，而非这个或那个学派；唯一阻挠我这如此热切之心的，是其中没有基督的名字。因为这个名字，按照祢的仁慈，主啊，这我的救主、祢圣子的名字，我温柔的心自幼便虔诚地吸入，深深蕴藏，甚至与我母亲的乳汁一同吸收；凡没有这个名字的任何东西，无论多么博学、精美、真实，都无法完全抓住我。\par

% structure: p0015 source=h2 target=subsection reason=relative-heading-rank; base=h2

\subsection{第五章 他拒绝接受圣经，认为它太简单，无法与\Person{图利}的尊严相比。}

因此，我决定专注我的心于圣经，以便看看它们是怎样的。看，我看到一种不被骄傲的人所领会，也不向儿童揭示的东西，而是卑微者近前，崇高者前行，却笼罩在奥秘之中；而我并非属于能够进入其内，或肯低头跟随其步伐的那类人。因为我当时转向圣经时的感受，并非像我现在说话时这样，而是它们在我看来不配与\Person{图利}的尊严相比；因为我膨胀的骄傲拒绝它们的风格，我机智的锐利也无法刺透它们的内在含义。然而，它们确实是会使小孩子成长的书；但我耻于做小孩子，而且，由于骄傲自大，我把自己看作一个大人。\par

% structure: p0017 source=h2 target=subsection reason=relative-heading-rank; base=h2

\subsection{第六章 因自己的过错而受骗，他陷入了\OriginalTerm{摩尼教徒}{Manichæans}的错谬中，他们夸耀对天主的真知识和透彻的探究。}

10. 因此，我陷入了一群人中间，他们骄傲自大，极尽肉欲而喋喋不休；他们的口是魔鬼的陷阱——那粘鸟胶是将祢的名、我们的主耶稣基督，以及那位\OriginalTerm{护慰者}{Paraclete}、\OriginalTerm{圣神}{Holy Ghost}、\OriginalTerm{护慰者}{the Comforter}的音节混合而成的。这些名字从不离开他们的口，但只限于声音和舌头的喧响，因为他们的心缺乏真理。他们仍呼喊：\Quote{真理，真理，}并就此事向我谈论许多，\BibleQuote{然而真理不在他们内}；\BibleRef{若一 2:4}。他们不仅对祢撒谎——祢确实是真理——也对这世界的元素，祢的受造物撒谎。而我，实在应当为了祢——我的圣父，至善者，一切美之美的爱——而忽略那些哲学家，即使他们论及这些事时说了真理。啊，真理，真理！那时，我灵魂的骨髓是如何深切地为祢而喘息啊，当他们经常地、以各种方式、在众多而庞大的书籍中向我吹嘘祢的名时，虽然那不过是一种声音！这些便是他们摆在我面前的菜肴，我渴望着祢，他们却以太阳和月亮代替祢，呈上祢美丽的工程——但仍只是祢的工程，不是祢自己，甚至也不是祢的初造之物。因为在这些有形体的工程之前，先有祢的灵性工程，它们虽是天上光明的。但我所饥渴的，甚至不是那些祢的初造物，而是祢自己，真理，\BibleQuote{在祂内没有变化或转动的阴影}；\BibleRef{雅 1:17}。而他们仍然在那些菜肴中给我端上灿烂的幻象，与其这样，不如去爱这太阳本身（至少对我们的视觉是真实的），也不要去爱那些通过眼睛欺骗心灵的幻影。然而，因为我以为它们就是祢，我吞食了它们；但并不贪婪，因为祢尝在我的口中不是祢自己，祢并非这些空虚的虚构；我也没有因它们得到滋养，反而更加困乏。睡梦中的食物看起来像我们醒时的食物；然而睡觉的人并未因此得到滋养，因为他们睡着了。但这些事一点不像祢，正如祢现在对我所说的，因为那些都是形体的幻象，虚假的身体，比起这些真实的身体——无论是天上的还是地上的，我们以肉体的视觉所看到的——要确实得多。这些东西，连野兽和飞鸟也和我们一样能见到，而它们比我们想象它们时更确实。而且，我们想象它们，比通过它们去设想其他那些更伟大、无限而不存在的形体，要确实得多。那时我便以这样的空壳为食，却并未得食。可是祢，我的爱，我为了寻求祢而衰败，以便我能刚强，祢既不是我们见到的那些形体，虽在天上，也不是我们在那里见不到的事物；因为祢创造了它们，也并不把它们算作祢最伟大的工程。那么，祢离我那些幻象有多远啊，那些绝非实体的形体之幻象，比起这些幻象，那些存在之物的影像更为确实，而存在之物本身更为确实，然而祢却不是它们；甚至那灵魂，即形体之生命，祢也不是。那么，形体之生命比形体本身更美好、更确实。而祢是灵魂的生命，生命中的生命，祢在自己内拥有生命；祢永不改变，啊，我灵魂的生命！\par

11. 那么，那时祢对我而言在哪里呢？离我有多远呢？我确实远离祢漂泊而去，甚至被拒于那些猪所吃的豆荚之外，而我曾以豆荚喂养它们。因为语法学家和诗人们的神话故事，比起这些罗网，要好得多！因为诗文、诗歌，以及那飞翔的美狄亚，确实比这些人的五种元素更有益——这些元素被描绘得五花八门，以对应五个黑暗的洞穴，其实无一存在，却扼杀信者。因为诗文和诗歌我可以转化为真正的食粮，但\Quote{飞翔的美狄亚}，即使我吟唱，我也不坚持其为真；即使我听到人吟唱，我也不信其为真；但那些东西我却真信了。祸哉，祸哉，我是怎样一步步被拖下\BibleQuote{地狱的深处}的啊！\BibleRef{箴 9:18}——我因缺乏真理而劳苦烦乱，当我寻找祢、我的天主时——我向祢忏悔，在我尚未忏悔时祢已怜恤了我——寻找祢，不是按照理智，即祢愿我以此超越禽兽的那种理智，而是按照肉体的感觉！祢比我最内在的部分更内在于我，比我最高的部分更高。我遇到了那个大胆的女人，她\BibleQuote{是愚蠢的，一无所知，}\BibleRef{箴 9:13} 那是撒罗满的隐语，她\BibleQuote{坐在家门口的座位上}，并说：\BibleQuote{偷来的水是甜的，暗中吃的饼是美味的。}这女人引诱了我，因为她发现我的灵魂离开了自己的内室，居住在我肉体的眼中，并思想着那样我通过眼睛吞食的食物。\par

% structure: p0020 source=h2 target=subsection reason=relative-heading-rank; base=h2

\subsection{第七章 他抨击摩尼教徒关于恶、天主及圣祖们正义的学说}

12. 因为我对那真实存在者一无所知，并且几乎是被强迫地去支持那些愚蠢的骗子，当他们问我：\Quote{恶从何来？}——以及，\Quote{天主是否受形体限制，是否有头发和指甲？}——以及，\Quote{那些同时拥有多个妻子、屠杀过人、并宰杀动物献祭的人，能算为义人吗？}\BibleRef{列上 18:40}。对于这些事，我因无知而深感不安，当我从真理退却时，我自以为正走向真理；因为我尚不知恶无非是\OriginalTerm{善的缺乏}{privatio boni}，直至最后完全不存在；我怎能看见这点呢？我眼之所见不出于形体，我心灵之所见不出于幻象！我也不知天主是神\BibleRef{若 4:24}，并不是具有长度和宽度延伸的肢体者，也不是具有体积的存在；因为任何体积在部分中小于整体，而如果它是无限的，它在被一定空间限制的部分中也一定小于其无限性；它也不能如神、如天主那样无所不在。至于在我们内有什么使我们肖似天主，并能在圣经中正当地被称为\BibleQuote{天主的肖像}，我完全无知。\par

13. 我也不认识那\OriginalTerm{真正的内在义德}{true inner righteousness}，它不随从习俗判断，而是依据全能天主的至善法律；这法律使各地各时的风俗适应其地方与时代，而它自身却恒常不变、处处如一，绝不因地而异。亚巴郎、依撒格、雅各伯、梅瑟、达味，以及所有蒙天主亲口称赞的人，都依照这法律而成为\OriginalTerm{义人}{righteous}\BibleRef{希 11:8}，但愚昧之人却凭人的判断，断定他们不义\BibleRef{格前 4:3}，并以自己风俗的狭小尺度，衡量全人类的习俗。如同在军械库中，有人不知各种护具当配何种肢体，竟把胫甲戴在头上，或用头盔当靴穿，而后抱怨不合；又如某日午后禁止营业，有人因不能像午前那样合法买卖而恼怒；或见一仆人在屋中手持某物，而这物管家不得触碰，或在马厩后作某事，而这在餐厅内是禁止的，便愤愤不平，以为在同一屋宇、同一家庭中，不当各处规矩不一。这些人正属于此类：他们无法容忍听说昔时对义人合法之事，今时已不合法；也不能接受天主因某些暂时的理由，命令前人作此，而今人作另一事，二者却同遵一义德。然而，他们在同一人、同一天、同一屋中，却能见到不同肢体适宜不同之事，先前合法的，过时便成不合法——在某处许可或命令的事，若行于他处，便公正地受到禁止和惩罚。正义岂是多变且变化的吗？绝非如此，而是其统驭的时代各不相同，因它们本是时代。世人岁月短少\BibleRef{约 14:1}，凭自己的识力，无法调和以往时代和其他民族（他们对此并无经验）的缘由，与他们所经验的缘由；虽然在同一身体、同一天或同一家庭中，他们很容易看出什么适合每个肢体、时辰、部分和人物——对某些事他们予以排斥，对另一些事则服从。\par

14. 这些事我当时既不知晓，也未观察。它们从四面呈现于眼前，我却视而不见。我作诗时，不允许在任何地方随意安放每一音步，而是在一种格律中这样做，在另一种中那样做，甚至在任一诗行中，同一音步也不可处处使用。然而，我赖以作诗的艺术本身，并非对这些不同情形有不同的原理，而是总括一切于一。我仍看不到：那善人与圣人遵从的\OriginalTerm{义德}{righteousness}，何等卓越而崇高地将天主所命的一切总括为一，毫不变化，尽管在不同的时代，它并不一次颁布所有事，而是按各时代分配并命令合宜的事。我既盲目，便责怪那些虔诚的先祖，不但因他们遵照天主的命令和启示利用当前事物，也因他们预告未来之事，正如天主所启示的。\par

% structure: p0024 source=h2 target=subsection reason=relative-heading-rank; base=h2

\subsection{第八章 他同样驳斥关于罪过理由的论点}

15. 无论在何时何地，人若以全心、全灵、全意爱天主，并爱近人如己，这岂能是不义的事？因此，那些违反自然的罪过，在任何时地都应受憎恶和惩罚；索多玛人的罪即属此类：即使所有民族都犯了这样的罪，他们也都按天主的法律同样被定为有罪，因为天主的法律并未使人得以如此彼此滥用。因为，当天主所创的\OriginalTerm{本性}{nature}被邪欲的乖张所玷污时，我们与天主之间应有的\OriginalTerm{共融}{fellowship}也就被破坏了。至于那些违反人间习俗的罪过，则当依照各处通行的习俗加以避免；如此，任何城邦或民族由习俗或法律所确认的协议，便不致因任何人（无论公民或外邦人）的任意妄为而遭破坏。因为任何与整体不协调的部分都是不合宜的。但当天主命令做某事，与任何民族的风俗或协约相悖时，即使该民族从未做过，也必须执行；若已中断，则当恢复；若从未建立，则当建立。因为，一位国王在其治下国家，能命令他本人及前人都未曾命令之事，而服从他并不被视为有损公益——甚至，若不顺从，反有损害（因为服从君主是人类社会普遍的契约）——何况我们更当毫不迟疑地服从统治一切受造物的天主！在人类社会权柄中，较大的权柄先于较小的权柄得到服从，那么更应当超越一切地服从天主。\par

16. 同样，在暴力行为中，若存在伤害的欲望，无论是通过侮辱还是加害；而这两者又或出于报复，如仇敌相待；或为了从他人处谋取利益，如同强盗之于行人；或为避免某种祸患，如同有人惧怕他人；或出于嫉妒，如不幸者对幸福者；或如某事一帆风顺者，对那恐怕会与他平起平坐、或因对方之平起平坐而忧伤的人；或纯粹以他人的痛苦为乐，如观赏角斗士者，或嘲笑讥讽他人者。这些便是从肉身的贪欲、眼目的贪欲和权势的贪欲中生发的主要罪恶，无论是单独一处，还是两相结合，甚或三者同时发作。如此，人们的生活便与那三与七，即那\Quote{十弦的}琴瑟——你的十诫——背道而驰，\OriginalTerm{至高至甘饴的天主}{God most high and most sweet}啊。然而，对于不可玷污的你，能有什么污秽的冒犯呢？对于不可伤害的你，又能有什么暴力的行为呢？但你却报复那些人对自身所犯的恶行，况且，他们得罪你时，也是邪恶地对待自己的灵魂；罪恶乃是自欺，或通过败坏、扭曲你创造并统序的本性，或通过无节制地使用许可之物，或在于对违禁之物的\BibleQuote{燃烧}，即那种违反本性的用途；\BibleRef{罗 1:24} 或在被定罪后，心怀并声讨你，用脚踢刺；\BibleRef{宗 9:5} 或冲破人类社会的藩篱，胆大妄为地根据个人好恶，在私自结党分裂中欢欣。凡此种种，皆发生在人离弃你之时，\OriginalTerm{生命之泉}{Fountain of Life}啊，你是宇宙唯一真实的创造者与统治者；人凭私意骄矜，从中拣选并贪恋任何一桩虚妄之事。如此，我们便以谦卑的虔敬归向你；你涤除我们的恶习，怜悯那些向你告明之人的罪恶，并\Quote{垂听囚徒的呻吟}，松解我们自造自铸的枷锁，只要我们不在你面前高举虚假自由的角——因贪婪而丧失一切，因爱我们自身的私善胜过爱你，万善之善。\par

% structure: p0027 source=h2 target=subsection reason=relative-heading-rank; base=h2

\subsection{第9章 天主与世人对人类暴力行为的判断，迥然不同}

17. 然而，在如此多污秽、暴力和罪恶的过犯中，也有那些总体而言正在进步之人的罪；按正确判断并依循成全之准则的人看来，这些罪虽遭到责难，却也因有望结实而受褒扬，犹如那生长中的谷物的青青嫩苗。还有一些行径，虽似污秽或暴力的过犯，却并非罪过，因为它们既不冒犯你——我们的主天主，也不违背社会习俗：例如，人因时制宜地预备生活所需之物，我们却无法确定是因为贪求占有；或者合法当局为纠改而施行的惩罚，我们却无法确定是因为贪求伤害。因此，许多在人眼中不被赞同的行为，却因你的见证而蒙悦纳；许多受人称赞的行为，有你作证，却受谴责；因为行为的表相、行事者的心态以及那时代的隐秘急需，往往各不相同。可是，当你出人意料地命令一件非同寻常、难以想象的事情时——是的，即便你从前曾予以禁止，且暂时将命令的缘由秘而不宣，甚至该命令相反某些人类社会的规条，谁能怀疑是否当行呢？因为那服事你的社会才是正义的。然而，明了你命令的人是有福的！因为你的仆人所作的一切，或是为显示当时所必要之事，或是为预示将来的事。\par

% structure: p0029 source=h2 target=subsection reason=relative-heading-rank; base=h2

\subsection{第10章 他斥责摩尼教徒有关地上果实的无稽之谈}

18. 由于对此事一无所知，我竟嘲笑你那些圣善的仆人和先知。我嘲笑他们，又从中得了什么益处，无非是被你嘲笑罢了，不知不觉、一点一滴地被引向那些愚妄之事，竟致相信，无花果树被采摘时会哭泣，母树也会流下乳般的泪。然而，那颗无花果若不是由于自身的邪恶，而是由于别人的邪恶被采摘，若被某位\Quote{圣者}吃了，融入他的五内，他便会从中呼出天使来；是的，他在祈祷中定会呻吟吐纳出\OriginalTerm{天主的微粒}{particles of God}，而这些至尊真天主的微粒，若非借着某位\Quote{被选的圣者}的牙齿与肚腹获得释放，本会永远束缚在那无花果中！我这可怜人，竟相信对地上的果实，比对人——它们正是为人而被造——更应慈悲；因为若一个饥饿的人——又不是摩尼教徒——乞求一点点果实，所给他的那一口，便仿佛被判了极刑一般。\par

% structure: p0031 source=h2 target=subsection reason=relative-heading-rank; base=h2

\subsection{第11章 他提到母亲的眼泪，以及天主赐给她关于儿子的那个难忘的梦}

19. 祢从高天伸出祢的手，把我灵魂从深重的黑暗中拉出来，那时我的母亲，祢的忠信者，为我向祢痛哭，胜过母亲们为子女肉身的死亡所惯常哭泣的。因为她看出我死于那从祢得来的信德和心神，祢就俯听了她，主啊。祢俯听了她，没有轻视她的眼泪；当她的眼泪倾流而下，在她祈祷的每一处，浇灌了脚下的土地时，祢确已俯听了她。因为祢用来安慰她的那个梦是从何而来的呢？那个梦使她准许我与她同住，并在家中同桌进餐，而这正是她先前所回避的，因她憎恨并厌恶我错误的亵渎。因为她看见自己站在一条木尺上，有一个光明的青年向她走来，欢欣喜笑地望着她，而她却忧愁悲伤，俯身哀恸。但他询问她为何愁苦，日日哭泣（他如常所为，愿教导人，而不愿受教），她回答说，是在哀悼我的丧亡，他便叫她安心，并要她注视观看：\Quote{她所在之处，我也在那里。}当她看时，看见我靠近她站在同一木尺上。这事从何而来，岂不是因为祢的耳朵倾向她的心？哦，至善全能者，祢如此眷顾我们每一个人，仿佛只眷顾他一人；又如此眷顾众人，仿佛他们只是一人！\par

20. 这事又从何而来呢？当她向我叙述这异梦，我试图这样解释：\Quote{她倒不必绝望，有朝一日会成为我这样的人；}她立即毫不犹豫地回答说：\Quote{不；因为告诉我的不是\InnerQuote{他在哪里，你也将在那里}，而是\InnerQuote{你在哪里，他也将在那里}}？我向祢忏悔，主啊，尽我记忆所及（我常谈起这事），祢通过我警醒的母亲所做的回答——她并未因我虚妄解释的似是而非而不安，瞬间就看清了当看的事，而我在她说话之前确实未曾领会——这件事当时比那异梦本身更感动我；那异梦为减轻她眼前的忧虑，早在许久以前就预言了这位虔诚妇女将要实现、却要在许久之后才实现的福乐。因为将近九年的时间过去了，我仍在那深坑的淤泥和谎言的黑暗中打滚，常常努力奋起，却摔得更重。然而那位贞洁、虔诚、节制的寡妇（祢所喜爱的那种），虽然希望倍增，但哭泣哀悼的热诚丝毫不减，在一切祈祷的时刻，从未停止向祢痛哭我的境况。她的祈祷进入祢的临在，而祢仍然任我反复卷入那黑暗之中。\par

% structure: p0034 source=h2 target=subsection reason=relative-heading-rank; base=h2

\subsection{母亲为儿子的皈依求问主教时，主教出色的回答。}

21. 同时，祢赐予了她另一个答复，我如今还记得；因为我急于奔向那些更强烈催促我向祢忏悔的事，便略过了许多，也有许多我已不记得。那时祢通过祢的一位司铎，就是某位主教，赐予了她另一个答复；他在祢的教会中成长，精通祢的书籍。当这位妇女恳求他肯与我交谈，驳斥我的谬误，破除我的恶习，教导我善道（这是他遇到合适的人时常做的事）时，他非常明智地拒绝了，正如我事后明白的。因为他回答说，我仍不堪受教，被那种异端的新奇所鼓胀，而且我已经用一些刁难的问题困惑了许多无经验的人，正如她告诉他的。他说：\Quote{暂时由他去吧，只为他祈求天主；他自己会通过阅读发现那是怎样的错误，其邪恶有多大。}同时，他也向她透露，自己年幼时，曾被他那误入歧途的母亲交给\OriginalTerm{摩尼教徒}{Manichæans}，不仅阅读过，甚至几乎誊写过他们所有的书籍，并且（无需任何人的论证或证明）看出那教派应如何远避，便远避了它。他说了这话之后，她仍不满足，反而更恳切地流泪祈求，要他来见我与交谈；他因她的纠缠而稍感烦恼，便大声说：\Quote{去吧，愿天主降福你，因为流这些眼泪的儿子决不会丧亡。}她接受了这个答复（如她在与我的谈话中常常提起的），仿佛这是从天而来的声音。\par

\SourceRecord{Translated by J.G. Pilkington.；Nicene and Post-Nicene Fathers, First Series；Vol. 1.；Edited by Philip Schaff.；Buffalo, NY: Christian Literature Publishing Co.,；1887.；<http://www.newadvent.org/fathers/110103.htm>.}

% structure: p0001 source=h1 target=section reason=page-title

\section{\WorkTitle{忏悔录（第四卷）}}

接着是他从十九岁开始的九年时期，在此期间他失去了一位朋友，皈依了\OriginalTerm{摩尼教徒}{Manichæans}，并撰写了关于美善与相称的著作，还发表了一部关于自由技艺以及亚里士多德范畴论的著作。\par

% structure: p0003 source=h2 target=subsection reason=relative-heading-rank; base=h2

\subsection{第一章 论及他受骗且骗人的至不幸时期，并论及嘲笑他认罪的人。}

1. 在这九年的时间里——从我十九岁到二十八岁——我们彼此诱惑且诱惑人，欺骗且被欺骗，陷入各种情欲：公开地，借着他们所谓\Quote{自由}的学问；私下地，借着那名为宗教的虚假。在这里骄傲，在那里迷信，处处虚妄！在这里，追逐着庸俗荣名的虚空，直至剧场喝采、诗歌竞赛、为草冠的无谓争执，以及演出之愚行与贪欲之无度。在那里，试图借由给那些被称为\Quote{\OriginalTerm{选民}{elect}}与\Quote{\OriginalTerm{圣者}{holy}}的人送去食物，以期从我们的败坏中得洁净——在他们的胃脏工场中，这些食物将为我们制造出天使和神明，好使我们借之获得解救——我热切地随从这些事，并与友人们一同实践：皆是由我欺骗，且与我一同被欺骗。任那高傲之辈，以及尚未因祢而被救援地倾覆和击打的人，我的天主啊，讥笑我吧；然而我仍要向祢承认我的耻辱，以赞美祢。我恳求祢，容忍我，并赐我恩宠，在我现在的记忆中回溯我过往错误的周折，并\Quote{向祢献上感恩的祭献}。离开了祢，我对自己而言，若不是一个导向自己堕落的向导，又是什么呢？或者即便在我最好的时候，我不过是\BibleQuote{吸吮祢的奶水}\BibleRef{伯前 2:2}，并以\BibleQuote{祢那不会朽坏的食物为生}\BibleRef{若 6:27}？然而一个人算什么呢，他终归只是个人。那么，让强壮者和权贵们讥笑我们吧，但我们这些\Quote{贫穷和困乏的}人要向祢认罪。\par

% structure: p0005 source=h2 target=subsection reason=relative-heading-rank; base=h2

\subsection{第二章 他教授他唯一喜爱的修辞学，并蔑视那答应让他获胜的占卜者。}

2. 在那些年间，我教授修辞术，因被贪欲所胜，便将那用于胜人的辩才出售。然而——主啊，祢知道——我更愿意收诚实的（如他们所认为的）学生；对于这些学生，我不用诡诈，却教他们诡诈之术，并不是要用来陷害无辜者的生命，有事时却可用来为有罪者的生命辩护。我的天主啊，祢从远处看见我在那条滑路上跌撞，且在浓烟中发出些许忠信的火花，这便是我在引导那些爱虚荣、寻求谎言的人时——我既是他们的同伴——所表现出来的。那些年间，我有过一个女人（我不是按照所谓的合法婚姻认识她，而是由于我那没有明智的放荡激情而发现的），但仅此一人，并且我甚至忠贞于她；在我自己的经验中，我确实发现了结婚盟约的约束（为了生育后代而缔结）与情欲爱情的契约之间的区别——在后一种情况下，子女违背父母之意而生，尽管出生以后，他们便会迫使父母产生爱心。\par

3. 我也记得，当我决定去角逐一个戏剧奖项时，有一个占卜者问我，要给他什么报酬，以便我能获胜；可是，我因憎嫌这类污秽的秘术，便回答说：\Quote{即使那冠冕由不毁的黄金制成，我也不愿为赢得它而害死一只苍蝇。}因为他原本要在祭献中杀某些活物，并借着这些祭仪邀请魔鬼来援助我。然而，我拒绝这恶事，也并非出于对我心的天主祢的纯洁爱情；因为我还不知道如何爱祢，也不知道如何想象任何超越形体明亮之美的事物。一个灵魂若渴慕这类虚构之事，难道不是对祢犯下奸淫，信赖虚假的东西，并\BibleQuote{牧养风}吗？\BibleRef{欧 12:1}然而，我确实不愿让人为了我向魔鬼献祭，虽然我本人却借着那种迷信向他们献上祭品。因为牧养风不就是牧养他们吗？也就是说，借着我们的放荡成为他们的享乐与笑柄。\par

% structure: p0008 source=h2 target=subsection reason=relative-heading-rank; base=h2

\subsection{第三章 即便是最有经验的人也无法使他相信他所热衷的占星术是虚妄的。}

4. 那时，对于那些被称作\OriginalTerm{星命学家}{Mathematicians}的骗子，我毫不迟疑地前去请教，因为他们不用祭献，也不呼求任何神灵来行占卜；这种技艺正是基督信仰和真虔敬所摈弃和谴责的。向祢认罪，并说\Quote{求祢怜悯我，医治我的灵魂，因为我得罪了祢}，这是美好的；而不是滥用了祢的慈善作为犯罪的许可，却是要记住主的话：\BibleQuote{看哪，你已痊愈，不要再犯罪，免得遭遇更不幸的事。}\BibleRef{若 5:14}他们竭力破坏这一切救人灵魂的劝告，说什么\Quote{你犯罪的原因不可免地在天上命定了}，又说\Quote{这乃是\OriginalTerm{金星}{Venus}，或\OriginalTerm{土星}{Saturn}，或\OriginalTerm{火星}{Mars}造成的}；如此，人为的血肉之躯和骄傲的朽腐便可无罪，而那创造和安排天地星辰的造主却要担当罪责。而这造主不是别人，正是祢，我们的天主，甘饴和正义的泉源，祢\Quote{按各人的行为报应他}，且不轻看\Quote{一颗破碎和谦卑的心}！\par

5. 在那些日子里，有一位智者，精通医术，且享有盛名。他曾以自己曾任总督的手，将竞技的花环戴在我混乱的头上；然而，他不是作为医生；因为这病惟有祢医治，祢拒绝骄傲者，赐\OriginalTerm{恩宠}{grace}给谦卑者。但是，祢难道藉着那位老人辜负了我，还是停止了医治我的灵魂？因为我与他变得熟识，专注而恒常地倾听他的谈话（虽然用词简单，却充满活力、生气与诚挚），当他从我的谈话中察觉我沉迷于占星家的书籍时，他便以友善和慈父般的态度劝我将它们丢弃，不要将用在有益事物上所需的精力和劳苦，空耗在这些虚幻之事上；他说他自己早年也曾研究这门技艺，想以此为业谋生，既然他能理解\Person{希波克拉底}，本也能很快掌握这个，但他放弃了，转而从事医学，唯一的原因是他发现它全属虚妄，而他作为正直的人，不愿靠欺骗他人维生。他说：\Quote{但你却有修辞学可以维生，所以你是出于自愿而非受迫从事此道——那么，你就更应当信任我了，因为我曾努力精通它，本想单靠它谋生。}当我问他为何它能预言许多真实的事时，他尽可能回答我：\Quote{那是由弥漫于整个自然秩序中的偶发之力所导致的。因为，若有人偶然翻开某位诗人的书页，而诗人所咏原意大相径庭，却常有诗句奇妙地恰合当前之事，那就不足为奇，}他继续说，\Quote{若从人的灵魂中，藉着某种更高的本能，并不自知其内部所发生之事，而由偶然而非技艺给出一个答案，恰与询问者的事情和行动相合，那也不足为怪。}\par

6. 因此，的确，无论是藉着他，还是通过他，祢都在关照着我。祢在我的记忆中绘出了我后来可以自己探寻的东西。但在那时，无论是他，还是我最亲爱的\Person{内布里迪乌斯}——一位最为良善且最为审慎的青年，他嘲笑那一整套占卜之术——都无法说服我放弃它，因为那些作者的权威更影响着我；而且，我还没有找到任何确切的证据——如我所寻求的——能够毫无疑问地表明，那些被咨询者所真实预言的事，是由于偶然或巧合，而非星象家的技艺。\par

% structure: p0012 source=h2 target=subsection reason=relative-heading-rank; base=h2

\subsection{因痛哭亡友而深感悲痛，他为自己寻求安慰。}

7. 在那些年里，当我最初在家乡开始教授修辞学时，我结交了一位非常亲爱的朋友，因同学之谊，与我同岁，正如我一样，刚刚步入青春年华。他与我自幼一同长大，我们既是同窗，又是玩伴。但是，他那时并非我的朋友，后来也不是，按照真正友谊而言；因为真正的友谊，只存在于祢所结合的人之间，他们因着那藉着所赐给我们的\OriginalTerm{圣神}{Holy Ghost}，倾注在我们心中的爱，而依附于祢。\BibleRef{罗 5:5} 然而，这份情谊由于相似的志趣热忱而成熟，太过甜美。因为我曾使他偏离了真正的信仰（对于这信仰，他作为一个青年，尚未健全而彻底地掌握），转向那些迷信而有害的寓言，正是这些寓言使我母亲为我哀伤。这个人的心思如今与我一同迷失，而我的灵魂没有他便无法生存。但是请看，祢紧随祢的逃亡者之后——祢是既为报复的天主，又为慈悲的源泉，以奇妙的方式将我们转向祢自己。当我们的友谊还未满一整年时，祢就将那人从今生带走，他对我来说，比我那生命中的所有甜蜜更为甜蜜。\par

8. \Quote{谁能述说祢一切的赞美}，就是他单在自己身上所经历的呢？我的天主，那时祢做了什么？祢的判断何其不可测度！因为，他患热病，沉重痛苦，长时间昏迷不醒，汗如雨下，濒临死亡，所有人都对他康复绝望时，他在不知情中领受了\OriginalTerm{圣洗}{baptism}；而我那时并不在意，还认为他的灵魂所保留的，将更多是从我这里所吸收的，而非对他无知觉的身体所做之事。然而，事情远非如此，因为他苏醒并康复了。我一能与他交谈（因为我一能就立刻与他交谈，因为我从未离开过他，我们彼此黏得太紧了），我试图和他开玩笑，仿佛他也会与我一起嘲笑那次他在心神和感官丧失时所领受的、但现在已知道曾领受了的圣洗。但他对我颤栗，如同我是他的仇敌；并且以一种惊人而意想不到的坦率劝诫我，如果我还希望继续做他的朋友，就不要再以这种方式对他说话。我惊惶困惑，隐藏了所有的情绪，直到他痊愈，身体健壮到足以让我随心所欲地应对他。但他却被从我的狂乱中带走，好使他与祢同在，为我的安慰而被保存。几天后，在我外出期间，他热病复发，便去世了。\par

9. 面对这忧伤，我的心全然黑暗，我所望见的一切都是死亡。我的家乡对我是种折磨，我父亲的家是种奇特的痛苦；我与他所共享的一切，因缺少了他，都变成了可怕的折磨。我的双眼四处寻找他，却寻不见；我憎恶一切地方，因为他不在其中；它们如今再也不能像他活着而外出时那样，对我说：\Quote{看，他来了。}我成了自己的一个大谜，我问我的灵魂为何如此忧愁，为何如此极度地扰乱我；但她不知如何回答我。我若说：\Quote{仰望天主，}她也极其合宜地不听从我；因为她所失去的那位最亲爱的朋友，作为一个真实的人，比她受命去仰望的那个幻影，更为真实，更为美好。惟有眼泪使我感到甘饴，它们在我最深的爱意中取代了我的朋友。\par

% structure: p0016 source=h2 target=subsection reason=relative-heading-rank; base=h2

\subsection{为何哭泣对不幸者是愉悦的。}

10. 主啊，如今这些事都已过去，时间已经治愈了我的创伤。愿我向祢——真理——学习，将我心灵的耳贴近祢的口，好让祢告诉我，为何哭泣对于不幸的人竟是如此甘甜。祢虽处处临在，难道却将我们的痛苦远远抛离？祢安居于己内，我们却因种种考验而骚动不安；然而，若非我们在祢耳中哭泣，我们便毫无希望存留。那么，从生命的苦涩中，从呻吟、眼泪、叹息和哀哭中，为何能摘取如此甘甜的果实？是因为希望祢垂听我们，便使之变甜吗？这在祈祷中确实如此，因为在祈祷中，有着亲近祢的渴望。但是否也在为失物而悲痛，并被当时压倒我的哀伤中呢？因为我既没有希望他复活，也没有用眼泪去寻求此点；我只是悲伤哭泣，因为我可怜，失去了我的欢乐。或者，哭泣本是一件苦事，是因为厌恶我们从前所享受的事物，甚至就在我们厌恶它们之时，它却给我们带来快感？\par

% structure: p0018 source=h2 target=subsection reason=relative-heading-rank; base=h2

\subsection{第六章 他的朋友被死亡夺去，他想象自己只剩下半个。}

11. 但我为何讲这些事呢？因为现在不是提问的时候，而是向祢忏悔的时候。我那时是悲惨的，每一个灵魂若被可朽之物的友谊所捆绑，当他失去它们时，便会被撕碎，那时才感受到他早在失去它们之前便已存在的悲惨。我当时就是那样；我痛哭流涕，在苦涩中找到安息。我如此悲惨，竟将那段悲惨的生命看得比我的朋友更宝贵。因为我固然愿意改变它，却更不愿失去它，胜过不愿失去他；我甚至不知道，我是否愿意为他而失去它，就像流传给我们的（若非虚构）\Person{皮拉得斯}和\Person{俄瑞斯忒斯}的故事那样，他们情愿为彼此而死，或一同死去，因为他们看来，不能共活比死亡更痛苦。然而，我内心里也生出一种相反的感受：因为生活对我来说极其厌倦，死亡又极其可怕；我想，我越爱他，便越憎恨并畏惧那将他从我夺走的死亡，视它如最残忍的敌人；我想象这死亡既然能征服他，也会突然消灭所有的人。我记得我当时就是这样。请看我的心，我的天主！请看并审视我，因为我清楚记得，我的希望啊！是祢洗净我这类情感的不洁，引导我的眼目仰望祢，并把我的脚从罗网中拉出来。我感到惊讶，其他凡人竟还能活着，因为我所爱的那位，仿佛永不死亡者，却已死去；我更诧异，那个身为他第二个的我，竟能在他死后还活着。有人论到自己的朋友说得好：\Quote{你，我灵魂的一半}，因为我感到我的灵魂和他的灵魂只是两个身体内的一个灵魂；因此，我的生命对我而言成了恐怖，因为我不愿只活在一半中。正因如此，我或许也害怕死亡，唯恐我深爱的他会完全死去。\par

% structure: p0020 source=h2 target=subsection reason=relative-heading-rank; base=h2

\subsection{第七章 被不安和悲伤所困扰，他再次离开家乡前往\Place{迦太基}。}

12. 哦，疯狂啊，竟不知如何以人当受的爱去爱人！我那时是多么愚蠢的人啊，竟如此不能忍受人类的命运！于是我烦躁、叹息、哭泣、折磨自己，既不得安息，也无从接受劝导。因为我背负着一个撕裂而污秽的灵魂，无法忍受它被我背负，却找不到安放它的地方。不在怡人的丛林中，不在玩耍或歌声里，不在芬芳之处，不在丰盛的宴席上，不在床榻的欢娱里，最终也不在书籍和诗歌中找到安息。一切都显得可怖，甚至光明本身也是如此；凡不是他那样的，都令人厌恶且可憎，唯有呻吟和眼泪，只有在这些中，我才找到些许安息。但每当我的灵魂被拉离它们，沉重的痛苦重担便压垮我。主啊，这本应被举向祢，让祢减轻并移去。这一点我是知道的，却既不愿也不能；更因为在我的思想中，祢对我而言并不是任何坚实实在的东西。因为祢并非祢自己，而是一个空虚的幻影，我的错误便是我的神。如果我试图将我的重担卸在其上，让它得享安息，它便沉入虚空，又迅速弹回压在我身上，于是我自己便成了一处不幸之地，既不能停留，也无法离去。因为我的心怎能逃离自己的心？我怎能逃离我自己？又怎能不处处跟随自己？然而我逃离了家乡；因为这样，我的眼睛在不再惯常看见他的地方，便不会那么急切地寻找他。于是，我离开了\Place{塔加斯特}城，来到了\Place{迦太基}。\par

% structure: p0022 source=h2 target=subsection reason=relative-heading-rank; base=h2

\subsection{第八章 他的悲伤因时间和朋友的安慰而消逝。}

时光从不虚度，也不在我们的感官中徒然流转。它们在人心中运行奇妙的工作。看，它们日复一日地来去，借着来去，在我心中播撒下别的意念和别的回忆，一点一点地，又以先前那种种欢愉把我修补起来，我那悲伤也俯首屈从了。但随之而来并非定是别的悲伤，却是别的悲伤的缘由。试问从前那悲伤何以如此轻易刺骨？岂不是因为我曾把灵魂倾注于尘土中，爱着一个必死之人，仿佛他永不会死？然而，最让我复苏且获得安慰的，乃是其他友人的慰藉，我与他们一同爱着那替代祢而爱的对象。此乃一个巨大的虚构故事，一个持久的谎言，借着其私通的接触，我们那发痒的耳朵所卧之灵魂，正在被染污。然而，每当我有朋友死去，这虚构故事并不与我一同死去。在他们身上，还有其他更能抓住我心的事物：一同谈论戏谑；彼此善待交流；同读怡情书籍；一同戏耍，一同认真；间或意见不一却无愠色，如人与己相异时那样；甚而借着这少见的不合，反为我们更频繁的合意增添滋味；有时教导，有时受教；切切思念不在场的，欣喜欢迎接到的。这类以及类似的表露，从那些爱人且被爱者的心灵发出，经由面容、舌、眼以及千般怡人举止，便如同燃料一般，将我们的灵魂融合在一起，使众人合而为一。\par

% structure: p0024 source=h2 target=subsection reason=relative-heading-rank; base=h2

\subsection{第九章 人之爱，虽恒于爱且回报爱，终归逝去；而爱天主者永不失友。}

14. 这便是人在友爱上所爱的；爱得如此深，以致人的良心若不爱那爱己者，或不以爱回报那爱己者，不求对方任何回报，只求其爱的明证，便要自责。由此而有丧友之哀恸，忧苦之阴郁，以泪水浸透的心，一切甜蜜变为苦涩，随着逝者生命之失去，生者亦经历死亡。爱祢、且在祢内爱朋友、为祢的缘故而爱仇敌的人，是有福的。因为惟有他，在不可丧失的那一位内，对所有人都是宝贵的，故而不会失去任何珍爱的人。这位是谁？岂不是我们的天主，那创造了天地\BibleRef{创 1:1}并充满天地\BibleRef{耶 23:24}，且因充满而创造的天主吗？惟离开祢的人，才失去祢。离开祢的人，无论走向何方，逃往何处，无非是从慈颜的祢转向震怒的祢而已。在他受的惩罚中，何处寻不见祢的法律？\BibleQuote{而祢的法律就是真理，}祢就是真理。\BibleRef{若 14:6}\par

% structure: p0026 source=h2 target=subsection reason=relative-heading-rank; base=h2

\subsection{第十章 万物存在趋于消亡，若非天主眷顾，我们便无安全。}

15. \BibleQuote{求你使我们回转，万军的上主天主，使祢的慈颜照耀，我们便会得救。}因为人的灵魂无论转向何方，除非转向祢，便与忧苦相连；哪怕它依附于在祢之外、在己之外的美好事物。然而，这些事物若不是从祢而来，根本无有。它们升起又落下；升起时，仿佛开始存在；成长，以求达至完美；及至完美，便会衰老消亡；虽非所有都衰老，但所有皆消亡。因此，当它们升起并趋向存在时，它们越是迅速成长以成其所是，便越是迅捷地趋于不复存在。这便是它们的道路。祢赐予它们的，仅此而已，因为它们是万物的部分，并非同时全部存在，而是经由离去与接续，共同构成宇宙整体，而它们便是其部分。我们的话语亦是如此，借着发出声音的符号得以完成；然而，这也非至善，除非一个词在发完其声部后消逝，好让后一个接续上来。天主啊，万有的创造者，愿我的灵魂因这万物赞颂祢；但不要让我的灵魂借着身体的感官，以爱的胶液粘附于这些事物。因为它们去向它们当去之处，不复存在；它们以有害的欲望撕裂灵魂，因为灵魂渴望存在，却又爱在所爱中歇息。然而，在这些事物中，并无地方可寻；它们不停留——它们飞逝；谁又能以血肉的感官追随它们呢？即便它们近在咫尺，谁又能抓住它们？血肉的感官是迟钝的，正因其是血肉的感官，其界限便是自身。它足以完成其受造之目的，却不足以留住那些从其指定的起点奔向指定终点的流逝之物。因为，在祢借着造它们的圣言中，它们听到了谕令：\Quote{由此至彼。}\par

% structure: p0028 source=h2 target=subsection reason=relative-heading-rank; base=h2

\subsection{第十一章 勿爱世间部分；造作者天主恒常不变，其圣言永恒。}

16. 我的灵魂啊，不要愚昧，也不要让你心中的耳朵因纷乱的愚昧而聋聩。也当静听。那圣言亲自呼唤你回归；那里才是不可扰乱的安息之所，爱若不自弃，便不会被舍弃。看，这些事物逝去，好让其他事物接续，为使这尘世宇宙在各部分上得以完整。但天主圣言却说：难道我会离开任何地方吗？在那里，你当安置你的居所。无论你从那里得到什么，我的灵魂，都该托付于那里；至少如今你已厌倦了欺骗。将从真理得来的一切，托付于真理，你便一无损失；你的枯萎将重焕生机，你的一切疾病将得医治，你那些可朽坏的部分将被改造更新，并汇聚到你身边；它们不会沉沦到它们下降之处，反而将与你同住，并在那永存永续的天主面前永远延续。\BibleRef{伯前 1:23}\par

17. 那么，你为何悖逆而随从你的肉身呢？不如让它归向你而跟随你。你借着她所感觉的，都只是部分；而你并不知道这些部分所属的全体，但这些部分却令你欣喜。但你肉身的感官倘若能够理解全体，并且它本身并非为了惩罚而被公义地局限于全体的一部分，你就会希望现在存在的一切都消逝，这样全体便能令你更为喜乐。因为我们所讲的话，你也是借这肉身的感官聆听的，然而你并不愿音节停留，却愿它们消逝，好让其他音节接续而来，使整段话能被听见。当任何单一事物由众多部分组成，而所有这些部分并非同时存在时，若一切都能同时被感知，那么全体就会比单个部分更令你喜乐。但比这些都更好的是那创造万有者；祂是我们的天主，祂并不消逝，因为没有事物能接替祂。如果形体令你喜悦，你就为它们赞美天主，并将你的爱回转到它们的创造者那里，免得在令你喜悦的事物中，你反倒令祂不悦。\par

% structure: p0031 source=h2 target=subsection reason=relative-heading-rank; base=h2

\subsection{第十二章 爱并不受谴责，但在天主内、借着耶稣基督而得享安息的爱，方为更可取。}

18. 如果灵魂令你喜悦，就当在天主内爱它们；因为它们也是可变的，但在祂内得到坚固，否则它们就会流逝而消亡。因此，要在祂内爱它们；并尽你所能带领更多灵魂与你一同归向祂，且对他们说：\Quote{让我们爱祂，让我们爱祂；是祂创造了这些，祂并不遥远。因为祂并非创造了它们后便离去；而是它们源于祂，在祂内。看，哪里有真理被人认识，祂就在那里。祂就在心灵深处，但心灵却已偏离祂。悖逆的人啊，归回你们的心灵吧！\BibleRef{依 56:8}要紧紧依附那创造了你们的祂。与祂同在，你们就能立稳；在祂内安息，你们就会得享安宁。你们为何走上崎岖的道路？往何处去？你们所爱的美善都来自祂；当它们被导向祂时，便是美好甘饴的，而一旦你们为它们离弃了祂，任何源于祂的事物都会公义地变为苦涩，因为它们被不义地爱了。那么，你们又为何在这些艰难费力的道路上越走越远？你们所寻求之处并无安息。去寻找你们所寻求的，但它不在你们寻找的地方。你们在死亡之地寻找有福的生命；那里没有。因为若是生命本身不在，哪里会有有福的生命呢？}\par

19. 然而我们的生命亲自降到这里，承受了我们的死亡，并以祂生命的丰盛将死亡杀死；祂大声疾呼，召唤我们从这里回到祂那里，回到祂从那里来到我们中间的那个隐秘之处——首先进入童贞玛利亚的胎中，在那里人性与祂结合——我们可朽的肉身，好使它不再永远可朽——从那里出来，\BibleQuote{像新郎出离洞房，又像壮士欣然奔赴赛程。}祂并未耽搁，而是奔驰着，通过言语、行为、死亡、生命、降下、升天，向我们大声疾呼，要我们回归于祂。祂从我们眼前离去，好让我们回归内心，在那里寻见祂。祂离去了，看，祂就在这里。祂不愿长久与我们同在，却并未离弃我们；因为祂离去的那个地方，祂从未离开过，因为\BibleQuote{世界是由祂造成的。}\BibleRef{若 1:10}祂曾在这世界上，祂来到这世界，是为拯救罪人，\BibleRef{弟前 1:15}而对这罪人，我的灵魂向他认罪，好使祂能治愈它，因为它冒犯了祂。人子们啊，你们的心竟然如此迟钝，要到几时呢？\BibleRef{路 24:25}如今，生命既已降下到你们这里，你们还不肯上升而生活吗？但你们既已在高处，且向天开口，你们要上升到何处呢？你们必须下降，才能上升，上升到天主那里。因为你们曾经\Quote{悖逆祂而上升}，从而跌倒了。把这些话告诉他们，好让他们在涕泣之谷中哭泣，这样你就带着他们和你一同归向天主，因为如果你怀着爱火说话，你便是借着祂的圣神这样对他们说的。\par

% structure: p0034 source=h2 target=subsection reason=relative-heading-rank; base=h2

\subsection{第十三章 爱源于吸引我们的恩宠与美。}

20. 那时我并不知道这些事，我只爱那些低下的美丽，并沉入深渊；我对我的朋友们说：\Quote{除了美，我们还爱什么？那么，什么是美？什么是美丽？是什么吸引我们并使我们与所爱之物结合？因为若事物中没有一种恩宠与美丽，它们绝不可能吸引我们。}我观察并体会到，在形体自身中，有一种美来自它们构成的某种整体，另一种美则来自彼此相宜，就像身体的某部分与整体相宜，或鞋与脚等等。这个想法从我内心的深处涌现，我写了几本书（我想是两三本）\WorkTitle{论美与适宜}。主啊，祢知道，这件事我已忘记；因为我没有保存它们，它们已不知如何地散失了。\par

% structure: p0036 source=h2 target=subsection reason=relative-heading-rank; base=h2

\subsection{第十四章 论他所著的\WorkTitle{论美与适宜}，并题献给\Person{希厄里乌斯}。}

21. 主，我的天主，是什么促使我将这些书题献给\Person{希厄里乌斯}，一位\Place{罗马}的演说家，我并未亲眼见过他，却因他闻名于世的学识而喜爱他，他以此著称，我也听过他的一些话，那些话使我愉悦。但更令我愉悦的是，他也为他人所喜爱，那些人对他推崇备至，他们惊异于一个\Place{叙利亚}人，起初学习希腊辩术，后来竟成为一位出色的拉丁语演说家，并且精通与智慧相关的研究。一个人不在面前时，就这样受人称赞和喜爱。这种爱是从称赞者的口进到听者的心吗？不是的。而是借着那位爱慕者，另一人被点燃了爱火。因为，当人们相信那称赞者是出于真诚的心来赞美他时，那被称赞的人便因此受爱戴；也就是说，当爱慕他的人称赞他时。\par

22. 因此，我那时是按照人的判断去爱人，而不是依照你的判断，我的天主，在你的判断中没有人受骗。但为什么不像称赞著名的赛车手、或凭俗众的赞誉声名远播的猎人那样呢？——而是大不相同地、严肃地、如同我渴望自己被称赞的方式？因为我不愿人们像称赞演员那样称赞我、爱我——尽管我自己倒是称赞并爱慕演员——我宁愿默默无闻，也不愿如此出名；宁愿被人憎恨，也不愿如此被爱。如今，这些各式各样的爱的影响，在同一个灵魂中是如何分布的呢？我在别人身上所爱的究竟是什么？如果我不恨这东西，我就不会从自己身上厌恶并排斥它，因为我们都是同样的人。这并不意味着，一个人喜爱一匹好马，即使他能成为那匹马也不愿成为它，因此一个演员也具有相同情形，因为演员与我们人性相同。那么，我这个人，是否在我所恨恶成为的人身上，爱着某种东西呢？人本身是一个巨大的深渊，\BibleQuote{主啊，连他的头发你都数过了，没有你的许可，一根也不会掉在地上。}\BibleRef{玛 10:29} 然而，他头上的头发，比他的情感和内心的活动更容易数算。\par

23. 然而，那位演说家正是我所喜爱的那种人，甚至希望自己也能成为那样的人；我由于骄傲的膨胀而步入歧途，\BibleQuote{被一切异风所飘荡}\BibleRef{弗 4:14}，却仍由你暗中引领。我从何而知，又怎能坦然向你承认：我之所以爱他，更多是因那些赞扬他的人的喜爱，而非因他们所赞扬的那些事迹本身？因为，倘若他遭遇非难，而这些人非但不赞扬他，反而以贬斥和轻蔑的语气说出同样的事迹，我绝不会如此动情并受激励去爱他。然而，事迹并无不同，他本人也无不同，变的只是叙述者的情感。看哪，尚未得到真理稳固支撑的灵魂何等软弱！它随着臆测者胸中吹出的言语之风，被抛来抛去，前行后退，于是光明遮蔽，真理不辨。而真理却明摆在眼前。对我来说，我的文采和学识若能为他所知，便是件大事；若他赞许，我便倍受激励；若他否定，我这颗缺乏你坚实支撑的虚荣之心便会感到被冒犯。至于我写给他的\Quote{美与适合}，即便无人与我共鸣，我仍欣然反思，静观，赞叹。\par

% structure: p0040 source=h2 target=subsection reason=relative-heading-rank; base=h2

\subsection{第十五章 写作时，为形像所蔽，未能认出天主的精神性本质。}

24. 但是，全能者啊，我尚未领悟，在你的智慧中，这无能之事的关键所在；你是\Quote{独自行了奇异之事}的主。我的心灵徘徊于有形的形式之间，我将本身即为美的定义为\Quote{美}，将因与他物相配而显为美的定义为\Quote{适合}；并以有形之物为例加以论证。我也曾转而关注心灵的本性，但我对精神事物所抱的错误见解，却阻碍我认识真理。然而，真理的力量本身强行闯入我的眼帘，我却将自己悸动的灵魂从无形的实体转向线条、色彩和庞大的量。由于不能在心灵中感知这些，我以为自己也感知不到心灵。我既在美德中爱慕和平，在邪恶中憎恶纷争，便在前者中辨认出统一，在后者中则是一种分裂。我设想，理性的灵魂、真理的本性以及至善，都包含在那统一之中；而不幸的我，却在这分裂中想象着某种非理性生命的实体，以及至恶的本性——它不仅是实体，而且还是真实的生命，却又不是源自你——我的天主，万物的根源。我将前者称为\OriginalTerm{一元}{Monad}，视其为无性别的灵魂；而将后者称为\OriginalTerm{二元}{Duad}——即在暴力行为中的愤怒，在情欲行为中的血气与贪婪——竟不知自己所言为何。因为我还不知道、也未曾学过：恶并非实体，我们的灵魂也非那至高不变之善。\par

25. 正如在暴力行为中，若那激发冲动的心灵情感堕落，并傲慢而悖逆地发作；在情欲行为中，若那放纵肉身享乐的心灵欲望毫不节制；同样，若理性灵魂本身堕落，错误与谬见便会玷污生命——那时的我便是如此，我不知道灵魂必须由另一种光辉照耀，方能分享真理，因为它本身并非真理的本体。因为你必点亮我的灯，主我的天主必照明我的黑暗；且\BibleQuote{从他的满盈中，我们都领受了}\BibleRef{若 1:16}，因为\BibleQuote{那光是真光，普照来到世上的每一个人}\BibleRef{若 1:9}；因为在你内\BibleQuote{没有变化，也不见转动的阴影}\BibleRef{雅 1:17}。\par

26. 然而我向祢冲去，却被祢推回，好使我品尝死亡，因为祢\BibleQuote{拒绝骄傲的人。}\BibleRef{雅 4:6} 可是，还有什么比我在惊人的疯狂中，竟敢宣称自己本性上就是祢所是，更为骄傲的呢？\OriginalTerm{可变的}{mutable}——这一点我很清楚，因为我渴望获得智慧，正是出于要从较差变为较好的愿望——却宁愿想象祢是可变的，而不愿承认我自己不是祢那样的本性。因此我被祢推回，祢抗拒了我顽梗悖逆的颈项；我虚构出各种\OriginalTerm{物质的形式}{corporeal forms}，自身是血肉，却责怪血肉，自己是\BibleQuote{一阵吹过的风，}\BibleRef{咏 78:39} 却不肯归向祢，反而徘徊游荡，趋向那些既不在祢内，也不在我内，亦不在物体内的虚无之物。它们并非由祢的真理为我创造，而是从物质事物之中由我\OriginalTerm{虚妄的幻想}{vain conceit}所构思出来的。我曾常常以轻率而愚昧的态度询问祢那些忠信的小孩子，我的同乡公民——我自己尚不自觉已被充军流放——我问说：\Quote{那么，天主所创造的灵魂为什么会错误呢？}我却不愿任何人问我：\Quote{那么，天主为什么错误呢？}我宁可坚持说，祢\OriginalTerm{不变的本体}{immutable substance}是迫不得已而错误，也不肯承认我\OriginalTerm{可变的}{mutable}本体是出于\OriginalTerm{自由意志}{free will}而迷失，并因惩罚而错误。\par

27. 当我写那几卷书时，约二十六七岁——那时我沉思着物质性的虚构，它们在我心耳中喧嚷。这些虚构，甘饴的真理啊，我引向祢内心的旋律，默想着\Quote{美善与合宜}，并切望停留，聆听祢，因新郎的声音而欢欣踊跃\BibleRef{若 3:29}，我却不能；因为我被自己错误的喧嚣驱逐，被我骄傲的重压沉入深渊。因为祢并没有\BibleQuote{使我听到欢乐和愉快的声音}\BibleRef{咏 51:10}，那尚未谦卑的骨骸也没有欢跃。\par

% structure: p0045 source=h2 target=subsection reason=relative-heading-rank; base=h2

\subsection{第十六章 他极容易领悟了自由学术和亚里士多德的范畴，却没有真正的收获。}

28. 当我还不满二十岁时，有一本\Person{亚里士多德}的著作，题为\WorkTitle{十大范畴}，落到我手中——我一听到这个书名，就像听到伟大而神圣的事物一般，当时我的雄辩术老师（在\Place{迦太基}），以及其他被视为有学问的人，总是带着充满骄傲的面颊提到它——我独自阅读并理解了，这对我有什么益处呢？后来我与别人讨论，他们说，即使有非常能干的老师（这些老师不仅口头讲解，而且在尘土中画了许多图形），他们才勉强理解，而告诉我的东西，竟不比我独自阅读时获得的更多，这又有什么益处呢？这本书在我看来，十分清楚地讲论了\OriginalTerm{本体}{substances}，比如人，以及本体的属性——比如人的形象，是怎样的；他的身高，几尺；他的关系，是谁的兄弟；或处在何处，何时所生；或他站立或坐着，或穿鞋或武装，或做什么，或遭受什么；以及无数可分列于这九个\OriginalTerm{范畴}{categories}之下的事物（我已举出数例），或者那个首要的本体范畴。\par

29. 这一切对我有什么益处呢？它甚至反而阻碍了我，因为我想象凡存在的一切都包含在那些十个范畴之内，便试图这样来理解，我的天主啊，祢那奇妙而不可改变的单一性，仿佛祢也隶属于祢自己的伟大或美丽之下，以致它们如同在物体中一样，以祢为主体而存在于祢内，然而祢本身就是祢的伟大和美丽。可是，一个物体并不因为它是物体而成为伟大或美丽，因为即使它较小或较丑，它仍然是个物体。但我对祢所构想的，却是虚假，不是真理——是我患难的虚构，而不是祢真福的支柱。因为祢曾吩咐，这事便在我身上应验：\BibleQuote{地就给我生出荆棘和蒺藜}\BibleRef{依 32:13}，且\BibleQuote{我须劳苦，才得食粮}\BibleRef{创 3:19}。\par

30. 我当时是卑劣情欲的卑贱奴隶，竟能独自阅读和理解我能找到的一切所谓自由学术的书籍，这对我有什么益处呢？我固然乐在其中，却不知其中一切真实和确定的东西从何而来。因为我那时背向光明，面朝被光照的事物；因此我用来辨识受光照之物的面孔，自身却未受光照。凡论及修辞学、逻辑学、几何学、音乐或算术的著作，我都毫不费力地，也无须任何人的教导，便能理解，主，我的天主，祢知道，因为敏捷的理解力和敏锐的领悟力都是祢的恩赐。但我仍未因此而向祢献祭。所以，这不但对我没有用处，反而更招致毁坏，因为我竟企图\BibleQuote{将我的产业如此丰厚的一份归于自己掌握}\BibleRef{路 15:12}；我不为祢保守我的力量，反而\BibleQuote{离开祢往远方去，在那里浪费在娼妓身上}\BibleRef{路 15:13}。好才能若不用于善行，对我有什么益处呢？因为我并未意识到，即使好学而天资聪颖的人，掌握这些技艺也是极其艰难的，直到我设法向这样的人讲解时才发现；而最能领会这些技艺的，是那些能不太迟钝地跟随我讲解的人。\par

31. 但这于我何益呢？即便假设祢，\OriginalTerm{主天主，真理}{O Lord God, the Truth}，是一具光明而庞大的躯体，而我乃此躯体之一部分？悖逆甚矣！但我当时正是如此。我的天主啊，我毫不羞愧地向祢承认祢所施于我的仁慈，并呼求祢——我，那时竟在人前宣扬我的\OriginalTerm{亵渎}{blasphemies}，向祢狂吠，且不以为耻。在那一切学科和所有那些盘根错节的书卷上，我无需人间导师之助，便能条分缕析，机敏才智于我何益？我竟在\OriginalTerm{虔敬的道理}{doctrine of piety}上犯下如此可憎的错误，行事如此\OriginalTerm{渎神的卑劣}{sacrilegious baseness}。反之，\OriginalTerm{祢的小子们}{Your little ones}，纵然才智远为迟钝，又有何妨？他们并未远离祢，在\OriginalTerm{祢教会的巢}{the nest of Your Church}中安然长成羽翼，以\OriginalTerm{健全的信德}{sound faith}为食粮，滋养\OriginalTerm{爱德的翅膀}{wings of charity}。主，我们的天主，让我们在祢翼荫下寄望，求祢保护我们，怀抱我们。\BibleQuote{祢将怀抱我们，从小直到白发，祢仍怀抱我们；} \BibleRef{依 46:4} 因为我们的\OriginalTerm{稳固}{firmness}，若出乎祢，方为稳固；若出乎我们自己，便是\OriginalTerm{虚弱}{infirmity}。\OriginalTerm{我们的善}{our good}永与祢同在，一旦背离，便趋堕落。主啊，让我们现在回转，免得倾覆，因为我们的善与祢同在，无丝毫阴翳，这善就是祢自己。我们亦无需忧虑，恐因堕落离去而找不到归去之处；因当我们不在时，我们的家——祢的\OriginalTerm{永恒}{Eternity}——并未崩塌。\par

\SourceRecord{Translated by J.G. Pilkington.；Nicene and Post-Nicene Fathers, First Series；Vol. 1.；Edited by Philip Schaff.；Buffalo, NY: Christian Literature Publishing Co.,；1887.；<http://www.newadvent.org/fathers/110104.htm>.}

% structure: p0001 source=h1 target=section reason=page-title

\section{\WorkTitle{忏悔录}（卷五）}

他记述自己二十九岁时，发现了\OriginalTerm{摩尼教徒}{Manichæans}的谬妄，在\Place{罗马}和\Place{米兰}教授修辞学。在听过\Person{安波罗修}之后，他开始醒悟。\par

% structure: p0003 source=h2 target=subsection reason=relative-heading-rank; base=h2

\subsection{第一章 灵魂理当赞颂天主并向祂认罪。}

1. 接纳我以口舌的代劳献上的忏悔之祭，这口舌是祢造成并赋予生命的，好使它向祢的圣名认罪；求祢医治我所有的骨头，让它们说：\BibleQuote{主啊，谁能与祢相比？}因为向祢认罪的人并不告诉祢他内心发生的事，因为关闭的心瞒不过祢的眼，人心的刚硬也推不开祢的手，但祢随己意消融它，或出于怜悯，或为施行惩罚，\BibleQuote{没有一人能在祢的心前隐藏自己。}但愿我的灵魂赞颂祢，好能爱慕祢；让它向祢宣认祢自己的仁慈，好能赞美祢。祢所有的受造物都不停息，也不沉默于对祢的赞美——不论是人的精神，借由朝向祢的声音，还是动物与有形之物，借着默想它们的人的声音；这样，我们的灵魂可以从困乏中起身向祢上升，倚靠祢所造的万物，并经由它们走向祢，祢奇妙地造了它们，那里有更新与真正的力量。\par

% structure: p0005 source=h2 target=subsection reason=relative-heading-rank; base=h2

\subsection{第二章 论那些企图逃避全能天主者的虚妄。}

2. 愿不安分和不义的人离开并逃避祢吧。祢既看见他们，也分辨阴影。看，与他们相关的一切都是美好的，他们自己却是丑陋的。他们怎样伤害了祢呢？或在何事上侮辱了祢那从天到地最低处公正且完全的治理？因为他们逃避祢的圣面时，逃到哪里去了？何处祢找不到他们？但他们逃避，是为了不要看见祢在看着他们，并因失明而撞上祢；见：\BibleRef{创 16:13}因为祢不遗弃任何祢所造的——不义的人便撞上祢，公正地受伤害，自己远离了祢的温良，撞上了祢的正直，跌在自己刚硬上。诚然，他们不知道祢无处不在地无所不包，不知道惟独祢亲近那些远离祢的人。所以，让他们回头并寻求祢吧；因为他们虽离弃了造物主，祢却没有离弃受造物。让他们回头并寻求祢吧；看，祢就在他们心里，在那些向祢认罪、扑倒在祢怀中、在顽梗之后在祢胸前哭泣的人心里，祢温柔地擦干他们的眼泪。他们哭得更甚，并在哭泣中喜乐，因为主啊，不是人，不是血肉，而是主，祢造了他们，再造他们，安慰他们。而我寻找祢时，我身在何处？祢本在我面前，我却甚至远离了自己；我既找不到自己，更找不到祢！\par

% structure: p0007 source=h2 target=subsection reason=relative-heading-rank; base=h2

\subsection{第三章 听了摩尼教徒中最有学问的主教\Person{福斯图斯}讲论后，他领悟到天主是生命与无生物的作者，尤其眷顾谦卑的人。}

3. 让我在天主面前揭示我二十九岁那年的事。那时，有一个摩尼教的主教来到\Place{迦太基}，名叫\Person{福斯图斯}，是魔鬼的大网罗，许多人因他流畅动听的言辞的魅力而被缠住；我虽赞赏他的口才，却能将它与我所渴求学习的那些真理分开。我并不看重他所端上的那盘小小的演讲术，而是看重他所赞扬的福斯图斯摆在我面前让我享用的知识。其实，名声早已向我传扬他，说他在一切适宜的知识上精通，尤其在自由学术上卓越。我读过并记住了许多哲学家的学说，于是常将他们的某些道理与摩尼教徒冗长的神话以及他们所说的较早的事相比较，那些人最多只能估量这下界的世间，却怎么也找不到它的主宰，见：\BibleRef{智 13:9}我觉得前一种更可信。因为主啊，祢是伟大的，垂顾卑微的人，却远远地认识骄傲的人。祢不亲近破碎的心以外的，骄傲的人也找不到祢——即使他们用巧技能数算众星和海沙，测量星宿的区域，追踪行星的轨迹。\par

4. 他们用祢所赐予的理智与才能去探究这些事，发现了很多东西，并在许多年前就预言了那些发光体——太阳和月亮的亏蚀，在哪一天，哪一时，从某几分的点上会发生。他们的计算没有失误，事情正如他们所预言的那样发生了。他们写下了所发现的规则，至今仍在诵读；从这些规则，其他人能预言在哪一年、哪一月、哪一天、哪一时、光亮被遮住几分，月亮或太阳将发生亏蚀，事实也恰如所预言。而不懂这些的人惊奇赞叹，知道的人则欢欣自豪；由于不敬虔的骄傲，他们背离了祢，离弃了祢的光，他们能提前这么久预言那将要发生的日光亏缺，却看不见自己现在已有的亏缺。因为他们不以宗教的精神追问自己从哪里得到这种探究的能力。他们发现是祢造了他们，却不把自己交给祢，好让祢保存祢所造的，也不把自我作为祭品献给祢——就像他们已经把自己塑造成的那样；他们不宰杀自己的骄傲，如同空中的飞鸟，不宰杀自己的好奇心（他们凭此如同海中的鱼，漫游深渊的未知路径），也不宰杀自己的放纵，像\BibleQuote{田间的野兽}那样，好让主，祢这\BibleQuote{吞噬的烈火}，见：\BibleRef{申 4:24}，能烧尽他们无生命的忧虑，并使他们不朽地更新。\par

5. 但是那道路——你的\OriginalTerm{圣言}{Word}\BibleRef{若 1:3}，你曾藉着祂造了这些他们所计算之物，又造了计算者，并他们用以察觉所计算之物的感官，以及他们赖以计算的判断能力——他们却不认识；也不认识你的智慧无可测度。然而，那\OriginalTerm{独生子}{Only-begotten}\BibleQuote{成了我们的智慧、正义和圣化}\BibleRef{格前 1:30}，被列入我们中间，并向\Person{凯撒}纳税\BibleRef{玛 17:27}。他们不认识这条道路，就是他们能从自身下降到祂面前的道路；也不认识藉着祂，他们能上升到祂那里。他们不认识这道路，却自以为被高举，与星辰同列\BibleRef{依 14:13}，光辉闪耀；看！他们却坠落到地上\BibleRef{默 12:4}，而\BibleQuote{他们愚昧的心就昏暗了}\BibleRef{罗 1:21}。关于受造物，他们说出了许多真实的事；但对真理，即受造物的创造者，他们却不以\OriginalTerm{虔敬}{piety}去寻求，因此他们找不到祂。或者即使他们找到祂，既知道祂是天主，却不以祂为天主而光荣祂，也不感谢祂\BibleRef{罗 1:21}，反而在他们的思想上成为虚妄，自称是智者\BibleRef{罗 1:22}，将属于你的事物归于自己；以此，凭着极其悖谬的盲目，他们企图将属于他们的事物归咎于你，编造谎言反对身为真理的你，并将不可朽坏的天主的光荣，改换为可朽坏的人的形像，以及飞禽、走兽和爬虫的形像\BibleRef{罗 1:23}——将你的真理变为谎言，崇拜并事奉受造物，超过事奉造物主\BibleRef{罗 1:25}。\par

6. 然而，关于受造物，我从这些人那里保留了许多真理，并且从计算、季节的交替，以及星辰的可见现象中，我得以明了原因；我将这些与\Person{摩尼}的言论相比较，他狂热地大量撰写了关于这些题目的著作，却发现其中没有任何关于至点、分点、日月食，或类似我从世俗哲学书籍中学到的这类事物的论述。但在这事上，我却被命令相信，然而它与计算所公认的法则以及我自己的观察并不相符，反而大相径庭。\par

% structure: p0012 source=h2 target=subsection reason=relative-heading-rank; base=h2

\subsection{第四章：对地上和天上事物的知识并不带来幸福，唯独认识天主才带来幸福。}

7. 那么，主，真理的天主，凡知道那些事的人，难道就因此而蒙你悦纳吗？因为那知道这一切，却不认识你的人，是不幸的；但那认识你的人，纵然不知道这些，却是幸福的。但那既认识你又认识这些的人，并不因这些而更幸福，而是单单因你而幸福，只要他在认识你时，光荣你为天主，感谢你，又不思想虚妄的事\BibleRef{罗 1:21}。正如一个人，知道如何拥有一棵树，并为它的使用而感谢你，虽然他可能不知道这树有多少肘高，或树冠有多宽，却比那测量它、数算所有枝条，却不拥有它、也不认识或爱慕它的创造者的人，更为幸福；同样，一个义人，整个世界是他的财富，他好像一无所有，却因紧紧依附你而拥有一切\BibleRef{格前 6:10}，万物都服从你，他即使连大熊星座的轨迹也不知道，但若怀疑他是否确实比那能量度诸天、计算星辰、衡量元素，却忘记了你的人更好，那是愚昧的；\BibleQuote{你曾以数目、重量和尺度处置了万物}\BibleRef{智 11:20}。\par

% structure: p0014 source=h2 target=subsection reason=relative-heading-rank; base=h2

\subsection{第五章：\Person{摩尼}顽固地教导虚假的道理，并骄傲地将圣神归于自己。}

8. 但究竟是谁命令\Person{摩尼}写这些事呢？这些事的技能对虔敬并非必需的。因为你曾告诉人要注重虔敬和智慧，即使对这些其他事物了如指掌，也可能对虔敬和智慧一无所知；但他既然不知道这些事，却极其无耻地胆敢教导它们，显然他与虔敬毫无关系。因为即使我们知晓这些世上的事，自夸这些事也是愚昧；惟独向你\OriginalTerm{宣认}{confession}才是虔敬。因此，这个迷途者之所以谈论许多这类事，意图在于：当被那些真正了解这些事的人证明为假时，他在那些更艰难的事上实际拥有的理解力便可以显明出来。因为他不想被人轻看，而是四处设法说服人，\Quote{圣神——你信徒的\OriginalTerm{护慰者}{Comforter}和\OriginalTerm{富源}{Enricher}——带着完全的权柄亲自居住在他内。}因此，当发现他关于天穹和星辰、日月的运行的教导是虚假时，尽管这些事与宗教的道理无关，他那亵渎神明的傲慢便充分暴露出来，因为他不仅肯定了自己一无所知的事，而且还扭曲了它们，并怀着如此极度的虚荣和骄傲，竟企图将它们归诸于他自己，仿佛他是神明一般。\par

9. 当我听见有基督徒兄弟对这些事无知，或在这些事上犯错，我能耐心容忍他持守自己的意见；我也看不出对这次受造的物质世界的方位或性质有所不知，会对他有什么损害，只要他不抱有任何不符合祢——万物之创造者、主——的任何信念。然而，他若认为这属于虔诚教义的道理，并满怀顽固地肯定自己所不知的事，那才是损害。但即使是在信德初萌时的这般软弱，我们的慈母爱德也会容忍，直到新人\BibleQuote{成长为一个完人}，不再\BibleQuote{随各种教义之风而飘荡}\BibleRef{弗 4:13}。至于那个如此胆大妄为，竟马上自命为所有能被他诱惑而相信此事的人的教师、创始人、首领和领袖，使得所有追随他的人都相信自己跟随的不是一个凡人，而是祢的圣神——当这种极为疯狂的错误教义一旦被证实是假的，谁不认为理当深恶痛绝并将之彻底弃绝呢？但我尚未明确断定，关于昼夜长短的变化、日夜本身、以及较大光体的亏蚀，以及我在其他书籍中读到的类似现象，是否与他的言论一致。即使我发现自己能做到，内心仍会存疑是否果真如此；尽管凭着他所号称为虔诚的名声，我或许可仰赖他的权威，安于信仰。\par

% structure: p0017 source=h2 target=subsection reason=relative-heading-rank; base=h2

\subsection{第6章 福斯图斯确是娓娓动听的演说家，但对\OriginalTerm{自由学术}{liberal sciences}一无所知。}

10. 在那些心神不定地追随他们的差不多整整九年里，我一直过分急切地期待着这位福斯图斯的到来。因为我偶然遇到的其他该派成员，当我提出那些问题而他们无法回答时，总是嘱咐我期待他的来临，说一旦与他交谈，这些困难以及更大的困难都会迎刃而解。他终于来了，我发现他是个谈吐悦耳愉快的人，所讲的内容与他们如出一辙，只是更流畅、辞令更佳罢了。但我的递酒人言辞雅致，于我又有何益？他并未献上我渴求的更珍贵的饮品。我的耳朵早已听厌了类似的东西；它们并不因表达得更好而显得更确凿，也不因辞藻华丽而成为真理；心灵不因面貌英俊、语言流利便必然睿智。可是那些向我夸耀他的人并不是称职的评判者；正因他言语温雅，在他们看来便显得精明智慧。然而我知道另有一类人，即使真理本身，若用流利优美的语言说出来，他们也会怀疑。但我的天主，祢早已藉奇妙隐秘的方式教导了我，我相信是祢教导了我，因为那是真理；而真理之外，没有其他教师——无论它在何时何地从何处照耀我们——唯有祢。所以从祢那里，我已学会：凡事不因说得娓娓动听就必须视为真实；也不因结结巴巴说出就是虚假；同样，不因表达得笨拙就必然真实，也不因措辞优美便一定不真。智慧与愚昧犹如有益与有害的食物，而文雅的或朴素的言辞则像城里的碟子或乡间的瓦器——两种食物都可以盛在两种器皿中。\par

11. 因此，我长久等候这个人的那份热切，确实因他辩论时的神态感受，以及他表达思想时的流利恰当辞令而感到满意。于是我满心欢喜，并与其他人一起（甚至超过他们）夸耀和颂扬他。然而，令我烦恼的是，在不便与他私下交流讨论的听众聚会上，我不被允许提出那些困扰我的问题。等到我可以发言，并与朋友们一起在不妨碍他与我讨论的时候引起他的注意，提出那些使我困惑的问题时，我发现他首先对\OriginalTerm{自由学术}{liberal sciences}除了文法学之外一无所知，而且那也只是平常的水平。然而，因他读过\Person{西塞罗}的几篇\WorkTitle{演说集}、\Person{塞内卡}极少数的著作、一些诗人的作品，以及他自己教派中用连贯拉丁文写成的几卷书，并且每日不断练习演讲，从而练就了一种辩才，这种辩才因他得心应手的机智和某种天生的优雅而愈加赏心悦耳、引人入胜。我主我的天主，我良心的审判者，我的回忆岂不正如此吗？我的心和我的记忆都摆在祢面前，祢那时以深不可测的上智安排引导着我，把我那些可耻的谬误摆在我眼前，好叫我看见并憎恶它们。\par

% structure: p0020 source=h2 target=subsection reason=relative-heading-rank; base=h2

\subsection{第7章 他看穿了\OriginalTerm{摩尼教徒}{Manichæans}的谬误，毅然退出，并得到天主奇妙的助佑。}

12. 原来，当我清楚看出，他对于我素来相信他精通的那些学问，其实一无所知时，我便不指望他能澄清并解释那些折磨我的种种疑难了。不过，他即使对这些一窍不通，若不是\OriginalTerm{摩尼教徒}{Manichæan}，仍可能持守虔诚的真道。因为他们的书里满是关于天空与星辰、太阳与月亮的长篇寓言，我早已不再认为他有能力令我满意地判断出我热切渴望知道的事——那就是，把这些说法与我在别处读到的计算相比较，\OriginalTerm{摩尼}{Manichæus}著作中的解释是否更胜一筹，或至少同样合理？但当我提议深思熟虑并辩论这些题目时，他却十分谦逊，不敢承担这一重任。因为他自知对这些事毫无知识，并且不耻于承认。他并非那种夸夸其谈的人，我曾为许多这样的人所困扰，他们答应教导我这些事，却什么也说不出来；但此人却有一颗心，这颗心虽然对祢并不正直，对自己却并非全然虚假。因为他绝非完全不知道自己无知，他也不会不加考虑便卷入一场争辩，使他既无法抽身，也不能公正地摆脱困境。正因如此，我对他更为满意，因为一颗真诚心灵的谦逊，比起我所渴望的知识的获得，是更为美好的。我发现，在一切更为深奥和精微的问题上，他都是如此。\par

13. 我对摩尼著作的热心既因此受到了挫折，又对其他教师们更加失望——因为在许多使我困惑的事上，这位在他们中间如此闻名的人物，竟也表现得如此——我便开始与他一起研习他也很热衷的那种文学，即我那时作为修辞学教授，正在\Place{迦太基}给青年学生们讲授的那门学问，并且和他一起阅读他想要听的，或我认为适合他心智倾向的作品。然而，通过结识此人，我原以为能在那个教派中有所精进的一切努力，至此完全落空。这并非说我完全离开了他们，而是就像一个找不到更好东西的人，决定暂时安于我以任何方式偶然遇见的，除非碰巧有更值得向往的东西出现。如此，那位\Person{浮士德}，既非出于本意，亦非出于预料，曾诱使许多人走向死亡的他，如今开始松开我已落入的网罗。因为，我的天主啊，祢的双手，依照祢隐秘的\OriginalTerm{上智安排}{Providence}，并未抛弃我的灵魂；而我母亲心头的血，藉着她日日夜夜倾流的眼泪，化为祭献，呈于祢前，为了我的缘故；祢以奇妙的方式待我。\BibleRef{岳 2:26} 我的天主啊，是祢成就了这事，因为人的脚步由主安排，他必为自己的道路定夺。若非从祢手中，我们又怎能获得救恩？是祢的手再造了它所造的。\par

% structure: p0023 source=h2 target=subsection reason=relative-heading-rank; base=h2

\subsection{第八章 他出发前往\Place{罗马}，他的母亲徒然哀哭}

14. 因此，祢这样安排，使我被说服前往\Place{罗马}，去讲授我当时在\Place{迦太基}讲授的课程，而非其他。我如何被说服这样做，我必不隐瞒于祢，因为在这事上，也必须深思并承认祢的智慧那最深邃的运作，以及祢对我们常存的慈悯。我之所以想去\Place{罗马}，并非因为说服我前往的朋友们向我保证，那里有更大的好处和尊荣——虽然当时这些考虑也影响了我——但我主要且几乎是唯一的动机，乃是我听说那里的青年求学更为安静，且受更严格纪律的约束，因此他们不会任性而放肆地闯入并非他们的老师的学校，未经许可，他们是被禁止进入的。在\Place{迦太基}，恰恰相反，学生中盛行一种可耻而放纵的自由。他们粗暴地冲进来，几乎以狂暴的手势，打断任何人为学生的利益而设立的制度。他们以令人吃惊的冷漠犯下许多暴行，这些行为若非习俗的纵容，本应受法律惩罚。这习俗正显出他们更加卑劣，因为他们现在所做的，按照惯例，仿佛合法，而按照祢不变的法律，却永远是不合法的。他们自以为这样做可以不受处罚，岂不知他们这样做时的盲目本身就是他们的惩罚，他们所遭受的远甚于他们所行的。于是，我作为一个学生时不愿沾染的习气，当了教师后却不得不忍受他人的这种习气；因此我十分乐意前往一个地方，所有知情人都向我保证那里没有类似的事。然而，祢——\Quote{我在活人之地的避难所和我的产业}——在\Place{迦太基}鞭策我，好使我由此被抽离，为保全我的灵魂而交换世俗的居所；同时，在\Place{罗马}，祢又借着那些迷恋这必死生命的人，向我展示诱惑，以吸引我到那里去——一方行事疯癫，另一方则保证虚幻之事；而且，为了纠正我的脚步，祢暗中利用了他们的和我的败坏。因为那些扰乱我平静的人，被可耻的疯狂蒙蔽了眼睛，而那些在别处引诱我的人，则散发着世俗的气息。而我，既憎恶此地的真实不幸，却在那地寻觅虚幻的幸福。\par

15. 然而，我离开那里前往彼处的原因，\Emph{祢}，天主，是知道的，却未曾向我或我的母亲透露；她为我的旅程悲痛不已，一直陪我走到海边。但我欺骗了她，她竭力挽留我，要么想留住我，要么想陪我同行，我则假装有一个朋友，我必须等到他遇到顺风才能启航。我向母亲——那是怎样的母亲啊！——说了谎，然后逃走了。即便如此，祢也已仁慈地宽恕了我，保全了我这满身可憎污秽的人，免于葬身海水，而用祢恩宠之水，使我洁净之后，我母亲眼中的泉源得以干涸，她曾日日为我把眼泪洒在她脚下的土地上。然而，她不肯没有我就回去，我好不容易才说服她那夜留在离我们船只很近的一个地方，那里有一座纪念\Person{圣西彼廉}的小堂。那夜我悄然离去，她却不停止祈祷和哭泣。主啊，她那时泪如泉涌，向祢所求的，不正是求祢不要让我起航吗？但祢以神秘的方式安排，并垂听了她内心真正的渴望，没有应允她当时的祈求，是为了成就她一直所祈求的——使我成为怎样的人。风起帆满，海岸从眼前退去；次日清晨，她悲痛欲狂，就在那里向祢哀号呻吟，祢却不予理睬；同时，祢借着我的种种渴望，催迫我走向一切渴望的终点，并以公义的痛苦鞭笞，涤除了她对我那出于尘世的爱。然而，像所有的母亲——甚至比她们更甚——她喜欢我在身边，却不知我离开后，祢正为她预备何等喜乐。她对此一无所知，所以哭泣哀伤，在痛苦中显露出\Person{厄娃}的遗绪——在痛苦中寻找她曾在痛苦中生下的孩子。但她指责我的背信与残酷之后，又开始为我向祢代祷，然后回到她惯常的地方，而我则前往\Place{罗马}。\par

% structure: p0026 source=h2 target=subsection reason=relative-heading-rank; base=h2

\subsection{第九章 他患病发烧，生命垂危。}

16. 看，我遭受了肉身疾病的鞭笞，身负一切罪恶——就是我对祢、对自己、对他人所犯下的一切重大且众多的罪——外加那使我们在\Person{亚当}内都死去的那原罪的桎梏，我正走向地狱。\BibleRef{格前 15:22} 因为这些罪，祢没有一样在基督内宽恕我，祂也未借着十字架\Quote{摧毁}那因我的罪而与祢\Quote{为仇}的关系。因为我曾误以为祂只是一具幻影，祂怎能借着幻影的十字架来摧毁呢？因此，我灵魂的死亡确实是真实的，正如祂肉身的死亡在我看来是不真实的那样；而祂肉身的死亡是真实的，正如同我那不置信的灵魂的生命是虚假的一样。热度加剧，我正滑向死亡和毁灭。如果那时我就此离去，除了因祢正义的安排、与我恶行相称的火刑之外，我能去何处呢？这一切，我母亲并不知道，虽然身在远方，她仍在为我祈祷。然而，无所不在的天主，垂听了她所在之处的祈祷，并怜悯了在此处的我，使我重获肉身健康，尽管我亵圣的心依然疯狂。因为那场危险并未使我希望受洗，我少年时曾恳求母亲虔诚为我求得的事，如我已述说并忏悔过的，比当时我所行的还要好。但我已长大自取其辱，疯狂嘲笑祢医治的一切计划；而祢却不允许我这样一个罪人遭受双重死亡。若是那创伤重击了我母亲的心，将永远无法治愈。因为我无法充分表达她对我的爱，也无法描述她如今在圣神内为我所受的生产之痛，那比当初肉身生我时更为剧烈。\par

17. 因此，我不能想象，如果我这样死去，刺透了她那充满母爱的心肠，她怎能得治愈。那时，她那如此恳切、频繁且不间断地只向祢呈上的祈祷将何去何从？难道至仁慈的天主，祢会轻看那贞洁又明智的寡妇\Quote{哀痛和谦卑的心}吗？她多么恒常地施舍，对祢的圣徒们多么亲切殷勤，从不间断地每日到祢祭台前祭献，每日早晚两次，必到祢的教会——不为闲谈，也不是老妇人的\Quote{无稽之谈，}\BibleRef{弟前 5:10} 而是为能在祢的讲道中听祢，在祈祷中让祢垂听。难道祢——她之所以能如此全凭祢的恩赐——会轻看并漠视这样一个人的眼泪，而不施救援吗？她向祢哀求，不求金银或任何易变易逝的好处，而是求她儿子的灵魂获得救恩。决不，主。祢确实就在近处，正在垂听，并按祢预定的方式予以应允。祢决不至于在那些神视和祢所赐的答复上欺骗她——其中有些我已提及，有些未提。她把这一切\BibleRef{路 2:19} 珍藏于忠信的心内，不断恳求，当作祢亲手写下的凭据呈给祢。因为祢，\Quote{因为祢的仁慈永远长存，}对于蒙祢宽恕了罪债的人，屈尊就卑，并以祢的许诺成为负债者。\par

% structure: p0029 source=h2 target=subsection reason=relative-heading-rank; base=h2

\subsection{第十章 离开摩尼教徒后，仍对罪恶及救主的起源持有错误看法。}

18. 祢那时使我从那次疾病中康复，使祢婢女的儿子身体恢复健康，好让他能为了祢而生活，赐予他一种更崇高和更持久的健康。而且那时在\Place{罗马}，我加入了那些既欺骗人又受人欺骗的\Quote{圣徒}；不仅是他们的\Quote{聆听者}——我在他家中生病并康复的那人就是其中之一——而且还加入了他们称为\Quote{选民}的人。因为在我看来仍觉得\Quote{犯罪不是我们，而是我不知道的什么别的本性在我们内犯罪。}这满足了我的骄傲，可以不受责备，在我犯了任何过错之后，却不承认自己做了任何事——\BibleQuote{好叫祢医治我的灵魂，因为我犯罪得罪了祢；}但我却喜欢为自己开脱，并归咎于与我同在但并非我自己的另一个什么（我不知是什么）。但确实整个是我，是我的不虔使我与自己分裂；而那罪更不可治愈，因为我不认为自己是个罪人。哦，全能的天主，那可憎的邪恶就是：我宁愿祢在我内被征服而毁灭我，也不愿我自己靠祢而获救恩！所以，祢还没有在我口前设立守望，保守我嘴唇的门户，免得我心偏向恶言，与作孽的人一同为自己的罪过文过饰非——因此，我仍与他们的\Quote{选民}联合在一起。\par

19. 但是现在，对在那错误的教义上取得进步已绝望，就连那些我曾决定在找不到更好的情况下用以自满的东西，我现在也更松懈和马虎地对待。因为我半倾向于相信，他们称为\Quote{学园派}的那些哲学家，比其余的人更明达，因为他们主张我们应当怀疑一切，并断定人没有能力领悟任何真理；因为那时我还不理解他们的意思，我也完全相信他们的想法正如人们常以为的那样。我也没有不坦诚地去抑制我的居停主人，在我观察中，他对那些摩尼著作中充斥的虚妄故事所抱有的确信。尽管如此，我同他们比同不属这异端的人，保持着更亲密的友谊。我也不再以从前的热情为它辩护；然而我与那教派的熟稔（他们许多人在\Place{罗马}隐藏着）使我更迟迟不去寻求另一条路——尤其因为我已对在祢的教会中寻获真理感到绝望，哦，天地的主，一切有形无形之物的创造者，他们曾使我从这真理偏离开了——而且在我看来，相信祢具有人的肉身的形状，并受我们肢体的形体轮廓所局限，是极不合适的。又因为每当我渴望默想我的天主时，我除了想到一堆物体之外，不知该想什么（因为凡不是这样的，在我看来就不存在），这就是我不可避免的错误的最大、也几乎是唯一的原因。\par

20. 因此，我也相信恶是一种相似的实体，拥有它自身污秽而畸形的团块——无论是他们称为土的稠密的，或是像空气体那样稀薄而精细的，他们想象某种恶灵爬行于那土中。又因为一种虔诚——尽管是那样的虔诚——迫使我相信，善良的天主从未创造任何邪恶的本性，我就设想出两个团块，一个与另一个对立，两者都是无限的，但恶的较为收缩，善的较为扩张。从这个有害的开端起，其他的亵渎思想便随之而来。因为每当我的心灵试图回归公教信仰，我就被抛了回来，因为我原先所持守的并非公教信仰。在我看来，视祢——我的天主，我向祢承认祢的仁慈——至少在其他方面是无限的，更为虔诚，尽管在那恶的团块与祢对立的那一侧，我被迫承认祢是有限的，免得我设想祢在各方面都受人身体的形状所局限。而我以为，相信祢未曾创造任何恶——在我无知看来，恶不仅是某种实体，而且是形体性的，因为我对心灵的理解，除了将它看作一种精细的、弥漫于空间中的物体外，别无概念——这比相信从我所以为的恶的本性那样的东西能从祢发出来，似乎更好。至于我们的救主，祢的独生子，我也相信祂仿佛是为了我们的救恩，从祢那极其光辉的团块中伸展出来，以至于关于祂，我只相信那些我在虚妄中所能想象的。那么，我认为这样的本性，若不同肉身相混合，便不能由童贞玛利亚所生；而我自己所构想的这东西，如何能在混合中不被玷污，我无法明白。因此，我不敢相信祂降生成人，免得我被迫相信祂被肉身所玷污。如今，祢的属神的人若读到我的这些忏悔，定会温柔而慈爱地嘲笑我；然而，我从前就是如此。\par

% structure: p0033 source=h2 target=subsection reason=relative-heading-rank; base=h2

\subsection{\Person{赫尔皮迪乌斯}就\BibleBook{新}{约}的真实性有力地驳斥了摩尼教徒}

21. 再者，无论他们在祢的圣经中指责什么，我都认为无法辩护；然而，有时我确实渴望就这些各点，与某位精通那些经书的人商讨，并试试他如何看。因为那时，在\Place{迦太基}，有个名叫\Person{赫尔皮迪乌斯}的人，当面对那些摩尼教徒讲论和辩论，已开始打动我，他从圣经中引证了不易反驳的事，而在我看来，他们的回答显得软弱无力。而这个回答他们并不公开提出，只私下对我们说——他们说\BibleBook{新}{约}的经书已被某些我不知道的人篡改过，那些人想把犹太法律嫁接到基督信仰上；但他们自己并未提出任何未经窜改的抄本。但我，念念于有形之物，深受羁绊，且在某种程度上被窒息，被那些团块所压迫；在其下渴望着祢的真理的气息，我却无法吸入纯洁而未被玷污的真理。\par

% structure: p0035 source=h2 target=subsection reason=relative-heading-rank; base=h2

\subsection{在\Place{罗马}教授修辞学，他发现学生的欺诈行为。}

22. 那时我开始努力实践我来\Place{罗马}的目的——教授修辞学；首先在家里聚集一些人，通过他们我开始为人所知；这时，看，我得知在罗马还发生着另一些坏事，这些在\Place{非洲}我无须忍受。因为那些被遗弃的年轻人在这里并不实行我听说的那种颠覆；然而，他们说，为了逃避支付老师的费用，许多青年合谋，转投他人——背信之徒，他们因贪爱金钱而轻视正义。我的内心也\BibleQuote{憎恶他们}，却不是\BibleQuote{完全的憎恶}；因为，或许我恨他们更多是由于我将因他们受苦，而非因他们所做的恶行。这些人实在卑鄙，他们对祢不忠，贪爱这些暂世事物的转瞬虚幻和肮脏的利诱，那利益会弄脏握住它的手；他们拥抱这飞逝的世界，蔑视祢——祢常存不移，召唤人回头，并宽恕堕落的灵魂，当它归向祢时。现在我憎恶这般邪曲悖谬的人，尽管若他们肯改过，宁愿看重所学胜于金钱，并透过学问认识祢——天主啊，你是真理，是圆满的、确定的善与至洁的平安——我便会爱他们。但那时，我更强烈的愿望是为自己免于受害，而非为天主的缘故让他们变好。\par

% structure: p0037 source=h2 target=subsection reason=relative-heading-rank; base=h2

\subsection{第13章 他被派往\Place{米兰}，为要教授修辞学，并由\Person{安博罗削}所认识。}

23. 因此，当米兰人派人到罗马见城市长官，请求为他们城市选派一位修辞学教师，并由公费派遣时，我就通过那些同样沉醉于\OriginalTerm{摩尼教}{Manichæan}虚幻的人——我正要离开他们以摆脱其影响，但我们双方都未意识到这点——去疏通，使得当时的城市长官\Person{叙马库斯}在命题考试后便派遣了我。我来到了\Place{米兰}，去见主教\Person{安博罗削}，他是举世知名的贤士中的佼佼者，是祢虔诚的仆人；他的雄辩那时正殷勤地将祢麦子的细面、祢\BibleQuote{油}的\BibleQuote{喜乐}，以及祢\BibleQuote{酒}的清醒的陶醉分施给祢的子民。我被祢不知不觉地引领到他面前，为要藉着他自觉地被引向祢。那位属天主的人像父亲一般接待了我，并以一种慈祥而主教式的和蔼看待我的迁居。我开始爱戴他，起初并非因他是真理的导师——我在祢的教会中对真理早已绝望——而是因他是一个待我友善的人。我用心听他向民众讲道，并非怀着应有的动机，而像是在试探他的口才是否名不虚传，或比传闻所言更为丰富或逊色；我全神贯注于他的言辞，对其内容却漫不经心、满心轻蔑；我欣赏他言说的美妙，比\Person{福斯图斯}更有学问，但不及后者那般轻松抚慰。然而，就内容而言，无法相提并论：因为后者游移于摩尼教的欺骗之中，前者却极其健全地教导救恩。但\BibleQuote{救恩远离恶人}，而我当时站在他面前正是如此；不过我正逐渐而不自觉地被拉近。\par

% structure: p0039 source=h2 target=subsection reason=relative-heading-rank; base=h2

\subsection{第14章 聆听主教后，他领悟到天主教信仰的力量，但仍像现代\OriginalTerm{学园派}{Academics}那样怀疑。}

24. 因为尽管我无心去了解他讲的内容，只想听他如何讲（我既已对通往祢的道路绝望，便只留下这种空虚的关注），然而，连同我所欣赏的言辞，那我不经心的内容也进入我的思想；因为我无法将它们分开。当我敞开心扉接纳他\Quote{讲得多么巧妙}时，同时也逐渐但随之进入了\Quote{而且他讲得多么真实！}因为首先，这些事也开始向我显得是可以辩护的；而天主教信仰，对于它我曾幻想无法回应摩尼教徒的攻击，现在却认为它可以不武断地加以维护；尤其在我听到旧约的一两处解释，且往往是寓意解释之后——过去按字面接受时，我在灵性上被\BibleQuote{杀死}了。许多这些书卷的章节既已向我阐明，我便责怪自己从前绝望地以为无法回应那些憎恨和嘲笑法律与先知的人。然而，我那时并未看出，就此便该采取天主教之路，因为它有其学识渊博的辩护者，他们能够并非不合理地回答反对意见；也未看出我所持守的就该被定罪，因为双方似乎都有辩护能力。因为在我看来，那条路并未被打倒；但也未显得胜利。\par

25. 于是我开始认真思考，是否可能以某种方式证明摩尼教徒是错误的。若我能想象出一种灵性实体，他们的一切堡垒便会被击垮，彻底逐出我的思想；但我做不到。然而，对于这个世界的形体及全部有形的自然，即肉体的感官所能触及的，我通过越来越多地思考和比较事物，断定大多数哲学家持有远为更可能的主张。因此，按照学园派的方式（如人们所假定的），怀疑一切并在万有之间摇摆不定，我决定放弃摩尼教徒；判定即使在怀疑时期，我也不能留在一个我偏爱某些哲学家的教派；然而，对于这些哲学家，由于他们没有救恩的基督之名，我全然拒绝将我那昏厥灵魂的治疗托付给他们。因此，我决心在天主教会中作一名\OriginalTerm{慕道者}{catechumen}，这是我父母曾嘱咐于我的，直到某些确定之事向我显明，使我知所航向。\par

\SourceRecord{Translated by J.G. Pilkington.；Nicene and Post-Nicene Fathers, First Series；Vol. 1.；Edited by Philip Schaff.；Buffalo, NY: Christian Literature Publishing Co.,；1887.；<http://www.newadvent.org/fathers/110105.htm>.}

% structure: p0001 source=h1 target=section reason=page-title

\section{忏悔录（第六卷）}

年届三十，他在\Person{盎博罗削}讲道的劝诫下，日益发现天主教教义的真实性，并思量着如何更好地规范自己的生活。\par

% structure: p0003 source=h2 target=subsection reason=relative-heading-rank; base=h2

\subsection{第一章 母亲随他来到\Place{米兰}，宣称若不见儿子接受天主教信仰，自己绝不瞑目。}

1. 我自幼的希望啊，祢在何处？祢往哪里去了？因为事实上，难道不是祢创造了我，使我与田野的走兽和天空的飞鸟有别吗？祢使我比它们更有智慧，我却仍在黑暗险滑之处徘徊，向自身之外寻求祢，却没有找到我心中的天主；且沉入大海深处，怀疑并绝望于找到真理。这时，我的母亲因她的虔诚而变得坚强，渡海越陆，跟随我而来，在一切危险中因祢而觉得安全。因为在海上遇险时，她反而安慰那些水手（不惯航行的旅客受惊时，通常反倒去向水手寻求安慰），确保他们能安全抵达，因为她在神视中已由祢确知此事。她见我处于严重的危险中，因为我对找到真理已感绝望。但当我向她透露，我已不再是摩尼教徒，虽然还不是天主教信徒，她并没有像遇到意外喜事那样欢喜雀跃；虽然此刻她对我那部分苦难已感安心，她曾为我哀哭，视我如死人，但相信我会被举扬到祢面前，她用思念的担架将我抬到祢面前，好使祢对寡妇的儿子说：\BibleQuote{青年人，我对你说：起来吧！}他便复活，开口说话，祢将他交还他母亲。\BibleRef{路 7:12}因此，当她知道她每日流泪祈求的事已大部分成就时——虽然我尚未把握真理，但已脱离谬误——她的心并未因任何剧烈的兴奋而波动。的确，正因她完全信赖那应许了全部的祢会赐予余下的，她极其平静，满怀信心地回答我说：\Quote{她相信基督，她离世之前必会见到我成为天主教信徒。}这是她对我说的话；但对祢，仁慈的泉源啊，她更频繁地倾流祈祷和眼泪，求祢加速祢的助佑，光照我的黑暗；她更殷勤地赶往教堂，紧靠着\Person{盎博罗削}的讲道，祈求那涌到永生的水泉。\BibleRef{若 4:14}因为她爱这位人如同天主的使者，她知道是藉着他，我才暂时被引入如今这种困惑不安的状态，她完全相信我会经历一次仿佛更剧烈的突发过度后，从疾病转入健康，医生称之为\OriginalTerm{转机期}{crisis}。\par

% structure: p0005 source=h2 target=subsection reason=relative-heading-rank; base=h2

\subsection{第二章 母亲因\Person{盎博罗削}的禁令，停止敬礼殉道者的纪念。}

2. 因此，有一次，我的母亲按照她在\Place{非洲}的习惯，带着一些糕饼、面包和酒来到为纪念圣人而建的圣堂，却被堂役阻止；当她得知这是主教的禁令时，便如此虔敬而顺从地接受了，以致我自己都惊讶，她竟如此容易地归咎自己的习惯，而不质疑他的禁令。因为嗜酒并未占据她的心灵，对酒的爱好也没有刺激她憎恨真理，不像许多男女，一听到\Quote{节制之歌}便作呕，如同醉酒之人饮净水一般。但她每次带着盛有节日食物的篮子，自己先尝一些，再把其余的施散，却从不让自己饮用超过一小杯的酒，按照她清淡的口味掺水稀释，只是出于礼貌才沾唇。如果有多座安眠圣人的圣堂需同样敬礼，她仍随身携带同一个杯子，到处使用；而这酒不仅掺了大量的水，而且因携带来去变得十分温热，她只以极小口分给周围的人；因为她寻求的是他们的热心，而非享乐。因此，当她发现这习惯被那位著名的讲道者和最虔诚的主教所禁止，即使是针对那些愿意有节制地使用的人，以免给醉酒之人提供放纵的机会，而且因为这些所谓敬礼亡者的节日，极似外邦人的迷信，她便极甘愿地戒绝了。她学会了，代替盛满地上果实的篮子，给殉道者的圣堂带去充满更纯洁祈求的心，并尽其所能施舍给穷人；好使主圣体的共融能在此处妥善举行，因为在那些地方，殉道者们效法基督的苦难，被祭献并蒙受冠冕。然而，主，我的天主，在我看来——我的心在祢面前也是如此想——若这禁令来自她并非如同喜爱\Person{盎博罗削}那样喜爱的其他人，她或许不会如此轻易地放弃这习惯；她因顾念我的救恩，极其深爱\Person{盎博罗削}；而他确实也爱她，因她极为虔诚的品行，在善工上如此\BibleQuote{心神火热}\BibleRef{罗 12:11}而常去教堂；因此他每逢见到我，常会突然对她大加赞扬，庆幸我有这样一位母亲——却不太清楚她在我身上有怎样的儿子，一个对这一切都怀疑、且不认为生命之道可以寻获的儿子。\par

% structure: p0007 source=h2 target=subsection reason=relative-heading-rank; base=h2

\subsection{第三章 \Person{盎博罗削}忙于事务与研读，奥古斯丁很少能就\WorkTitle{圣经}问题向他请教。}

3. 那时，我甚至在祈祷中也不叹息求祢帮助我；我的心思完全专注于求知，热衷于辩论。而\Person{安博罗削}本人，按世俗的标准，我视他为幸福的人，因为那样显赫的人物都尊敬他；只是他的独身生活在我看来是一件痛苦的事。但他怀有何等的希望，为战胜伴随他各种优长而来的试探作了何等的挣扎，在逆境中有何等的安慰，当他的内心隐藏的口在默思祢的食粮时，这食粮给予他怎样的甘饴喜乐，我既不能猜想，也未曾体验。他也不知道我的窘迫，不知道我所陷入的危险深坑。因为我不能如愿地向他请教我所想知道的事，由于他被一群忙碌的人所隔绝，使我无法与他交谈，那些人将自己的软弱全都交托给他。当他不和他们在一起时（这时间很少），他不是在给身体补充必要的营养，就是在阅读中给心灵补充营养。但在他阅读时，他的眼睛扫描着书页，内心探寻着意义，但他的声音和舌头却保持沉默。很多时候，当我们去见他时（因为没有人被禁止进入，而且他的习惯是不通报来访者的到来），我们都看到他这样默读，从来不见他做别的。我们长久地默默坐着（因为谁敢打断如此专注的人呢？），便只好离开，猜想他在那短暂的时间里，为了恢复心智，免受别人事务的喧嚣干扰，是不愿被打扰的。也许他担心，如果他研读的作者在表达上有些模糊，某个好疑而专心的听者会要求他解释，或讨论一些较深奥的问题，这样占用了他的时间，他就不能按所愿的翻阅那么多卷册了；尽管保存他那很容易变得微弱的声音，可能是他默读的更真实原因。但不管他这样做的动机是什么，在这样一个人身上，无疑是一个好的动机。\par

4. 然而，我确实找不到机会从祢的那位如此圣洁的代言人——他的胸怀——那里探知我所渴望的事，除非能简要地谈论一下。但我内心翻腾的思绪，需要他完全闲暇时，好让我向他倾诉，却始终无法找到他如此闲暇的时候。我每个主日都听他在民众中\BibleQuote{正确地讲解真理的言语}\BibleRef{弟后 2:15}，并且我越来越确信，那些欺骗我们的人所编织的、针对圣经的各种诡诈诽谤的结，是可以解开的。然而，当我同时理解了\Quote{是照那创造他者的肖像造的}这句话，并不是像祢属神的子女们（祢借着恩宠由慈母天主教会所重生的人）所理解的那样——仿佛他们相信并想象祢被人的形状所限制——尽管我当时对\OriginalTerm{属神实体}{spiritual substance}的性质连最模糊的概念也没有——但我还是欣然地感到羞愧，羞愧自己这么多年来一直狂吠，不是反对\Emph{公教}信仰，而是反对肉性想象的虚构。因为在这件事上，我既是不敬的，又是轻率的：我本应通过探询去学习的事，我却定断并加以谴责。因为祢，至高而又至近，最隐秘而又最临在的，没有或大或小的肢体，而是全然无所不在，却不在空间中的任何地方，也不是具有这样形体的；然而祢却照祢自己的肖像造了人，看啊，人从头到脚却被空间所限制。\par

% structure: p0010 source=h2 target=subsection reason=relative-heading-rank; base=h2

\subsection{第四章 他认识到自己见解的错误，并铭记\Person{安博罗削}的话语。}

5. 那么，由于我不知道祢的肖像应该如何存在，我本应叩门，提出疑问，问明应如何信仰，而不是侮辱性地反对，仿佛信仰是这样的。因此，对什么是应该坚持的真理的焦虑，越是猛烈地啃啮我的灵魂，我越是感到羞愧，因为我曾长久地被确定性的许诺所欺骗和迷惑，竟以幼稚的错误和狂妄，把这么多不确定的东西当作确定的东西来唠叨。因为它们是虚假的，这是我后来才明白的。然而，我确定的是它们是不确定的，而我从前却以盲目的争辩，指控祢的\Emph{天主教会}教导这些东西，尽管我尚未发现她教导真理，但她确实没有教导我所猛烈指控她的那些东西。于是，我感到混乱和转变，我的天主啊，我欣喜地发现，那唯一的教会，祢唯一圣子的身体（我在婴儿时，在这教会内被冠以基督的名字），并不看重那些幼稚的琐事，在她的健全教义中，并不主张任何将祢——万物的创造者——限制在空间中的论点，不论这空间多么巨大和广阔，却仍被人的形状的限制所四面围住。\par

我也欢欣，因为那些古老的律法和先知的圣经摆在我面前，供我细读；如今不再用以前那种眼光去看，那时我觉得它们极其荒谬，并指责你的圣人们如此想，但其实他们并非如此想；我也愉快地听到\Person{安博罗削}在向民众讲道时，常常极恳切地推荐这句经文作为准则——\BibleQuote{文字叫人死，圣神却叫人活}；同时，他揭开那神秘的帷幔，以\OriginalTerm{属神的方式}{spiritualiter}揭示了那些若按\BibleQuote{文字}理解则似乎教导荒谬道理的经文——他对此的讲解毫不冒犯我，尽管他所教导的那些事，我还不知道是否真实。在这整段时间里，我抑制自己的心，不轻信任何事，生怕冒失跌倒；但由于悬而不决，我更加痛苦。因为我渴望对那些看不见的事，能有如我知道七加三等于十那样确定的把握。因为我并非疯狂到认为这些事无法理解；但我希望其他事物也能如此清楚，无论是那些不呈现于感官的形体之物，或是那些属神的事物——对后者，我除了用形体的方式外，不知如何设想。借着相信，我原可以得医治，好使我灵魂的视力得以澄清，能以某种方式导向你的真理，那真理永远常在，绝不缺失。但犹如一个人若曾尝试过不良的医生，便不敢再信任好医生，我灵魂的健康也是如此；它只能借着相信而得医治，但为了避免相信谬误，它拒绝被治愈——抗拒你的双手；你为我们准备了\OriginalTerm{信仰的药剂}{medicamenta fidei}，并将它们施用于整个世界的病痛，且赋予它们如此大的权威。\par

% structure: p0013 source=h2 target=subsection reason=relative-heading-rank; base=h2

\subsection{信德是人生活的基础；人不能发现圣经所启示的真理。}

然而，由此我逐渐被引导更倾向\OriginalTerm{公教教义}{doctrinam catholicam}，我感觉到它以一种更为温和与诚实的态度，命令人相信那些未经证明的事（无论是因这些事虽可证明，并非人人能够明白，抑或根本不能证明），这远胜于摩尼教徒的方法；在摩尼教徒那里，我们的轻信被他们狂妄地应许知识所嘲弄，随后又把许多极其荒诞无稽的事强迫我们相信，就因为它们无法被证明。此后，主啊，你以最温柔和最慈悲的手，一点一点地牵引并平静我的心，使我思量：有多少事情我从未见过，也不曾在发生时在场，诸如世俗历史中的许多事，以及许多我未见过的地方与城市的叙述；还有许多关于朋友、医师，有时是这些人、有时是那些人的事，除非我们相信，否则此生将一事无成；最后，我以何等不可动摇的确信，相信自己的生身父母是谁——这事我除了道听途说外，绝无可能知道——考虑到这一切，你说服我：那些相信你的圣经（你以如此大的权威，将之建立于几乎万邦之中）的人不应受责备，反而那些不信的人才该受责备；而且那些会对我说\Quote{你怎么知道这些圣经是借着唯一真实且至真天主的圣神传授给人类的？}的人，也不该听从。因为这正是最应相信的同一件事：既然没有任何亵渎性问题的争论（我在那些自相矛盾的哲学家著作中读到过许多这样的争论），能一度从我心中夺去\Quote{你存在}的信念——无论你是怎样的，尽管我并不知道——或\Quote{人类事务的管理属于你}的信念。\par

这是我多少相信的，有时更坚定些，有时稍弱，但我始终相信你存在，且关怀我们，尽管我既不知该如何思虑你的\OriginalTerm{本体}{substantia}，也不知哪条道路能通往你或引回你那里。既然我们单凭理性过于软弱，无法寻获真理，因此需要圣经的权威，我已开始相信：你绝不会将如此卓越的权威赋予这些圣经，使之遍传全地，除非你的旨意正是要人藉着它相信你、寻求你。先前我在圣经中觉得不合理、常使我反感的地方，自从听到其中多处被合理地解释后，我都归之于\OriginalTerm{奥秘}{mysteria}的深度；圣经的权威在我看来，也变得更为令人敬畏，更配受信仰的尊崇；因为，它虽让人人都可阅读，却将其秘密的威严蕴藏在高深内涵之中，以最浅显的语言和谦卑的风格俯就一切人，同时又吸引那些心灵不轻浮的人深入钻研；这样，它便能将众人纳入它共通的怀抱，并通过窄路引领少数人到你这来——然而比起它若没有立在如此崇高的权威之巅，也未以其圣洁的谦卑吸引众人进入它怀抱的情况，这少数人仍要比那时更多。我默想这些事，你与我同在；我叹息，你俯听了我；我动摇，你引导了我；我漫游在这世界的广阔道路上\BibleRef{玛 7:13}，你却没有离弃我。\par

% structure: p0016 source=h2 target=subsection reason=relative-heading-rank; base=h2

\subsection{论真正喜乐的根源与原因——以欢喜的乞丐为例}

9. 我渴望荣誉、利益、婚姻；祢却嘲弄了我。在这些欲望中，我经受了极度的苦楚，祢越是仁慈，就越不容许任何非祢的事物在我心中生出甘甜。主啊，请看我的心，祢愿意我回忆这一切，并向祢认罪。现在，让我的灵魂依附于祢吧，祢已将它从那紧粘不放的死亡罗网中解救出来。它曾是何等可怜！祢激起了它的伤痛感，使它离弃一切，转向祢——祢超乎万有之上，没有祢，万有皆为虚无——转向祢而得医治。那时我是何等的可怜，而祢又是如何对待我，让我在那一天意识到自己的可怜；那一天我正准备为皇帝朗诵一篇颂词，其中我要说许多谎言，而知道我在说谎的人却要为之鼓掌；我的心为这些忧虑而喘息，因着那吞噬思想的狂热而沸腾。因为，当我走在\Place{米兰}的一条街上时，我看见一个穷乞丐——那时，我想他定是酒足饭饱——正在谈笑欢乐；我叹了口气，对身边的朋友们说起我们疯狂所导致的诸多愁苦；我们费尽心机——就像我当时在欲念的鞭策下，拖着我那不幸的重担，越拖越沉重——所追求的，无非是那乞丐已在我们之先得到的喜乐，而他或许永远也无法得到！因为他靠乞讨几个小钱所获得的，正是我费尽多少可怜曲折的心机所要图谋的——那短暂幸福的快乐。他确实没有真正的喜乐，然而我，怀着这些野心，所寻求的却是更加虚假的东西。其实他是欢乐的，我是忧愁的；他无忧无虑，我满是恐惧。但是，若有人问我，是愿意快乐还是恐惧，我会回答快乐。再问我，是愿意像他那样，还是像我当时那样，我仍会选择做我自己，尽管被忧虑和恐惧所困扰，但这却是出于乖谬；因为事实果真如此吗？我不应该因为自己比他更有学问，就自以为优于他，因为我对这学问并无乐趣，而只是想用它来取悦于人；并非为教导人，而只是为讨好人。因此，祢也用祢惩戒的杖打断了我的骨头。\BibleRef{箴 22:15}\par

10. 那么，让那些对我灵魂说\Quote{一个人的喜乐来自何处，这是有区别的。那乞丐沉醉于酒中取乐；你却渴望在荣耀中得乐。}的人离开我的灵魂吧。主啊，什么荣耀？那不是存在于祢内的荣耀。因为正如他的不是真喜乐，我的也不是真荣耀；而且它更颠覆了我的灵魂。他那一夜会消化他的醉意，而我却有许多夜与我的荣耀同眠，又与之同起，不知还要多少次与之同睡同起。的确，\Quote{一个人的喜乐来自何处，这是有区别的。}我知道是如此，也知道忠信望德的喜乐远超这等虚幻。的确，那时他胜过了我，因为他实实在在比我更幸福；不仅因为他完全沉浸在欢笑中，而我被忧虑撕扯，更因为他通过祝愿得到了酒，而我却通过谎言追逐骄傲。当时我对亲爱的朋友们大概说了这些话，我也常常观察到他们的情形如何反映我自己的情形；我发现我过得很糟，便烦恼，使那糟糕加倍。若有任何顺境向我微笑，我却厌烦去攫取它，因为几乎在我能抓住它之前，它就飞走了。\par

% structure: p0019 source=h2 target=subsection reason=relative-heading-rank; base=h2

\subsection{他引领被\OriginalTerm{竞技表演}{Circensian Games}狂热攫住的朋友\Person{阿黎庇乌斯}归向正道}

11. 我们这些像朋友一样生活在一起的人，共同哀叹这些事，但尤其经常和亲密地，我与\Person{阿黎庇乌斯}和\Person{内布里迪乌斯}讨论这些。\Person{阿黎庇乌斯}与我同乡，他的父母在那里地位最高，但他比我年轻。因为他曾在我门下求学，起初是在我们本城，后来在\Place{迦太基}；他非常敬重我，因为我在他眼中显得善良而有学问；而我则因他天生热爱德行而看重他，这在年纪不大的人身上已相当卓越。然而，\Place{迦太基}习俗的漩涡（那里的人们狂热地追逐这些轻浮的表演）诱使他陷入了\OriginalTerm{竞技表演}{Circensian Games}的狂热。当他可怜地在这狂热中颠簸时，我正在那里教授修辞学，并开设了一所公开的学校。由于我与他父亲之间产生了一些嫌隙，他那时还不听我的课。当时我发现他如何致命地痴迷于竞技场，并深感痛心，因为他似乎——倘若他还没有这样做过——要将这莫大的天赋弃之不顾了。然而，我既无法以朋友的善意，也无法以师长的威权，通过劝告或某种约束来挽回他。因为我想象他对我的感情与他父亲一样；但他并非如此。因此，在那件事上，他不顾父亲的意愿，开始来问候我，并进入我的讲堂，听一会儿便离去。\par

12. 但是我竟忘了与他商量，好教他不致因盲目固执的空虚娱乐而毁掉如此卓越的才智。然而，主啊，祢掌管祢所造万物的舵，并未忘记他，他终有一天将成为祢的儿女，祢圣事的主持者；为使他的改正明显归功于祢自己，祢便藉着我成就了此事，而我却毫不知情。一日，我如常坐下，学生们在我面前，他进来，向我致意，坐下，专心听我正讲的题目。恰巧我手头有一段文字，在解释时，我想起一个从\OriginalTerm{竞技场游戏}{Circensian games}借用的比喻，觉得它能使我要传达的意思更加生动明白，同时辛辣地讥讽那些被那种疯狂所俘虏的人。我们的天主啊，祢知道，当时我并无治愈\Person{阿利比乌斯}此疾的念头。但他却认为我是冲他而来，以为若非为他，我不会那样说。别人或许会因此对我怀恨，这位可敬的青年却以此作为对己恼怒的理由，并更热切地爱我。因为祢早已说过，并记在祢的书中：\BibleQuote{责备智慧人，他必爱你。}\BibleRef{箴 9:8} 但我并未责备他，而是祢，利用一切有意的和无意的，按祢自己所知的秩序（那秩序是正义的），从我的心灵和口舌中取出火炭，用以燃起并医治那垂危的有望心灵。愿那不思祢仁慈的人，在赞美祢时缄默不语——我这内心最深处正在向祢宣认这仁慈。就因那番话，他从深渊中冲出——那座他自己甘愿投入、并被其悲惨娱乐所蒙蔽的深渊；他唤醒心神，以坚决的节制；于是竞技场娱乐的一切污秽都从他身上脱落，他再也不去接近了。此后，他劝说心有不甘的父亲允许他做我的学生。他父亲让步了，同意了。\Person{阿利比乌斯}再次来听我讲课后，便陷入了与我相同的迷信，喜爱\OriginalTerm{摩尼教徒}{Manichæans}那种炫耀的节制，他以为那是真实而不虚伪的。然而，那是一种愚昧而诱人的节制，网罗宝贵的灵魂——他们尚不能达到德性的高度，易被那仅是影子般虚假德性的浮饰所欺骗。\par

% structure: p0022 source=h2 target=subsection reason=relative-heading-rank; base=h2

\subsection{第八章 同一人，在\Place{罗马}时，被他人引入圆形剧场，为角斗比赛所迷。}

13. 他并未放弃父母所引诱他追求的那条世俗之路，先我一步前往\Place{罗马}学习法律；在那里，他以异乎寻常的方式，怀着难以置信的热切投入角斗表演。本来，他对这类表演深恶痛绝，极为反感；可有一天，他偶然遇到几位相识和同学用完餐归来，他们以友善的暴力，将他强拉硬拽到圆形剧场——那天正上演残忍而致命的表演——他一边激烈反对、奋力挣扎，一边抗议说：\Quote{就算你们把我的身体拖到那个地方，按在那里，难道你们还能强迫我用心看吗？我将身虽在场，魂却缺席，这样便能胜过你们和表演。}他们听了这话，仍拖他同行，或许是想试试他是否能说到做到。当他们抵达那里，各就其位，整个场地因无人性的比赛而骚动起来。但他紧闭双眼的门户，禁止心灵游荡于这等邪恶之事；真希望他当时也堵上耳朵！因为，当格斗中一人倒下，全场观众发出震耳欲聋的呼喊，他受到强烈刺激，被好奇心击败，原以为可以轻蔑视之、超然其上，无论发生什么，他睁开眼睛，便受了比他所要看的那个人身体更重的灵魂之伤；他比那激发巨大呼喊的倒下者更可怜地倒下了。那呼喊声从耳朵进入，打开了眼睛，好让攻击和击倒灵魂，而他的灵魂一直以来只是大胆，并非真勇敢；正因如此，它更软弱，因为它信赖自己，原本应该依靠祢的。因为，他一看见那鲜血，便同时吸入了一种残忍；他没有转开头，而是定睛观看，不知不觉地疯狂汲取，陶醉于那罪恶的竞赛，沉醉于那血腥的娱乐。他不再是从前的他，而成了他所加入人群中的一员，成了带他去那里之人的真正同伙。何须多言？他观看，他呼喊，他兴奋，带着那疯狂的冲动离去，不仅与最初怂恿他的人一起，甚至比他们更甚，还吸引别人前往。然而，祢以最强大、最仁慈的手，从这一切中将他一举拔出，教导他不要信任自己，而要信赖祢——不过这是很久以后的事了。\par

% structure: p0024 source=h2 target=subsection reason=relative-heading-rank; base=h2

\subsection{第九章 无辜的\Person{阿利比乌斯}被当成窃贼抓住，因一位建筑师的机智而获释。}

但这一切都被储存在他的记忆中，作为将来的良药。正如那件事一样：当他还在\Place{迦太基}跟我学习时，有一天正午，他在市场上默想需要背诵的内容（一如学子们惯常的练习），祢竟容许他被市场管理人员当作小偷捉住。我们的天主啊，我料想祢之所以容许这事，没有别的原因，就是要让这个将来要成为如此伟大的人，现在就开始学习：在判断案件时，人不该轻率地凭著鲁莽的轻信就定别人的罪。当时他正独自拿着写字板和笔，在审判台前踱步，忽然有一个年轻人，是学生之一，真正的小偷，悄悄地带着一把斧头，乘\Person{阿利比乌斯}没看见时，走到保护银匠铺的铅栅栏前，开始砍削铅皮。但斧头的声音被听见了，下面的银匠们骚动起来，派人去捉拿他们能找到的任何可疑之人。那贼听见声音，丢下斧头就逃，怕被抓住时人赃俱获。\Person{阿利比乌斯}虽没看见他进来，却在他出去时看到了他，并注意到他跑得多么快。出于好奇，想知道缘由，他就走进了那地方，发现了斧头，正站在那儿惊奇思索时，看，那些被派来的人抓住了他单独一人，手里还拿着斧头，正是那斧头的声音惊动了他们并引他们前来。他们抓住他，拖着他，召集市场上的商户们围拢过来，夸口说捉住了一个臭名昭著的小偷，接着就把他带走去见法官。\par

然而，他该受的教训也就到此为止。因为主啊，祢立刻来援救他的无辜，只有祢是这无辜的唯一见证。当他正被带往监狱或受刑时，遇见了一位建筑师，这人主管公共建筑。他们特别高兴碰见他，因为他们常被这人怀疑偷窃市场上失落的物品，好像这下终于可以让他确信这些偷窃都是谁干的了。而这位建筑师呢，曾屡次在一位元老院议员的家里见过\Person{阿利比乌斯}，他经常去那里拜见；因此立刻认出了他，就拉着他的手把他带到一边，询问如此大祸的缘由，听完整个事情的经过后，便命令当时在场的所有乌合之众（他们吵吵嚷嚷，威胁连连）随他同去。他们来到了那犯事的年轻人的家。门前有个少年，年纪很小，小到不会因害怕伤害主人而忍住不吐露全部实情。原来他曾跟着主人去过市场。\Person{阿利比乌斯}一认出他，就告诉了建筑师；建筑师把斧头拿给少年看，问他这是谁的。他立刻说：\Quote{是我们的}；再行盘问，他就全部供了出来。这样一来，罪行转嫁到了那家，那群开始向\Person{阿利比乌斯}耀武扬威的乌合之众蒙受了羞辱，而他——将来要成为祢圣言的宣讲者、祢教会中无数案件的审断者——离去时获得了更好的经验与教训。\par

% structure: p0027 source=h2 target=subsection reason=relative-heading-rank; base=h2

\subsection{\Person{阿利比乌斯}在审判中的惊人正直。\Person{内布里迪乌斯}与\Person{奥古斯丁}的持久友谊。}

因此，我在\Place{罗马}遇到了他，他以最牢固的情谊依恋着我，陪我去到\Place{米兰}，既不愿离开我，也可运用他所学的法律知识，与其说是为他自己，不如说是为取悦他的父母。他在那里曾三次担任陪审官，其廉洁不为他人所称奇，他反倒惊奇那些人竟能重黄金而轻正直。他的品性不仅受到贪婪的诱饵考验，也受到恐惧的刺激考验。在\Place{罗马}，他任意大利财政伯爵的陪审官。当时有一位极有权势的元老，许多人仰他的恩惠，也有许多人怕他。他妄想依仗惯常的权势，求得一件法律所禁止的事。\Person{阿利比乌斯}抗拒了；有人许以贿赂，他全心鄙视；动用威胁，他踩在脚下——众人无不惊奇如此罕有的气概，既不贪图一个人的友谊，也不畏惧他的敌意，而此人既如此有权有势，又因其无数行善或作恶的手段而声名赫赫。就连\Person{阿利比乌斯}担任其参议的那位法官，虽然也不愿做这件事，却未公开拒绝，而是把事情推给\Person{阿利比乌斯}，声称是他不许做；因为事实上，若法官做了，\Person{阿利比乌斯}就会做出相反的裁决。在学习方面，有一件事曾差点将他引偏——那就是他可以按\OriginalTerm{裁判官价格}{prætorian prices}替自己抄写书籍。然而，他征询了公正之后，便改弦易辙，选择了更好的路，他认为那阻碍他的公平，比那允许他的权势更为有益。这些都是小事，但\BibleQuote{在小事上忠信的，在大事上也忠信。}\BibleRef{路 16:10} 从祢真理之口所出的一切，都绝不可能落空。\BibleQuote{如果你们在不义的钱财上不忠信，谁还把真实的钱财委托给你们呢？如果你们在别人的财物上不忠信，谁还把属于你们的交给你们呢？}\BibleRef{路 16:11} 他既是如此的人，那时便依恋着我，并且也像我一样，犹豫不决，该选择什么样的生活道路。\par

奈布里迪乌斯也一样，他离开了家乡\Place{迦太基}附近的地方，也离开了通常居住的\Place{迦太基}城，撇下美好的祖业、房屋和不愿跟随他的母亲，来到\Place{米兰}，别无所求，只为能与我一同最热切地追寻真理和智慧。他像我一样叹息，像我一样彷徨，是真实生命的热切寻求者，最精深的隐微问题的钻研者。于是我们成了三个嗷嗷待哺的人，彼此叹诉着各自的渴望，等候祢按时赐下食粮。在祢因仁慈而让我们世俗追求中所受的一切辛酸里，当我们思考这一切苦难的结局为何临到我们时，黑暗笼罩了我们；我们转身叹息呼喊说：\Quote{这些事要到几时呢？}我们常常这样说；虽然这样说，却并未舍弃它们，因为我们尚未发现任何确定的东西，可供我们在舍弃后去投靠。\par

% structure: p0030 source=h2 target=subsection reason=relative-heading-rank; base=h2

\subsection{深陷严重错误而苦恼，默想开始新生活。}

我反复思索和回顾这些事，最使我惊讶的，是从十九岁那年开始，便燃起了追求智慧的愿望，决意寻获后便抛下一切虚妄希望的虚妄妄想。而今，我已年近三十，仍旧陷在同一泥沼中，贪享着眼前飞逝而毁灭我的事物，一面自语：\Quote{明天我必寻得；看，事情将清楚地显现，我必抓住它；看，\Person{福斯图斯}会来并解释一切！啊，你们学园派的大人物们，那为生活定秩序而有所确知的事，真的无法获得！不，让我们更用心探求，不要绝望。看，教会书籍中那些从前似乎荒谬的事，如今却不显得荒谬了，且可能有其他端正的解释。我要把脚立在儿时父母安置我的那个台阶上，直到发现清楚的真理。但是，何时何地去寻求呢？\Person{盎博罗削}没有空闲，我们也没有空闲去读。我们从哪里找到这些书？何时何地购得？向谁借阅？让我们定下时间，分出特定的时辰用于灵魂的健康。心中燃起了大希望，天主教信仰并不教导我们所设想并加以诬蔑的东西。她的博学之士认为，相信天主受限于人形是可憎的。我们岂可怀疑\BibleQuote{敲门}，为使其余的\BibleQuote{敞开}吗？\BibleRef{玛 7:7}上午要教学生；其余的时间我们怎样利用？为什么不着手这事？但是，何时去拜访大人物朋友呢？我们需要他们的恩惠。何时预备学生向我们购买的东西？何时休息，使心灵从忧劳中放松？}\par

\Quote{让一切灭亡吧，让我们抛开这些虚空的浮华，专心于寻求真理！生命悲惨，死亡无常。倘若死亡突然袭来，我们将以怎样的状态离开，又将去哪里学习在此地所忽视的？岂非更要为此疏懒受罚？如果死亡果真彻底结束一切忧虑和感觉呢？这一点也须探究。但求天主莫使如此。基督信仰的权威，其如此崇高的高度传遍全世界，并非没有理由，也非虚妄。倘若肉身的死亡就毁灭了灵魂的生命，就绝不会为我们成就如此伟大的事。那么，我们为何迟延舍弃世俗的希望，而全心寻求天主和真福生活呢？且慢！那些事物也令人愉悦，且有几分不小的甘甜。我们不可轻易抛舍，因为再回头是可耻的。看，现在获取某个显赫职位是件大事！我们还能贪求什么？我们有一群有势力的朋友，即使别无所长，如能抓紧，一个裁判官职位也会到手；再娶一个带来一些财产的妻子，使她不增添我们的开支，这便达到愿望的极致了。许多伟大而堪为模范的人，已在婚姻状况中致力于追求智慧。}\par

我口中叨念着这些，风向转变，我的心随风飘荡摇动，时光流逝；我却迟迟不肯转向主，日复一日拖延生活在祢内，却又每日死亡在己身之中。我贪慕幸福的生活，却又怕它所在之处，逃离它，又去寻求它。我以为若被剥夺女子的怀抱，我便会太为不幸；而对于祢那医治此病的仁慈良药，我未曾想到，也未曾尝试。至于\OriginalTerm{节德}{continency}，我想象它是受我们自身力量控制的（虽然我未在自己身上发现它），竟愚昧地不知经上写的：若非祢赏赐，无人能节德；而只要我由心底叹息叩祢的耳朵，并以坚定的信德将我的忧虑托付于祢，祢必会赏赐。\par

% structure: p0034 source=h2 target=subsection reason=relative-heading-rank; base=h2

\subsection{与阿利比讨论独身生活。}

21. 事实上，是\Person{阿利比乌斯}阻止了我结婚，他说这样我们绝不可能像我们长久以来所渴望的那样，在爱智慧中拥有不受干扰的闲暇而共同生活。因为他自己在这事上如此贞洁，以至令人惊奇——更令人惊奇的是，他在早期的青年时代曾涉足那条路，却未沉溺其中；相反，他对此感到悲伤和厌恶，从此以后直到如今都过着极其节制的生活。但我用那些已婚者仍爱慕智慧、得蒙天主悦纳、并与朋友忠信友爱的榜样来反驳他。我从这些人精神的伟大中远落其后，为肉体的疾病及其致命的甘甜所束缚，拖曳着我的锁链，害怕被解开；而且，仿佛它挤压我的伤口，我拒绝了他善意的规劝，如同拒绝那想为我解开锁链的人之手。更有甚者，那蛇正是借着我对\Person{阿利比乌斯}说话，用我的舌头在他的路上编织并铺设了悦人的罗网，使他那尊贵而自由的脚被缠住。\par

22. 因为当他感到惊奇，我这个他并非不看重的人，竟如此深陷那欢愉的粘鸟胶中，以至每当讨论此事，我都断言独身生活对我来说是不可能的；而当我看到他惊奇时，我为自己辩护说，他暗中偷偷尝试过的生活（对此他现在只有模糊的记忆，因此可以轻易地毫不惋惜地轻视）与我长期熟悉的生活之间有着巨大的差别，倘若加以婚姻的尊荣之名，他对我不能鄙视那种生活就不会感到惊讶了——于是他也开始想要结婚，并非被这种欢愉的贪欲所压倒，而是出于好奇。因为，他说他急切地想知道那东西到底是什么，没有它，我那令他如此喜悦的生活在我看来就不是生活，而是刑罚。因为他的心灵，摆脱了那种束缚，对我的奴役感到震惊，并且通过这种震惊，正走向尝试它的愿望，从尝试的愿望走向实际的尝试，然后，或许会陷入他如此震惊的奴役之中，因为他正准备\BibleQuote{与死亡立了盟约}\BibleRef{依 28:15}；\BibleQuote{而爱危险的人必陷于危险之中}\BibleRef{德 3:27} 因为，无论婚姻的尊荣在妥善安排已婚生活和养育子女的职务中有多大，对我们影响甚微。但最折磨我、已使我成为其奴隶的，是一种满足贪得无厌的欲望的习惯；而对于即将受奴役的他，则是钦佩的惊奇吸引着他。我们就这样陷入这种状态，直到祢，至高者啊，不舍弃我们的卑微，可怜我们的困苦，以奇妙而隐秘的方式前来拯救我们。\par

% structure: p0037 source=h2 target=subsection reason=relative-heading-rank; base=h2

\subsection{第十三章 在母亲的催促下娶妻，他寻找一位令他中意的少女}

23. 人们积极努力为我物色妻子。我追求，我订婚，母亲为此事费尽心机，希望我一旦结婚，带来救恩的圣洗便能洗净我；她因我日渐准备妥当而欣喜，说她的愿望和祢的许诺在我的信德中得以实现。那时，我们确实按照我的请求和她的愿望，每天以强烈的心声向祢呼求，恳求祢藉异象向她启示一些关于我未来婚姻的事。但祢并未应允。她确实看到一些虚幻而离奇的事，如同专注于此事的人心神热切所幻象出来的；她将这些事讲给我听，却不像祢向她显示什么时那样带有惯常的确信，而是轻视它们。因为她宣称，她能凭着某种她无法用言语表达的感觉，分辨祢的启示与她自己的心神梦境。然而，婚事仍在推进，向一位距适婚年龄还差两年的少女求婚；因她令人中意，便等她成年。\par

% structure: p0039 source=h2 target=subsection reason=relative-heading-rank; base=h2

\subsection{第十四章 与朋友建立共同家庭的设计迅速受阻}

24. 我们许多朋友商议并憎恶人生扰攘的烦恼，考虑并几乎决定安闲生活，远离人群的喧嚣。要实现这一点，方法如下：我们各自尽其所能提供财物，组成一个共同的家园，以便透过我们友谊的真挚，没有什么东西更属于某个人而非他人；而是所有物品，从众人而来，整体属于每个人，也整体属于所有的人。我们认为这个团体可以由十个人组成，其中有些人非常富有，尤其是\Person{罗曼尼安}，他是我们的同乡，我自幼的亲密朋友，当时因重大的事务被召到朝廷；他对此计划最为热心，而且他的意见在推荐这项计划时极有分量，因为他的产业远比其余的人丰厚。我们还安排每年选出两位管事，负责供应一切所需，而其他人则免受干扰。但当我们开始考虑我们中有些人已有妻室，另一些人希望娶妻，不知妻子们是否会容许这样做时，那安排得如此妥善的计划，在我们手中化为碎片，彻底破灭并被弃置一旁。于是我们又陷入叹息呻吟之中，我们的脚步再次踏上了世界\BibleQuote{那宽阔且被人踏惯的道路}\BibleRef{玛 7:13}；因为我们心中有许多想法，但祢的谋略永远常存。祢以祢的谋略嘲笑我们的计划，并准备祢自己的安排，决意按时赐给我们食粮，张开祢的手，使我们的灵魂充满祝福。\par

% structure: p0041 source=h2 target=subsection reason=relative-heading-rank; base=h2

\subsection{第十五章 他遣走一位情妇，又选择另一位}

25. 在此期间，我的罪不断加增。我那情妇，作为我结婚的障碍，被从我身边扯离；我那依恋她的心，备受折磨，受伤流血。她返回了\Place{非洲}，向祢立誓不再认识其他男人，将我与她所生的儿子留给了我。然而我，不幸的人，无法效法一个女人，因迫不及待——由于我须等待两年才能娶到我所追求的女子，我并非热衷婚姻，实为情欲的奴隶——便另觅了一个（虽非妻子），好使我的灵魂之疾，借恒常习惯的束缚得以滋养、维持其活力，甚或加剧，直至进入婚姻的国度。我与先前情妇分离所造成的创伤仍未痊愈，反在灼烧和最剧烈的痛苦之后变得坏死，疼痛虽已麻木，却更加绝望。\par

% structure: p0043 source=h2 target=subsection reason=relative-heading-rank; base=h2

\subsection{第十六章 对死亡与审判的恐惧，以及相信灵魂不死，使他从先前信奉\Person{伊壁鸠鲁}见解的恶行中转回}

26. 赞颂归于祢，荣耀归于祢，仁慈的泉源啊！我愈是悲惨，祢便愈亲近。祢的右手常准备将我从泥沼中拉出，并将我洗净，我却茫然不知。没有任何事物能将我从肉欲享乐的更深渊薮中唤回，唯有对死亡与祢未来审判的恐惧；这恐惧，在我一切意见的摇摆中，从未离开我的心。当我与朋友们，\Person{阿里匹乌斯}与\Person{内布里迪乌斯}，争论善恶的本质时，我认为\Person{伊壁鸠鲁}在判断上已赢得了胜利——如果我不相信死后灵魂仍有生命，且有赏报之地，而\Person{伊壁鸠鲁}不信这些。我问道：\Quote{假定我们不死的，且生活在永恒肉身快乐的享受中，全无失去这快乐的恐惧，那我们为何并不幸福，为何还要寻求其他呢？}——我竟不知道，这本身就是我大不幸的一部分：我既如此沉溺与盲目，竟不能分辨那为自身之故当被追求的光荣与美善之光，这光非肉眼所能见，唯内在的人可得而睹。我也不幸地不去思考，这些话的矿脉从何而来，我竟能愉快地与朋友们谈论这些可厌之事。按照我当时的幸福观念，无论我多么丰盈地拥有肉欲的享乐，没有朋友，我也不能使自己幸福。我确实为朋友本身的缘故才爱他们，也知道自己同样为朋友本身的缘故被爱。邪曲的道路啊！那狂妄的灵魂有祸了，它竟希望若离弃了祢，能寻得更好的事物！灵魂辗转反侧，仰卧、侧卧、俯伏，处处是坚硬，唯有祢是安息。看啊，祢就在近处，将我们从可悲的流浪中解救出来，使我们稳立于祢的道路上，安慰我们，并说：\Quote{奔跑吧；我将扶持你，是的，我将引领你，在那里我也将扶持你。}\par

\SourceRecord{Translated by J.G. Pilkington.；Nicene and Post-Nicene Fathers, First Series；Vol. 1.；Edited by Philip Schaff.；Buffalo, NY: Christian Literature Publishing Co.,；1887.；<http://www.newadvent.org/fathers/110106.htm>.}

% structure: p0001 source=h1 target=section reason=page-title

\section{\WorkTitle{忏悔录}（第七卷）}

他回顾自己青年时期的开始，即他三十一岁时，在那时，关于天主的本性及恶的起源等极为严重的错误被辨明，\WorkTitle{圣经}被更准确地认识，他终于对天主有了清晰的认识，尽管尚未正确理解耶稣基督。\par

% structure: p0003 source=h2 target=subsection reason=relative-heading-rank; base=h2

\subsection{第一章 他确实不以人形看待天主，而是将天主视为弥漫空间的形体物质}

1. 我邪恶而可憎的青年时期如今已死去，我正步入成年初期：随着年龄增长，我在虚荣中变得越发污秽，除了亲眼所见的实体外，我想象不出任何别的实体。天主啊，我并未把祢想象成人形。自从我开始听到些许智慧，便始终避免如此设想；并且我欣喜地在我们灵魂的母亲——祢的天主教会——的信德中找到了同样的智慧。但除此之外，我不知道该把祢想象成什么。而我，一个人，且是那样一个人，想方设法要构想祢——至高无上、唯一的真天主；我内心深处确实相信祢是\OriginalTerm{不可朽的}{incorruptible}、\OriginalTerm{不可侵犯的}{inviolable}、\OriginalTerm{不变迁的}{unchangeable}；因为尽管我不知缘由，也不明其理，却非常清楚地看出并确信：能被朽坏的，必定劣于不能朽坏的；毫不犹豫地，我宁取那不可侵犯的，而不取那能被侵犯的；我认为那不经受任何变化的，胜于那能变化的。我的心猛烈地呐喊，反对我所有的幻象，并力图用这致命一击，从我灵魂的眼目前驱走所有萦绕其周围的不洁人群。看啊，它们几乎未被驱散，眨眼之间，便成群地重新压迫我，撞击我的面容，使其朦胧不明；以至于，尽管我并未把祢想象成人的形体，却不得不将祢想象为空间中某种有形的物体，这物体或者灌注于世界之中，或者无限地弥漫在世界之外——即便是那\OriginalTerm{不可朽的}{incorruptible}、\OriginalTerm{不可侵犯的}{inviolable}、\OriginalTerm{不变迁的}{unchangeable}者，正是我宁取于那可朽、可侵犯、可变者之上的。因为凡我所构想的，一经剥夺了这种空间，在我眼中便成为空无，是的，完全虚无，甚至不是虚空，就好像一个物体从其位置上移走，而该位置不存有任何物体，无论是土质的、属地的、水性的、气态的或是天界的，只留下一个虚空之处——好似一个空旷的虚无。\par

2. 因此，我心肠愚顽，连对自己都不清晰，凡是未被延展于某些空间、未被扩散、未被聚集、未被膨胀，或者不能或不会承受这些度量中任何一种的事物，我都断定它全然是虚无。因为那时我的心所游弋的，正是我的双眼惯常掠过的那些形式；我并未察觉，我用来形成那些意象的观察本身，并非此类，然而它若本身不是某种伟大的事物，便无法形成那些意象。同样，我设想祢——我生命的生命——，如同充满无限空间的庞大实体，从各处穿透整个世界的质量，并在世界之外，以无可无量的空间向四面八方延伸；以至于大地拥有祢，天拥有祢，万物拥有祢，它们被限定在祢内，而祢却不在任何地方。因为，就像大地之上的空气，并不阻隔太阳的光穿过它、渗透它，不是通过爆裂或切割，而是完全充满它一样，我便设想不只是天、气、海的形体，甚至大地的形体，对祢都是通透的，在它最大与最小的部分中，都能渗透，以接纳祢的临在，藉着一种隐秘的感应，内在地、外在地治理着祢所创造的一切。我之所以如此推测，是因为我无法想到别的事物；但这并非真实。因为按照这种想法，大地更大的部分将包含祢更多的部分，较小的部分则包含较少的部分；万物都将如此充满祢，以至于一头大象的身体所包含的祢，要比一只麻雀所包含的更多，因其体积更大、占据更多空间；这样，祢便将祢自身的各部分，以碎块的方式，临在于世界的各个部分，大的碎块给大的部分，小的碎块给小的部分。然而祢并非如此；况且那时祢尚未照亮我的黑暗。\par

% structure: p0006 source=h2 target=subsection reason=relative-heading-rank; base=h2

\subsection{第二章 \Person{尼勃里底乌斯}与\OriginalTerm{摩尼教徒}{Manichæans}的辩论：论\Quote{天主是可朽的还是不可朽的}这一问题}

3. 主啊，对于那些受骗的骗子与哑口的饶舌者（既然你的言语不从他们口中发出，便为哑口），很久以前我们在\Place{迦太基}时，\Person{尼勃里底乌斯}常提出的一个论点，便足以为我驳斥他们了，我们所有听到的人当时都为之震惊：\Quote{那个所谓的黑暗之邦——\OriginalTerm{摩尼教徒}{Manichæans}惯于将其立为与祢对峙的一团——如果祢决意与之争战，它会对祢造成什么伤害呢？倘若回答说：\InnerQuote{它会给祢造成些许伤害}，那么祢就会遭受暴力与朽坏；但如果回答是：\InnerQuote{它不能给祢造成任何伤害}，那么祢与它争战便毫无理由可言；然而战争的结局却是：祢的某个部分、肢体，或祢本质的子嗣，与并非祢所创造的敌对力量和本性混合在一起，被它们败坏、腐化，以至于从幸福转为不幸，需要援助才能得到解救和净化；而这个祢本质的子嗣便是灵魂，它在受奴役、受污染、被朽坏时，祢那自由、纯洁、完整的言语便会来救援；可是，这言语本身也是\OriginalTerm{可朽的}{corruptible}，因为它源自同一本质。因此，他们若肯定祢——不论祢是什么，即祢的本质——是\OriginalTerm{不可朽的}{incorruptible}，那么所有这些说法便都是虚假而可憎的；但若肯定祢是可朽的，那么从一开始这话便是虚假的，且该受憎恶。}所以，这个论点足以反驳那些理应被从饱胀的胃中呕吐出去的人，因为他们无法逃脱，除非犯下心灵与口舌的可怕亵渎，竟这样思想和谈论祢。\par

% structure: p0008 source=h2 target=subsection reason=relative-heading-rank; base=h2

\subsection{第三章 恶的原因是意志的自由抉择}

4. 然而我，虽然曾说并坚信，祢，我们的主，真实的\Person{天主}，不仅创造了我们的灵魂和肉身，而且不止于我们的灵魂和肉身，还创造了所有受造物和万有，是不可玷污、不可转变、丝毫不可改变的；但我仍未轻易清晰地明了恶的原因是什么。然而无论它是什么，我发觉必须寻求它，却不能由此强迫我相信不可改变的天主是可改变的，免得我自己成为我正寻求的对象。因此，我无忧无虑地寻求它，确信这些人所说的虚妄不实，我全心远避他们；因为我发觉，他们在寻求恶的起源时，充满了恶意，因为他们更愿意认为祢的\OriginalTerm{本体}{substance}遭受恶，而不愿承认是他们自己犯下了恶。\par

5. 我集中注意，要辨明我当时听到的：自由意志是我们行恶的原因，而祢公义的审判是我们受苦的原因。但我无法清楚辨明。于是，我试图将心灵的眼目从那个深坑中拉出来，却又被投入其中；多次尝试，就多次被抛回。但有一件事使我朝向祢的光上升：我既知道自己有意志，也知道自己有生命；因此，当我愿意或不愿意做某事时，我最肯定的是，除我自己之外，没有人是愿意或不愿意的；我立即意识到，这就是我犯罪的原因。至于我身不由己所做之事，我看作是受苦而非主动为之，我认为这不是我的过错，而是我的惩罚；由此，信赖祢是最公义的，我很快承认自己受罚并非不公。但我又说：\Quote{谁创造了我？难道不是我的天主，祂不仅美善，而且是美善本身吗？那么，我何以愿意行恶，不愿意行善，以致有了我应受公义惩罚的理由呢？是谁把这放在我里面，在我内种下了\OriginalTerm{苦根}{root of bitterness}，既然我完全由我最甘饴的天主所造？如果魔鬼是创始者，那么魔鬼从何而来？如果魔鬼也是因着自己的邪恶意志，从善天使变成了魔鬼，那么他得以变成魔鬼的邪恶意志又从何而来，既然那天使是完全由至善的造物主所造的善体？}这些思绪又使我沮丧窒息；但却没有坠入那错误的地狱（在那里无人向祢认罪），去认为祢容许恶，而不是人自己行恶。\par

% structure: p0011 source=h2 target=subsection reason=relative-heading-rank; base=h2

\subsection{第四章 天主是不可朽坏的，若祂可朽坏，就根本不是天主。}

6. 我如此努力寻求其余问题，正如已经发现不可朽坏之物必然优于可朽坏之物；因此，祢，无论祢是什么，我都承认是不可朽坏的。因为从未有过，也绝不会有灵魂能设想出比祢更优越的事物，祢是至高至善的善。既然最真实确切的是，不可朽坏者确实优于可朽坏者（正如我现在宁愿选择一样），那么，如果祢不是不可朽坏的，我就能在思想中达到比我的天主更优越的某物了。因此，我既看到不可朽坏者优于可朽坏者，我就该在那里寻求祢，并在那里观察\Quote{恶本身从何而来}，即那绝不能玷污祢的\OriginalTerm{本体}{substance}的朽坏从何而来。因为朽坏确实不会以任何方式损害我们的天主——无论出于任何意志，任何必然，任何意外——因为祂是天主，祂所意愿的是善，祂自己就是那善；而被朽坏则不是善。祢也不被迫做任何违背祢意志的事，因为祢的意志并不大于祢的能力。如果祢的意志更大，那祢自身就要大于自身了；因为天主的意志和能力就是天主自己。有什么是祢所不能预见的呢？祢通晓一切。没有任何一种本性是祢不知晓的。我们还需要多说什么，\Quote{为什么天主所是的那\OriginalTerm{本体}{substance}不可朽坏}，因为若非如此，它就不能是天主？\par

% structure: p0013 source=h2 target=subsection reason=relative-heading-rank; base=h2

\subsection{第五章 关于恶的起源与天主相关的问题，祂既是至善，就不能是恶的原因。}

7. 我曾探问\Quote{恶从何来？}，却以一种邪恶的方式探问；我在探问中未看出自己的恶。我把我心神的视线前的整个受造界，以及我们能在其中分辨的一切，如地、海、空气、星辰、树木、生物，都整理排列；是的，甚至包括其中我们看不见的，如穹苍、所有的天使以及天上的一切居民。但是这些存在物，就像是有形体似的，我的想像将它们安排在这样那样的地方；我把祢的一切受造物造成一个巨大团块，按形体种类分别开来——有的是真正的形体，有的则是我自己为精神体假想出来的形体。我使这团块成为巨大的——并非它本来如此，这我无法知道，而是以我为合宜的那么大，但无论如何仍是有限的。然而祢，主啊，我设想祢在各方面环绕渗透它，尽管祢是无限无量的；好像处处有海，而且在无边无际的四面八方，除了一片无限的海洋，别无他物；这海洋自身内含有一块巨大却有限的海绵，这海绵的每一部分都被无垠的海洋所充满。我就这样设想着祢的受造界是有限的，并被祢——\OriginalTerm{无限者}{the Infinite}——所充满。我说：看，天主，以及天主所创造的；天主是善的，是的，至大、无可比拟地超越这一切；然而祂——善者——创造了它们，使之成为善，看，祂如何环抱并充满它们。那么，恶在哪里，从何而来，又是如何潜入此处的呢？它的根是什么，种子又是什么？抑或它根本不存在？那么，我们为什么惧怕并逃避那不存在的东西呢？若我们是无谓地惧怕，那这惧怕本身就是恶，是心灵无端地受刺伤和折磨——而且更是一个大恶，因为我们没有可惧怕的，却仍然惧怕。所以，要么我们惧怕的对象是恶，要么惧怕行为的本身是恶。既然善的天主创造了这一切美善的事物，那么恶从何而来？祂——那\OriginalTerm{至大至高的善}{the greatest and chiefest Good}——创造了这些较小的善；但造物主和受造物，全都是善的。恶从何而来？或许有一种\OriginalTerm{恶的物质}{evil matter}，祂用它制造、塑造、安排了世界，却在其中留下了一些祂没有转化为善的东西？但为什么这样呢？祂是全能者，难道无力改变整个泥团，使其中不留丝毫的恶吗？最后，祂为什么一定要用这恶的物质制造什么，而不以同样的全能使其根本不存在呢？或者这恶的物质能违背祂的意愿而存在吗？又或者，如果它从永恒就存在，为什么祂在过去无尽的时间里容许它如此，而又在此后这样久才愿意用它来制造什么？再或者，如果祂突然之间愿意做些什么，那么全能者倒不如实现这样的事：使这恶的物质根本不存在，而只有祂自己是\OriginalTerm{全体、真实、至高、无限的善}{the whole, true, chief, and infinite Good}。再或者，如果善者不去做善的制造者与创造者，这便不是善，那么那恶的物质被除灭、归于虚无之后，祂可以形成善的物质，以此创造万物。因为祂若不能无需自己未造的物质之助而创造善的事物，祂便不是全能的了。这样的事情我反复在我悲惨的胸怀中思索，被剧烈的忧虑所压倒，深恐在我发现真理之前死去；然而，对祢的基督——我们的主和救主——的信德，一如在天主教会内所保持的，已牢牢固定在我心中，固然在许多要点上尚未成形，且偏离教义准则，但是我的心神并未彻底离弃它，反而日日愈加汲取它。\par

% structure: p0015 source=h2 target=subsection reason=relative-heading-rank; base=h2

\subsection{第六章：他驳斥占星家由星座推出的占卜}

8. 此时，我也弃绝了那些占星家的虚妄占卜和不虔的荒谬。我的天主啊，愿你的仁慈由我心灵深处为此也称谢你。因为是你，全是你——除了那不知死亡的生命，和那无需光亮而能开启需要光照之心灵的\OriginalTerm{上智}{Wisdom}，宇宙藉着它被掌管，甚至达于飘动的树叶——又有谁能把我们从一切错谬的死亡中召回呢？你还对付了我的顽固，我曾与机敏的老人\Person{Vindicianus}和才华横溢的青年\Person{Nebridius}为此争执；前者激烈地断言，而后者虽带几分怀疑，却常说 \Quote{并不存在预见未来的方术，人们的猜测常靠运气相助，他们预言许多事时，由于多说而偶然言中。} 因此，你为我提供了一位朋友，他并非草率地求问占星家，但也并不精通这些技艺，而是如我所说，怀着好奇心去咨询他们；他知道一些事，据他说从父亲那里听来，至于这会在多大程度上颠覆对方术的评价，他却不知道。此人名叫\Person{Firminius}，受过良好教育，精通修辞学，把我当作挚友，就他的一些事务——他的世俗希望因所谓的星象而高涨——来请教我的看法；而我在这方面已开始倾向\Person{Nebridius}的意见，确实没有拒绝推敲此事，并告诉他我那犹豫不定之心中所浮现的想法，但仍旧加上一句，说我几乎已确信，这些不过是空洞可笑的愚妄。听罢，他告诉我，他的父亲曾热衷于这类书籍，且有一位朋友，对此兴趣不亚于他；他们一同研究、商讨，煽起了他们对此把戏的热爱之情，以至于连家中饲养的哑巴畜生何时产仔，他们也要观察时辰，然后观测当时天象的位置，以收集这门所谓技艺的新证据。他更说，他父亲曾告诉他，当他的母亲即将生下他（\Person{Firminius}）时，他父亲那位朋友的一名女仆也怀了孕，这事瞒不了她的主人，那位主人连自己狗的产期都以极其精密的准确性来记录。于是，经过最精心的观察，为妻子和女仆计算日子、时辰以及更细的时刻，两人竟在同一时刻分娩，以至于他们不得不为儿子和年轻奴隶记载下完全相同、丝毫不差的星象。因为妇女们刚一开始阵痛，他们便互相通报各人家中发生的情形，并准备好报信的人，一得到确实分娩的消息便立即派往对方，而且各自轻易地安排好了即时通报的方式。他说，就这样，双方的报信者在两家的等距离处相遇，以致于无论星宿的位置还是其他最细微之处，他们都无法辨认任何差异。然而，\Person{Firminius}出生在父母家中的高贵地位，走过世上的亨通之路，财富增加，荣升高位；而那奴隶——他命运的轭未曾放松——则继续服事主人，正如认识他的\Person{Firminius}告诉我的。\par

9. 听到并相信这位如此可靠之人所讲述的事，我一切的抗拒都瓦解了；首先，我试图使\Person{Firminius}本人脱离那种好奇心，告诉他说，我在查看他的星象时，若要说真话，就应从中看出，父母在邻人中显赫，在本城是贵族高门，出身高贵，教养良好，学问博雅。可是，若那奴隶以同样的星象来咨询我——因为这些星象也是他的——我同样要说真话，就应告诉他在其中看出出身卑微、境遇卑下，以及所有与前者完全相反、毫无共同之处的一切。那么，观看同样的星象，我若说真话，就会说出迥异的话；若说相同的话，就是说假话。由此确凿地推断出：凡据星象所说而应验的，不是靠方术，而是碰巧；凡不验的，不是因方术不精，而是碰巧出错。\par

10. 既开了这个窍，我心中寻思这类事，使那些迷恋此事（我渴望抨击他们并加以嘲笑驳斥）的妄人当中，没有一个能反驳我说\Person{Firminius}，或他的父亲，提供了假消息。我把思绪转向了双胞胎，他们通常先后出离母胎，时间相隔甚近，以致无论他们如何争辩这在自然中有多大的力量，人与人之间的观察无法察觉，也无法用占星家为宣说真理而检查的那些图形表达出来。其实它们也不可能是真实的。因为如果查看同样的图形，他就必定对\Person{厄撒乌}和\Person{雅各伯}预言同样的事情，然而同样的事并未发生在他们身上。所以，他必定说假话；若说真话，则查看同样的图形时，他就不能说同样的话。因此真话不是靠方术，而是碰巧说的。主啊，宇宙最公义的主宰，不问者和被问者都不知晓，凭你隐密的启示，使他们按照灵魂隐秘的功过，听你公义审判的深渊中所当听的话。人不要向你说：\Quote{这是什么？}或\Quote{为什么那样？} 他不要这样说，因为他只是人。\par

% structure: p0019 source=h2 target=subsection reason=relative-heading-rank; base=h2

\subsection{第七章 为恶之根源备受熬炼}

11. 如今，我的助佑者啊，祢已解除了我的那些束缚；我追问：\Quote{恶从何来？} 却毫无结果。但祢不容我因思想上的任何摇摆而远离信仰，这信仰使我相信祢存在，祢的本体永不变易，祢眷顾并审判世人；且在祢的圣子、我们的主基督内，以及在祢天主教会权威敦促我接受的圣经中，祢已为人类筹划了救恩之路，通向死后那未来的生命。这些事既在我心中稳固而不可动摇地确立了，我便急切地追问：\Quote{恶从何来？} 那时，我痛苦的心忍受了何等的折磨！我的天主啊，我发出了多少叹息！然而就在那里，祢的耳朵是敞开的，我却不知道；当我安静地竭力寻求时，我灵魂那些无声的痛悔，向祢的慈悲发出了高亢的呼声。没有人知道，唯有祢知道我所忍受的。因为那时从我口中倾注到我最亲密朋友们耳中的是什么？我灵魂的全部喧嚣，是时间和言语都无法容纳的，难道它达到了他们那里吗？然而，它全部进入了祢的耳朵，我从心底发出的哀号，祢都听见了；我的渴望呈现在祢面前，而我眼中的光明却不在我内；因为那光在内里，我却在外边。那光并不居于空间，我却专注于空间内的事物；在那里我找不到安息之所，它们也不以让我能说\Quote{这足够了，这就好了}的方式来接纳我；它们也不容我返回，回到那可能让我满足的地方。因为对于这些事物，我居于其上，但与祢相比，我则在其下；当我服从祢时，祢是我真正的喜乐；祢也将祢所造低于我的一切，使之服属于我。这便是真正的中和与中道，是我得享安全的境域：即保持祢的肖像，并通过侍奉祢而统治肉身。但当我骄傲地起来反抗祢，并以我盾牌厚厚的凸起，\BibleQuote{挺著颈项，持著坚盾攻击全能者}\BibleRef{约 15:26}时，甚至那些低于我的事物也置于我上，压迫着我，毫无喘息与轻松之处。它们成群结队地在四面八方与我的目光相遇，当我正要归向祢时，物质的形象在思想中闯入，彷佛要对我说：\Quote{你这卑劣不堪的家伙，往哪里去？} 这些事都从我的创伤中涌出；因为祢使骄傲的人卑微，如同受伤的人，我因自己的肿胀而与祢隔离；是的，我那胀得太大的脸遮蔽了我的眼睛。\par

% structure: p0021 source=h2 target=subsection reason=relative-heading-rank; base=h2

\subsection{第八章 赖天主助佑，他逐渐达到真理}

12. \Quote{但是，主啊，祢永远常存，} 但祢不会永远向我们发怒，因为祢怜悯我们的尘土灰烬；祢乐意改造我的畸形，用内心的刺痛搅扰我，使我不安，直到祢在我心灵的目光中成为确实。祢藉着祢良药奥秘的手，减轻了我的肿胀，以有益的悲伤为尖刻的眼药，使我心灵混乱而昏昧的眼目日益痊愈。\par

% structure: p0023 source=h2 target=subsection reason=relative-heading-rank; base=h2

\subsection{第九章 他将柏拉图学派关于\OriginalTerm{圣言}{\GreekText{Λόγος}}的学说与卓越得多的基督宗教教义相比较}

13. 祢愿首先向我显示，祢如何\BibleQuote{拒绝骄傲人，却赏赐恩宠于谦逊人}，并以何等伟大的慈悲之举，为人类指出谦卑的道路，即祢的\BibleQuote{圣言成了血肉}，寄居在我们中间——祢便通过一个因极度骄傲而膨胀的人，为我获得了一些柏拉图学派的著作，由希腊文译成了拉丁文。在其中我读到，虽非同样的字句，而意义相同，且有许多不同论证支持，说：\BibleQuote{在起初已有圣言，圣言与天主同在，圣言就是天主。这圣言在起初就与天主同在。万物是藉着他而造成的；凡受造的，没有一样不是由他而造成的。在他内有生命，这生命是人的光。光在黑暗中照耀，黑暗决不能胜过他。}\BibleRef{若 1:1} 并读到人的灵魂，虽为光作证，但本身不是那光；而是天主圣言，就是天主，才是那普照每人的真光，正在进入这世界。\BibleRef{若 1:9} 又读到：\BibleQuote{他已在世界上，世界是由他造成的，世界却不认识他。他来到了自己的领域，自己的人却没有接受他。但是，凡接受他的，他给他们，即给那些信他的名字的人权能，好成为天主的子女。} 这些我并未在那里读到。\par

14. 同样，我在那里读到\OriginalTerm{天主圣言}{God the Word}不是由血肉，也不是由情欲，也不是由人欲，而是由天主所生。但是\BibleQuote{圣言成了血肉，寄居在我们中间}我却没有在那里读到。因为我在那些书中发现，它以多种方式说了，子具有父的形像，且\BibleQuote{不以自己与天主同等为僭越}，因为祂本性就是同一\OriginalTerm{实体}{substance}。但是祂却空虚自己，\BibleQuote{取了奴仆的形体，成了与人相似的形状；既有了人的样子，就谦抑自下，听命至死，且死在十字架上。为此，天主极其举扬祂}从死者中，\BibleQuote{赐给祂一个名字，超越所有的名字，好使上天下地和地下的一切，一听到耶稣的名字，无不屈膝；一切唇舌无不宣认耶稣基督是主，以光荣天主圣父；}\BibleRef{斐 2:6} 那些书上却没有。因为在你万古以前、万世之上，你的独生子永远不变地与你同为永恒；并且人从\BibleQuote{祂的满盈}中领受恩宠，\BibleRef{若 1:16} 为能蒙受祝福；并因分有常存的智慧而更新，为能成为智者，这些是有的。但是\BibleQuote{基督在指定的时期为不虔敬者死了}，\BibleRef{罗 5:6} 并且你没有怜惜你的独生子，反而为我们众人将祂交出，\BibleRef{罗 8:32} 这些是没有的。\BibleQuote{因为你将这些事向智慧和明达的人隐藏起来，而启示给了小孩子}；\BibleRef{玛 11:25} 为使那些\BibleQuote{劳苦和负重担的}能\BibleQuote{到}祂那里去，祂要使他们安息，因为祂是\BibleQuote{心谦良善的}。\BibleQuote{良善的人，祂必按公道引导；良善的人，祂必指示自己的道路}；垂顾我们的卑微与苦难，宽恕我们的一切罪恶。但那些因自封的高深学问而骄傲自大的人，却听不到祂说：\BibleQuote{跟我学吧！因为我是心谦良善的：这样你们必要找得你们灵魂的安息。}\BibleRef{玛 11:29} \BibleQuote{因为他们虽然认识了天主，却不以他为天主而光荣他或感谢他；他们的思想反而成为虚幻，他们愚昧的心就昏暗了。他们自负为智者，反成了愚拙。}\BibleRef{罗 1:21}\par

15. 因此，我还在那里读到，他们曾将你那不朽本性的光荣变成偶像和各种形像——\BibleQuote{变成可朽坏的人、飞禽、走兽和爬虫的形状}的像，也就是厄撒乌为了一碗羹汤而失去长子名分的那种埃及食物；\BibleRef{创 25:33} 因为你的首生子民，心里回向埃及，不崇拜你，却崇拜四足走兽的头，并在\BibleQuote{吃草的牛}的像前，叩拜你的肖像——即他们自己的灵魂。我在那里发现了这些事，但我并没有以它们为食粮。因为主啊，你乐意除去雅各伯的次子耻辱，使长子要服事幼子；\BibleRef{罗 9:12} 你并且召叫了\OriginalTerm{外邦人}{Gentiles}承受你的嗣业。而我正是从外邦人中来到你面前，我竭力追求你愿意你的子民从埃及带走的黄金，因为无论它在哪里，都是属于你的。你曾借你的\OriginalTerm{宗徒}{apostle}向雅典人说，在你内\BibleQuote{我们生活、行动、存在}，正如他们自己的一位诗人也说过。\BibleRef{宗 17:28} 的确，这些书卷也是从那里来的。但我并没有留意埃及的偶像，他们用你的黄金事奉它们，\BibleRef{欧 2:8} \BibleQuote{他们将天主的真实变为虚妄，去崇拜事奉受造物，以代替造物主。}\BibleRef{罗 1:25}\par

% structure: p0027 source=h2 target=subsection reason=relative-heading-rank; base=h2

\subsection{第十章 凡退入心灵深处者，天主之事更为显明}

16. 由此受到警示，要我回归自己，我便在你的引导下，进入了我的心灵深处；我所以能如此，是因为你作了我的助佑。我进入后，便以我灵魂的眼目（虽然黯淡）看见，在我灵魂的眼目之上，在我心智之上，有那\OriginalTerm{不变之光}{Unchangeable Light}。这光不是凡人都能看见的普通之光，也不是同一类型而更为强大的光，仿佛这光的光辉更为灿烂，且以它的巨大充满一切。这光并非如此，而是完全不同，的确，与所有这些全然相异。它也不是像油在水上，或天在地之上那样高于我的心智；它之所以在上，是因为它创造了我，我之所以在下，是因为我由它而受造。认识真理者，认识这光；认识这光者，认识永恒。爱认识它。啊，永恒的真理，真实的爱情，可爱的永恒！你是我的天主；我日夜向你叹息。当我初次认识你时，你将我提升，使我看见我所能看见的，并且看见那不是我凭自己所能看见的。你击退了我视力的软弱，将你强烈的光芒照射在我身上，我因爱与畏惧而战栗；我发现自己远离你，处于\OriginalTerm{不相像的区域}{the region of dissimilarity}，仿佛我从高处听到你的声音说：\Quote{我是强者的食粮；你长大，就能以我为食粮；你并不会像转化你肉体的食粮那样将我转化成你，而是你将转化成为我。}我便知道你因人的罪过而纠正人，并使我的灵魂如蜘蛛网般消耗。我说：\Quote{真理，难道因为既不弥散于空间，也不为有限或无限，就成了虚无吗？}而你从远处向我呼喊说：\Quote{的确，\InnerQuote{我是自有者}。}我在心中聆听了，如同在内心听到的事一样，毫无怀疑余地；我宁愿怀疑我是否生活，也不怀疑真理的不存在，因为真理是\BibleQuote{借着受造之物，明显可见，并可理解的}。\BibleRef{罗 1:20}\par

% structure: p0029 source=h2 target=subsection reason=relative-heading-rank; base=h2

\subsection{第十一章 受造物皆可变，唯天主不变}

17. 我观察了在祢之下的其他事物，觉察到它们既不完全存在，也不完全不存在。它们确实存在，因为它们来自祢；但它们又不存在，因为它们不是祢之所是。因为那真正存在的是不变地恒存着的。所以，对我而言，紧紧依附天主是善的，因为如果我不在祂内恒存，我也将不在我自己内；但祂，恒存于自己内，更新万物。\BibleRef{智 7:27} 祢是主，我的天主，因为祢不需要我的善。\par

% structure: p0031 source=h2 target=subsection reason=relative-heading-rank; base=h2

\subsection{第12章 善的天主所创造的一切都好得很}

18. 我又清楚认识到，那些会败坏的事物仍是善的，因为除非它们是善的，否则就不能败坏；它们既非至善，若非善便不能败坏。若为\OriginalTerm{至善}{supremely good}，便不会败坏；若毫无善，便无可败坏。因为败坏造成伤害，除非它能减少善，否则不能造成伤害。那么，要么败坏不造成伤害——这是不可能的——要么最确定的是，所有被败坏的，都是被剥夺了善。但若它们被剥夺了一切善，它们就不再存在。因为如果它们存在，且完全不能败坏，它们会变得更好，因为它们会\OriginalTerm{不朽}{incorruptible}地恒存。还有什么比宣称那些丧失了一切善的事物变得更好更荒谬的呢？因此，若它们被剥夺一切善，便不复存在。所以，只要它们存在，便是善的；因此，凡是存在的，都是善的。那么，我所探求其来源的\OriginalTerm{恶}{evil}，并非任何\OriginalTerm{实体}{substance}；若它是实体，便是善的。因为要么它是\OriginalTerm{不朽}{incorruptible}的实体，即\OriginalTerm{至善}{chief good}；要么是\OriginalTerm{可朽}{corruptible}的实体，若非善便不能败坏。因此，我认识到，并清楚地知道，祢创造了万物，都是善的，没有任何实体不是祢所造；又因为祢所造的并非都同等，所以万物都存在；因为各自都是善的，总体而言便是好得很，因为我们的天主创造的一切都是好得很。\par

% structure: p0033 source=h2 target=subsection reason=relative-heading-rank; base=h2

\subsection{第13章 应当赞美造物主，因祂创造了天地间的美善事物}

19. 对祢而言，绝无丝毫之恶；不仅对祢，对祢的整个受造界也是如此；因为没有任何外在之物能闯入并破坏祢所安排的秩序。但在其各部分中，某些事物因与其他部分不和谐而被视为恶；然而这些事物却与其他部分和谐，且是善的，在其自身亦是善的。所有这些彼此不和谐的事物，却与我们称之为地的低下部分相和谐，这地有其相应的阴云密布、狂风大作的天。因此，我绝不说：\Quote{这些事物不该存在。}因为即便我只见到这些，我确实会渴望更好的事物；但即便如此，单为这些，我也应当赞美祢；因为应当赞美祢，这话显明于\BibleQuote{大地、蛟龙和一切深渊，火焰、冰雹、雪和云雾，遵行祢命令的狂风，山岳和一切丘陵，果树和一切香柏，野兽和一切家畜，爬虫和飞鸟，世上的君王和万民，首领和世上一切审判官，少年人和童贞女，老人与儿童，赞美祢的名。}但当这些受造物，\BibleQuote{从天上赞美祢，在高处赞美祢！众天使赞美祢，众军旅赞美祢！太阳和月亮赞美祢，一切星辰和光明赞美祢！天上的天和天上的水赞美祢的名！}那时，我不再渴望更好的事物，因为我思考着万有；并以更佳的判断反思到，在上的事物比在下的更美善，但万有整体却比仅有在上的事物更美善。\par

% structure: p0035 source=h2 target=subsection reason=relative-heading-rank; base=h2

\subsection{第14章 对天主的部分受造物感到不悦，他便构想了两个原初实体}

20. 那些对祢的任何受造物感到不悦的人，毫无完整可言，其程度不亚于当年的我，那时祢所造的许多事物都令我不悦。因为我的灵魂不敢对我的天主感到不悦，它便不愿承认任何令它不悦的事物是属祢的。于是它陷入两种实体的见解，不抗拒，反而胡言乱语。又从那里回来，它为自己造了一位\OriginalTerm{神}{god}，遍布无尽空间；想象这就是祢，将其置于心中，再次成为自己偶像的殿宇，这在祢眼中是可憎的。但后来，祢抚慰了我无知觉的头，合上我的眼睛，免得它们\BibleQuote{看虚假}，我暂时止息自己，我的疯狂得以安眠；我在祢内醒来，看见祢是无限的，却是另一种方式；这景象并非来自肉体。\par

% structure: p0037 source=h2 target=subsection reason=relative-heading-rank; base=h2

\subsection{第15章 万有的存在都归功于天主}

21. 我回顾其他事物，觉察到它们的存在都归功于祢，它们都在祢内被限定；却是另一种方式，并非如同存在于空间内，而是因为祢以真理之手掌握万物：一切事物，只要存在，便是真实的；除非那不存在者被认为存在，否则便无虚假。我又看见万物彼此和谐，不仅与处所，也与时节相和谐。祢，唯有祢是永恒的，并非在无数的时间之后才开始工作；因为一切时间的间距，无论已过或将过，若非通过祢的运作与恒存，都不会来去。\par

% structure: p0039 source=h2 target=subsection reason=relative-heading-rank; base=h2

\subsection{第16章 恶并非源于实体，而是源于意志的扭曲}

22. 我辨识并发现，这不奇怪：不合病人口味的面包，对健康者却是美味；刺伤痛眼的光芒，对健全的眼睛却是愉悦。祢的\OriginalTerm{正义}{righteousness}使恶人不悦；至于毒蛇和小虫，祢造它们是好的，与祢受造物的低下部分相合；恶人自己也与之相合，他们越不肖似祢，就越与之相合；而与上等的受造物相合，则与肖似祢的程度成比例。我探问\OriginalTerm{不义}{iniquity}是什么，发现它并非实体，而是\OriginalTerm{意志}{will}的扭曲，从那\OriginalTerm{至高的实体}{Supreme Substance}——天主，祢，转向这些低下的事物，抛弃自己的内心，向外膨胀。\par

% structure: p0041 source=h2 target=subsection reason=relative-heading-rank; base=h2

\subsection{超越他可变的理智，他发现真理的不变作者。}

23. 我惊讶于我现在竟然爱慕祢，而非一个代替祢的幻象。然而我仍不配享受我的天主，而是被祢的美所吸引奔向祢，但很快又被我自身的重量从祢那里拉开，带着忧伤沉沦于这些低等事物之中。这重量就是肉欲的习惯。然而对祢的记忆仍存留在我内；我绝不怀疑有一位我可以依附的，只是我尚未成为能依附于祢的人；因为这可朽坏的肉身重压着灵魂，这尘世的寓所使多虑的心智下沉。\BibleRef{智 9:15} 我极其确信：祢那自创世以来不可见的美善，\BibleQuote{都可凭祢所造的万物得以辨识，连祢永恒的大能和天主性也都可以看见。}\BibleRef{罗 1:20} 因为，当我探问我是从哪里得以赞赏物体的美，无论是天上的还是地上的，又是凭什么使我能正确判断可变的事物，并说出\Quote{这该当如此，那不该如此，}——当我探问我如此判断的缘由，既然我确实如此判断，我便发现了真理的不变而真实的永恒，超越于我易变的心智之上。于是，我逐步地从物体提升到那运用身体感官来感知的灵魂；又从那里到那内在的官能，身体感官将外物呈现给它，而这也正是禽兽所能达到的；然后，我又从那里进到那推理的官能，身体感官所接受的一切都交给它来审断，这官能也发觉它自身在我内是变动不居的，便提升自己到达自己的理智，并拖带我的思想脱离习惯，从众多矛盾的幻象中抽身，好能发现那曾洒落其上的光，当它毫无怀疑地呼喊出\Quote{不变者当优先于可变者；}它由此也认识了那不变者，除非它以某种方式早已认识，否则它绝无确切的理由去优先选择不变者。就这样，随着一道战栗的目光一闪，它抵达了那自有者。于是，我看见祢那不可见的美善，从祢所造的万物中得以辨识。\BibleRef{罗 1:20} 但我未能将目光固定于此；我的软弱被击退，我又被抛回我惯常的习惯中，随身携带的，除了对此的爱慕记忆，以及一种渴望，如同嗅到了香气，却还未能品尝。\par

% structure: p0043 source=h2 target=subsection reason=relative-heading-rank; base=h2

\subsection{耶稣基督，这位中保，是唯一的救援之路。}

24. 我寻求获得足够力量以享受祢的方法；但直到我拥抱了那位\BibleQuote{在天主与人之间的中保，就是降生成人的基督耶稣，}\BibleRef{弟前 2:5}\BibleQuote{祂是在万有之上，永远可赞颂的天主，}\BibleRef{罗 9:5}祂呼唤我，并说\BibleQuote{我是道路、真理、生命，}\BibleRef{若 14:6}且将我那无法领受的食粮与我们的血肉相混合——因为\BibleQuote{圣言成了血肉，}\BibleRef{若 1:14}好使祢用以创造万物的智慧，能为我们的孩提时代提供奶汁。因我未领悟我的主耶稣——我虽谦卑，却未领悟那谦卑的；我也不知祂的软弱将教导我们什么功课。因为祢的圣言、永恒的真理，高居祢创造的高层之上，提升那些顺从祂的人；但在这低层世界里，祂为自己用我们的泥土建造了卑微的居所，借此祂意图使那些愿顺从的人自我贬抑，并带领他们归向祂，平息他们的骄傲，培养他们的爱；好叫他们不再依赖自己，反而变得软弱，看见在他们脚前，\OriginalTerm{天主性}{Divinity}因穿上了我们的\Quote{皮衣}而变得软弱；并且疲倦了，他们便投靠在祂身上，而祂在复活时，便会将他们提升起来。\par

% structure: p0045 source=h2 target=subsection reason=relative-heading-rank; base=h2

\subsection{他尚未完全理解若望的话：\BibleQuote{圣言成了血肉。}}

25. 但我那时不这么看，只把我的主基督视为一位卓绝智慧的人，无人能及；特别是因为他由童贞女奇妙诞生，似乎在我们身上彰显了天主的眷顾，从而获得了如此伟大的领导权威——作为轻看暂世事物以获得永生的榜样。但\BibleQuote{圣言成了血肉}这句经文蕴藏的奥迹，我却无法想象。我只是从有关他的文字记载中得知，他吃、喝、睡、行走、心神喜乐、忧伤、与人交谈；不只是肉躯与祢的圣言结合，而是结合了人的灵魂与肉身。所有认识祢圣言永恒不变的人，都知道这点；这我现在已尽我所能地认识，并且毫不怀疑。因为，时而随己意活动身体的肢节，时而不动；时而受某种情感激动，时而不动；时而以记号表达智言慧语，时而保持沉默——这都属于可变化的灵魂与理智的特性。倘若这些关于他的记载是虚假的，其余的一切便有受指摘的危险，那些经卷中便不会为人类留下任何带来救恩的信德。既然这些记载真实不虚，我便承认在基督内有一个完人——不只是人的身体，也不只是身体加上一个无理性的感觉魂，而是一个真正的人；我断定他不只作为真理的典范，更因人性某种卓越和更完全地分沾智慧，而应凌驾众人之上。但\Person{阿利比乌斯}以为天主教信友相信天主这样地取肉躯，以致在基督内除了天主与肉体，没有灵魂，也不认为他有人的理智。因为他深信，记载中有关他的行为，若非由一个有生命和理智的受造物，便不可能完成，所以他向基督信仰的迈进较为缓慢。不过，他后来得知这是阿波林派异端者的错误，便欣然接受天主教信仰，并与之相符。然而，我承认稍微后来我才明白，在\BibleQuote{圣言成了血肉}这句话上，天主教的真理如何能与\Person{佛提努斯}的谬论区分开来。因为异端者的悖逆，反使祢教会的训导和纯正的教义更为彰显。诚然，必须要有异端，好叫那些经得起考验的人在软弱者中显明出来。\BibleRef{格前 11:19}\par

% structure: p0047 source=h2 target=subsection reason=relative-heading-rank; base=h2

\subsection{他庆幸自己是由\Person{柏拉图}走向圣经，而非相反。}

26. 然而，当时我已读了那些柏拉图主义者的书卷，并且受它们提醒去寻求无形的真理，我看见了\BibleQuote{祢的不可见之物，借着所造之物得以理解}；\BibleRef{罗 1:20} 尽管被阻挡，我却感知到那因我心智的幽暗而不得凝望的是甚么——我确信祢是，祢是无限的，却不散布于有限或无限的空间；祢真实地是，祢永远如一，在部分或动态上毫无变异；而其他万物都从祢而来，其存在的唯一确据，就是它们存在。这些事我确实确信，但太过软弱而不能享受祢。我像一个老练的人喋喋不休；但如果我没有在救主基督内寻求祢的道路，我非但不熟练，且将趋向灭亡。因为那时，我满受惩罚，已开始渴望显为聪明；却不哀哭，反而因知识而自高自大。\BibleRef{格前 8:1} 因为那建立在谦卑\Quote{根基}上的爱德，\Quote{这根基就是耶稣基督}，又在何处？\BibleRef{格前 3:11} 这些书卷又何时能教导我这一点？因此，我相信，祢乐意让我在研究祢圣经之前，先碰上这些书，好让这些感受深印我心；然后，当我被祢的经卷驯服，当我的创伤被祢医治的手指触摸时，我就能分辨和区别，自恃与宣认之间的差别——那些看见当往何处去，却看不见道路的人，与那条不仅引人观看、更引人安居于福地之路之间的差别。因为，倘若我先在祢的圣经中被塑造，倘若祢在熟悉运用它们时变得甘甜于我，而后我碰上了那些书卷，它们或许会将我从虔诚的稳固地基上夺走；或者，即使我仍坚定于从那里吸收来的健康气质，我也可能以为单靠研究这些书卷便可达到这种境界。\par

% structure: p0049 source=h2 target=subsection reason=relative-heading-rank; base=h2

\subsection{他在圣经中发现而\Person{柏拉图}著作中却找不到的东西。}

27. 于是，我极其热切地抓住了祢圣神的那部可敬著作，尤其是宗徒\Person{保禄}的书信；先前那些曾使我以为他自相矛盾，其讲论与法律和先知的见证不相符的困难，便都消散了。那纯洁言语的面貌向我显现，是唯一而相同的；我学会了\BibleQuote{欢跃而颤栗}。于是我着手阅读，发现我以前在那里读到的一切真理，在这里都是凭着祢恩宠的推荐而宣告；使那看见的人不可夸耀，好像他不仅所看见的不是领受的，就连他能看见也不是领受的（因为他有什么不是领受的呢？）；并且不仅劝勉他看见祢——祢是永恒不变的——也得以治愈，好能把握住祢；而那从远处不能看见的人，仍可行在那条道路上，藉以抵达、仰瞻并拥有祢。因为虽然一个人\BibleQuote{按内心喜爱天主法律}，\BibleRef{罗 7:22} 但他肢体中另有一条法律，与他内心的法律交战，将他掳去，附属于那在他肢体中的罪恶法律，他能做什么呢？因为祢是公义的，主，但我们犯了罪，行了不义，作了恶事，祢的手沉重地压在我们身上，我们被公义地交于那远古罪人，死亡的主宰；因为他引诱我们的意志，效法他的意志，因而他未存留在祢的真理中。\BibleQuote{可怜的人}能做什么呢？\BibleQuote{谁能救我脱离这死亡的肉身呢？} 除非藉着祢的恩宠，\BibleQuote{靠我们的主耶稣基督}。\BibleRef{罗 7:24} 这基督，祢所生的\OriginalTerm{同永恒者}{co-eternal}，又在祢行径的起始创造了祂；世界的首领在祂身上找不到任何可死的罪行，\BibleRef{若 18:38} 却把祂杀害了，那反对我们的罪状也被涂抹了？\BibleRef{哥 2:14} 那些书页不包含这虔诚的表达——忏悔的眼泪、祢的祭献、忧苦的精神、\BibleQuote{破碎痛悔的心}、人民的得救、许配的城、\BibleRef{默 21:2} 圣神的\OriginalTerm{质押}{earnest}、\BibleRef{格后 5:5} 我们的救赎之杯。在那里没有人咏唱：\BibleQuote{我的灵魂岂不该归顺天主吗？因为我的救恩是从祂而来，祂是我的天主，我的救恩，我的堡垒，我决不再动摇。} 在那里没有人听到祂呼唤：\BibleQuote{你们凡劳苦和负重担的，都到我这里来吧。} 他们不屑于向祂学习，因为祂是良善心谦的；\BibleRef{玛 11:28} 因为 \BibleQuote{祢将这些事瞒住了智慧和明达的人，而启示给了小孩子。}\BibleRef{玛 11:25} 因为从树木茂密的山顶望见和平之地，\BibleRef{申 32:49} 却找不到通往那里的道路——徒然尝试那些不可通行的路径，遭到由叛逆和逃亡者组成的伏兵敌对和阻截，他们的首领是那 \BibleQuote{狮子} \BibleRef{伯前 5:8} 和 \BibleQuote{巨龙}；\BibleRef{默 12:3} 而另一件事则是沿着通往那里的道路前行，那条路由\OriginalTerm{天上大将}{heavenly general}的军队守护，在那里，那些逃离天军的人——他们像躲避酷刑一样躲避天军——不再进行劫掠。当我读到祢的\BibleQuote{最小的宗徒}，默观祢的工程，大为战栗时，这些话语奇妙地深入我的肺腑。\par

\SourceRecord{Translated by J.G. Pilkington.；Nicene and Post-Nicene Fathers, First Series；Vol. 1.；Edited by Philip Schaff.；Buffalo, NY: Christian Literature Publishing Co.,；1887.；<http://www.newadvent.org/fathers/110107.htm>.}

% structure: p0001 source=h1 target=section reason=page-title

\section{\WorkTitle{忏悔录}（卷八）}

最终，他描述了他三十二岁这一年，这是他整个生命中最值得铭记的一年。在这一年里，他受\Person{西姆普利齐亚努斯}指点，得知他人的皈依及行事方式，经过一番严酷的挣扎后，他的整个心灵得到更新，并皈依了天主。\par

% structure: p0003 source=h2 target=subsection reason=relative-heading-rank; base=h2

\subsection{他现在已献身于神圣之事，但仍被爱欲所缠累，就心灵的更新之事请教\Person{西姆普利齐亚努斯}}

哦，我的天主，让我怀着感激之情，记忆并向你宣认你赐予我的仁慈。愿我的骨骸沉浸于你的爱中，并让它们说：上主，谁能相似你？\BibleQuote{你解开了我的羁绊，我要向你献上感恩祭。}我要述说你如何解开了它们；所有敬拜你的人听到这些事，必将说：\BibleQuote{天上地上的上主是应受赞美的，他的名大而可畏。}你的话语已深深扎入我的胸中，我四面被你围护起来。\BibleRef{约 1:10}对于你永恒的生命，我现在已确信无疑，尽管我曾\BibleQuote{对着镜子观看，模糊不清}\BibleRef{格前 13:12}然而，我不再怀疑有一个\OriginalTerm{不朽的实体}{incorruptible substance}，一切别的实体都由它而来；我也不再渴求更确知你，而是渴求在你内更坚定不移。至于我现世的生命，一切都不确定，我的心尚需清除\OriginalTerm{旧酵母}{old leaven}\BibleRef{格前 5:7}那\BibleQuote{道路}\BibleRef{若 14:6}即救主自己，令我喜悦，但我尚不愿穿越它的窄狭。你又把去见\Person{西姆普利齐亚努斯}的念头放进我心里，这念头在我看来是好的；他在我看来是你忠信的仆人，你的恩宠在他身上闪耀。我也听说他从青年时代便极虔诚地生活于你内。如今他年事已高，因在如此高龄中，仍以如此热忱追随你的道路，在我看来他必已积累了丰富的经验；事实也确是如此。我愿他以这些经验，（向他倾诉我的苦楚）告诉我像我这样痛苦的人，最宜如何行走你的道路。\par

因为我见教会内充满了人，有的走这条路，有的走那条路。然而我自己仍度着世俗的生活，令我不悦；是的，如今我的情欲已不再像往日那样以荣誉和财富的希冀来激动我，承受如此沉重的奴役实在令我痛苦不堪。因为，比起你的甘饴，及我所爱的你殿宇的美，那些事物已不再令我喜乐。但我仍极其顽固地被恋女之情所束缚；宗徒并未禁止我结婚，尽管他劝我追求更好的境界，尤为希望众人都如同他一样。\BibleRef{格前 7:7}但我体质软弱，便选择了这较为惬意的生活方式，并仅因此而在其他一切事上动荡不安，因憔悴的忧虑而萎靡不振，因为在其他事上，尽管不情愿，却被迫适应婚姻生活，我已将自己交付并受其束缚。我曾听真理亲口说过：\BibleQuote{有些\OriginalTerm{阉人}{eunuchs}，却是为了天国而自阉的；}\BibleRef{玛 19:12}但他说：\BibleQuote{这话不是人人所能领悟的，只有领悟的人才能领悟。}\BibleRef{玛 19:12}凡不认识天主的人，确实都是虚幻的，他们未能从可见的美物中，寻见那至善者。\BibleRef{智 13:1}但我已不再处于这种虚幻之中；我已超越它，并藉你所有受造物的一致见证，找到了你，我们的创造者，以及你的圣言，与你同在的天主，并与你及圣神同为唯一天主，你借着他创造了万物。另有一种不虔敬的人，\BibleQuote{他们虽然认识了天主，却不以他为天主而予以光荣，也不感谢他}\BibleRef{罗 1:21}。我也曾陷入这种状况；但你用右手扶持我，把我带走，并将我安置于能复原之处。因为你对人说过：\BibleQuote{看，敬畏上主，就是智慧；}\BibleRef{约 28:28}并且不要自以为聪明，\BibleRef{箴 3:7}因为\BibleQuote{他们自负为智者，反而成为愚妄。}\BibleRef{罗 1:22}但我已寻到了那颗宝贵的\OriginalTerm{珍珠}{pearl}，本应卖掉所有的一切去购买它；\BibleRef{玛 13:46}我却犹豫了。\par

% structure: p0006 source=h2 target=subsection reason=relative-heading-rank; base=h2

\subsection{虔诚的老人喜悦于他阅读了\WorkTitle{柏拉图}与\WorkTitle{圣经}，并告诉他修辞学家\Person{维克托里努斯}如何因阅读\WorkTitle{圣经}而皈依了信仰}

于是我去见\Person{西姆普利齐亚努斯}——在领受你恩宠方面，他堪称是\Person{盎博罗削}（当时为主教）的灵性之父，\Person{盎博罗削}也真诚地爱他如父。我向他诉说了我错误的曲折历程。当我向他提及，我读过某些\OriginalTerm{柏拉图学派}{Platonists}的著作，是由\Person{维克托里努斯}——曾为\Place{罗马}\OriginalTerm{修辞学教授}{rhetorician}（我听说他后来以基督徒身份去世）——译成拉丁文的，他便祝贺我，因为我没有落在那些充满谬误与欺诈、\BibleQuote{按世俗的粗浅学说}\BibleRef{哥 2:8} 的其他哲学家的著作中，而这些柏拉图学派的著作却在许多方面引导人相信天主及其圣言。然后，为了劝我修习基督的谦卑，就是那瞒住了智慧和明达的人，而启示给小孩子的\BibleRef{玛 11:25} 的，他谈起了\Person{维克托里努斯}本人；他在\Place{罗马}时曾与他十分熟识。关于他，他讲述的事，我不会保持沉默。因为这包含了对你的恩宠的极大赞美，应当向你宣认：这位极有学问的老人，精通一切\OriginalTerm{自由学科}{liberal sciences}，读过、评论并讲解过许多哲学家的著作；是众多高贵元老的老师；作为他杰出履行其职责的标志，他甚至还（此世之人视为莫大荣誉）在\Place{罗马广场}赢得并树立了一座雕像——他直到这般年纪，仍是个\OriginalTerm{偶像崇拜者}{worshipper of idols}，参与那些几乎全体罗马贵族都热衷的\OriginalTerm{亵圣仪式}{sacrilegious rites}，并激发了人们对\par

犬首的\OriginalTerm{阿努比斯}{Anubis}，与一群怪神，\newline{}与\OriginalTerm{尼普顿}{Neptune}对阵，\newline{}对抗\OriginalTerm{维纳斯}{Venus}和\OriginalTerm{弥涅耳瓦}{Minerva}，铁甲的\OriginalTerm{玛尔斯}{Mars}，\par

罗马曾征服而如今敬拜的那些神灵，老\Person{维克托利努斯}曾以雷霆般的口才为之辩护了这么多年——他现在不以成为祢的基督的孩子为耻，在祢的泉源旁成为婴孩，将颈项置于谦卑的轭下，使额头屈服于十字架的耻辱。\par

4. 主，主啊，祢曾使天低垂而降临，触摸群山而冒烟，祢用何法将自己送入他的心怀中？他常——据\Person{辛普利齐亚努斯}说——勤读圣经，极其渴望地寻求并钻研所有基督徒著作，并对\Person{辛普利齐亚努斯}说——不是公开地，而是私下、如朋友般——\Quote{你要知道，我是一名基督徒。}\Person{辛普利齐亚努斯}答道：\Quote{我不会相信，也不会将你列入基督徒中，除非我在基督的教会中见到你。}他嘲弄地反问：\Quote{那么，是墙壁造就基督徒吗？}他常这样说，说自己已是基督徒；而\Person{辛普利齐亚努斯}总是同样回答，那\Quote{墙壁}的观念便也常由他反复提起。因为他害怕得罪他的朋友们，那些骄傲的崇拜邪魔者，从他们巴比伦般的尊高位份上，如同尚未被主折断的\Place{黎巴嫩}香柏，他以为仇恨的暴风雨会向他倾泻。但后来，借着阅读与探求，他获得了力量，便惧怕若如今在人前羞于承认基督，将来基督要在圣天使面前否认他\BibleRef{路 9:26}，并自觉犯了大罪，竟以祢圣言的谦卑圣事为耻，却不以那些骄傲邪魔的亵渎仪式为耻——他曾效法他们的骄傲，采用他们的仪式——于是他变得对虚荣厚颜，对真理羞惭，忽然出人意料地对\Person{辛普利齐亚努斯}说——这是他亲口告诉我的——\Quote{我们到教会去吧；我愿意成为基督徒。}\Person{辛普利齐亚努斯}喜不自胜，便陪同了他。接受了初步教导的圣事后，他不久便报名登记，为使能藉圣洗重生——\Place{罗马}惊奇，教会欢跃。骄傲的人看见，勃然大怒；他们咬牙切齿，却消融无踪！但主天主是祢仆人的希望，祂不看顾虚幻和谎言的狂妄。\par

5. 最后，到了他宣认信仰的时刻（在\Place{罗马}，凡即将领受祢恩宠的人，惯常于高台上，在信众面前，背诵预先记熟的固定格式的言辞），据他说，司铎们建议\Person{维克托利努斯}更为私下地宣认，如同惯常对待那些可能因羞怯而胆怯的人；但他却宁愿在神圣会众面前宣认他的救恩。因为他在修辞学中所教导的并非救恩，然而却曾公开宣认那（修辞学）。那么，当他宣讲祢的圣言时，他又怎应畏惧祢温良的羊群呢？他宣讲自己言语时，都不曾畏惧那些疯狂的群众！于是，当他登上台去宣认时，众人一认出他，便互相低语他的名字，发出祝贺之声。他们中间谁不认识他呢？欢欣的群众口中传出一阵低语：\Quote{维克托利努斯！维克托利努斯！} 一见到他，突然爆发出欢呼；又突然安静下来，以便听他发言。他以超凡的勇气宣认了真实的信仰，众人渴望将他拥入心底——是的，他们以爱与喜乐将他拥入那里；这便是拥抱他的双手。\par

% structure: p0012 source=h2 target=subsection reason=relative-heading-rank; base=h2

\subsection{第三章 论天主与天使因一个罪人的回头，比因众多义人更欢喜}

6. 美善的天主，在人心中发生了什么，使他对那绝望而得救、脱离更大危险的灵魂的救恩，比对他若常存希望或危险较小之时更为欢喜呢？因为祢，仁慈的父啊，也正是如此：\BibleQuote{对于一个悔改的罪人，比那九十九个无须悔改的义人，更觉欢喜。} 每当我们听到，如何有亡羊被善牧扛在肩上带回，天使为之欢欣，达玛被找回放入祢的宝库，邻人与那妇女一同欢乐时，我们总是满心喜悦地聆听\BibleRef{路 15:4}；当在祢的殿中诵念祢的小儿子\BibleQuote{死而复生，失而复得}时，祢殿中的庄严礼仪之喜乐催人泪下\BibleRef{路 15:32}。因为祢既在我们内也在祢的圣天使内欢跃，他们因神圣的爱德而成为圣的。因为祢是永恒不变的；对于那些既非恒常亦非永存的事物，祢总以同一方式知晓。\par

7. 那么，灵魂中发生了什么，使得它对于寻获或恢复所爱之物，比一直拥有它们时更为喜悦呢？是的，其它事情也为此作证；万物都充满了见证，呼喊道：\Quote{正是如此。} 获胜的统帅举行凯旋；然而，他若不战斗便不能得胜，战斗的危难愈大，凯旋的欢乐也愈大。风暴颠簸航行者，威胁颠覆船只，众人面临死亡，脸色苍白；但天空与大海归于平静，他们便因曾深惧而深喜。一位亲爱者患病，脉搏显示危险；所有渴望他平安的人立时心中忧戚：他康复了，虽尚未能如往昔健步行走，却带来一种他健壮行走时所没有的喜乐。的确，人生之乐——不仅是那些不期而至、违反我们意愿的，也包括那些自愿与安排的——人们都是通过艰难而获得。若无饥渴之苦在先，饮食便毫无乐趣可言。醉汉们吃某些咸肉，旨在制造一种令人不适的热渴，再以酒浇熄而得快感。习俗亦使已订婚的新娘不应立即过门，免得丈夫因在订婚时未生渴望而轻看她。\par

这条规律适用于卑贱可憎的喜乐，也适用于那许可与合法的喜乐；适用于诚实友谊的真诚，也适用于那曾死而复生、失而复得的那一位\BibleRef{路 15:32}。更大的喜乐总是以更大的痛苦为先导。主，我的天主，这究竟意味着什么？祢本身就是永恒的喜乐，而环绕祢的万物也常在祢内欢跃。这万物中的一部分如此消长交替，时而触怒，时而和解，又意味着什么？难道这就是它们的样式？是祢从诸天最高处到大地最低处，从世界的开始到终结，从天使到虫豸，从最初的动作到最后的活动，将一切各归其位，各定其时，而使万物各按其类成为美好时，所分配给它们的命运吗？我真是苦啊！在至高之处祢如此崇高，在至深之处祢又如此深邃！祢从不远离，而我们却几乎难以\Emph{回归}于祢。\par

% structure: p0016 source=h2 target=subsection reason=relative-heading-rank; base=h2

\subsection{第四章 他以\Person{维克托林}为例，说明贵族归主更令人喜乐。}

9. 主，快快行动吧；唤醒我们，召回我们；点燃我们，吸引我们归向祢；激发我们，使我们尝到祢的甘甜；让我们现在爱祢，让我们\BibleQuote{追随祢。}\BibleRef{歌 1:4} 难道不是有许多人，从比\Person{维克托林}更深的地狱般的盲目中，归向祢，靠近祢，并被光照，领受那光，凡领受的人就从祢得到权能，得以成为祢的子女吗？\BibleRef{若 1:12} 但是他们若在民众中较不为人知，那么即使认识他们的人，也为他们欢喜较少。因为当许多人一同欢乐时，每个人的喜乐因彼此的激发和点燃而更为充足。再者，因为那些为众人所知的人，能影响更多的人走向救恩，并带领众人跟随他们。因此，那些先于他们的人也大大为他们欢喜，因为他们不单为他们自己欢喜。切愿在祢的帐幕里，不按富人的身份而偏待穷人，也不按贵族的身份而偏待卑贱者；因为祢反倒\BibleQuote{拣选了世上软弱的，为羞辱那强壮的；世上卑贱的和受人轻视的，祢也拣选了，以及那无有的，为废除那有的。}\BibleRef{格前 1:27} 然而，就连那位\BibleQuote{宗徒中最小的，}\BibleRef{格前 15:9} 祢借他的舌头发出这些话，当方伯\Person{保禄}\BibleRef{宗 13:12}——他的骄傲为宗徒的战斗所克服——负起祢基督那轻松的轭\BibleRef{玛 11:30}，成为伟大君王的臣民之后，他也——除了他原有的名字\Person{扫禄}外——愿称为\Person{保禄}，以见证如此伟大的胜利。因为敌人在他掌握得多的人身上，才被更彻底地战胜，也借着他掌握更多的人。然而他借着贵族身份掌握骄傲者较多；借着权威，借他们掌握的人更多。因此，\Person{维克托林}的心——那曾被魔鬼视为不可攻克的堡垒——和\Person{维克托林}的舌头——他曾用这强大锋利的武器杀死了许多人——是多么更受欢迎；祢的子女们更应大大欢喜，因为我们的君王已捆绑了那壮士，\BibleRef{玛 12:29} 他们看到他的器皿被夺去，被洗净，成为合乎祢荣耀的器皿，准备为天主所用，行各种善事。\BibleRef{弟后 2:21}\par

% structure: p0018 source=h2 target=subsection reason=relative-heading-rank; base=h2

\subsection{第五章 论使我们远离天主的原因。}

10. 但是，当祢的人\Person{辛普利齐亚努斯}向我讲述了\Person{维克托林}的这件事时，我热切渴望效法他；他正是为此目的才讲述的。然而，当他再补充说，在\Person{尤里安}皇帝时代，曾颁布一项法律，禁止基督徒教授文法与修辞学，而\Person{维克托林}为服从这法律，宁愿放弃他那夸夸其谈的学校，也不舍弃祢的圣言——祢以此使哑巴的舌头变得流畅，\BibleRef{智 10:21}——他便在我看来，不仅是勇敢的，更是有福的，因为他由此寻得了一个专事祢的机会，这正是我所叹息渴望的，而我却被捆锁，不是被别人的铁链，而是被我自己的铁制意志所捆锁。敌人控制了我的意志，并由此为我制造了一条锁链，将我束缚。由于邪僻的意志，产生了私欲；放纵私欲，成为习惯；习惯不加抗拒，便成必然。这些环节，仿佛彼此相连（因此我称之为\Quote{锁链}），形成了一种严酷的奴役，紧紧捆住了我。但是，那在我内开始萌芽的新意志——即自由地朝拜祢，并渴望享有祢，噢天主，唯一确实的享受——尚不能战胜我先前因长久放纵而变得强大的顽梗。就这样，我的两个意志，一旧一新，一属肉欲，一属神的，在我内交战，它们的冲突使我的灵魂紊乱。\par

11. 这样，我从亲身经验中明白了所读到的话：\BibleQuote{肉欲相反圣神，圣神相反肉欲。}\BibleRef{迦 5:17} 我确实两边都有欲求，但在我内心所赞同的事上，欲求更多，超过我所不赞同的事。因为对于后者，已不再是\BibleQuote{我}\BibleRef{罗 7:20}在作，因为我在很大程度上是逆己意而受苦，而非心甘情愿去做。然而，习惯因我而变得更与我交战，因为我曾自愿地走向我所不愿之处。那么，当罪人遭受公义的惩罚时，谁还能合理地说什么呢？我也不再有惯常的借口，即我之所以迟迟不愿超越世俗、事奉祢，是因为我对真理的认识还不确定；因为现在真理已经确定了。但我，仍被世俗捆绑，拒绝作祢的战士；我害怕摆脱一切纠缠，正如我们应该害怕被纠缠一样。\par

12. 我就这样为世俗的行李甜蜜地压着，犹如在睡梦中一般；我思念祢的种种思绪，好比那些切望醒来的人的努力，他们虽仍被浓重的睡意所笼罩，却力图挣醒，旋即又沉入其中。正如没有一个人愿意永远睡着，按一切人的清醒判断，觉醒总比沉睡好，然而一个人在四肢都感到沉重的昏倦时，往往迟迟不肯摆脱睡意，虽感不满，甚至到了该起来的时候，仍乐意再躺一会，我也如此，虽然确信，把自身交付于祢的爱德，远比随从我的私欲更好；前者说服了我并战胜了我，后者却使我悦乐并束缚了我。当祢呼唤说：\BibleQuote{你这睡眠者，醒起来吧！从死者中起来吧！基督必要光照你！}\BibleRef{弗 5:14} 我没有什么可回答祢。祢从各方面给我指出祢的话是真的，我被真理折服了，一句话也答不出来，只吐出了慢吞吞的、昏沉的言语：\Quote{马上，看，马上；}\Quote{再等一会儿。} 可是这 \Quote{马上，马上} 却没有马上；我的 \Quote{再等一会儿} 却延长了很久。我虽然 \BibleQuote{照内心的人是喜悦天主的法律}，可是 \BibleQuote{另有一条法律，与我理智所赞同的法律交战，并把我掳去，叫我隶属于那在我肢体内的罪恶的法律。} \BibleRef{罗 7:22-23} 因为罪恶的法律是习惯的暴力，它强拉着人的心志，违其意愿而被禁锢，这原是心志自愿投入其中的，所以应当受这禁锢。我真是 \BibleQuote{我这个人真不幸呀！谁能救我脱离这该死的肉身呢？}\BibleRef{罗 7:24} 只赖祢的恩宠，通过我们的主耶稣基督，才能。\par

% structure: p0022 source=h2 target=subsection reason=relative-heading-rank; base=h2

\subsection{第六章 蓬提提亚努斯讲述隐修制度创立者安当以及效法他的一些人的事迹}

13. 至于祢如何解救我于肉情的束缚，就是那束缚我极紧的，又解救我于世俗事务的劳苦之中，我现今要向祢的名宣告并忏悔：\BibleQuote{上主，我的力量！我的救主！}\BibleRef{咏 19:14}。在日益增多的忧虑中，我料理着日常的事务，每日向祢叹息。我尽可能常到祢的圣堂中去，只要那令我悲叹的重担下的事务，给我些许空闲。\Person{阿利比乌斯}与我在一起，他经过三次开庭之后，他的法律事务暂告一段落，正等待再次出售他的法律建议的机会，如同我惯于出售讲演之能——如果这种能力能由教授供给的话。\Person{内布利提乌斯}则由于我们的友谊，同意在\Person{韦雷昆都斯}门下教书，\Person{韦雷昆都斯}是\Place{米兰}公民，文法学者，也是我们大家极亲密的朋友；他迫切地需要，并且凭友谊之权，要求我们团体给予他极需的忠实援助。\Person{内布利提乌斯}并非被任何私利引诱而这样做的（因他若愿意，他的学识本可获利更多），但他既是最为和善可亲的朋友，便不愿不尽友爱的义务而轻视我们的请求。然而他行事非常审慎，小心避免让世俗所推重的那些大人物认识他；这样便能避免心神纷扰，他极愿心神自由，并拥有尽可能多的闲暇时间去探索、阅读或聆听有关智慧的事情。\par

14. 于是某一天，\Person{内布利提乌斯}不在场（原因我记不起了），看，有一位名叫\Person{蓬提提亚努斯}的同乡来访我和\Person{阿利比乌斯}，就他是\Place{非洲}人而言；他在皇帝的朝廷中担任要职。他找我们有什么事，我不得而知，但我们坐下来交谈时，他注意到在我们面前的一张游戏桌上有一本书；他拿起书来，打开一看，出乎他意料，发现是宗徒\Person{保禄}的书信——因为他原以为是我在教学中弄到精疲力竭的某类书。他微笑着看我，表示又喜出望外又诧异，居然在我眼前意外地发现了这本书，而且只这本书。原来他既是基督徒，又领过洗，常跪在圣堂中向祢、我们的天主祈祷，恒心又日行不断。于是，当我告诉他，我对这些著作下了很大功夫时，便展开了交谈，他谈起\Person{安当}这位埃及的隐修士，其名在祢的众仆人中间极为闻名，但那时我们却未听说过。他得知这情况后，便详述此人，把他这位杰出人物的事迹讲给我们听，并对我们的无知感到惊讶。而我们则惊讶不已地听着，因为祢奇妙的作为，在如此晚近、几近我们这个时代，在纯正的信仰与\Emph{天主教会}内，如此充充足足地显示了出来。我们大家都惊奇——我们惊奇这些事如此伟大，他则惊奇我们竟从未听说过。\par

谈话由此转向了隐修院中的团体，以及他们那芬芳馥郁、令祢喜悦的品行，又谈到了那丰美荒野中的旷野，对此我们一无所知。在米兰有一座隐修院，满有好弟兄，位于城墙之外，在\Person{圣安博罗削}的照料之下，我们竟也不知道。他继续讲述，我们专注而静默地聆听。他接着告诉我们，某天下午，在\Place{特里尔}，皇帝正忙着观看赛马竞技，他和另外三个同伴外出散步，来到城墙附近的花园中。他们在那里偶然两人一行地走着，其中一人陪他走开了，另外两人则单独行走；这两人在漫步时，来到一间小屋，里面住着祢的一些仆人——\BibleQuote{神贫的人}，因为\BibleQuote{天国是他们的}——在那里他们找到了一本书，上面记载着\WorkTitle{圣安多尼传}。其中一人开始阅读，惊叹不已，心中燃起热火；在阅读中，他思索着要接受这样一种生活，放弃世俗的职务来事奉祢。他们都是被称作\Quote{公事代理人}的人。随后，他突然被一种圣爱和清醒的羞愧所压倒，对自己感到愤怒，转眼看向他的朋友，喊道：\Quote{请告诉我，我恳求你，我们通过这一切劳苦努力追求的是什么目的？我们的目标是什么？我们服侍的动机何在？我们在宫廷中的希望难道能超越成为皇帝的仆役吗？而身处这样的位置，有什么不是脆弱且充满危险的？要经历多少危险才能达到更大的危险？我们何时才能到达那里呢？但我若渴望成为天主的朋友，看啊，我现在就已经成就了。}他这样说着，在新生命分娩的阵痛中，他再次转眼看向书页，继续阅读，内心在祢的注视下被改变，他的心思脱离了世俗，很快就显明了出来；因为当他阅读时，心潮起伏，他一度挣扎，最终辨明并决心走上更好的道路，如今他已属于祢，就对朋友说：\Quote{现在我已挣脱了我们那些希望，决意事奉天主；从此刻、在此地，我便开始实行。你若不愿效法我，也不要拦阻我。}另一个人回答说，他要与他联合，共担如此大的赏报和如此大的事奉。于是，他们二人都属于祢了，便以必要的代价建造一座塔——\BibleRef{路 14:26}——舍弃一切所有，跟随祢。之后，\Person{蓬提齐亚努斯}和那位与他同行到花园其他部分的人，寻找到他们所在的地方，发现他们后，提醒他们天已近晚，该回去了。但他们向他说出他们的决心和意向，以及这决定是如何产生并坚定的，恳求他们若不愿加入，就不要打扰他们。但那两人，丝毫未改变从前的自己，却（据他所说）为自己感到悲哀，并虔诚地祝贺他们，托付自己代祷；他们的心仍倾向于尘世的事物，便回宫去了。但另外两人，将他们的爱慕放在天上的事上，留在了小屋里。他们二人都曾订有未婚妻，当未婚妻们听说这事后，也将自己的童贞献给了天主。\par

% structure: p0026 source=h2 target=subsection reason=relative-heading-rank; base=h2

\subsection{第七章 他哀叹自己的悲惨，年届三十二岁，仍未寻获真理。}

这就是\Person{蓬提齐亚努斯}讲述的故事。但是，主啊，当他说着的时候，祢使我转向我自己，将我从我的背后提出来，我本把自己放在那里，不愿进行自我省察；祢又使我面对面地面对自己，让我看见自己是多么污秽，多么弯曲、肮脏，布满污点和溃烂。我看见并厌恶自己；却不知要逃到哪里离避自己。如果我试图从我身上移开目光，他就继续他的叙述，而祢再次使我面对自我，将我推到自己的眼前，使我发觉自己的邪恶，并憎恨它。我从前认识这邪恶，却假装不认识——对它视而不见，并将它遗忘。\par

然而现在，我越热烈地爱慕那些我听说的、拥有健全情感的人，他们完全把自己交给祢，以获得治愈，相比之下，我就越发恨恶自己。因为自从我十九岁那年读了\Person{西塞罗}的\WorkTitle{霍尔滕西乌斯}，被激发起对智慧的渴望以来，又已经过去了我的许多年月（或许有十二年）；而我仍然迟迟不抛弃单纯尘世的幸福，不致力于寻求那样一种东西——光是寻找它，更不用说找到它，都理应比这个世界的财富和王权，即使已经得到，也比身体的享乐，即使随我心意环绕我，更受喜爱。但是我，可怜的青年，在我青春伊始便最为可怜，曾向祢恳求贞洁，并说：\Quote{求祢赐我贞洁和节制，但不要立刻。}因为我害怕祢会很快俯听我，很快治愈我那私欲偏情的疾病，而我宁愿它得到满足，也不愿它被消灭。我曾在一个亵渎的迷信中行走邪路；虽不确信于它，却宁愿选择它，而非其他，我不是虔诚地寻求，而是恶意地反对。\par

18. 我曾以为，我之所以日复一日拖延抛弃世俗希望、唯独跟随祢，是因为没有显出任何确定的方向指引我的道路。如今，日子到了，我在自己面前赤露敞开，我的良心将要责备我。\Quote{你在哪里，我的舌头啊？你确实说过，为了一种不确定的真理，你不愿抛弃虚荣的包袱。看哪，如今它是确定的，但那重担仍然压着你；而那些既没有如此耗尽心力寻求它，也没有耗费十年甚至更久思考它的人，却已卸下肩头的重负，并长出翅膀飞去了。}当\Person{蓬提提亚努斯}讲述这些事时，我内心备受煎熬，被一种可怕的羞耻重重压倒。他讲完故事，办完了来的事，便走了。对我自己，我内心什么没说呢？我用怎样谴责的鞭子抽打我的灵魂，叫它跟从我，努力追随祢而去！但它却后退；它拒绝，不肯振作。它所有的推诿都耗尽了，被驳倒了。只剩下无声的战栗；它惧怕被制止于那风俗之流，仿佛惧怕死亡，而它正是在那风俗中衰败，走向死亡。\par

% structure: p0030 source=h2 target=subsection reason=relative-heading-rank; base=h2

\subsection{第八章 与\Person{阿利比乌斯}的谈话结束后，他退入花园，他的朋友跟随其后。}

19. 当时，在我内心居所的巨大斗争之中——我在心室内极力煽起对抗我灵魂的斗争——我心神不宁，面容不安，一把抓住\Person{阿利比乌斯}，喊道：\Quote{我们到底怎么了？这是怎么回事？你听到了什么？无学问的人奋起夺取天国，\BibleRef{玛 11:12}而我们虽有学问，却缺少心志，看我们在血肉中打滚！因为别人已经走在我们前头，我们是羞于跟随，还是更应羞于不跟随呢？}我说了这类的话，激动地猛然离开他，而他则默默惊讶地凝视着我。因为我说得不像平常的语调，我的额头、脸颊、眼睛、面色、声调，都表达出我的激动，更胜言语。我们住处有一个小花园，我们可以使用，如同整间房屋一般；因为主人，也就是我们的房东，不住在那里。我心中的风暴把我赶到那里，在那里没有人能阻碍我正与自己进行的激烈较量，直到这较量达到祢知道的结局，而我并不知道。但我是在疯狂中求全愈，在死亡中求生命，知道自己是多么邪恶，却不知道自己不久将成为什么样的善。于是，我退入花园，\Person{阿利比乌斯}步步跟随。因为他的在场并不妨碍我的独处；或者他怎能离弃如此不安的我呢？我们在尽可能远离房屋的地方坐下。我心神不宁，对自己极度不耐，因为我尚未进入祢的旨意和盟约，我的天主啊，我全身的骨骼都向我呼喊，要我进入，并将它赞美到天上。我们进入彼处，不是借着船只，或车辆，或双脚，不，甚至也不是像我那样从房屋走到我们坐着的地方那么远的距离。因为，不仅要行走，还要进入那里，那无非是意愿去行，并且是坚决彻底的意愿；而不是这样那样摇摆不定，一个变化无常、半受伤的意志，扭斗不休，一部分落下，另一部分又抬起。\par

20. 最终，在犹豫不决的狂热中，我做出许多身体的动作，就是人们有时候想做却做不到的，要么因为他们没有肢体，要么因为肢体被绳索捆绑、因病衰弱，或其他任何方式受阻。这样，如果我揪头发，打额头，或者手指交缠，抱住膝盖，我这样做是因为我愿意。但我可能愿意却做不到，如果我肢体的运动能力没有响应。于是，我做了许多事，这时意愿并不等于能力，而我却没有做那件我怀着无与伦比的渴望更想做的事，并且我不久将愿意时就可能有能力做的事；因为不久我若愿意，我必彻底愿意。在这样的事上，能力与意愿是合一的，愿意就是行动，然而却未做成；身体更容易服从灵魂最轻微的愿望，在心灵的指挥下活动肢体，而灵魂却不服从它自己，去在单纯的意志中完成它这伟大的意愿。\par

% structure: p0033 source=h2 target=subsection reason=relative-heading-rank; base=h2

\subsection{第九章 心灵命令心灵，但并不完全愿意。}

21. 这种怪异的事情从何而来？原因何在？愿祢的仁慈照耀我，让我能探询，看看人类惩罚的隐秘之处，以及\Person{亚当}子孙最黑暗的痛悔，或许能回答我。这种怪异的事情从何而来？原因何在？心灵命令身体，身体立即服从；心灵命令自己，却遭到抗拒。心灵命令手动作，命令如此灵便，几乎分不清命令与执行。然而，心灵是心灵，手是身体。心灵命令心灵去意愿，然而，尽管它是它自己，却不服从。这种怪异的事情从何而来？原因何在？我再说，它命令自己去意愿，若非已愿意，就不会发出命令；然而它所命令的却未达成。但它并不完全愿意；因此它不完全命令。因为它命令到什么程度，就愿意到什么程度；而所命令的事未完成到什么程度，就不愿意到什么程度。因为意志命令要有意志；——不是另一个，而是它自己。但它不完全命令，因此所命令的就不存在。因为若是完全的意志，它甚至不必命令自己存在，因为它已经存在了。因此，部分愿意、部分不愿意，并不是什么怪异的事，而是心灵的软弱，因它不能完全奋起，虽有真理支持，却被习惯压伏。于是便有了两个意志，因为其中一个不是完整的；而一个所缺的，恰好由另一个所需来补足。\par

% structure: p0035 source=h2 target=subsection reason=relative-heading-rank; base=h2

\subsection{第十章 他驳斥\OriginalTerm{摩尼教徒}{Manichæans}关于两种心灵——一善一恶——的意见。}

22. 愿他们从你面前消灭，天主啊，如同那些灵魂的\BibleQuote{空谈者和骗子}\BibleRef{铎 1:10}消灭一样；他们观察到在考虑中有两个意志，便断言我们有两种心灵——一善一恶。他们自己确实是恶的，当他们持有这些邪恶见解时；他们若持有真理，并将赞同真理，就会成为善的，好让你的宗徒对他们说：\BibleQuote{你们从前是黑暗，但现在你们在主内却是光明。}\BibleRef{弗 5:8} 然而，他们渴望成为光明，却不是\Quote{在主内}，而是凭自己，揣想灵魂的本性与天主一样，结果成了更严重的黑暗；因为由于可憎的傲慢，他们离你更远了，你乃是\BibleQuote{那普照每人的真光，正在进入这世界。}\BibleRef{若 1:9} 注意你们所说的话，并要羞愧脸红；接近他而\Quote{蒙光照}，你们的面容就不会\Quote{羞愧}。我，当我正在考虑现在事奉主我的天主，如同我长久以来所定意的——是我愿意，是我不愿意。是我，正是我自己。我既不全心愿意，也不全不愿意。因此我与我自身交战，并被我自己所毁灭。这毁灭违背我的意愿发生在我身上，然而并非显示有另一种心灵，而是对我自己心灵的惩罚。\BibleQuote{所以现在，已不再是我做这事，而是住在我内的罪}\BibleRef{罗 7:17}——这是更不受拘束之罪的惩罚，因为我是\Person{亚当}的子孙。\par

23. 因为如果有多少相反的意志，就有多少相反的本性，那么现在就不会只有两种本性，而是许多种。若有人考虑是去他们的集会所，还是去剧场，那些人立刻喊道：\Quote{看啊，这里有两种本性——一种善，引人往这边；一种恶，拉人往那边；否则为何会有这对立意志间的犹豫不决呢？} 但我回答说两者都是恶的——那引人到他们那里的，和那引人回到剧场的。但他们不信那引向他们的是善的意志。那么，假设我们中有一人考虑，因两种意志的冲突而犹豫是否去剧场还是我们的教会，这些人岂不也不知如何回答吗？因为要么他们必须承认，他们虽不愿，那引向教会的意志是善的，如同那些领受并持守他们奥秘之人的意志是善的一样；要么他们必须想象在一个人内有两种恶的本性和两种恶的心灵，彼此交战；而他们所谓的一种善一种恶就不再正确了；要么他们必须转向真理，不再否认，在任何考虑中，有一个灵魂在相互冲突的意志之间摇摆不定。\par

24. 那么，当他们看到同一个人内有两种彼此对立的意志时，不要再说是两种对立心灵、两种对立实体、来自两种对立本原的斗争，一善一恶。因为你，真实的天主啊，驳斥、阻止并说服他们；比如，当两个意志都是恶的时，一个人考虑是否用毒药杀人，还是用剑杀；是否侵占别人的这一处或那一处产业，而他不能两者兼得；是否通过奢侈挥霍购买快乐，还是因贪婪守住钱财；是否去竞技场还是剧场，如果两者在同一天开放；或者，第三，是否抢劫别人的房屋，如果有机会；或者，第四，是否犯奸淫，如果同时有办法——所有这些事在同一时刻发生，同样渴望，却不可能同时施行。因为心灵被四个，乃至（在人们渴望的众多事物中）更多对立的意志所撕裂，但他们还不承认有那么多不同的实体。在善的意志中也是如此。因为我问他们，喜欢阅读\WorkTitle{宗徒书信}是善吗？喜欢庄重的\BibleBook{}{圣咏}是善吗？喜欢讨论\BibleBook{}{福音}是善吗？对每一个他们都会回答：\Quote{这是善的。} 那么，如果所有同样令我们喜悦，且都在同一时间呢？当一个人考虑更应选择哪一个时，不同的意志岂不分散心灵吗？然而它们都是善的，且彼此争斗，直到选定一个，那时整个统一的意志被带向那里，而之前它被分为许多。因此，同样，当永恒的喜乐自上吸引我们，而暂时的快乐自下羁绊我们时，是同一个灵魂不愿全心要这个或那个，因此被严重的困惑所撕裂，它因着真理更爱那一个，却因着习惯不能放弃这一个。\par

% structure: p0039 source=h2 target=subsection reason=relative-heading-rank; base=h2

\subsection{神魂如何与肉身交战，为能脱离虚幻的束缚。}

25. 我如此病苦，折磨不已，比往常更严厉地自责，在我的锁链中辗转反侧，直到那锁链完全断裂；我先前只是轻微地被它束缚，却仍被缠住。主啊，你以严厉的慈爱压迫我的内心，加倍增添恐惧和羞愧的鞭打，免得我再退让，而那条脆弱的残余束缚若未断裂，便会恢复力量，更紧地捆绑我。因为我在心里说：\Quote{看，现在成就吧，现在成就吧。} 我这样说着，几乎作出了决定。我几乎做了，却没有做。我既没有退回到旧日的光景，却站在旁边，喘了口气。我又尝试，离成功只差极小的一点，又差更少的一点，然后几乎触到并抓住了它；然而我仍未达到，未触到，未抓住，犹豫着向死亡而死，向生命而活；那习惯已久的恶事比那从未尝试过的善事更有力地控制我。那将使我成为另一个人的时刻越是临近，它越是令我恐惧；但它既未击退我，也未使我转向别处，只是让我悬而不决。\par

26. 玩具中的玩具，虚幻中的虚幻，我旧日的情妇们仍然束缚着我；她们摇动我的肉性外衣，柔声低语说：\Quote{你要离弃我们吗？从此刻起，我们就永远不再与你同在了吗？从此刻起，这个或那个也永远不再许可给你了吗？} 她们用\InnerQuote{这个或那个}向我暗示什么呢？我的天主啊，她们暗示的是什么呢？愿祢的仁慈使它远离祢仆人的灵魂。她们暗示了何等的不洁！何等的羞耻！现在我只听到她们不足一半的声音了，她们不再公然显现来反驳我，而是在我背后喃喃低语，在我离去时悄悄拉我，要我回头注视她们。然而她们仍然拖延我，使我迟疑不肯挣脱她们、摆脱她们，跳到我蒙召前往的地方——一个不驯的积习对我说：\Quote{你以为没有它们，你还能活下去吗？}\par

27. 但如今，这话说得很微弱；因为在我面向的一方，在我战栗前往之处，\OriginalTerm{节德}{Continence}的贞洁尊严显现给我，愉悦而不放荡轻浮，诚然吸引我前来，毫不疑惑，并伸出她圣洁的双手，充满众多良好的榜样，要接纳拥抱我。那里有如此多的青年男女，众多各年龄层的群众，庄重的寡妇和年老的贞女，节德在这一切之中，并非不生育的，而是喜乐的儿女们的多产母亲，借着祢，主啊，她的丈夫。她以鼓励的讥嘲向我微笑，仿佛在说：\Quote{这些青年男女能做到的，你不能做到吗？或者，他们能做得到，是靠他们自己，而不是靠主他们的天主吗？主他们的天主将我赐给了他们。你为什么靠自己的力量站立，却又站不稳呢？把自己投靠给祂吧！不要怕，祂不会退缩使你跌倒；毫无恐惧地投靠给祂，祂会接纳你，治愈你。} 我羞愧难当，因为我仍听见那些玩物的低语，悬而未决。她又仿佛说：\Quote{要闭耳不听你们在地上那些不洁的肢体，好使它们被治死。\BibleRef{哥 3:5} 它们向你讲论逸乐，却不是按主你的天主的法律。} 我心中的这场争论，无非是自我对抗自我。而阿利比乌斯，静坐我身旁，默然等待我这不寻常的激情的结果。\par

% structure: p0043 source=h2 target=subsection reason=relative-heading-rank; base=h2

\subsection{第十二章 向天主祈祷后，他倾流如雨的泪水，并受一个声音的提醒，打开书卷诵读\BibleRef{罗 13:13}的经文；由此，他整个人灵发生变化，向他的朋友和母亲宣示了天主的恩宠。}

28. 然而，当一次深刻的反省，从我灵魂的秘密深处，将我所有的悲苦聚集堆积在我心目之前时，兴起了一场强大的风暴，伴随着同样强大的泪雨。为了能充分地倾流这泪雨，又以自然的方式表达出来，我悄悄离开阿利比乌斯；因为我觉得独处更适合哭泣之事。于是我退到一段距离之外，以致他即使在场，也不至于令我压抑。当时我就是如此，他也觉察到了；因为，我相信，我曾说了些什么，声音似乎哽咽在泪水中，我就是在那种状态下起身的。于是他留在我们坐过的地方，万分惊讶。我投身倒在一棵无花果树下，怎样倒下的，我不知道，任凭眼泪流淌，我眼中的泪水涌出，成为蒙祢悦纳的祭献。\BibleRef{伯前 2:5} 我并非用这些话，但大意如此，向祢说了许多——\Quote{但是祢，主啊，要到几时呢？} \Quote{主啊，要到几时呢？祢要永远发怒吗？求祢不要记念我们从前的罪债；} 因为我感到我被它们束缚住了。我发出这些哀号——\Quote{要到几时，要到几时呢？明天，又是明天？为什么不是现在？为什么不在这一时刻结束我的不洁呢？}\par

29. 我正说着这些，心中怀着最辛酸的痛悔而哭泣，看呐，我听见邻近屋中传来一个声音，像是男孩或女孩，我不知道是谁，反复诵唱着：\Quote{拿起来，读吧；拿起来，读吧。} 我的脸色立刻变了，开始最认真地思量，孩子们在某种游戏中是否常有唱诵这类词句的习惯；我也想不起曾听到过类似的话。于是，我抑制住泪水的洪流，站起身来，不作他解，只当作这是上天给我的命令，要我打开书卷，诵读我首先翻到的一章。因为我曾听说过安当，他在福音诵读时偶然入内，领受了劝谕，仿佛所读的是直接对他说的：\BibleQuote{去，变卖你所有的，施舍给穷人，你必有宝藏在天上；然后来跟随我。} \BibleRef{玛 19:21} 因了这样的神谕，他立即归向了祢。所以我急速回到阿利比乌斯坐着的地方；因为我起身时，就把宗徒的书卷放在了那里。我抓起、翻开，默默地读了我的目光首先落下的那一段：\BibleQuote{不可狂宴豪饮，不可淫乱放荡，不可争斗嫉妒；但该穿上主耶稣基督；不应只挂虑肉性的事，以满足私慾。} \BibleRef{罗 13:13} 我不再往下读，也没有必要了；因为就在这句话结束的瞬间，一道安稳的光明，仿佛注入我的心中，一切疑惑的阴霾顿时消散了。\par

于是，我合上书，把手指——或其他什么东西——夹在书页间，然后带着平静的神情，将这事告诉了\Person{阿利比乌斯}。他也向我吐露了他内心所发生的变化，那是我所不知道的。他要求看看我读的那段经文。我指给他看；他看的地方更超过我读到的部分，而我不知道下文是什么。那正是：\BibleQuote{信德软弱的，你们要接纳；}\BibleRef{罗 14:1} 他把这话应用到自己身上，并向我说出了。这劝勉使他坚强起来；他怀着与他本性完全相合的善志和决心（在这方面，他向来远胜于我），没有片刻不安地迟延，便和我站在了同一阵线。接着，我们进去见我的母亲。我们把这消息告诉她——她欢欣鼓舞。我们叙述事情如何发生——她欢喜雀跃，得胜欢呼，并赞美祢，因为祢\BibleQuote{能成就一切，远超我们所求所想的。}\BibleRef{弗 3:20} 她看出，为了我，祢赐给她的，远超过她曾凭着可怜而极其痛苦的呻吟所求的一切。因为祢使我这般归向祢，以至我不再寻求妻子，也不再寻求这世界的任何其它希望，而是站立在那信德的准则上，就是许多年前祢曾在神视中向她预示我必会站立的准则。祢把她的哀伤化为喜乐，那喜乐远比她所渴望的更丰盛，远比她曾因我有肉身的子孙所渴求的更宝贵、更纯洁。\par

\SourceRecord{Translated by J.G. Pilkington.；Nicene and Post-Nicene Fathers, First Series；Vol. 1.；Edited by Philip Schaff.；Buffalo, NY: Christian Literature Publishing Co.,；1887.；<http://www.newadvent.org/fathers/110108.htm>.}

% structure: p0001 source=h1 target=section reason=page-title

\section{忏悔录（第九卷）}

他谈及自己放弃修辞学职业的计划；谈及朋友\Person{内布里迪乌斯}和\Person{维雷昆都斯}的死亡；谈及自己三十三岁时领受了圣洗；并谈及母亲\Person{莫尼加}的德行与逝世。\par

% structure: p0003 source=h2 target=subsection reason=relative-heading-rank; base=h2

\subsection{第一章 他赞美天主——安全的创作者，及救赎者\Person{耶稣基督}，承认自己的邪恶。}

\BibleQuote{主啊，我确实是祢的仆人；我是祢的仆人，祢婢女的儿子：祢解开了我的束缚。我要向祢献上感恩的祭。} 愿我的心和我的口赞美祢，愿我一切的骨骸说：\BibleQuote{主啊，谁能像祢？} 让他们这样说吧，求祢应允我，\BibleQuote{向我的灵魂说：我是你的救恩。} 我是谁？我的本性又是什么？我的行为何其邪恶；若非行为，则是言语；若非言语，则是意志。但祢，主啊，是美善和仁慈的，祢的右手垂顾了我死亡的深渊，从我心底清除了那腐化的深渊。结果便是，我不愿行我所愿的，而愿行祢所愿的。但是，在这许多年中，从何等深邃隐秘之处，我的\OriginalTerm{自由意志}{free will}竟在片刻之间被召唤出来，使我将颈项置于祢的\BibleQuote{柔和的轭}下，将双肩负上祢的\BibleQuote{轻松的担子}，\BibleRef{玛 11:30}，耶稣基督啊，\BibleQuote{我的力量，我的救赎主}？突然之间，脱离虚幻之乐，对我来说是多么甘饴！昔日我所害怕失去的，如今却乐于将其抛弃。因为祢，真实而至高的甘饴，将它们从我这里抛弃了。祢抛弃了它们，而取而代之进入了祢自己，——比一切享乐更甘饴，却非血肉所能感受；比一切光明更明亮，却比一切奥秘更隐晦；比一切尊荣更崇高，却非那自高之人所能体得。如今我的灵魂摆脱了追求与获取的忧苦啮噬，以及沉溺并煽动情欲之痒的苦楚。我向祢，我的光明、我的财富、我的健康，主，我的天主，絮叨不已。\par

% structure: p0005 source=h2 target=subsection reason=relative-heading-rank; base=h2

\subsection{第二章 因肺部不适，他考虑退出公众的赏识。}

在祢面前，我以为善的，是不要喧嚷地猛然抢夺，而是轻轻地撤回我口舌从空谈行业中的服务；好使那些不思念祢的法律，也不思念祢的平安，而只顾谎诞的谬见与法庭争论的青年，不再从我的口中购买用以助长其气焰的装备。恰逢葡萄收获假前几天，我决意忍耐到那时，以便照常离开，并且因为我已被祢赎身，不再回去出卖。我们的意向，祢是知道的；但除了我们自己的朋友外，人并不知道。因为我们在私下议定，不让这件事传到任何外人那里；尽管祢已赐予我们，在我们从泪谷上升、歌唱\OriginalTerm{登阶歌}{song of degrees}时，\BibleQuote{利箭}和\BibleQuote{毁灭的炭火}，来对付\BibleQuote{欺诈的舌头}，那在进言时攻击，示爱时吞噬的舌头，正如对待它的食物那样。\par

祢的爱德穿透了我们的心，我们将祢的言语，仿佛铭刻在肺腑中；祢仆人的榜样——他曾由黑变明，由死复生——簇拥在我们的心怀，焚烧并吞噬了我们沉重的麻木，使我们不致坠入深渊。他们极大地激励了我们，以致反对者的欺诈舌头的每一呼气，反而更焚烧我们，却不能熄灭我们。然而，因为祢的名，就是祢在全地所圣化的名，我们的誓愿与决意，也可能找到赞许者，若不等候这已近在眼前的假期，却提前离开一个众目睽睽下的公开职业，便显得像是自我炫耀；以致所有目睹我这行为，并看到近在咫尺的葡萄收获期我竟有意提前的人，都会大加议论，以为我试图显得是个大人物。众人思量并争论我的意向，而我们的善被毁谤，这又有什么益处呢？\BibleRef{罗 14:16}\par

此外，就在这年夏天，因过度的文字辛劳，我的肺部开始虚弱，难以深呼吸；胸中的疼痛显示肺部已受损伤，且不能承受过响或过久的讲话。起初这对我是一种磨难，因为它几乎迫使我辞去那教学的重担；或者，如果我能够痊愈并康复，至少也要暂停一阵。可是，当那想要安闲，以便看到祢就是主的渴望充分升起，并在我的内心巩固时，我的天主，祢知道我甚至开始庆幸自己已有这个现成的藉口——而且不是假造的——这颇能减轻那些为了子女的利益，巴不得我永无子女自由之权者的不满。于是，我满怀这种喜乐，忍耐到那段时日过去——或许是二十天左右——然而它们却是英勇地忍受了；因为那惯常承受这重担部分的贪欲已经离去，倘若没有忍耐取而代之，我早已被压垮了。祢的一些仆人，我的弟兄们，或许会说，在这事上我犯了罪，因为我既然已经全心全力投入祢的军伍，却又容许自己在那谎言的座位上安坐一刻。我并不想争辩。但是，至仁慈的主啊，祢岂不曾也连同我其它那些可怖致命的罪恶，在圣水中将这罪一并赦免且宽宥了吗？\par

% structure: p0009 source=h2 target=subsection reason=relative-heading-rank; base=h2

\subsection{第三章 他退隐到朋友\Person{维雷昆都斯}的庄园，\Person{维雷昆都斯}当时还不是基督徒，并提及他的皈依与死亡，以及\Person{内布里迪乌斯}的皈依与死亡。}

5. \Person{Verecundus}因我们的幸福而焦虑憔悴，因为他被自己的束缚牢牢捆住，看到他将失去我们的友谊。他当时还不是基督徒，尽管他的妻子是一位信友；然而，正是由于她，他比受其他任何事物的束缚都更紧，无法踏上我们已开始的道路。而且，他表示，除了那些不可能的条件外，他不愿成为基督徒。不过，他非常客气地邀请我们使用他的乡间别墅，只要我们留在那里。主啊，祢必为这事在\Quote{义人复活的时候}\Quote{报答}他，\BibleRef{路 14:14}因为祢已经赐给了他\Quote{义人的一份}。虽然我们不在时，他在罗马身患疾病，并因此成为基督徒，成为一名信友，然后离世，但祢怜悯了他，不仅怜悯了他，也怜悯了我们；\BibleRef{斐 2:27}免得我们想起这位朋友对我们无微不至的关怀，却又不能把他算在祢的羊群中，而遭受无法忍受的痛苦。感谢祢，我们的天主，我们是属于祢的。祢的劝慰、安慰和忠信的应许使我们确信，祢现在为\Person{Verecundus}在\Place{Cassiacum}的那座乡间别墅报答了他，在那里我们得以在祢内摆脱世俗的狂热而安息，用祢乐园的永恒清凉，因为祢赦免了他尘世的罪过，在那座流奶的山上，那座果实累累的山——祢自己的山。\par

6. 那时他满心忧伤；但\Person{Nebridius}却充满喜乐。虽然他当时也还不是基督徒，曾陷入那最有害的错误，相信祢的圣子只是一个幻影，但他摆脱出来后，就持守与我们相同的信仰；尽管尚未领受祢教会任何圣事的入门礼，但他是一位最热切的真理探求者。在我们因祢的圣洗而皈依和重生之后不久，他也成了天主教会忠实的一员，在他的故乡阿非利加，在自己的族人中，以完全的贞洁与节制侍奉祢，又使他的全家都归向了基督。之后，祢让他摆脱了肉身；如今他安息在\Person{亚巴郎}的\Quote{怀中}。无论那\Quote{怀中}意味着什么，我的\Person{Nebridius}，我亲爱的朋友，主啊，祢的儿子，一个自由人的养子，就活在那里；他活在那里。因为像他那样的灵魂，还能有什么别的地方呢？他就活在那里，关于那里，他曾多次问我——问那个毫无经验、软弱无能的我。如今，他不再把耳朵贴近我的口，而是把他属灵的唇贴近祢的泉源，尽情畅饮他渴慕的智慧——享有无尽的幸福。我也不相信他会因沉醉其中而忘记我，因为祢，他所畅饮的主，是眷顾我们的。因此，我们就这样安慰忧伤的\Person{Verecundus}（我们的友谊未受影响），关于我们的皈依，并按照他的状况——我指的是他的婚姻状况——劝勉他归向信德。我们还在等待\Person{Nebridius}跟随我们，他即将行动，那时，那些日子终于过去了；因为在我看来，它们漫长而多，缘于我对安逸自由的热爱，好能从我的骨髓深处向祢歌唱。我的心对祢说——我寻求祢的面；\Quote{主啊，我要寻找祢的面。}\par

% structure: p0012 source=h2 target=subsection reason=relative-heading-rank; base=h2

\subsection{第四章 在乡间，他专心于文学，并就\Person{Alypius}的幸福皈依阐释\BibleBook{}{第四篇圣咏}。他备受牙痛困扰。}

7. 终于，那一天来到了，我真的要卸下修辞学教授的职务，其实我在思想上早已卸下。这完成了；祢解救了我的口舌，如同祢先前已解救了我的心灵。我满心喜乐地赞美祢，然后与所有家人退居到别墅。在那段日子里，我的写作完全献身于侍奉祢，尽管在这个间歇期，我仿佛还在从骄傲的学府中喘息，我的著作可作见证——那些我与朋友们辩论的著作，以及我与祢独处时默思的著作；还有我与不在身边的\Person{Nebridius}的通信，可作见证。我何时才能有时间一一细数祢在那时所赐给我们的种种大恩呢？尤其是我正急忙要讲述更大的仁慈。我的记忆召唤我，主啊，向祢忏悔是多么甘怡：祢用什么内心的刺鞭驯服了我，祢如何使我谦卑，推倒我思想的崇山峻岭，修直我的弯曲，铺平我的坎坷之路；\BibleRef{路 3:5}祢又用什么方法驯服了我心爱的兄弟\Person{Alypius}，使他归向祢的独生子、我们的主和救主耶稣基督的圣名，这圣名他起初拒绝写入我们的著作。因为他宁愿让它们带有学府中\Quote{香柏}的气味，就是那已被主折断的香柏，而不是教会那些对蛇有害、有益健康的药草的气味。\par

8. 我的天主啊，当我阅读\Person{达味}的\WorkTitle{圣咏}——那些忠信之歌、虔诚之音，它们排除一切骄傲的狂傲——那时我初蒙你的真爱，正与\Person{阿利比乌斯}一同退居乡间别墅；他与我一样是\OriginalTerm{慕道者}{catechumen}，我的母亲伴在我们身边——外表虽是女子，却具男子般的信心，满有长者之宁静，充盈着母爱与基督徒的热忱。我向你献上何等的心声啊！我在那些圣咏中是如何被感动，如何为你燃烧，渴望如果能，便在全世界宣扬它们，以对抗人类的骄傲！然而，它们已在全世界被歌唱，没有人能躲过你的热力。我以何等激烈痛苦的悲伤愤恨\OriginalTerm{摩尼教徒}{Manichæans}；然而我又怜悯他们，因为他们不认识那些圣事、那些良药，反而疯狂抗拒那能治愈他们的解毒剂！我多么希望那时他们就在附近，且不知我在场，能够看见我的面容，听见我的话，当我在闲暇之时阅读\BibleBook{圣咏}{第四篇}时——那篇圣咏如何触动我。\BibleQuote{我呼求你时，你应允了我，我公义的天主；你从困苦中使我宽广；求你垂怜我，听我的祈祷！}\BibleRef{咏 4:1} 哦，但愿他们能听见我论这经文所说的话，而我却不知道他们在听，免得他们以为我是因为他们才这样说！其实，如果我察觉他们在听在看，我就不会说出同样的话，也不会以同样的方式说话；即使我说了，他们也不会像当我独自在你面前、出于灵魂的私密感受而说时那样领受。\par

9. 我时而恐惧战栗，时而充满希望，并在你的慈爱中欢欣，哦，圣父啊。这些情感都从我眼中和声音中流露出来，当你的善神转向我们，说：\BibleQuote{你们世人啊，你们的心迟钝要到几时呢？你们喜爱虚幻，追求诈欺要到几时呢？}\BibleRef{咏 4:2} 主啊，你已经显扬了你的圣者，使他从死者中复活，坐在你的右边\BibleRef{弗 1:20}，从那里他要在至高之处派遣他所预许的\BibleRef{路 24:49}\OriginalTerm{护慰者}{Paraclete}，\BibleQuote{就是真理之神。}\BibleRef{若 14:16}其实他已经派遣了他\BibleRef{宗 2:1}，只是我并不知道。因为他从死者中复活，升天之后，现在已被显扬，才派遣他来。果然，直到那时，\BibleQuote{圣神还没有赐下，因为耶稣还没有受到光荣。}\BibleRef{若 7:39}先知呼喊道：\BibleQuote{你们的心迟钝要到几时呢？你们喜爱虚幻，追求诈欺要到几时呢？该知道：上主已经显扬了他的圣者。}\BibleRef{咏 4:2-3} 他高呼：\BibleQuote{要到几时呢？}\BibleRef{咏 4:2} 他又高呼：\BibleQuote{该知道}\BibleRef{咏 4:3} 而我，长久以来蒙昧无知，曾\BibleQuote{喜爱虚幻，追求诈欺。}\BibleRef{咏 4:2} 所以我听了战栗，因为这些话正是对像我这般记起自己从前所作所为的人说的。在我曾误认为真实的那些\OriginalTerm{幻象}{phantasms}中，确有\BibleQuote{虚幻}和\BibleQuote{诈欺}。我在回忆的痛苦中大声热烈地讲了许多话，但愿那些现今依然\BibleQuote{喜爱虚幻，追求诈欺}\BibleRef{咏 4:2}的人能听见！他们也许会震动了，把它呕吐出来，而你必俯听他们向你呼求；因为那位如今在天上为我们转求\BibleRef{罗 8:34}的，的确为我们以血肉之躯真实地死了。\par

10. 我又读到：\BibleQuote{你们可以发怒，但是不可犯罪。}\BibleRef{弗 4:26} 我的天主啊，我多么激动！我已经学会为过去的事向自己\Quote{发怒}，以便将来不再犯罪。是的，要正义地发怒，因为并不是有另一个黑暗种族的本性替我犯罪，就如那些不肯对自己发怒、为自己积蓄在震怒之日和天主公义审判显示之时的震怒\BibleRef{罗 2:5}的人所声称的那样。我的善事也已不在身外，也不是以肉眼在那个太阳中寻求；因为那些想从身外得享欢乐的人，很容易陷入遗忘，浪费在可见和暂时的事物上，并用饥饿的思虑去舔舐自己的影子。唉，但愿他们因斋戒而困乏，能说：\BibleQuote{谁能使我们看到善？}\BibleRef{咏 4:6} 而主啊，我们会回答，且让他们听见。你的慈颜的光辉已高举在我们之上。因为我们并不是那光照普世的人的光\BibleRef{若 1:9}，而是被你所照亮的，我们从前是黑暗，如今在你内却是光明\BibleRef{弗 5:8}。哦，但愿他们能看见那内在的永恒者！我尝到了它，而痛恨自己不能将它明示给他们；那时他们却把自己的心灵放在眼睛里，带到外边远离你，说：\BibleQuote{谁能使我们看到善？}\BibleRef{咏 4:6} 但是就在那里，在我的内室里，我对自己发怒，在那里我内心被刺痛，在那里我献上了\BibleQuote{祭献}\BibleRef{咏 4:5}，把我的旧人钉死，开始决心度新生活，将我盼望放在你身上——正是在那里，你开始使我感到甘饴，\BibleQuote{使我的心欢乐。}\BibleRef{咏 4:7} 当我外面念这经文，心里感觉这经文时，我大声呼号。我也不再贪图世俗的财富，不再浪费时间，也不再被时间所消耗，因为我在你永恒的单一中已拥有另一种谷物、酒和油。\par

11. 我从内心深处大声呼喊，用以下诗句呼号：\Quote{啊，在平安中！}以及\Quote{永不变者！}啊，他所说：\Quote{我躺下安睡！}因为当 \Quote{那时就要应验经上所记载的：\BibleQuote{死亡在胜利中被吞灭了}？} \BibleRef{格前 15:54} 而祢是至高无上的\Quote{永不变者}，祢从不改变；在祢内才有安息，使人忘却一切劳苦，因为在祢以外别无真神，我们也不应追求那些与祢不同的众多事物；然而，主，唯独祢使我安居于望德之中。我读了这些，便心火炽燃；但我却不知道怎样对待那些既聋且死的人，我曾是他们中一个有危害的成员——一个刻薄盲目的毁谤者，毁谤那些浸润着天蜜、闪耀着祢光辉的著作；我也因这圣经的仇敌而忧心如焚。\par

12. 我将何时忆起那些假期中所发生的全部事情？但我既没有忘记，也不会对祢的鞭笞之严酷和祢慈悯之迅速绝口不言。那时祢以牙痛折磨我；当牙痛加剧，至使我无法言语时，我心里想到要敦促所有在场的朋友，为我向祢——那位全能医治之天主——祈祷。我便将此事写在蜡版上，交给他们阅读。随即我们以顺服的心屈膝下跪时，疼痛便离开了。但这究竟是怎样的疼痛？又如何消失的？我承认，我的主、我的天主，我极其害怕，因为自我幼年以来，从未经历过这样的疼痛。那时，祢的旨意深深烙印在我心中；我因信德而欢欣，并颂扬了祢的名。而这份信德，不让我安于过往的罪过，这些罪过尚未因祢的圣洗而获得赦免。\par

% structure: p0019 source=h2 target=subsection reason=relative-heading-rank; base=h2

\subsection{第五章 在\Person{盎博罗削}的推荐下，他阅读\BibleBook{依撒意亚}{先知书}，却不理解。}

13. 葡萄收获假期结束后，我通知\Place{米兰}城的居民，他们可以为自己的学生另觅一个辞藻贩子；因为我已经决定事奉祢，而且由于呼吸困难和胸痛，我不能再继续担任教职。我并通过书信，向祢的主教、圣人\Person{盎博罗削}，陈述了我此前的谬误和如今的决意，以便他建议我最好阅读祢的哪一部经典，从而使我更好地准备并适合领受如此大的恩宠。他推荐了\BibleBook{依撒意亚}{先知}；我相信，这是因为他比其他先知更清楚地预示了福音和外邦人的蒙召。但我因不理解这部书的前一部分，并以为全书都是如此，便把它搁置一旁，打算等我更熟悉主的话语之后再读。\par

% structure: p0021 source=h2 target=subsection reason=relative-heading-rank; base=h2

\subsection{第六章 他在\Place{米兰}与\Person{阿利比乌斯}及其子\Person{阿得奥达图斯}一同受洗。\WorkTitle{师说}一书。}

14. 此后，当我要报名\EditorialNote{候洗}的时刻到来时，我们离开乡村，返回\Place{米兰}。\Person{阿利比乌斯}也乐意与我一同在祢内重生，他此时已穿戴好与祢的圣事相称的谦逊，并且他是一个如此勇敢的克己者，竟以不同寻常的坚毅，赤足踏在意大利冻土之上。我们让那个孩子\Person{阿得奥达图斯}加入我们的行列，他是我肉身所生，生于我的罪中。祢造他造得甚好。他年仅十五，但才智却超过许多稳重博学的人。我向祢承认祢的恩赐，我的主、我的天主，万有的创造者，且有大能使我们的缺失得以重整；因为在那孩子身上，除了罪过，没有什么属于我。我们按照祢的训导教养他，是因为祢——绝无他人——启示了我们。我向祢承认祢的恩赐。我们有一本书，题名为\WorkTitle{师说}。那是他与我之间的对话。祢知道，书中论辩对手口中所说的一切，都是他十六岁时的思想。我还发现他有多项更奇妙的才能。那种才智令我敬畏。除祢之外，谁能是这些奇事的创造者呢？祢迅疾将他从世上接走；如今我忆念他时心中平安，因为我对他童年或青年，或对他的整个生命，都一无所惧。我们使他在祢的恩宠内与我们同龄，在祢的训导中受教育；我们接受了圣洗，过往生活的焦虑便离开了我们。那些日子里，我沉思祢关于人类救恩的深远计划，那种奇妙的甘饴滋味，我怎么尝也尝不够。我在祢的圣诗和圣歌中如何尽情流泪，被祢那言词甘美的慈母教会的歌声深深感动！歌声涌入我的耳中，真理灌注进我的心田，敬虔之情便由此沸腾，眼泪夺眶而出，我为此而感到无比幸福。\par

% structure: p0023 source=h2 target=subsection reason=relative-heading-rank; base=h2

\subsection{第七章 论\Place{米兰}教会所建立的圣诗；\Person{犹斯提纳}引发的对\Person{盎博罗削}的迫害；以及两位殉道者遗骸的发现。}

15. \Place{米兰}教会开始运用这种慰藉和劝勉的方式，弟兄姊妹们以极大的声音和心灵热诚一同咏唱，为时并不长。大约一年，或稍久一点，那时，童子皇帝\Person{瓦伦提尼安}的母亲\Person{犹斯提纳}，因被亚略派所诱惑而堕入异端，为了她的异端，迫害祢的仆人\Person{盎博罗削}。虔诚的信众在圣堂内守护，准备与他们的主教、祢的仆人同死。在那里，我的母亲、祢的婢女，在那些挂虑和守夜中担当了重要的角色，恒心祈祷。我们，当时尚未因祢圣神的热度而熔化，但仍被这惊惶扰攘的城池所触动。就在那时，教会仿效东方教会的做法，规定要咏唱圣诗和圣咏，免得民众因长时间的忧伤而消沉；这一自那时起一直保留到如今的习惯，如今已被遍及世界各地祢的众多会众，甚至几乎所有会众所效法。\par

16. 然后，你藉着异象指示你荣名卓著的主教，在你隐秘的库藏中保存了那么多年不朽的殉道者\OriginalTerm{热尔瓦削}{Gervasius}和\OriginalTerm{普罗塔削}{Protasius}的遗体所在之处；时机一到，你便将他们公诸于世，以抑制那妇人般的、却属王族的暴怒。当他们被发掘、起出，并以应有的敬意迎至\OriginalTerm{盎博罗削}{Ambrosian}大殿时，不仅那些被不洁之魔缠绕的人（魔鬼自己承认）得到痊愈，还有一位失明多年的、本城知名的公民，询问了民众欢腾的原因后，便冲上前去，让向导领他到那里。到了之后，他恳求允许用他的手帕触摸你圣者的灵柩；他们的死亡在\Emph{你}眼中极为珍贵。他照做，并将手帕敷在眼上，眼睛立即复明。于是声名远播；于是对你的颂扬燃烧——闪耀；于是那个敌人的心，虽尚未敞开全心相信，却也被抑制了迫害的狂怒。感谢你，我的天主。你究竟从何处、引向何处，使我的记忆如此引导，让我向你忏悔这些事——伟大的事，虽然我这健忘的人曾略过不提？然而那时，当你的\Quote{香液}的\Quote{芬芳}如此馨香时，我们岂不\Quote{追随你}\BibleRef{歌 1:3}。因此，在咏唱你的圣诗时，我更泪流不止；从前我渴慕你，如今终能在你内呼吸，至少这气息能回荡在这草芥之屋中。\par

% structure: p0026 source=h2 target=subsection reason=relative-heading-rank; base=h2

\subsection{第八章 \Person{埃沃迪乌斯}的归化，及其母亲在随他返回非洲途中去世；他深情叙述母亲的教育。}

17. 你使人和睦同居在一家之内，你又让我们的同乡青年\Person{埃沃迪乌斯}亦与我们结伴；他曾在处理公务时，先于我们归化于你并领了洗；他舍弃了世俗职务，预备自己专务于你。我们在一起，并打算怀着圣洁的意向共同居住。我们寻找最能有助于事奉你的地方，正一同返回非洲。当我们抵达\Place{台伯河口的奥斯蒂亚}时，我母亲去世了。许多事我略去不提，因急于赶路。请收纳我的忏悔和感谢，我的天主，为那无数我缄默不言的事。但关于你那婢女——即是生养我肉身的她，叫我能生于这短暂的光明，并在她心中叫我生于永生的那位——凡是我心灵所涌现的，我必不忽略。我不讲论她的长处，而是你在她身上的恩赐；因为她既非自造，亦非自教。你造了她，她的父母也不知将生出何等样人。是你基督的杖、你唯一圣子的管束，在她敬畏你的氛围中，于你一位忠实信徒的家中把她教养成人；那人是你的教会中健全的肢体。然而她并不将这种良好的管教，归于她母亲的勤勉，倒是归于某位年老衰迈的侍女；这位侍女曾在她父亲还是幼童时背负他，如同小小孩童惯于由较年长的女孩背负那样。为此，并因她年迈而品德极佳，深得那个基督徒家庭主人的敬重。于是便将照料主家女儿们的责任交给了她；她勤谨执行，必要时以圣洁的严厉克制她们，并以清醒的明智教导她们。因为，她们在父母餐桌前非常节制用餐的时刻以外，即使她们渴得难受，她也不许她们喝水；以此防备恶习，并加上忠告：\Quote{你们只喝水，因为你们不能支配酒；但到了你们出嫁，成为储藏室与酒窖的女主人时，水便被你们轻视，但饮酒的习惯却会存留。}藉着这样的训导和权威的命令，她约束了她们稚龄的贪欲，把女孩子们的渴望调节到合宜的限度，以至于凡不端方的事，她们便不渴求。\par

18. 尽管如此——如祢的婢女、我的母亲向我讲述的那样——一种对酒的嗜好曾悄然袭上她的心头。当她还是个端庄的少女时，一如往常被父母吩咐从桶中汲酒，拿着器皿放在出口下，在将酒倒入瓶之前，她会用一点点酒沾湿嘴唇，再多她的嗜好就不允许了。她这样做，并非出于任何饮酒的渴望，而是出于她那个年纪满溢的活泼，这种活泼因戏谑而冒泡，在年轻的心灵中，通常被长者的严肃所压制。于是，在那一点点之上，日复一日地增加一点点（因为\BibleQuote{凡轻忽小事的，必渐渐跌倒}），她便养成了一种习惯，竟至急切地将她那几乎满斟的小杯一饮而尽。那么，那位明智的老妇和她那恳切的约束在哪里呢？主，如果没有祢的药剂看顾我们，还有什么能战胜隐秘的疾病呢？父亲、母亲和养育者不在场，祢却在场，祢创造了我们，召唤了我们，也借着那些管理我们的人，为我们的灵魂的救恩行了些许善事，我的天主啊，那时祢做了什么呢？祢如何医治了她？祢如何使她痊愈？祢岂不是从另一个妇人的灵魂中引出苛刻而刻薄的辱骂，如同外科医生的刀从祢的秘密库房中取出，一刀切除所有腐烂吗？因为那个经常陪她去酒窖的女仆，恰好与她的年轻的女主人争吵，在单独相处时，用非常刻薄的话辱骂她，当她的面指控她的恶习，称她为\Quote{酒鬼}。她被这嘲讽刺痛，察觉了自己的污秽，立即谴责并戒绝了它。就如朋友以奉承使人败坏，敌人也常以嘲弄纠正我们。然而，祢并不按照他们所做的报应他们，而是按照他们所意图的。因为她发怒，是想激怒她的年轻的女主人，而不是要医治她；而且是在暗中进行，要么是因为争吵的时机和地点正好如此，要么是担心自己因这么晚揭露此事而陷入危险。但是，主啊，祢是天地万物的统治者，祢将最深的洪流转变为祢的目的，安排时代的狂乱潮流，祢借着一个灵魂的不健全来医治另一个灵魂；免得任何人注意到这事时，若另一个人因他的话语而悔改，便归功于自己的权能。\par

% structure: p0029 source=h2 target=subsection reason=relative-heading-rank; base=h2

\subsection{他描述了他母亲值得称赞的习惯；她对丈夫和儿子们的慈爱}

19. 她就这样被端庄而清醒地教养，与其说是由她的父母使她服从于祢，不如说是由祢使她服从于她的父母。当她到了适婚年龄，被嫁给了一个她当作主来服侍的丈夫。她努力使他归向祢，以她的举止向他传扬祢；祢使她因此显得美丽、可敬可爱，令她的丈夫赞叹不已。因为她在婚床被冒犯时如此忍受，以致从未因此与丈夫发生任何争执。因为她等候祢的怜悯临到他，使他因信靠祢而变得贞洁。此外，他在友谊上热心，在愤怒中却暴烈；但她学会了，愤怒的丈夫不应被顶撞，无论是行动还是言语。但等他一旦平静下来，她看到合适的时机，就会为她的行为给出理由，如果他此前是无缘无故发怒的话。简而言之，当许多主妇——她们的丈夫更为温和——在她们受羞辱的脸上带着殴打的痕迹，私下交谈时责备丈夫的生活，她却责备她们的舌头，严肃地告诫她们，好像在开玩笑：\Quote{从她们听到所谓的婚约书被宣读给她们的那一刻起，她们就应想到它们是使她们成为婢女的文书；因此，常记念自己的身份，就不应与她们的主人对立。}当她们知道她忍受着一个多么暴烈的丈夫，并且惊奇地发现从未有传言或任何迹象表明\Person{巴特利丘斯}殴打妻子，或者他们之间有过哪怕一天的争吵，她们私下向她询问原因时，她教给她们我上面提到的规则。那些遵守的人体验了这条规则的智慧，并欢欣快乐；那些不遵守的人则被辖制、受苦。\par

20. 她的婆婆起初因心怀恶意的仆人的谗言而对她存有偏见，她却以顺从得胜，以忍耐和温良坚持，以致婆婆主动向儿子告发那些搬弄是非的仆人的舌头，这些仆人曾扰乱她与媳妇之间的家庭和睦，并要求儿子为此惩罚他们。因此，当他——遵照他母亲的意愿，为使家庭纪律严明，并确保未来家庭成员的和谐——按照揭发者的意愿，用鞭打惩罚了那些被发现的人之后，她应许给予同样的奖赏，给任何为了讨好她而向她讲她媳妇坏话的人。此后，再没有人敢这样做，她们便以相互善意的奇妙甜美一起生活。\par

21. 我的天主，我的仁慈，祢也将这伟大的恩赐赐给了祢的那位好婢女，祢从她的胎中创造了我。她能在任何时候，在任何有分歧和不和的人之间，表现出自己是一个和平的调解者。当她从双方听到极其刻薄的话时——就如同膨胀而未消化的不和所惯于发泄的那样——当仇恨的粗劣以刻薄的话语向在场的朋友唾骂不在场的敌人时，她从不向一方透露另一方的任何事，只透露有助于他们和解的事。这在我看来似乎是件小事，但令我痛心的是，我认识无数的人，他们由于某种可怕且蔓延的罪恶感染，不仅向彼此发怒的敌人透露对方在激动时所说的攻击对方的话，而且还加上一些从未说过的话。而对于一个大度的人，不通过恶言挑起或加深人与人之间的仇恨，这似乎应是件小事，除非他同样努力以善言消除它们。她就是这样的人——祢，她最亲密的导师，在她内心的学校里教导她。\par

22. 最终，她也将自己的丈夫，在他尘世生命即将结束时，争取到了祢面前；她不必抱怨他作为信徒的任何事，那是她在他成为信徒之前所忍受的。她也是祢仆人的仆人。认识她的任何人，都在她身上大大地颂扬、尊敬并爱慕祢；因为通过圣善生活的果实见证，他们感知到祢临在于她的心中。因为\Quote{她曾是一个丈夫的妻子}，孝敬了父母，虔诚地治理家务，\Quote{因善行而有好名声}，\Quote{养育儿女}，每当她看到他们偏离祢时，就为他们再受产痛\BibleRef{迦 4:19}。最后，主啊，对所有我们这些在她于祢内安眠\BibleRef{得前 4:14}之前领受了祢圣洗的恩宠、共同生活在一起的人，她都献出了关怀，仿佛她是我们所有人的母亲；她服事我们，仿佛她是我们所有人的孩子。\par

% structure: p0034 source=h2 target=subsection reason=relative-heading-rank; base=h2

\subsection{他与母亲关于天国的谈话}

23. 她离世的日子临近了（这日子祢知道，我们却不知道），正如我所相信的，祢以祢隐秘的方式安排——有一天，我和她两人独自倚靠着一扇窗户，从那窗口可以看到我们在\Place{奥斯蒂亚}居住的房子的花园；在那里，远离人群，我们在长途旅行后的疲惫中为航海而休息。我们当时单独愉快地交谈着；\BibleQuote{忘尽我背后的，只向在我前面的奔驰}\BibleRef{斐 3:13}，我们在真理——即祢自己——面前，彼此探求圣人们的永生将是怎样的，那是眼未曾见，耳未曾闻，人心也未曾想到的。但我们仍敞开心灵的口，朝向祢那自高天涌流的泉源，就是\BibleQuote{生命的泉源}，\BibleQuote{在祢那里}；好让我们按自己的能力蒙受它的滋润，得以稍许揣摩这高深的奥秘。\par

24. 当我们的谈话达到这一点时，即使是最高的肉欲快感，在最明亮的物质光明中，也因那生命的甘美而显得不仅不值得比较，甚至不值得提及；我们怀着更炽热的情感举心朝向\Quote{那自同者}，逐渐穿越一切有形之物，甚至穿越那太阳、月亮和星辰照耀大地的天穹；是的，我们更以内心的默想、谈论和赞美祢的工程而向上飞升；我们到达了自己的心灵，并超越了它们，好能攀升到那永不匮乏的丰盛之境，在那里祢以真理之粮永远牧养以色列，在那里生命就是那\OriginalTerm{智慧}{Wisdom}，万物都是藉着她而造成的，无论是已有的，还是将来的；她并非受造，而是她如何存在，就将永远如何存在；是的，\Quote{曾经}和\Quote{将来}并不在她内，只有\Quote{现在}，因为她是永恒的，而\Quote{曾经}和\Quote{将来}并非永恒。当我们这样谈论，并努力趋赴她时，我们以整个心灵的努力，已稍稍触及了她；我们叹息，并在那里留下\BibleQuote{圣神的初果}\BibleRef{罗 8:23}，然后回落到我们自己唇舌的声响中，在那里言语有始有终。有什么能相似祢的\OriginalTerm{圣言}{Word}，我们的主呢？祂永存不变，却\BibleQuote{更新一切}\BibleRef{智 7:27}？\par

25. 当时我们说：如果一个人的肉体的喧嚣得以静默——大地、水和空气的幻象得以静默——天极也得以静默；甚至灵魂自身也向自己静默，并通过不思自己而超越自己——幻影和虚假的启示得以静默，一切唇舌、一切标记，以及一切暂存的事物都得以静默，因为，若有人能聆听，这一切都会说：\Quote{我们并非自造，而是由那永存者所造}。如果它们在说出这话后便静默了，只唤醒我们的耳朵去聆听它们的创造者，而祂不通过它们发言，却亲自发言，使我们能听见祂的圣言，不是经由肉体的唇舌，不是经由天使的声音，不是经由雷鸣，也不是经由比喻的隐晦，而是能听见祂——我们在这些事物中所爱的祂——脱离这一切，如同我们二人此刻努力伸展自己，并以神速的思维触及那凌驾万物之上的永恒智慧。如果这能持续下去，其他截然不同的景象都退去，而这景象将观者攫取、吸纳并笼罩在这些内心的欢乐中，以至他的生命能永远如同我们此刻所渴慕的那一刹那的明悟，那岂不就是\BibleQuote{进入你主人的福乐吧}\BibleRef{玛 25:21}？而何时才能实现呢？就在我们众人都复活之时，但并非所有人都要改变。\par

我正说着这类话；即使不是以这种方式、用这些言语，主啊，祢知道，那天当我们这样交谈时，这个世界及其一切欢乐，就在我们说话之际，变得令我们鄙夷。那时我母亲说：\Quote{儿子，至于我自己，此生任何事物我都不再觉得有乐趣。在此我还想要什么，为何还在此，我已不知，既然我在这世上的希望已经满足。确实曾有一件事，我愿在此生稍作停留，那就是在我去世之前能见你成为一位天主教基督徒。我的天主已远超此望，使我看到你鄙视一切尘世幸福，成为祂的仆人——我还在此做什么呢？}\par

% structure: p0039 source=h2 target=subsection reason=relative-heading-rank; base=h2

\subsection{第11章 他的母亲，因热病倒下，在\Place{奥斯蒂亚}去世。}

我怎样回答她这些话，我记不太清楚了。然而，不到五天，或稍多一点时间，她因热病倒下；在病中，有一天她陷入昏厥，短时间对可见之物失去知觉。我们急忙赶到她面前；但她很快恢复了神志，看着站在她身旁的我和我弟弟，探询地对我们说：\Quote{我在哪里？}然后专注地看着我们因悲伤而茫然的样子，她说：\Quote{在这里，你们将埋葬你们的母亲。}我沉默不语，忍住不哭；但我弟弟说了些话，希望她更幸福地在自己的家乡而不是异乡去世。她听到这话，面带忧虑之色，用目光止住他，似乎察觉他的话语中有此俗念，然后看着我，说：\Quote{看，他说的什么话；}稍后她又对我们俩说：\Quote{将这身体随便埋在什么地方，不要为照料它而烦恼。我只请求你们一件事：无论你们在哪里，都要在主的祭台上纪念我。}当她用能说出的话语表达了这番意思后，就沉默了，因病情加重而痛苦着。\par

然而，当我默想祢的恩赐，啊，不可见的天主，祢将这些恩赐注入祢忠信者的心中，从中生出如此奇妙的果实，我便欢喜，向祢感恩，回想从前知道的，她曾怎样急切地操心自己的葬身之地，并为自己在丈夫遗体旁预备了地方。因为他们生前非常和睦地一起生活，她的愿望也曾是（人心能把握神圣之事的能力何其微小）将此加到那幸福之上，并流传于人间，即在她渡海漂泊之后，竟蒙准许，使他们二人在世上如此结合，最后同葬一穴。但当这种无谓的念头，因着祢的慈恩，开始从她心中消失时，我并不知道，我满心欢喜，惊叹她向我透露的这番话；虽然确实在我们窗前的谈话中，当她说\Quote{我还在此做什么呢？}时，她似乎并不渴望死在自己的家乡。后来我也听说，我们在\Place{奥斯蒂亚}的时候，有一天我不在场，她以母亲般的信赖，与我的一些朋友谈论轻看此生和死亡之福；当他们——惊讶于祢赐给她、一个妇人的勇气——问她是否不害怕把遗体留在离自己故乡如此遥远的地方时，她回答说：\Quote{没有什么对天主是遥远的；我也不必害怕在末世时，祂不知道从哪里将我复活。}于是，在她患病的第九天，在她五十六岁，我三十三岁的那年，那颗虔诚而热心的灵魂脱离了肉身。\par

% structure: p0042 source=h2 target=subsection reason=relative-heading-rank; base=h2

\subsection{第12章 他如何哀悼他的亡母。}

我合上她的眼睛；一股巨大的悲伤流入我心，正要化为泪水，但我的双眼同时，因心灵的强力克制，把那泉源吸得干涸，在如此挣扎中我何其痛苦！然而，她刚一咽气，男孩\Person{阿德奥达图斯}突然放声大哭，但被我们大家制止后，他便安静下来。同样，我自己孩童般的情感，正通过心中的年幼之声在泪水中寻求宣泄，也被约束而沉默了。因为我们认为以哀哭和呻吟来举行那葬礼是不合适的；如此哀悼的方式，通常是为那些死得不幸福、或完全死去的人而设。但她既非不幸地死去，也非完全死去。我们对此确信无疑，乃因她良好品行的见证，她\BibleQuote{无伪的信德}，\BibleRef{弟前 1:5}以及其他充分的理由。\par

30. 那么，是什么在我内里痛切地折磨我呢，难道不是那新受的创伤，即那最甜蜜、最亲爱的共同生活习惯被突然中断吗？我从她的见证中确实满心欢喜，因为在她最后患病时，她夸赞我的孝敬，称我为\Quote{善良}，并以极大的爱意回忆说，她从未从我口中听到任何苛刻或责备的话语。然而，啊，我的天主，创造我们的天主，我对她的敬爱怎能与她为我所受的辛劳相比呢？于是，当我失去在她内如此大的安慰时，我的灵魂遭受打击，那由她和我的生命共同融为一体的生活，仿佛被撕裂了。\par

31. 那时，那孩子被制止了哭泣，\Person{埃沃迪乌斯}拿起\BibleBook{圣咏}{集}，开始歌唱——全家应和——那首圣咏：\BibleQuote{我要歌颂仁慈与公义，主啊，我要向祢歌唱。}但当他们听到我们所作的事时，许多弟兄和虔敬的妇女都聚集起来；那些按惯例负责料理丧事的人正在准备葬礼时，我在房子的一个方便处，同那些认为不该留下我独自一人的人一起，谈论着适合当时情形的话；并借着这真理的慰藉减轻了祢所知道的痛苦——他们并未察觉，专心地听着，以为我毫无忧伤的感觉。但在祢的耳中，他们谁也听不到，我责备自己情感的软弱，抑制涌出的悲泣，它稍许向我屈服；但一阵剧痛又涌现，虽不至于使我放声大哭，或面容改变，然而我深知我在心中所压抑的是什么。我极其烦恼，这些人性的事物竟有如此力量支配我，这些事按照我们自然条件的正常秩序和命运是必然要发生的，我怀着新的悲伤为自己的悲伤而悲伤，被双重的哀伤所折磨。\par

32. 于是，当遗体被运出去时，我们往返都没有流泪。因为既在我们将救赎的祭献为她的缘故向祢呈上时——按照当地的习俗，在安葬之前，遗体正被放在坟墓旁边——我们所倾吐的那些祈祷中，我没有流泪；在那些祈祷中我也没有流泪；但我整个白天藏于内心的悲哀极重，并以不安的心灵，尽我所能地恳求祢治愈我的忧伤，但祢没有应允；我相信，祢藉此教训将一切习惯捆锁的力量，铭刻在我记忆里，即使是对一个不再以虚谎之言为食的心灵。我也觉得去沐浴是件好事，因为我听说浴室（\OriginalTerm{浴场}{balneum}）其名出自希腊文\OriginalTerm{浴场}{\GreekText{βαλανεῖον}}，因为它能驱除心中的烦恼。看啊，我同样向祢的仁慈承认，\BibleQuote{孤儿的父亲}，我沐浴了，感觉与沐浴前一样。因为我悲痛的苦涩并未从心中渗出。然后我睡去，醒来时发现悲痛减轻了不少；我独自躺在床上时，祢的\Person{安布罗削}的那些真实的诗句浮上心头，因为祢是——\par

万物创造者天主，\newline{} 穹苍的主宰，祢以灿烂的光辉装饰白昼，\newline{} 以恩宠的睡眠赋予夜晚；\newline{}\par

让安息使松弛的肢体\newline{} 恢复工作的精力，\newline{} 安慰疲惫的心神，\newline{} 解除焦虑的忧伤。\newline{}\par

33. 然后，我一点一点地又想起祢的婢女从前的种种，想起她对祢的虔敬生活，想起她对我们的圣善温柔和无微不至的关怀，而这一切骤然从我这里被夺去；在祢的鉴临之下，为她、为我、关于她、关于我自己而流泪，对我成了甜蜜的事。我任由先前压抑的眼泪，任其随意流淌，让它们在我心底下蔓延；我的心便安息其中，因为祢的耳贴近我——不是人的耳朵，他们会对我的哭泣加以轻蔑的解释。但现在，我正以文字向祢，主啊！忏悔。谁愿意读，就让他读去，任凭他怎样解释；如果他发现我为母亲哭泣了如此短的一段时光——这位母亲暂时于我眼中死去，却曾多年为我哭泣，好让我能在祢眼前活着——便算有罪，那么，让他不要嘲笑我，相反，如果他是心怀崇高爱德的人，让他为我对祢、祢的基督所有弟兄的父所犯的罪而哭泣吧。\par

% structure: p0050 source=h2 target=subsection reason=relative-heading-rank; base=h2

\subsection{第十三章 他为她的罪恳求天主，并劝勉读者以虔敬之心记念她。}

34. 但是——我如今的心已从那个伤痛中痊愈，就它所能被指为具有肉性的\BibleRef{罗 8:7}情欲而言——天主，我们的天主啊，我为了祢的那位婢女，向祢倾流一种迥然不同的眼泪，这眼泪是从一个因想到每一个死于亚当内的灵魂所面临的危险而破碎的心灵中流出来的。尽管她在脱离肉躯之前，已经\BibleQuote{在基督内活过来}，并且度着藉信德和言行赞美祢圣名的生活，但我仍不敢说，从祢藉圣洗重生她时起，从未有一句违背祢诫命的话出自她的口。\BibleRef{玛 12:36} 祢的圣子、真理曾宣告说：\BibleQuote{谁若向自己的弟兄说\InnerQuote{笨蛋}，就要受地狱的火刑。}\BibleRef{玛 5:22} 在人值得称赞的生命中，假若祢撇开仁慈而去究查，也有祸了。但因为祢不会苛刻地追究罪过，我们便满怀信心地期望在祢那里寻得某种宽恕之所。然而，谁若向祢历数他真正的功劳，他除了向祢历数祢自己的恩赐外，还有什么呢？啊，但愿人人都知道自己不过是人，并且\BibleQuote{夸耀的，应当因主而夸耀！}\BibleRef{格后 10:17}\par

35. 我，我的赞美、我的生命、我心灵的天主啊，暂且放下她那些使我欢喜感谢祢的善行，我现在为母亲的罪恳求祢。请通过那位被悬于木架之上、我们创伤的良药，并现今坐在祢的右边\BibleQuote{为我们转求}者的转求，俯听我。\BibleRef{罗 8:34} 我知道她怀着仁慈行事，并从心里\BibleRef{玛 18:35}宽恕了她的欠债者的债务；求祢也宽恕她的债，无论她在领受救恩之水后的这么多年间欠下了什么。主啊，求祢宽恕她，求祢宽恕她，我恳求祢；不要\BibleQuote{与她严加审判}。\BibleRef{雅 2:13} 愿祢的仁慈超越祢的公义，因为祢的言语是真实的，祢曾向\BibleQuote{仁慈的人}许诺了仁慈；\BibleRef{玛 5:7} 因为祢怜悯谁就怜悯谁，恩待谁就恩待谁。\BibleRef{罗 9:15}。\par

36. 我深信祢已应允了我的祈求；但\Quote{主啊，请收纳我口自愿的祭献。}因为她在临终之日临近时，并不思虑将她的遗体华丽地裹殓，或用香料涂尸；也不贪求精选的坟茔，或渴望葬在祖茔。这些事她没有托付我们，只愿在祢的\OriginalTerm{祭台}{altar}前记念她的名字，她曾在那里无一日间断地事奉；她知道在那祭台上分发那\OriginalTerm{圣祭献}{holy sacrifice}，藉此涂抹了那反对我们的字据；\BibleRef{哥 2:14} 藉此战胜了仇敌，这仇敌数算我们的过犯，搜寻可告发我们的把柄，却在祂——我们因祂而得胜的那一位——身上，一无所获。\BibleRef{若 14:30} 谁能将那无辜的血归还给祂？谁能偿还祂买赎我们的代价，好从祂手中将我们夺去？祢的婢女将她的灵魂，用\OriginalTerm{信德}{faith}的锁链，系于我们赎价的\OriginalTerm{圣事}{sacrament}。谁也不能将她从祢的护佑中夺去。不要让\Quote{狮子}和\Quote{龙}以强力和诡计介入。因为她不会说她不欠什么，以免被狡猾的欺骗者定罪并战胜；但她会回答说她的\Quote{罪赦了}\BibleRef{玛 9:2}，是那位无人能偿还代价的所赦免的，祂本不欠什么，却为我们交付了这代价。\par

37. 愿她因此与她的丈夫安息于和平中，她在丈夫之前之后都未嫁他人；她以坚忍服从丈夫，为祢结出果实，\BibleRef{路 8:15} 好使她也为祢赢得他。我的主，我的天主，请启发——启发祢的仆役、我的弟兄，祢的子女、我的师长，就是我在心声、行为和著作中所服事的那些人——使他们中凡读到这\WorkTitle{忏悔录}的人，都在祢的祭台前记念\Person{莫尼加}，祢的婢女，以及她的丈夫\Person{帕特里丘}，她曾藉他的肉躯，以我所不知的方式，带我进入此世生命。愿他们以虔诚的爱情，记念我今世转瞬之光中的父母，记念我那些同居于祢——我们的父——和我们的天主教慈母内的弟兄们，以及永世\Place{耶路撒冷}中的我的同胞——那耶路撒冷，祢的子民自离散以来直到归回，都在漂泊中叹慕。这样，我母亲对我最后的恳求，藉着这\WorkTitle{忏悔录}比我的祈祷更有力地，通过众人的祈祷，更丰盛地应允于她。\par

\SourceRecord{Translated by J.G. Pilkington.；Nicene and Post-Nicene Fathers, First Series；Vol. 1.；Edited by Philip Schaff.；Buffalo, NY: Christian Literature Publishing Co.,；1887.；<http://www.newadvent.org/fathers/110109.htm>.}

% structure: p0001 source=h1 target=section reason=page-title

\section{忏悔录（卷十）}

他已经表明了他过去和现在的所是，并展示了他忏悔的巨大果实；在即将考察以何种方法可以找到天主及幸福生活时，他详述了记忆的本性和功能。然后，他从三重试探的划分角度检视自己的行为、思想和情感；并记念主，那位天主与人之间的唯一中保。\par

% structure: p0003 source=h2 target=subsection reason=relative-heading-rank; base=h2

\subsection{第1章 人的望德与喜乐唯在天主内。}

1. 使我认识祢，认识我的那一位；使我认识祢，如同我受认识一样。\BibleRef{格前 13:12} 我灵魂的勇力啊，请进入我的灵魂，为祢自己预备它，好使祢能够拥有并保持它，\BibleQuote{毫无瑕疵，毫无皱纹}。\BibleRef{弗 5:27} 这就是我的望德，\Quote{因此我说了话}；当我这样有节制地喜乐时，我就在此望德中喜乐。对于此生中的其他事物，人越是为它们忧伤，它们就越是少值得忧伤；人越是不为它们忧伤，它们就越是应当令人忧伤。因为看，\BibleQuote{祢喜爱真理}，而凡行真理的，\BibleQuote{必来就光明}。\BibleRef{若 3:20} 这就是我愿做的事：在祢面前用心灵忏悔，并在许多见证人面前用文字忏悔。\par

% structure: p0005 source=h2 target=subsection reason=relative-heading-rank; base=h2

\subsection{第2章 万物在天主面前都是敞开的；向祂的忏悔不是用血肉的言语，而是用灵魂的言语和沉思的呼喊。}

2. 主啊，在祢眼前，人良心的深处是赤露敞开的，\BibleRef{希 4:13} 即使我不愿向祢忏悔，我内里又能有什么可隐藏的呢？因为那样我只会把祢从我面前隐藏，而不是把我从祢面前隐藏。但现在，我的叹息证明我对自己的不满，祢就发出光芒，使我满足，并且祢被爱慕、被渴求；好叫我为自己羞愧，弃绝自己，选择祢，并且除了在祢里面，既不能取悦祢，也不能取悦我自己。所以，主啊，无论我是什么，我都在祢面前显露无遗，至于我能怀着何种果实向祢忏悔，我已经说过了。我这样做，不是用血肉的言语和声音，而是用灵魂的言语，以及祢耳朵所熟悉的沉思的呼喊。因为我邪恶时，向祢忏悔无非是厌恶自己；当我真正虔诚时，忏悔无非是不将此归于自己，因为祢，主啊，\BibleQuote{祝福义人}；但首先祢\BibleQuote{称不虔敬者为义}。\BibleRef{罗 4:5} 所以，我的天主，在祢眼前的忏悔，是默默向祢而发，却又不是默默。因为在声响中它静默，在情感中它高呼。因为我向人说出的任何正确的话，没有一句祢不曾先从我这里听到；祢从我这里听到的此类话，也没有一句不是祢自己先对我说的。\par

% structure: p0007 source=h2 target=subsection reason=relative-heading-rank; base=h2

\subsection{第3章 善向天主忏悔者，最认识自己。}

3. 我与人们有何相干，要他们来听我的忏悔，好像他们能治愈我所有的疾病似的？他们是热衷于知道别人生活，却懒于纠正自己生活的民众。他们不愿意从祢那里听到他们自己是谁，为什么却渴望从我这里听到我是谁呢？而且，当他们从我这里听到关于我自己的事，他们怎能分辨我是否在说真话？因为除了在人内的灵，没有谁认识在人内的事。\BibleRef{格前 2:11} 但如果他们从祢那里听到关于他们自己的事，他们就不能说：\Quote{主在说谎}。因为从祢那里听到关于自己的事，不就是认识自己吗？人若认识自己，却说：\Quote{这是假的}，除非他本人说谎，否则有谁这样说呢？然而，因为\BibleQuote{爱德凡事相信}（至少对那些它通过合一而联结为一体的人）\BibleRef{格前 13:7}，所以，主啊，我也这样向祢忏悔，让人可以听到。对于他们，我无法证明我的忏悔是否真实，但那些耳朵被爱德开启的人，却相信我。\par

4. 但是，祢——我最隐秘的医生——请让我明白，我这样做能收获什么果实。因为我过去诸罪的忏悔——祢曾\Quote{赦免}并\Quote{遮盖}，为能在祢内使我获得幸福，用信德和祢的圣事改变了我的灵魂——当这一切被读到和听到时，就激励人心，不使它沉溺于失望而说\Quote{我不能}；但能唤醒它，在祢仁慈的爱和祢恩宠的甘饴中清醒过来，软弱的人因之而刚强，\BibleRef{格后 12:10}，如果他因此意识到自己的软弱。至于善人，他们乐于听闻如今已得解脱的人从前的过错；他们喜乐，不是因为那些事是过错，而是因为它们曾经是过错，如今已不再是了。那么，主、我的天主，我的良心每日向祢忏悔，信赖祢的仁慈远超乎自身的无罪——那么，我祈求祢，我借着这本书，在祢面前向人忏悔我此刻的所是，而非我曾是怎样的，究竟是为着什么果实呢？关于那种果实，我既已看见，也已讲述；但对于我此刻，就在我作此忏悔的这一刻，各种人都渴望知道我是怎样的人——无论是认识我的，或不认识我的，无论是听说过我，或从我这里听说过什么的——但他们的耳朵却不在我的内心，而我真正之所是，就在那里。因此，他们渴望听我忏悔我内里的所是，那是他们既无法用眼、也无法用耳、更无法用思想触及的地方；他们是抱着信仰的意愿来渴望此事——但他们能理解吗？因为爱德——他们藉以成为善人的爱德——对他们说，我在忏悔中没有说谎；这爱德在他们内相信我。\par

% structure: p0010 source=h2 target=subsection reason=relative-heading-rank; base=h2

\subsection{第4章 他在忏悔中为能行善而顾及他人。}

但他们希望这样得到什么果实呢？他们希望我幸福吗？当他们得知由于你的恩赐我多么接近你，并为我祈祷，当他们得知我多么被自己的重量拖累。对于这样的人，我愿表白自己。因为这并非小果实，我主天主，使我们为你而献上许多感谢\BibleRef{格后 1:11}，且有许多人因我们而向你恳求。愿弟兄之灵魂爱在我内你所教导应爱的东西，并哀叹在我内你所教导应哀叹的东西。愿这样做的不是陌生人的灵魂，也不是那些\Quote{外邦之子，口说虚妄，右手是虚谎的右手}的灵魂，而是弟兄之灵魂，当他赞同我时，为我欢喜；当他批评我时，为我难过；因为无论赞同还是批评，他都爱我。对于这样的人，我愿表白自己；让他们因我的善行而自由呼吸，因我的恶行而叹息。我的善行是你的安排和你的恩赐，我的恶行是我的过失和你的审判。让他们因善行而自由呼吸，因恶行而叹息；让赞美诗和眼泪从弟兄的心---你的香炉---升到你面前\BibleRef{默 8:3}。主啊，你喜悦你圣殿里的香，求你按你的大仁慈怜悯我，\Quote{为了你的名}。绝不要放弃你在我内已经开始的工程，求你完成我身上不完美之处。\par

这就是我的忏悔所结的果实，不是我曾经如何，而是我现在如何，我这样忏悔不仅在你面前，在暗中战栗的欢欣和暗中怀着希望的忧伤中，也在信仰者的耳中---他们是我的喜乐的分享者，我必死性的共担者，我的同胞和朝圣的旅伴，那些先行者、后来者，以及我旅途的同伴。他们是你的仆人，我的弟兄，你愿意他们成为你的子女；他们是我的主人，你命令我服侍他们，如果我希望与你共处并靠你而活。但如果你的话只以言语命令，而不以行动先行，那这话对我意义甚微。所以我在行为和言语上这样做，在你的翅膀荫下这样做，若非我的灵魂在你翅膀荫下顺服于你，我的软弱为你所知，这实在太危险。我是个小孩子，但我的父永远活着，我的护卫者\BibleQuote{足够}\BibleRef{格后 12:9}为我。因为他既是生我的，也是护卫我的；你就是我的一切善；全能的你，与我同在，并且在我与你同在之前就已经如此。因此，对于你命令我服侍的那些人，我愿表白，不是表白我曾经如何，而是我现在如何，还有我依然如何。但我也不判断我自己\BibleRef{格前 4:3}。这样我才希望被人听见。\par

% structure: p0013 source=h2 target=subsection reason=relative-heading-rank; base=h2

\subsection{人并不完全认识自己。}

因为，主啊，是你在审判我；\BibleRef{格前 4:4} 因为虽然\BibleQuote{除了在人内里的心神外，没有人能知道人的事}，\BibleRef{格前 2:11} 然而在人内仍有连\BibleQuote{他内里的心神}自己也不知道的事。但你，主啊，你造了他，完全认识他。我虽然在你眼前鄙视自己，视自己\BibleQuote{不过是尘土和灰烬}，\BibleRef{创 18:27} 却对你略有所知，而这是对自己所不知的。确实，\BibleQuote{我们现在是借着镜子观看，模糊不清}，尚未\BibleQuote{面对面}。\BibleRef{格前 13:12} 因此，只要我\BibleQuote{离开}你，我便\BibleQuote{更与自己同在}，而非你；\BibleRef{格后 5:6} 但我却知道你无法遭受暴力；至于我自己，我不知我能抵抗什么诱惑，不能抵抗什么诱惑。但仍有希望，因为你是忠信的，必不允许我们受诱惑超出我们所能承受的，且在受诱惑时，也必给我们开一条出路，使我们能承担。\BibleRef{格前 10:13} 因此，我愿承认我对自己所知的一切；我也愿承认我对自己所不知的一切。因为我对自己的所知，是因你光照我；我对自己的所不知，尚需等候，直到我的\BibleQuote{黑暗如同正午} \BibleRef{依 58:10} 在你面前。\par

% structure: p0015 source=h2 target=subsection reason=relative-heading-rank; base=h2

\subsection{天主的爱，因其本质超越万有，是通过感官知识和理性的运用而获得的。}

主啊，我爱你，并非带着不确定的意识，而是带着确定的意识。你以你的话击伤了我的心，我便爱你。还有天、地及其中的一切，看哪，它们从四面八方都对我说，我应当爱你；并且不住地向众人述说，\BibleQuote{使他们无可推诿}。\BibleRef{罗 1:20} 但更深的是，你将对谁施予仁慈，便对谁施予仁慈，将对谁施予怜悯，便对谁施予怜悯，\BibleRef{罗 9:15} 否则天地只向聋聩之耳宣扬你的赞颂。然而，爱你时我所爱的是什么？不是形体之美，不是时间的光辉，不是那悦目的光线灿烂，不是种种歌曲的甜美旋律，不是鲜花、香膏和香料的芬芳，不是玛纳与蜂蜜，不是肢体适于肉体拥抱的快乐。我爱我的天主时，不是爱这一切；然而我在爱我的天主时，确实爱某种光、声音、芬芳、食物和拥抱，这光、声音、芬芳、食物和拥抱是在我\OriginalTerm{内在之人}{homo interior}内的---那光照耀我灵魂，不受空间限制；那声音响起，不被时间攫取；那芬芳香飘，不随微风飘散；那食物滋养，吃用不尽；那拥抱相契，永不厌腻。这就是我爱我的天主时所爱的。\par

9. 这是什么？我问大地，它回答说：\Quote{我不是祂}；大地内的一切都同样宣认。我问海洋和深渊，以及其中一切爬行的生物，它们回答：\Quote{我们不是你的天主，向比我们更高的地方寻求吧。}我问微风轻拂的空气，以及整个大气连同其中的万物，它们回答：\Quote{\Person{阿那克西美尼}被骗了，我不是天主。}我问苍天、太阳、月亮和星辰，它们说：\Quote{我们也不是你寻求的那位天主。}我于是对站在我肉体门户周围的这一切事物说：\Quote{你们已经对我谈了我的天主，说你们不是祂；请告诉我一些关于祂的事。}它们便高声呼喊：\Quote{祂创造了我们。}我的询问就是我对它们的观察；它们的美丽便是它们的回答。我转而反观自己，问道：\Quote{你是谁？}并答道：\Quote{一个人。}看啊，在我身上既有身体亦有灵魂，一个在外，一个在内。我该用哪一个来寻求我的天主呢？我曾凭身体从地到天寻求祂，尽我所能派遣使者——即我的视线；但更好的部分是内在的那一个，因为所有肉体的使者——上天、大地及其中的万物——都对这内在的部分（它犹如主席与法官）提交了答案，它们说：\Quote{我们不是天主，但祂创造了我们。}内在的人借着外在的服侍得知这一切；我，内在的人，就是灵魂，通过身体的感官知道了这一切。我向我天主的广大大地询问，它回答我：\Quote{我不是祂，但祂创造了我。}\par

10. 这种美，凡感官健全的人不都看得见吗？为什么它不对所有的人说同样的话呢？大大小小的动物也都看见，却不能向它提问，因为它们的感官没有赋予理智，不能对所报告的事物作出判断；但人能提问，致使\BibleQuote{祂那看不见的属性……都可凭受造之物辨识和洞悉；}\BibleRef{罗 1:20}然而人由于爱这些事物，反被它们辖制；被辖制者便无法判断。受造物也不回答那些提问的人，除非他们能判断；它们也不会改变自己的声音（即它们的美丽），假设一个人只看见，另一个人既看见又提问，以致在此人眼中显现为一个样，在彼人眼中显现为另一个样；但它们对两者显现的却是同一个样，对此人是沉默，对彼人则说话——的确，它们对所有人都说话，但只有那些将外面接收的声音与内心真理相对照的人才能领悟。因为真理对我宣告：\Quote{无论苍天、大地，或其他物体，都不是你的天主。}它们的本性向观者如此宣告。\Quote{它们是团块；团块在部分中比在整体中小。} 我的灵魂啊，你是我较好的部分，我向你发言；因为你使你身体的团块充满生机，赋予它生命，而这生命没有一个身体能赋予另一个身体，但是你的天主对你而言，却是生命中的生命。\par

% structure: p0019 source=h2 target=subsection reason=relative-heading-rank; base=h2

\subsection{第七章 既不能从身体的能力，也不能从灵魂的能力寻得天主。}

11. 那么，当我爱我的天主时，我爱的究竟是什么？那位超越我灵魂顶端的祂是谁？我要借着自己的灵魂上升到祂那里。我要飞越我用来依附身体、并以生命充满其整个结构的那种能力。我不能以那种能力找到我的天主；否则，\Quote{无知的}骡马也能找到祂了，因为使它们身体生活的正是同一种能力。但另有其他能力，不仅是我用以赋予肉体生机的那种，而且是我用以使肉体（主为我所造）具有感觉的那种；它命令眼不要听，耳不要看；让这眼为我观看，让这耳为我听闻；其他感官也各有适当的本位和职分；我，这唯一的理智，通过它们各司其职。我还要飞越我这种能力；因为这种能力骡马也有，它们也借身体分辨事物。\par

% structure: p0021 source=h2 target=subsection reason=relative-heading-rank; base=h2

\subsection{第八章 记忆的本性与惊人能力。}

12. 因此，我还要飞越我本性的这种能力，一步步上升，趋近那创造我的主。我进入记忆的旷野和广阔厅堂，那里储藏着由感官从各式事物收入的无数的影像宝藏。凡我们想到的，或通过对感官所触及的事物加以扩大、缩小，或任何方式的变化而想象的一切，也都储藏在那里；还有那些未被遗忘吞没和埋葬的一切，都托付并储存其中。当我进入这个库房时，我要求把我想找的取出来，有些立刻出现；有些则需要更长时间的寻觅，仿佛从某个隐藏的容器中拖出来；另有一些则成群涌出，趁我正在寻找和询问别的东西时，它们跳跃而现，仿佛说：\Quote{莫非是我们吗？}我以心灵之手，将这些从记忆眼前驱开，直到我所要的那一个从隐秘的深处显露出来。其他事物不请自来，依次鱼贯而出，应召响应——前方的为后面的让路，在让路时又储存起来，以备我再次需要时取用。这一切的发生，正是当我凭记忆复述一件事的时候。\par

13. 所有这些事物，各经由其自身的入口进入，都清晰地在普遍名目下贮存于彼：例如，光，以及物体的各种颜色与形状，经由双眼；各类声音经由双耳；一切气味经由鼻孔的通路；一切味道经由口；而全身的触觉则带来硬或软、热或冷、平滑或粗糙、重或轻，无论是身体之外的或内在的。那记忆的巨大府库，连同其众多且难以形容的贮藏室，接纳这一切，以备不时所需，加以回想并提取；每一事物由自己的门户进入，便隐藏其中。然而，事物本身并不进入记忆，只有被感知到的事物的影像存于其内，随时可供思想回想。谁能说出这些影像是如何形成的？尽管我们清楚知道每一影像是经由哪一感官摄入并储存于内的。因为即使我身处黑暗与静默之中，只要我愿意，仍能在记忆中带出颜色，辨别黑白以及我想要的其它颜色；声音也不会闯入并扰乱我所汲取并正在思考的视觉影像，尽管它们也存在于彼处，被隐藏、如被分离存放一般。只要我喜欢，我也可以召唤这些声音，它们便立即出现。我的舌头歇息，喉咙无声，我仍能随己意尽情歌唱；而那些颜色的影像，虽也在其内，却不会在我思考另一经由耳朵流入的宝藏时，插进来扰乱我。其余由其它感官载入并堆积起来的事物，我也可随意回想。我能分辨百合与紫罗兰的香气，尽管并未嗅闻；我能喜爱蜂蜜胜过葡萄浆，喜爱平滑胜过粗糙，尽管那时我不曾品尝也不曾触摸，只是回忆。\par

14. 我是在我记忆的广阔厅堂之内做这些事的。因为在我那里，有天、地、海，以及我能在它们之中思想的一切，除了我已忘记的那些。我在那里也遇见我自己，并回忆我自己——我做了何事，何时、何地做的，以及做时我是怎样的心情。一切我记得的，无论是我亲身经历的，还是因着对他人的信赖而得知的，都在那里。我从这同一宝库中，以过去的经历，为自己构建出各种事物的相似形象，或是我曾经历过的，或是因经历而相信的；由此又构想出未来的行动、事件和希望，并对这一切进行沉思，宛如它们就在眼前。我对自己说：\Quote{我要做这事或那事，}而我在那充满如此众多宏大事物影像的、我心灵的广阔腹地中，对自己说：\Quote{于是这事或那事便会随之而来。}\Quote{哦，但愿这事或那事能够发生！}\Quote{愿天主阻止这事或那事！}我便这样对自己说着；当我说话时，我所说的一切事物的影像都从同一记忆宝库中呈现出来；要是这些影像不在，我根本无法说出任何有关它们的话。\par

15. 这记忆的力量是多么伟大，极其伟大，我的天主——那是一间内室，既宽广又无边际！有谁曾探得其底呢？然而，它是我的一种能力，属于我的本性；我自身并不能完全把握我之所是。所以，心灵太狭窄，无法容纳自身。那么，那自身不能容纳的部分究竟在哪里？是在自身之外，而非在自身之内吗？既然如此，为何还不能把握自身呢？我心中涌起莫大的赞叹；惊愕攫住了我。人们外出，赞叹高山的耸峙、大海的巨浪、江河的浩荡、海洋的辽阔，以及星辰的运行，却忘记了赞叹他们自己；他们并不思索，当我说起这一切时，我并未用双眼凝视它们，但是，若非我曾在记忆之中内在地看见那些我见过的山峦、波涛、河流、星辰，以及我所相信的那片海洋，且它们彼此间的广阔距离正如我于外界所见那样，我便根本无法言说它们。然而，当我用双眼观望它们时，我并非通过观看便将它们占为己有；与我同在的并非事物本身，而是它们的影像。我知道每一影像是由何种身体感官刻印在我身上的。\par

% structure: p0026 source=h2 target=subsection reason=relative-heading-rank; base=h2

\subsection{第9章 不仅事物，就连文学和影像也取自记忆，并藉回忆活动而呈现。}

16. 然而，这些并非我那无边无际的记忆容量所留存的一切。在此，也有我所领悟的一切人文学科知识，尚未被遗忘——它们仿佛被移入了一个内在之处，但那并非一个地点；所留存的并非它们的影像，而是事物本身。因为，无论文学，论辩之术，以及我所知的各种各样的问题，都如此存在于我的记忆之中，并非我只是摄取了影像而将事物本身留在外面，也不像声音那样，它响起后便消逝，仅以痕迹印刻于耳，借此能被记录，仿佛它仍在回响，其实早已止息；也不像气味，于飘散逝去、散入风中时，触动嗅觉，由此将自身的影像传入记忆，我们在回忆时方能觉察；也不像食物，在胃中确实已无味道，却在记忆中存有某种滋味；也不像任何由身体触觉感受到的事物，即使离我们而去，仍能在记忆中被想象出来。因为这些事物本身并未被置入记忆，只是它们的影像被以惊人的迅捷捕获，并储存，仿佛在最奇妙的仓库之中，当我们回忆时又奇妙地被呈现出来。\par

% structure: p0028 source=h2 target=subsection reason=relative-heading-rank; base=h2

\subsection{第10章 文学并非经由感官进入记忆，而是由其更隐秘之处所呈现。}

因此，当我一听有三种问题：\Quote{一件事物是否存在？——它是什么？——它属于何种性质？}我的确牢牢抓住构成这些话的声音的意象，我知道那些声音曾随着响声穿过空气，而现在已无存。但是，这些声音所意指的事物本身，我却从未借助任何身体感官获得，也从未以心智以外的方式感知到；我在记忆中保存的，不是它们的意象，而是它们自身。它们是如何进入我内的，让它们自己若能告诉我的话。因为我检查了我肉身的所有门径，却找不到它们是由哪一扇门进入的。因为眼睛说：\Quote{如果它们有颜色，我们报告了。}耳朵说：\Quote{如果它们发出声音，我们通知了。}鼻孔说：\Quote{如果它们有气味，它们就由我们通过了。}味觉说：\Quote{如果它们无味，不要问我。}触觉说：\Quote{如果它没有形体，我就没有触碰过它；如果我从未触碰过它，我就未曾通知过它。}这些事物是从哪里、又是如何进入我的记忆的？我不知道是如何进入的。因为我学习它们时，我没有信赖另一个人的心，而是在我自己的心中感知到它们；我认可它们为真，并将它们托付给记忆，仿佛存放起来，以便我随愿取出。那么，它们在我学习之前就已经在那里了，但却不在我的记忆中。那么，它们当时在哪里？或者，当它们被说出时，我为何会认出它们并说：\Quote{是这样，这是真的，}除非它们已经存在于记忆中，不过像是被推到深处、隐藏在更隐秘的洞穴里，若非由他人的提醒而被引出，我或许根本无法构想它们？\par

% structure: p0030 source=h2 target=subsection reason=relative-heading-rank; base=h2

\subsection{第十一章 学习与思考是什么}

因此，我们发现，学习这些我们并非通过感官吸收其影像、而是在自身内按其本身无影像地感知的事物，无非是通过默想这般的专注，并通过观察，留心使那些先前在记忆中散乱纷杂的概念，得到整理就绪，仿佛就在那同一记忆中，那些先前隐伏、散乱而被忽略的概念得以呈现在眼前，使善于观察的心智更容易得到它们。我的记忆保存了许多这类已被发现的事物，而且如我所言，它们仿佛已备妥待用；这些我们说是学习和知道了的事物；但若我们在短时内停止回想，它们就会再次沉没，仿佛滑入更偏僻的密室，必须像新事物一样重新展开（因为它们没有其他领域），并必须再次\OriginalTerm{组合}{cogenda}才能被知晓；也就是说，它们必须像从涣散中再次\OriginalTerm{收集}{colligenda}起来；由此我们有了\OriginalTerm{思考}{cogitare}这个词。因为\OriginalTerm{我收集}{cogo}与\OriginalTerm{我再收集}{cogito}的关系，就如\OriginalTerm{我行动}{ago}与\OriginalTerm{我反复行动}{agito}、\OriginalTerm{我做}{facio}与\OriginalTerm{我反复做}{factito}的关系一样。但心智将这个词（思考）据为己有，以至于不是在任何地方收集起来的，而是在心智中收集、即整理的，才恰当称为被\Quote{思考}。\par

% structure: p0032 source=h2 target=subsection reason=relative-heading-rank; base=h2

\subsection{第十二章 关于数学性事物的回忆}

记忆也包含数字和维度的理则与无数法则，它们中没有一个受到身体感官的印象，因为它们既无颜色，也无声音、味道、气味或触感。我听过用于讨论这些事物时所指的那些字词的声音；但声音是一回事，事物本身是另一回事。声音在希腊语里是一回事，在拉丁语里是另一回事；但事物本身却既非希腊语，也非拉丁语，也不是任何别的语言。我见过匠人的线条，甚至是最细的，如同蜘蛛网；但那些线条是另外一种，并非我肉眼所见的那些影像；认识它们的人，无需任何关于形体的观念，是在自己内感知它们的。我也曾观察过我们用来计数一切身体感官之物的数目；但我们用来计数的那些数目是另一种，不是这些的影像，因此它们确实存在。让那看不见这些事的人嘲笑我这样说；当他嘲笑我时，我要怜悯他。\par

% structure: p0034 source=h2 target=subsection reason=relative-heading-rank; base=h2

\subsection{第十三章 记忆保存一切事物}

这一切我都保存在我的记忆里，我是如何学到它们的，我也保存在记忆里。我还保存了许多我听过的、极不正确地反对它们的话；这些话虽然虚假，但我记得它们这件事却不虚假；我也记得，我曾将这些真理与那些反对它们的虚假言论区分开来；我现在看到，区分这些事是一回事，记得我经常区分它们——当我常反复思考它们时——是另一回事。因此，我既记得我常常理解了这些事，又将我现在所区分和领悟的存放于我的记忆中，以便将来我能记得我此刻理解了它。因此，我也记得我曾记得；这样一来，若以后我想起我曾能够记得这些事，那将是借着记忆的能力我才想起它。\par

% structure: p0036 source=h2 target=subsection reason=relative-heading-rank; base=h2

\subsection{第十四章 关于喜乐与忧伤如何带回心灵与记忆的方式}

21. 这同一个记忆也包含我心灵的情感；不是以心灵本身经受它们时那种方式，而是以一种非常不同的、记忆特有的能力。因为我可以不喜乐却记得自己曾有过喜乐；不悲伤却回想过去的悲伤；曾经害怕的事，我不用害怕就能记起；曾经有过的欲望，我不用欲望就能回忆。有时，恰恰相反，喜乐时我记得过去的悲伤，悲伤时我记得喜乐。就身体而言这并不奇怪；因为心灵是一回事，身体是另一回事。因此，如果我快乐时回想起过去身体的疼痛，并不是什么怪事。但现在，既然这记忆本身就是心灵（因为当我们吩咐人记住某事时，我们说：\Quote{记住，你要把这事放在心里；}当我们忘记某事时，我们说：\Quote{这事没进入我心里}，以及\Quote{这事从我心中滑掉了}，这样把记忆本身称为心灵），既然如此，怎么当我喜乐时记起过去的忧伤，心灵有喜乐，而记忆有忧伤——心灵因其中的喜乐而喜乐，记忆却并不因其中的悲伤而悲伤？难道记忆不属于心灵？谁会这么说？记忆无疑可以说是心灵的肚腹，而喜乐和悲伤就像甜与苦的食物，当它们被交付给记忆时，就像被送进了肚腹，在那里可以被储存，却失去味道。想象这些完全一样是可笑的；然而它们也并非全不相似。\par

22. 但是，请注意，当我肯定心灵有四种\OriginalTerm{扰动}{perturbations}——欲望、喜乐、恐惧、忧伤，并就此进行辩论，将每种分为其特有类别并加以定义时，我是从记忆中引发出这些的；我在那里找到我要说的，并从那里引发出它们；然而，当我通过回忆把它们召来时，我并不被这些扰动所困扰；在我回忆和审视它们之前，它们已经在那里；因此通过回忆才得以从那里引出它们。也许，就像反刍时食物从肚腹中取出，同样，通过回想，这些从记忆中引出。那么，为何辩论者这样回忆时，在他的沉思之口中，感受不到喜乐的甘甜或忧伤的苦涩？难道这个比较在这方面不类似，因为并非在所有方面都类似？因为谁愿意谈论这些话题，如果每次我们提到忧伤或恐惧，就不得不感到忧伤或恐惧？然而，我们若要谈论它们，就必须不仅在记忆中找到由身体感官印刻的名字之声音的影像，而且要找到事物本身的概念，这些概念我们从未通过肉体的任何门户获得，而是心灵本身通过自己\OriginalTerm{情欲}{passions}的经验认识到，然后交付给记忆，或者记忆自己保留下来，而并未交付给它。\par

% structure: p0039 source=h2 target=subsection reason=relative-heading-rank; base=h2

\subsection{第十五章 记忆中也有不在场之物的影像}

23. 但是，究竟是不是通过影像，谁能明确断言？我命名一块石头，我命名太阳，这些事物本身并不临在于我的感官，但它们的影像却贴近我的记忆。我命名身体的某种疼痛，当没有疼痛时它并不临在；然而，如果它的影像不在我记忆里，我就不知道该说它什么，在辩论中也无法区分它与快乐。我身体健康时命名身体的健康；事物本身确实与我同在，但除非它的影像也存于我的记忆，我绝不可能想起这个名字的声音所指为何。病人当听到\Quote{健康}一词时，也不会知道所言何意，除非同样的影像被记忆的能力所保留，尽管事物本身并不在身体内。我命名我们用以计数的数字；不是它们的影像，而是它们本身在我的记忆里。我命名太阳的影像，这也在我的记忆里。因为我并非回忆起那个影像的影像，而是影像本身，因为当我记起它时，影像本身临在。我命名记忆，我也知道我命名的是什么。然而，我除了在记忆本身里知道它，还能在哪里呢？难道它也是通过影像临在于自己，而不是通过它自己？\par

% structure: p0041 source=h2 target=subsection reason=relative-heading-rank; base=h2

\subsection{第十六章 记忆的缺失就是遗忘}

24. 当我命名遗忘，并且也知道我命名的是什么，我若不记得它，又从何知道呢？我指的不是名字的声音，而是它所指的事物，如果我已忘记它，就不可能知道那个声音意指什么。因此，当我记起记忆时，记忆便通过它自己临在自己。然而，当我记起遗忘时，记忆和遗忘同时临在——记忆使我记起，遗忘是我记起的事物。但遗忘岂不是记忆的缺失？那么，当我必须记起遗忘时，它又如何临在？因为当它临在时，我便不能记起。但如果我们所记起的东西都保留在记忆里，而除非我们记得遗忘，否则在听到这个名称时绝不可能知道它所指的事物，那么遗忘便由记忆保留着。因此，它临在，以免我们忘记它；当它临在时，我们却在遗忘。这是否意味着，当我们记起遗忘时，遗忘并非通过自身临在于记忆，而是通过影像？因为如果遗忘通过自身临在，它就不会引导我们记起，而是引导我们遗忘。现在谁将深究此事？谁能理解这是怎样的？\par

25. 主啊，我确实在其中辛劳，也在我自己内辛劳。我已成为一块需要过度劳作的麻烦土壤。因为我们如今不是在探索诸天的轨迹，不是在测量星辰的距离，也不是在探究大地的重量。正是我自己——我，心灵——在记忆。这并不令人惊异，如果我自己所不是的远离于我。但有什么比我自身更亲近于我呢？看，我竟不能领悟我记忆的能力，尽管没有它我连自己也不能称呼。因为当我很清楚我记起\OriginalTerm{忘性}{forgetfulness}时，我该说什么呢？我该断言我所记起的并不在我的记忆里吗？或者我该说忘性存在于我的记忆里，为的是我不忘记吗？这两者都是最荒谬的。还有什么第三种观点呢？我怎能说当我的记忆保留着忘性的影像，而不是忘性本身时，我却记起它呢？我又怎能说这话，既然任何事物的影像印入记忆时，那事物本身必须先已临在，以便那影像得以印入？因为我就是这样记起\Place{迦太基}的；这样记起我所去过的所有地方；这样记起我所见过的人的面容，以及其他感官所报告的事物；这样记起身体的健康或疾病。因为当这些对象临在时，我的记忆从它们接受了影像，当它们临在时，我可以在心中凝视并回想，当它们不在时，我记起它们。因此，如果忘性是通过其影像，而不是通过本身被保留在记忆中，那么本身曾一度临在，以便获取其影像。但当它临在时，它又怎样才能把它的影像刻在记忆上，既然忘性藉其临在竟能抹掉甚至它发现已经记下的东西？然而，无论以什么方式，尽管它是不能领悟、不能解释的，我仍然十分确定我也记起忘性本身，那忘却我们所记忆之物的能力。\par

% structure: p0044 source=h2 target=subsection reason=relative-heading-rank; base=h2

\subsection{天主不能被那走兽和飞鸟也拥有的记忆能力所达到。}

26. 记忆的能力伟大；我的天主啊，它极为奇妙，幽深而无限的万象；这能力就是心灵，这就是我自己。那么，我的天主啊，我是什么呢？我由何本性构成？一种多样、万千而极其广阔的生命。看啊，在我记忆的数不尽的田野、洞穴和洞窟中，充满了无数类别的事物，或是通过影像，如同一切形体；或是通过事物本身的临在，如同各种技艺；或是通过某种概念或觉察，如同心灵的情感，虽然心灵并不受苦，记忆却加以保留，而凡在记忆中的也都在心灵中：我穿越这一切，往来奔飞；我能深入到哪一方，就深入到哪一方，却没有止境。记忆的能力如此伟大，人生存的能力如此伟大，即使这生命是有死的。那么，我该做什么呢，我真正的生命、我的天主？我将超越我这被称为记忆的能力——我将超越它，好能走向你，甘饴的光明。你对我说什么呢？看哪，我正以我的心灵高翔向你，你常居我之上。我也将超越我这被称为记忆的能力，切望到达你那能被到达的地方，并且依附你那能被依附的地方。因为甚至走兽和飞鸟也拥有记忆，否则它们绝不能找回自己的巢穴，以及它们所习惯的许多其他事物；若非靠记忆，它们也绝不能习惯任何事物。那么，我将超越记忆，好能到达那将我与四足走兽和天空飞鸟分开，使我比它们更有智慧的祂那里。我将超越记忆，但是我将在哪里找到你呢，真正美好和确实甘饴的\Emph{你}？我将在哪里找到你？如果我在记忆之外找到你，那么我就是忘却了你。如果我不记得你，我又怎能找到你呢？\par

% structure: p0046 source=h2 target=subsection reason=relative-heading-rank; base=h2

\subsection{失物若非被记忆保留，便不能寻获。}

27. 因为那位丢失了达玛的妇人，点着灯寻找，\BibleRef{路 15:8}，除非她记得那块钱，否则绝不能找到。因为当找到了的时候，她若不记得，怎能知道是同一块呢？我记起我曾经遗失和找到过许多东西；由此我知道，当我寻找其中任何一件时，人若问我：\Quote{是这个吗？} \Quote{是那个吗？} 我就回答说：\Quote{不是}，直到我所寻找的那件被递给我。我若不记得那件东西——无论它是什么——即使递给了我，我也不会找到，因为我不能认出它。每当我们寻找并找到任何遗失之物时，情形总是如此。尽管如此，如果有什么东西意外地从眼前遗失，而非从记忆中遗失——如同任何可见的形体——其影像仍保留在内，并被寻找，直至它重现眼前；当它被找到时，便藉着内在的影像被认出。我们若认不出所遗失的东西，就不说自己找到了；我们若不记得它，便不能认出它。但这件东西，虽然从眼前遗失，却仍保留在记忆中。\par

% structure: p0048 source=h2 target=subsection reason=relative-heading-rank; base=h2

\subsection{记忆是什么。}

28. 但当记忆本身失去某物时又如何呢？比如我们忘记某事并试图回想起来。我们最终在哪里寻找呢？岂不还是在那记忆本身中吗？在那里，如果偶然出示一个事物代替另一个，我们便拒绝它，直到遇见我们所寻找的；那时，我们就喊道：\Quote{就是这个！}若非我们重新认识它，我们便不会如此，若非我们记得它，我们也无法重新认识它。因此，我们确实曾经将它遗忘。或者，并非整个事物都溜出了记忆，而是我们抓住某部分，从而寻找其余部分；因为记忆察觉到它并不像往常那样完整地运转，由于惯常联系被割裂，记忆便滞碍不前，要求补全缺失的部分。例如，我们看见或想到一个认识的人，却忘了他的名字，努力想记起它；此时，任何其他出现的事物都与它无关，因为人们素来不把它与这人联系起来思考，因此就被拒绝，直到那个作为其惯常对象、能使知识能安顿下来的东西出现。这东西若非来自记忆本身，还能来自何处呢？因为即使当我们因他人的提醒而认出它时，它也是来自记忆。因为我们并不把它当作新事物来相信，而是在回忆中承认人所言的正确。但如果它已完全从心灵中抹去，那么即使有人提醒，我们也不会记起。因为我们并没有完全忘记我们记得自己已忘记的事。因此，一个我们已经完全忘记的、遗失了的概念，我们甚至无法去寻找。\par

% structure: p0050 source=h2 target=subsection reason=relative-heading-rank; base=h2

\subsection{第二十章 我们若不认识天主与幸福生活，便不会寻求它们。}

29. 那么，主啊，我如何寻求祢呢？因为我寻求祢、我的天主时，我是在寻求\OriginalTerm{幸福生活}{happy life}。我要寻求祢，使我的灵魂得以生活。\BibleRef{亚 5:4} 因为我的身体靠我的灵魂生活，而我的灵魂靠祢生活。那么，我如何寻求幸福生活呢？既然我要等到能说\Quote{这就够了！}的时候，它才属于我，而我应在何处说这话呢？我如何寻求它呢？是通过回忆，仿佛我曾忘记它，并且知道我曾忘记它吗？还是因为我渴望学习一件未知的事——我从未知道过，或遗忘得连自己曾经遗忘都不记得了？幸福生活岂不是人人追求的事物吗？又有谁是绝对不追求的？但他们从哪里获得对它的认识，而如此渴望它呢？他们在哪里见过它，而如此爱它呢？的确，我们有它，但我不知是怎么有的。是的，当有人拥有它时，便是幸福的，这是另一种方式；也有人在希望中是幸福的。这些人拥有幸福的方式，比那些实际上幸福的人要低一等；然而他们仍比那些既非实际也非在希望中幸福的人好得多。即使这些人，若不以某种方式拥有幸福，也不会如此渴望幸福，而他们确实渴望，这是最确定不过的。他们如何知道它，我说不上来，但他们以某种我所不知的认识方式拥有它；我仍在深深怀疑，它是否存于记忆之中。因为如果它在那里，那么我们曾经幸福过；至于是一切人分别曾幸福过，还是在那个首先犯了罪的人内——我们众人都死在他内，并且都由他生于苦难中——我现在不问这个；但我要问，幸福生活是否在记忆中？因为我们若不知道它，就不会爱它。我们听到这名，我们全都承认我们渴望这事；因为我们并非只喜欢那声音。因为当一个希腊人听到有人用拉丁文说这话时，他并不感到欣喜，因为他不懂说的是什么；但我们感到欣喜，正如他若听到希腊文也会欣喜一样；因为这事本身既不是希腊的也不是拉丁的，但希腊人、拉丁人以及所有其他语言的人，都如此热切地渴望获得它。因此，这是众人都知道的，倘若能用同一声音问他们是否愿意幸福，无疑他们全会回答愿意。而若非那名称所指的事物本身保留在他们的记忆中，这事便不可能发生。\par

% structure: p0052 source=h2 target=subsection reason=relative-heading-rank; base=h2

\subsection{第二十一章 幸福生活如何存留在记忆之中}

30. 但这是否如同一个见过\Place{迦太基}的人记得它呢？不。因为幸福生活不是肉眼可见的，因为它不是形体。那么，是否如同我们记得数字呢？不。因为拥有数字知识的人，不再努力追求更多；然而我们虽在某种认识中拥有幸福生活，因而爱着它，却仍然希望进一步获得它，以便能幸福。那么，是否如同我们记得口才呢？不。因为尽管有些人听到这个名称时，会想起那事物，而他们事实上并非口才出众者，也有许多人希望如此，由此可见它存在于他们的认识中；然而这些人通过身体的知觉留意到别人有口才，并为此感到欣喜，且渴望自己也如此——尽管若非有某种内在的认识，他们便不会欣喜，若非欣喜，他们也不会渴望如此——可是对于幸福生活，我们却无法通过任何身体的知觉在别人身上经验到。那么，是否如同我们记得喜乐呢？也许是的；因为我即使在悲伤时也记得我的喜乐，正如我在痛苦时记得幸福生活一样。我从未借身体的知觉看见、听见、嗅到、尝到或触到我的喜乐；但我欢跃时，是在我的心灵中经验到它；关于它的认识黏附在我的记忆上，以致我能根据我现在记得自己曾为之喜乐的事物之不同，时而以轻蔑、时而以渴望回想它。因为我曾从污秽的事物中沐浴到某种喜乐，如今回想起来，我憎恶并诅咒它；在另一些时候，我从善良而正直的事物中获得喜乐，我怀着渴望回想它们，尽管它们或许不在眼前，我却悲伤地回想往昔的喜乐。\par

31. 那么，我何时何地经验过我的\OriginalTerm{幸福生活}{beata vita}，使我忆起它、爱慕它并渴求它？不仅是我，或少数人，希望幸福，而是所有人确实如此；对此，除非我们以某种确切的知识得知，否则我们不会以如此确定的意志去愿望。然而，这又是怎么回事：若问两个人是否愿意当兵，或许一人回答愿意，另一人回答不愿意；但若问他们是否愿意幸福，二人都会毫不犹豫地说愿意；而那人愿意当兵，这人却不愿意，岂不是都只为了幸福吗？这岂不是因为，正如有人以此事为乐，有人以彼事为乐，所有人都一致希望幸福，如同他们若被询问，会一致希望拥有喜乐——并将这喜乐称为幸福生活？那么，虽然一人以此方式追求喜乐，另一人以彼方式追求，所有人都有同一个目标——竭力追求拥有喜乐。这种生活，无人能说自己未曾经验过，因此它在记忆中被寻获，每逢听到幸福生活这一名称时便被认出。\par

% structure: p0055 source=h2 target=subsection reason=relative-heading-rank; base=h2

\subsection{第二十二章 幸福生活在于在天主内并为天主而喜乐}

32. 主啊，这绝不可——绝不可进入向祢告罪的你的仆人之心；无论何种喜乐，都绝不可让我自以为幸福。因为有一种喜乐，并非赐给\BibleQuote{恶人}\BibleRef{依 48:22}，而是赐给那些感恩地敬拜祢的人，祢本身就是他们的喜乐。幸福生活正是这样——因着祢、在祢内、并为了祢而喜乐；这便是它，别无其他。但凡以为还有另一种幸福生活的人，其实追寻的是另一种喜乐，而非真正的喜乐。然而，他们的意志并非背离了某种喜乐的影子。\par

% structure: p0057 source=h2 target=subsection reason=relative-heading-rank; base=h2

\subsection{第二十三章 一切人都希望在真理中喜乐}

33. 那么，并非一切人都希望幸福，这并不确定，因为那些不愿在祢内喜乐的人——祢才是唯一的幸福生活——实则并不真正渴望幸福生活。或者，所有人都渴望幸福，却因为\BibleQuote{肉身的欲望相反圣神，圣神相反肉身}，以至于\BibleQuote{不能做他们所愿意的}\BibleRef{迦 5:17}，于是他们堕入自己力所能及的事，并以此为满足？因为那力不能及的事，他们并不如此愿意，以至使自己有能力做到。我问每一个人，他宁愿在真理中喜乐，还是在虚假中喜乐。他们会毫不犹豫地说\Quote{在真理中}，正如他们会说\Quote{希望幸福}一样。因为幸福生活就是在真理中的喜乐。这就是在祢内喜乐，祢就是\BibleQuote{真理}\BibleRef{若 14:6}，我的天主，\Quote{我的光明，我面容的救援，我的天主}。所有人都希望这幸福生活；所有人都希望这唯一幸福的生活；所有人都希望在真理中的喜乐。我经验过许多想欺骗人的人，却没见过一个想受人欺骗的人。那么，他们从何处得知这幸福生活，岂不是也从得知真理之处吗？因为他们也爱真理，因为他们不愿受骗。当他们爱幸福生活——那无非是在真理中的喜乐——时，无疑也爱真理；而他们若不记忆中对真理有某种认识，就不会爱它。那么，他们为何不在真理中喜乐？为何不幸福？因为他们更全神贯注于那些反使他们痛苦的事物，而非那能令他们幸福的、他们记忆得极为微弱的事物。在人内还有一点光明；让他们行走吧——让他们\BibleQuote{行走}，免得\BibleQuote{黑暗}抓住他们\BibleRef{若 12:35}。\par

34. 那么，为什么真理会招致仇恨，而你们那位传讲真理的人\BibleRef{若 8:40}竟成了他们的仇敌，然而幸福生活——无非是在真理中的喜乐——却为人所爱？除非人们对真理的爱是这样的：凡爱其它事物的人，希望他们所爱的事物成为真理，并且，因他们愿意受骗，就不愿意被说服自己受了骗。因此，他们为了所爱而取代真理的事物，而恨恶真理。当真理光照他们时，他们爱真理；当真理谴责他们时，他们恨真理。因为，他们既不愿受骗，又想欺骗人，所以当真理自我启示时他们爱之，当真理揭露他们时他们恨之。为此，真理将如此报应他们：那些不愿被她发现的人，她既违背他们的意愿揭露他们，又将自身向他们隐藏。人心，如此盲目、病弱、卑劣、不体面，确实如此渴望自己隐藏起来，却不愿有任何事物向它隐藏。但得到的恰恰相反——它自己不能向真理隐藏，而真理却向它隐藏。然而，即使如此可怜，它仍宁愿在真理中喜乐，而不愿在虚妄中喜乐。那么，何时没有烦扰介入，它只在那使万物为真的唯一真理中喜乐，它就幸福了。\par

% structure: p0060 source=h2 target=subsection reason=relative-heading-rank; base=h2

\subsection{第二十四章 寻获真理者，即寻获天主}

35. 看，主啊，我如何在记忆中广寻祢；但在记忆之外，我未曾寻获祢。凡有关祢的事，我所寻获的，无一不是我自从认识祢以来保存在记忆中的。因为自从我认识祢，我从未忘记祢。因为我在何处寻得真理，便在何处寻得我的天主，祂就是真理本身，自从我认识真理，我未曾忘记它。同样，自从我认识祢，祢就住留在我的记忆之中；在那里，每当我回想祢时，我便寻见祢，并在祢内欢欣。这些是我的\OriginalTerm{圣洁的欢愉}{sanctae delectationes}，是祢按祢的慈悲，垂顾我的贫穷而赐给我的。\par

% structure: p0062 source=h2 target=subsection reason=relative-heading-rank; base=h2

\subsection{第二十五章 他欣然得知天主居住在他的记忆里}

36. 但是，主啊，祢住在我的记忆中的何处？祢在那里住在何处？祢为自己建造了怎样的室？祢为自己立起了怎样的圣所？祢已赐予我的记忆这份尊荣，得以让祢居住其中；但我正在默想祢住在其中的哪一部分。因为当我记念祢时，我飞越了那些连兽类也拥有的部分，因为在那些形体事物的影像中我找不到祢；我到达了我存放心灵情感的那些部分，在那里也没有找到祢。我又进入了心灵在我记忆中的本座，因为心灵也记得它自己——在那里祢也不在。因为祢不是形体的影像，也不是活物的感情，如我们欢欣、哀悼、渴望、畏惧、记起、遗忘之类；同样，祢也不是心灵自身，因为祢是心灵的主天主；这一切都是变化的，而祢超越万有，永不变易，却自打我认识祢时起，就屈尊居住在我的记忆里。但我如今为何要寻找祢住在其中的哪一部分，仿佛那里真有地方似的？祢确然住在其中，因为我自认识祢时起就记着祢，当我记念祢时，便在其中找到祢。\par

% structure: p0064 source=h2 target=subsection reason=relative-heading-rank; base=h2

\subsection{第二十六章 天主处处回应向祂请教的人。}

37. 那么，我在哪里找到祢，以致能认识祢呢？因为在我认识祢之前，祢并不在我的记忆里。那么，我在哪里找到祢，以致能认识祢，岂不是在高于我的祢内找到吗？没有地方；我们或往\Quote{后退}，或往\Quote{前进}，\BibleRef{约 23:8}也没有地方。真理啊，祢处处指导所有请教祢的人，同时答复他们，尽管他们请教的是不同的事。祢的答复是清晰的，尽管并非所有人都清晰地听见。无论他们希望什么，都来请教祢，尽管他们听到的并非总是他们所希望的。祢最出色的仆人，不是专注于从祢听到自己所希望的，而是希望从祢听到的。\par

% structure: p0066 source=h2 target=subsection reason=relative-heading-rank; base=h2

\subsection{第二十七章 他悲痛自己竟这么久没有天主。}

38. 我爱祢爱得太晚了，万古常新的美啊，我爱祢爱得太晚了！看，祢原在我内，我却在外，在那里寻找祢；我这丑陋的人，贸然投身在祢所造的美丽事物中。祢和我在一起，我却没有和祢在一起。那些事物使我远离祢，它们若非在祢内，便不存在。祢呼唤，高声呼喊，冲破我的耳聋。祢放射光芒，照耀四方，驱散我的瞽盲。祢散发芬芳，我吸入，便渴慕祢。我尝到了祢，便饥渴不已。祢触摸了我，我便为祢的平安而燃烧。\par

% structure: p0068 source=h2 target=subsection reason=relative-heading-rank; base=h2

\subsection{第二十八章 论人生的不幸。}

39. 当我以整个的我与祢结合时，我便不再有痛苦和劳碌；我的生命将是一个真实的生命，完全充满祢。如今，因为祢所充满的，正是祢所提拔的人，我因未充满祢而成为自己的重负。悲中的乐与乐中的悲相争；我不知道胜利在何方。我有祸了！主，可怜我！我邪恶的悲愁与我善良的喜乐相争；我不知道胜利在何方。我有祸了！主，可怜我！我有祸了！看，我不掩藏我的创伤；祢是医生，我是病人；祢慈悲，我可怜。人生在世，岂非诱惑吗？谁愿意烦恼和困难呢？祢命令我们忍受它们，而不是喜爱它们。因为没有人喜爱他所忍受的，尽管他可能喜爱忍受。因为他虽然欢喜忍受，却宁愿没有可忍受的。在逆境中，我渴望顺境；在顺境中，我畏惧逆境。那么，在这两者之间，何处人生不是诱惑呢？尘世的荣华，由于畏惧不幸和欢乐的败坏，实是一而再的祸患！尘世的逆境，由于渴望顺境，且因逆境本身是难忍之事，使忍耐覆舟，实是一而再，三而四的祸患！人生在世，岂非持久的诱惑吗？\par

% structure: p0070 source=h2 target=subsection reason=relative-heading-rank; base=h2

\subsection{第二十九章 全部希望在于天主的仁慈。}

40. 我的全部希望仅仅在于祢极大的仁慈。赐予祢所命令的，并随祢的意愿命令吧。祢将\OriginalTerm{节德}{continency}加于我们，正如有人所说：\Quote{我明白，除非天主将她赐给我，我便不能获得她……知道她是谁的恩赐，这已是一桩智慧的事了。}，\BibleRef{智 8:21} 因为借着节德，我们得以收敛，合一，从此不再散乱于众多之中。因为谁若爱任何事物与祢一起，却不是为祢而爱，他爱祢便嫌太少。爱啊，祢常燃烧，永不熄灭！\OriginalTerm{爱德}{charity}啊，我的天主，燃着我吧！祢命令节德；赐予祢所命令的，并随祢的意愿命令吧。\par

% structure: p0072 source=h2 target=subsection reason=relative-heading-rank; base=h2

\subsection{第三十章 论他愿去除梦中的邪像。}

41. 的确，祢命令我节制\BibleQuote{肉身的贪欲，眼目的贪欲，以及人生的骄奢}。祢命令我戒绝姘居；至于婚姻本身，祢劝谕了比许可更好的事。因为祢恩赐了这事，它便成就了，且在我成为祢圣事的分施者之前。但在我记忆里——我已谈得很多——仍存留着那些习惯已深植其中的事物的影像；当我清醒时，这些影像虽无力，却涌进我的思想；而在睡眠中，它们不仅带来快感，甚至博得同意，几乎与事实无异。是的，那影像的幻相在我的灵魂与肉身中如此得势，以致于在梦中，虚假的说服我去做那真实清醒时不能做到的事。主，我的天主，那时我不是我自己吗？然而，在我自己与自己之间有着如此大的差别，就在我从清醒转入睡眠，或从睡眠返回清醒的那一刻！那么，清醒时抵抗这些诱惑的理性在哪里呢？即使那些事物强加于它，我仍不为所动。难道理性是随眼睛一同关闭了吗？还是随肉身的感官一同入睡了呢？然而，何以即使在梦中我们也常常抵抗，牢记我们的志向，极其贞洁地持守它，不认同这些诱惑呢？还有如此大的差别：若情况相反，醒来时我们恢复了良心的平安；就是透过这差异，我们发现那原不是我们做的，虽然我们仍因这事以某种方式在我们内发生而感到难过。\par

42. 全能的天主，祢的手不能治愈我灵魂的一切疾病，并用祢更丰厚的恩宠熄灭我睡眠中的淫欲冲动吗？主，祢必在我内日益增加祢的恩赐，好使我的灵魂追随我到祢面前，摆脱贪欲的粘鸟胶；使灵魂不再自相反抗，即使在梦中，不仅不因感官的影像而做出那些败坏的丑行，乃至肉身受玷污，而且连认同它们也不发生。因为对于那位\BibleQuote{能成就一切，远超我们所求所想的}\BibleRef{弗3:20}的全能者来说，使这样的影响——甚至连一个记号就能约束的轻微影响——也不给予睡梦中贞洁的意愿丝毫快感，并非难事；而且不但在今生，就在我现今的年龄也能做到。至于我在这类邪恶中仍是如何，我已向我慈善的主坦承了；对于祢所赐予我的，我战战兢兢地欢欣；对于我仍不完全之处，我自悲自叹；并信赖祢必在我内完成祢的仁慈，直至圆满的平安，到那时，我内里的与外在的都将在死亡被胜利吞灭了时，与祢共享平安。\BibleRef{格前15:54}\par

% structure: p0075 source=h2 target=subsection reason=relative-heading-rank; base=h2

\subsection{第31章 论及肉身的贪欲的诱惑，他先慨叹饮食的贪欲。}

43. 今日还有另一种恶，我愿它\BibleQuote{足够一天受的}。\BibleRef{玛6:34}因为我们藉着饮食补偿身体的日常损耗，直到祢毁灭了食物与肚腹，那时祢将以奇妙的饱饫灭绝我的需缺，并给这可朽坏的穿上永恒的不朽。\BibleRef{格前15:54}但现在，必需对我成了甘饴，我却与这甘饴作战，免得受它奴役；我藉着禁食进行每日的战斗，常常\BibleQuote{使我的身体为奴}\BibleRef{格前9:27}，而我的痛苦却被快乐驱走了。因为饥渴在某种意义上是痛苦；它们像热病一样消耗、摧毁我们，除非营养的药物来救治。而这药物由于祢恩赐的慰藉而使地、水、空气服务我们的脆弱，于是我们的患难竟被称为快乐。\par

44. 祢已教了我这么多，叫我应当以食物为药物。但当我要从缺乏的不适过渡到饱足的安宁时，就在这过渡之中，贪欲的罗网潜伏着等待我。因为这过渡本身就是快乐，此外别无它路可达那必需强迫我们前往的境地。虽然健康是饮食的理由，却有一个危险的乐趣如同婢女般依附上来，它多半试图抢先，好叫我为了它的缘故而做我自称或欲为健康而做的事。两者并没有同样的界限：因为对健康足够的，对快乐来说却太少。时常令人疑惑的是：究竟是身体必要的照料仍在要求营养，还是感官贪欲的罗网在提供服务。在这不确定中，我不幸的灵魂竟沾沾自喜，并由此准备托辞作为防卫，庆幸那为健康的节制所足够的分量尚未显露，于是可以假健康之名掩盖快乐的勾当。这些诱惑我天天努力抵抗，我呼求祢的右手救助，并将我的激动呈于祢前，因为我对这事至今尚无决断。\par

45. 我听到我的天主的声音命令说：\BibleQuote{你们的心不要为贪饕和沉醉所累。}\BibleRef{路 21:34} \Quote{沉醉}离我很远；祢必施以怜悯，使它不要靠近我。但是\Quote{贪饕}有时却潜至祢的仆人身上；祢必施以怜悯，使它远离我。因为除非祢赐予，无人能节制。\BibleRef{智 8:21} 我们祈求的许多事物，祢都赐予我们；我们领受的任何美善，即使是在祈求之前就已领受的，也是从祢手中领受的，而且我们后来能知道这是从祢领受的。我从未成为醉汉，但我认识醉汉因祢而成为清醒的人。因此，是祢的作为，使那些从未如此的人不会如此，也是因祢的作为，使那些曾经如此的人不再永远如此；更是因祢的作为，使他们都知道这能力是从谁而来。我听到祢的另一声音说：\BibleQuote{不要随从你的私欲，却要抑制你的嗜好。}\BibleRef{德 18:30} 并且我也因祢的恩宠听到了另一句我十分喜爱的话：\BibleQuote{我们即使吃，也无益；我们即使不吃，也无损。}\BibleRef{格前 8:8} 这就是说，那一样不会使我富足，另一样也不会使我困苦。我还听到另一个声音：\BibleQuote{因为我已经学会了，在所处的任何环境中都可以知足。我懂得受贫贱，也懂得享富裕……我赖加强我力量的那位，能应付一切。}\BibleRef{斐 4:11} 看，一位天上营帐的战士——不像我们是尘土。但是主啊，请记起，\BibleQuote{我们原是尘土，}祢由尘土造了人；\BibleRef{创 3:19} 他\BibleQuote{曾是失而复得的。}\BibleRef{路 15:32} 他也不能凭自己的能力做到这事，因为那位我如此深爱、藉着祢的\OriginalTerm{灵感}{afflatus}说出这番话的人，也是出自同样的尘土。他说：\BibleQuote{我赖加强我力量的那位，能应付一切。}\BibleRef{斐 4:13} 求祢加强我，使我能够做到。求祢赐下祢所命令的事，并命令祢所愿的事。他承认这一切都是领受的，当他夸耀时，他就因主而夸耀。\BibleRef{格前 1:31} 我听到另一个人恳求说——\BibleQuote{求祢除掉我，}他说，\BibleQuote{肚腹的贪欲；}\BibleRef{德 23:6} 由此显出，我的神圣的天主，当祢命令的事被执行时，就是祢的恩赐。\par

46. 祢教导了我，善父啊，\BibleQuote{为洁净的人一切都是洁净的；}\BibleRef{铎 1:15} 但\BibleQuote{谁若吃了而心中起疑，就有罪了；}\BibleRef{罗 14:20} \BibleQuote{凡受造的都是好的，若以感恩的心受之，一无所弃；}\BibleRef{弟前 4:4} 且\BibleQuote{食物不能使我们见悦于天主；}\BibleRef{格前 8:8} 人也不可\BibleQuote{在任何饮食上判断我们；}\BibleRef{哥 2:16} 那吃的，不要轻视不吃的；那不吃的，也不要判断吃的。\BibleRef{罗 13:23} 这些我都学会了，感谢赞美归于祢，我的天主，我的师傅，祢敲我的耳，光照我的心；求祢救我脱离一切诱惑。我所怕的不是食物的不洁，而是贪欲的不洁。我知道\Person{诺厄}获准吃一切可作食物的肉；\BibleRef{创 9:3} \Person{厄里亚}也靠肉食维持；\BibleRef{列上 17:6} \Person{若翰}，虽被赋予奇妙的节制，食那活物（就是蝗虫\BibleRef{玛 3:4}）也未受沾染。我也知道\Person{厄撒乌}因贪求扁豆而被欺骗，\BibleRef{创 25:34} \Person{达味}因想喝水而自责，\BibleRef{撒下 23:15} 我们的君王受试探，不是用肉，而是用饼。\BibleRef{玛 4:3} 旷野中的百姓因此也当受责，并非因为他们想吃肉，而是因为他们为了饮食之求而抱怨主。\EditorialNote{户 11}\par

47. 因此，我置身于这些诱惑之中，每日与对饮食的贪求搏斗。因为这并非那种我可以下决心一劳永逸地断绝、之后不再触碰的事物，像我能对姘居所做的那样。所以，喉舌的辔头必须保持在不松不紧的适中位置。主啊，有谁在某种程度上不曾被带出需求的界限之外呢？不论是谁，他若能做到，必是伟大的；他必彰扬祢的名。但我不是这样的人，\BibleQuote{因为我是个罪人。}\BibleRef{路 5:8} 然而我也彰扬祢的名；那\BibleQuote{战胜了世界}的\BibleRef{若 16:33}，正在为我代求，向祢转达我的罪过，\BibleRef{罗 8:34} 并将我算作祂身体中\BibleQuote{软弱的肢体}，\BibleRef{格前 12:22} 因为祢的眼看见了其中尚未成全的部分；在祢的册上一切都要被记录。\par

% structure: p0081 source=h2 target=subsection reason=relative-heading-rank; base=h2

\subsection{第三十二章 关于香味的诱惑，较易克胜。}

48. 对于气味的诱惑，我并不怎么困扰。不在时，我不追求；在时，我不拒绝；并准备随时没有它们。至少我自己看来如此；也许我是在欺骗自己。因为那也是可悲的黑暗，我内在的能力隐而不显，以致我的心灵反省自身的能力时，不敢轻易信任自己；因为大多数情况下，已在我内的，除非经验将其揭示，都隐藏着。此生被称为完全的诱惑，没有一个人该觉得安全，因为那能从恶改善的人，也可能从善变恶。我们唯一的希望、唯一的信心、唯一确实的应许，就是祢的仁慈。\par

% structure: p0083 source=h2 target=subsection reason=relative-heading-rank; base=h2

\subsection{第三十三章 他克胜了耳朵的愉悦，虽然在教会中常因歌唱而喜乐，过于所唱的内容。}

49. 听觉的欢愉曾更有力地诱骗并征服了我，但祢解开了束缚，释放了我。如今，当那些乐曲——祢的言语赋予其灵魂——以甜美而训练有素的歌声唱出时，我能在其中稍得安息；但并非依恋不舍，而是可以随意抽身。然而，这些乐曲连同赋予它们生命的言语，为了得以进入我心，力求在我内心取得某种尊位；我几乎无法为它们安排合适的位子。有时，我觉得自己似乎给了它们过多的尊重，因为我发觉，当神圣的词句以这样的方式歌唱时，比不唱时更能使我们的心灵虔诚而热切地燃起虔敬之火；而我们精神的每一种情感，按其各自的性质，在声音与歌唱中都有其相应的尺度，它们因某种我所不知的隐秘关联而受到激发。然而，肉身的享乐——心灵绝不应交付给它而变得软弱——常常欺骗我，因为感官并不如此服从理性，以至于耐心地跟随它；相反，感官仅因理性的缘故而得进入后，便力图奔在理性之前，充当其领导者。于是，在这些事上，我不知不觉地犯了罪，但事后我才知晓。\par

50. 有时，我又极力躲避这同样的骗局，却因过分拘谨而犯错；有时甚至想要将所有据以歌唱\Person{达味}的\BibleBook{}{圣咏集}的优美曲调从自己的耳中，也从教会本身的耳中逐出；我记得人们常向我提及\Place{亚历山大}城的主教\Person{亚大纳削}的做法，他要求诵读圣咏的人以极轻微的声音读出，使其更近于说话而非歌唱，那方式在我看来似乎更稳妥。尽管如此，每当我回忆起初恢复信仰时，在祢教会的歌声中所流的眼泪，又想到即使在如今，当那些圣咏以清澈而巧妙调谐的嗓音歌唱时，我感动的并非歌唱本身，而是所唱的内容，我便承认这种惯例具有极大的益处。我就这样在危殆的愉悦与可靠的健全之间摇摆不定；虽然我对此不发表不可更改的意见，但还是更倾向于赞同在教会中使用歌唱，好使软弱的灵魂藉着听觉的愉悦而被激发起虔诚的心境。然而，当歌唱比所唱的内容更打动我时，我便承认自己犯了大罪，那时我宁愿没有听到那歌唱。看哪，我现在的处境！你们这些善于约束内心情感，从而结出善果的人，请与我一同哭泣，为我哭泣吧。至于那些并不如此行事的人，这些事与你们无关。但是，祢，上主我的天主，求祢侧耳倾听，俯视垂顾，怜悯我，医治我——在祢眼中，我已成了自己的谜；\Quote{这便是我的软弱。}\par

% structure: p0086 source=h2 target=subsection reason=relative-heading-rank; base=h2

\subsection{第34章 论眼睛极其危险的诱惑；因形体之美，当赞颂造物主天主。}

51. 还有我肉身眼目的欢愉，我当在祢圣殿的耳中——那些弟兄般的虔诚之耳——聆听时，为此而忏悔；并以此结束那\BibleQuote{肉身的贪欲}的诱惑 \BibleRef{若一 2:16}，这些诱惑仍在攻击我，令我呻吟切望得穿上自天而来的住所 \BibleRef{格后 5:2}。眼目喜悦于美丽多样的形态，以及明亮怡人的色彩。请不要让这些占据我的灵魂；宁愿让天主占据，是祂造了这一切，确实\BibleQuote{很好} \BibleRef{创 1:31}；然而祂是我的善，而非这些。这些在白天我醒着时触动我；它们从不像音乐的旋律那样，有时能在寂静中完全止息，让我得到安宁。因为那色彩之后——光，洒遍我们所见的一切，无论白天我身处何处，它以万千形态从我眼前流过，当我忙于其他事而不加留意时，它便抚慰我。它如此强烈地潜入，以至于一旦突然消失，人便渴慕追寻，若久而不见，便使心灵忧郁。\par

52. 啊，祢是那光，是\Person{多俾亚}所见的光 \EditorialNote{多俾亚传 4}，那时他双目失明，却教给儿子生命的道路；他自己以爱德之脚前行，从未迷路。或者，那光是\Person{依撒格}所见的光，那时他肉身的\BibleQuote{眼目昏花，不能看见} \BibleRef{创 27:1} 因为年老的缘故；他蒙准许，不是在知情下祝福儿子们，而是在祝福中认识了他们。或者，那光是\Person{雅各伯}所见的光，他也因高龄而失明，却以受了光照的心，在儿子的身上，照亮了由他们所预象的未来民族的支派；又神秘地将手交错，按在\Person{若瑟}的两个孙儿头上，并非如他们父亲从外表看那样纠正他，而是如同自己分辨了他们一般 \BibleRef{创 48:13}。这就是那光，独一无二，凡看见并爱慕这光的人，都合而为一。但我之前所说的那物质之光，却以诱人而致命的甘甜，为这世界的盲目爱恋者调味人生。然而，那些知道因此而赞美祢的人，\Quote{啊，天主，世界的伟大建筑师}，便在赞颂祢的诗歌中将它高举，而不在昏睡中沉溺于它。我愿成为这样的人。我抵抗眼目的诱惑，免得我在祢道路上行走的脚被缠住；我向祢抬起无形的眼睛，恳求祢乐意 \BibleQuote{将我的脚从网罗中拔出}\BibleRef{咏 25:15}。祢不断将它们拔出，因为它们常被网住。祢从未停止拔出，而我，却始终困在四周遍设的网罗中；因为祢——\BibleQuote{那保护以色列的，不打盹也不睡觉}\BibleRef{咏 121:4}。\par

53. 人为了迷惑眼目，加上了多少由各种技艺和制造所出的数不清的东西——在我们的衣服、鞋子、器皿和每一种制品中，也在图画和各种形像中——这些都远远超出了必要和适度的用途及圣洁的象征；他们在外追求自己所造的东西，内心却抛弃了那位造他们的主，甚至毁灭了他们自身受造的本性！但我，我的天主，我的喜乐，我要为此向祢高唱赞美诗，向我的\OriginalTerm{圣化者}{Sanctifier}献上赞颂之祭，因为这些美丽的模型，通过人的灵魂传递到他们艺术家的手中，都是从超越我们灵魂的那美而来的——我的灵魂日夜为此美叹息。至于那些外在美的制造者和追随者，他们从其中得到了赞赏它的方式，却得不到使用它的方式。他们虽看不见祂，祂仍在那里，为使他们不致迷途，保持他们的力量为祢所用，而不致消散在逸乐的懒散中。至于我，我虽如此说，也如此领悟，仍被这些美妙的事物窒碍行程；但是主啊，祢拯救我，祢拯救我；\BibleQuote{因为祢的仁慈常在我眼前。}因为我可怜地被捉住，祢却慈悲地拯救我；有时我毫无察觉，因为我犹豫地走近它们；有时则痛苦不堪，因为我被它们紧紧束缚。\par

% structure: p0090 source=h2 target=subsection reason=relative-heading-rank; base=h2

\subsection{第三十五章 另一种诱惑是猎奇心，它由眼目的贪欲所激起}

54. 此外还有另一种形式的诱惑，它的危险更为复杂。因为除了那在于满足一切感官和快感的\OriginalTerm{肉身的贪欲}{concupiscence of the flesh}，其奴仆们在其中\BibleQuote{远离祢而丧亡}，还有一种属于灵魂的、通过身体的同样感官而来的、虚幻而好奇的渴望，披着知识和学问的外衣，其目的不在于在肉身中获得快乐，而在于通过肉身进行试探。这种渴望，由于源于对知识的嗜求，而在获取知识上视觉是感官中的主宰，因此在神圣的语言中被称为\OriginalTerm{眼目的贪欲}{the lust of the eyes}。\BibleRef{若一 2:16}因为看见严格说来属于眼睛；但我们将这个词也用于其他感官，当我们在求知中运用它们时。因为我们不说，听它如何发亮，闻它如何闪烁，尝它如何照耀，或摸它如何闪光，因为这一切都说成是看见。然而我们不仅说，看它如何照耀，这只有眼睛能察觉；而且也说，看它如何发声，看它如何闻起来，看它如何尝起来，看它如何坚硬。因此，如前所述，感官的一般经验被称为\OriginalTerm{眼目的贪欲}{the lust of the eyes}，因为视觉的功能，眼睛在其中拥有优越地位，其他感官以类似的方式加以借用，每当它们寻求任何知识时。\par

55. 但借此可以更清楚地区分，何时是感官追求快乐，何时是追求猎奇。因为快乐追求的是美丽、悦耳、芬芳、可口、柔和的事物；但猎奇，为了试探的缘故，寻求的是与这些相反的东西——不是有意去经受不适，而是出于试探和认识它们的冲动。因为在看到一具撕裂的尸体时，有什么快乐可言呢？那会令你颤栗。但若它横在那里，我们便蜂拥而至，使自己伤心，面色苍白。即使在梦中他们也害怕见到它。就好像在醒着时，有人强迫他们去看，或有任何关于它美丽的传言吸引了他们一样！其他感官也是如此，逐一追寻起来甚是繁冗。出于这种猎奇之症，所有那些奇观都在剧场中展出。由此我们进而去探究自然界的隐秘力量（这超出了我们的目的），知道这些对人并无益处，人们所求的不过是知道而已。由此，出于同一扭曲认知的目的，我们去求问魔术。由此，甚至在宗教本身中，天主也受到试探，当人们急切地向祂要求征兆和奇迹时——不是为任何拯救的目的，而是只想试探而已。\par

56. 在这如此广阔、布满陷阱和危险的荒原中，看，我已砍掉了许多，将它们从我心里除去，一如祢，我救恩的天主，使我能够做的。然而，既然在日常生活中有这么多此类事物嗡嗡作响，我几时才敢说我未曾被这类事物吸引去注目，或在我内产生虚妄的牵挂呢？诚然，剧场不再使我沉迷，我也不再在意星辰的运行，我的灵魂也从未求问过亡魂；一切亵渎神圣的誓言我都深恶痛绝。主我的天主，我应欠祢一切谦卑和纯朴的事奉，敌人用多么狡猾的暗示来煽动我向祢要求某些征兆！但凭我们的君王，并凭我们纯洁国土贞洁之城\Place{耶路撒冷}，我恳求祢，既然凡同意此类想法的都已远离我，也愿它永远离我越来越远。然而，当我为任何人的救恩恳求祢时，我的目的却全然不同，祢这行事全凭己意者，赐予我并乐意赐予我\BibleQuote{跟随}祢。\BibleRef{若 21:22}\par

57. 尽管如此，我们在多少极细微、极可轻视的事上，每天受好奇心的诱惑，而谁能数算我们屈服了多少次？多少次，当人们讲说闲谈时，我们起初容忍，免得开罪软弱的人；随后渐渐地，我们竟乐于听受！如今我不去竞技场看狗追兔子；但若偶然在田野间经过这样的追逐，它或许竟使我离开某种严肃的思想，随它而去——不是扭转骑乘的身体，而是心灵的方向。若非祢，借着向我显示我的软弱，迅速警醒我，或通过那景象本身借某种反思举心向祢，或完全鄙视它、越过它，我这虚妄的人便被它吞没了。当我坐在家中，一只壁虎捕蝇，或蜘蛛缠住撞入它网中的飞虫，又如何时常吸引我呢？难道因为这些都是微小生物，好奇心便不一样吗？我从它们进而赞美祢——万有的奇妙创造者与安排者；但最初吸引我注意的并非此意。迅速站起是一回事，不跌倒则是另一回事，而我的一生充满了这类事；我唯一的希望就是祢极大的仁慈。因为当我们的这颗心成了这类事物的接收器，并负载着大量汹涌的虚荣，我们的祈祷便时常因此中断扰乱；正当我们在祢面前，将心中的呼声转向祢的双耳时，这重大的事便被一些不知所由的闲思所打断。\par

% structure: p0095 source=h2 target=subsection reason=relative-heading-rank; base=h2

\subsection{第三十六章 第三种是取悦于人而非取悦于天主的\Quote{骄傲}}

58. 那么，我们也将这事算为应轻看的事物吗？或者除祢完全的慈爱之外，还有什么能使我们恢复希望呢？因为祢既已开始改变我们。祢知道祢已经改变我到了何种程度，祢先医治了我辩护己身的欲望，好这样赦免我其余一切的\Quote{邪恶}，医治我一切的\Quote{疾病}，从腐朽中救赎我的生命，以\Quote{慈爱和怜悯}为我的冠冕，并以\Quote{美物}满足我的心愿；祢以祢的敬畏抑制我的骄傲，并使我的颈项服在祢的\Quote{轭}下。现在我背起它，它于我乃是\Quote{轻省}的\BibleRef{玛 11:30}，因为祢曾如此应许，并使它如此，事实上它当初就是如此，只是我尚不知道，当时我还害怕肩负它。但是，主啊，——唯有祢才是毫无骄傲而统治的主，因为祢是唯一的真主，没有主——这第三种诱惑曾离开我吗？或者它能在今生离开我吗？\par

59. 愿意被人畏惧和受人爱戴，除了为此获得一种并非欢乐的欢乐之外别无目的，这便是一种可怜的生活，一种可耻的夸耀。尤其因此，我们不热爱祢，也不虔敬畏惧祢。因此祢拒绝骄傲人，却赐恩宠给谦逊人；\BibleRef{雅 4:6}祢向世界的野心图谋发出雷霆，\Quote{山岳的根基}也战栗。因为如今某些人类社会的职责使人不得不被人爱戴和畏惧，我们真福的仇敌便紧紧压迫我们，到处撒布\Quote{做得好！做得好！}的陷阱；使我们在热切获取它们时，不经意地陷入其中，将我们的欢乐与祢的真理分离，而系于人类的欺诳；乐于受人爱戴和畏惧，不是为了祢，而是取代祢，这样，他使那些与他相似的人，不是在爱的和谐中，而是在刑罚的共融中，成为他的从属；他曾企图将自己的宝座高举于北方，\BibleRef{依 14:13}让那些黑暗冷酷者服事他，以邪恶扭曲的方式模仿祢。但是，主啊，请看，我们是祢的\BibleQuote{小小羊群}；\BibleRef{路 12:32}求祢拥有我们，展开祢的双翼遮盖我们，让我们投奔其下。求祢作我们的荣耀；让我们为祢的缘故受人爱戴，并愿祢的圣言在我们内受人敬畏。谁若在祢谴责时渴望受人们赞美，在祢审判时将不受人们辩护；当祢定罪时，也不会被拯救。但是，当不是罪人因灵魂的欲望受称赞，也不是作不义者受祝福，而是人因祢赐予他的某种恩赐受赞美，且他对自己的赞美比对他拥有那受赞美的恩赐更为喜悦时，这样的人受赞美时，祢却在谴责。其实，称赞他的人，比那受称赞的人更为善。因为在前者，天主在人内的恩赐使他喜悦；而后者更喜悦人的恩赐胜于天主的恩赐。\par

% structure: p0098 source=h2 target=subsection reason=relative-heading-rank; base=h2

\subsection{第三十七章 他深受贪图赞美之心的驱使}

60. 主，我们每日受这些诱惑的试探；是的，我们不停地受试探。我们每日的\Quote{炉}就是人的口舌。在这方面祢也命令我们要节制。求祢赐下祢所命令的，并命令祢所愿意的。关于此事，祢知晓我内心的呻吟和我眼泪的河流\BibleRef{哀 3:48}。因为我无法确知我距离这瘟疫有多洁净，我极其恐惧我的\Quote{隐密的过错}，这些过错祢的眼目看见，我的却看不见。因为在其他种类的诱惑中，我有某种能力省察自己；但在此事上，几乎一无所有。因为，对于肉体的享乐和无聊的好奇心，当我自愿或因它们不在眼前而缺少它们时，我看出我能如何节制我的心灵；那时我自问，没有它们对我而言是多么地更少烦恼或更多烦恼。至于财富，人们寻求它是为了服务于这三种\Quote{私欲偏情}中的一种、两种或全部。如果心灵不能清楚地看出当它拥有财富时是否轻视财富，就可以将财富丢在一旁，好证明自己。但若我们想要试验自己能否忍受没有赞美，难道必须过着邪恶的生活，并且那样卑鄙放肆，以致所有认识我们的人都憎恶我们？还有什么比这更大的疯狂可言？但赞美既是善生和善工的常伴，也当是必备之伴，我们不当缺少它的伴随，正如不当缺少善生本身。但除非一件事物不在眼前，我不知我是否会因缺少它而满足或烦恼。\par

61. 那么，主，在这种诱惑中，我向祢承认什么呢？无非是我喜欢赞美，但更喜欢真理本身过于赞美吗？因为如果让我选择，我宁愿在发疯或在一切事上迷失时受众人赞美，还是在坚持并确信真理时受众人责备，我知道我会选择哪一个。然而我不愿意，因我拥有的任何善，别人的赞许竟增加我的喜乐。但我承认它确实会增加我的喜乐，而且不仅如此，贬斥也会减少我的喜乐。当我因自己的这种不幸而不安时，一个借口便呈现在我面前，其价值祢，天主，知道，因为它使我犹疑不定。因为祢不仅命令我们节制，即当从何事上收回我们的爱，也命令我们义德，即当将我们的爱投向何事上，并且祢愿意我们不仅爱祢，也要爱我们的近人---所以，当我因明智的赞美而喜悦时，我自觉乃是因近人的进步或向善而喜悦；同样，当我听到他贬斥他所不理解的事物或善事时，我便为他里面的恶而难过。因为我有时因自己的赞美而忧伤，或是当我在自己里面所不悦的事物竟在我身上受到赞美时，或是即使微小不足道的善也被过分看重时。但是，再说，我怎能知道我之所以有此感受，是因为我不愿意那赞美我的人对我的心意与我不同---并非出于对他的关心而受感动，而是因为那些令我喜悦的我之内在的善事，当它们也令他人喜悦时，就更令我喜悦？因为在某种意义上，当我对自己的看法不受赞美时，我便未受赞美；因为要么我里面不悦的事物受到赞美，要么那些不太令我喜悦的事物反倒更受赞美。那么，在这事上我难道对自己犹疑不定吗？\par

62. 看，真理啊，我在祢里面看见，我不应为了自己而因自己的赞美受感动，却应为了近人的益处。事实是否果真如此，我实在不知。因为关于这事，我对自己所知的远少于祢。我现在恳求祢，我的天主，也给我启示我自己，好让我能向我的弟兄们，那些将为我祈祷的人，承认在我内发现的软弱。让我再更勤奋地省察自己。若在我自己的赞美上，我是出于关心近人而受感动，那么为什么当别人受到不公正的贬斥时，我的感动不及当我自己受到时呢？为什么我对加于自身的侮辱比对他人在我面前所遭受的同等不义的侮辱更为恼怒？我连这事也不明白吗？还是我仍在自欺\BibleRef{迦 6:3}，且在祢面前心里和口舌上没有\Quote{真理}\BibleRef{若一 1:8}？主，求祢使这种疯狂远离我，免得我的口成为\OriginalTerm{罪人的油}{oil of sinners}，来膏我的头。\par

% structure: p0102 source=h2 target=subsection reason=relative-heading-rank; base=h2

\subsection{第38章 虚荣是最大的危险。}

63. \Quote{我贫穷可怜}，但我私下呻吟时，对自己不满，寻求祢的仁慈，直到我内所欠缺的得以更新并完成，达到那骄傲的眼目所不知的平安，这样倒更好。然而，从口中发出的话语，以及人所共知的行为，却因爱慕赞美而有极危险的诱惑，这种诱惑为了建立某种自己的优越，收集恳求来的赞赏。即使我内心因此责备自己，它仍以受责备为缘由而诱惑人；人往往因轻蔑虚荣而归虚荣于自己的轻蔑更甚；因此他归荣于自己的不再是轻蔑虚荣了，因为他内心归荣时并未真正轻蔑它。\par

% structure: p0104 source=h2 target=subsection reason=relative-heading-rank; base=h2

\subsection{第39章 关于那些自悦而不悦天主的人之恶。}

64. 内在也有，内在有另一种恶，来自同一种诱惑；那些在自身内自悦的人，虽然不取悦于人，或使人生厌，或企图取悦于他人，他们却因此成为空虚。但在自悦中，他们大大地不悦于祢，他们不仅以非善为善而自以为乐，而且以祢的善视为己有；或即使视为祢的，却仿佛出于自己的功绩；或即使视为出于祢的恩宠，却没有友好的欣喜，反而嫉妒那恩宠临于他人。在所有这些及类似的危险与劳苦中，祢洞悉我内心的战栗，我更觉得我的创伤是由祢治愈，而不是由我造成的。\par

% structure: p0106 source=h2 target=subsection reason=relative-heading-rank; base=h2

\subsection{第40章 灵魂唯一安息之所乃在于天主。}

65. 真理啊，祢何处不曾伴随我，教导我当避何事，当慕何物？当我将力所能及领会到的月下诸事呈献给祢，并向祢请教时，祢就教导我。我用外在的感官，尽我所能观视世界，留意身体从我所得的生命，以及我这感官本身。由此我向内进入记忆的深处——那些多种多样的殿堂，奇妙地充满了繁多丰裕；我思量而惊惧，发现所有这些事离了祢便一无所辨，也找不到其中有哪一个是祢。我并非发现这些事的那一位——我遍察一切，辛苦分辨，按各物的尊卑衡量，凭感官的报告接受一些，对自己感觉与己混杂之事加以询问，区别并数点报告者本身，在我记忆的巨大库房里研究某些事，收藏另一些，取出别些。做这事时，我本身（即我做此事的能力）并非是我，也不是祢，因为祢是那永不熄灭的光，我就这一切向祢请益，探究它们是否确实，价值几何；我听到祢教导并命令我。这事我屡屡行之；这是我甘甜之所在，且只要能从必需事务中稍得解脱，我便投身于此乐中。然而，当我向祢咨询而遍观这一切时，我找不到一个灵魂安稳之处，除非在祢内，在祢内我散佚的肢体方能团聚，而我无一部分离开祢。有时祢引领我进入内心一种至罕的情感，一种无法言喻的甜蜜，若这事在我内成全，我不知生命能达到何种地步。但我因这些可怜的重担（\BibleRef{希 12:1}）而再次陷入这些事中，被旧习惯所吞噬，被困住，虽大感忧愁，却仍被紧束。习惯的重压如此沉重地将我们压伏。如此，我能如此，却不愿；彼处我愿，却无能——两边皆悲。\par

% structure: p0108 source=h2 target=subsection reason=relative-heading-rank; base=h2

\subsection{第四十一 战胜三重欲望后，他获得救恩。}

66. 我便这样反思我诸罪的劳苦，即那三重的\Quote{贪欲}，并呼求祢的右手扶助。因为我以受伤之心看祢的光辉，被击退时呼喊：\BibleQuote{谁能达至？}\BibleQuote{我从祢眼前被隔绝。}祢是统御万有的真理，但我却因贪吝，不愿失去祢，却愿与祢同得谎言；正如无人愿说虚假，而自己竟不知真理。这样我便失去了祢，因为祢不屑与谎言同享。\par

% structure: p0110 source=h2 target=subsection reason=relative-heading-rank; base=h2

\subsection{第四十二 许多人如何寻求中保。}

67. 我能找谁使我与祢和好？要我哀求天使吗？用何祈祷？用何圣事？许多人勉力归向祢，却自身不能，据我所闻，曾尝试此法，竟陷入对怪异景象的渴望，并被认为配受欺蒙。因为他们高高在上，以学识之傲寻求祢，挺身向前而不捶胸，因此以心之相应，招来空中权势者（\BibleRef{弗 2:2}），就是那骄傲的同谋与同伴；因魔力之能，他们受了欺蒙，寻求一位可借以得洁净的中保，却一无所得。原来那是魔鬼，将自己伪装成光明的天使（\BibleRef{格后 11:14}）。他大大引诱骄傲的肉身，因他没有血肉之体。他们是必死的，有罪的；但祢，主啊，他们傲然欲与祢和好，祢却是不死不灭的，无罪的。然而，天主与人之间的中保，当有与天主相似之处，也有与人相似之处；免得若在两方面都像人，就远离了天主；若在两方面都像天主，就远离了人，这样就不能作中保。那欺骗人的中保，按照祢隐秘的判断，骄傲应受欺蒙，他与人有一样共同点，即罪；另一样他假作与天主同有，不以血肉之死亡为衣，便夸耀自己是不死的。可是，既然\BibleQuote{罪的薪饷就是死亡}（\BibleRef{罗 6:23}），他与人类另有一共同，即他们都要被判处死亡。\par

% structure: p0112 source=h2 target=subsection reason=relative-heading-rank; base=h2

\subsection{第四十三 耶稣基督，同时是天主又是人，是真确最有效的中保。}

68. 但那真正的中保，祢以隐秘的慈爱向谦卑者指出，并派遣他来，好叫他们借祂的榜样也学得同样的谦卑——那\BibleQuote{天主与人之间的中保，基督耶稣这人}（\BibleRef{弟前 2:5}）——出现于有死的罪人与不死的义者之间：与人同死，与天主同义；为叫义德之赏报是生命与平安，祂借与天主结合的义德，废除成义之罪人的死亡，这死亡祂甘愿与他们同当。因此，祂自古便被指示给圣者，好叫他们借信仰将来的受难，正如我们借信仰已过的受难，能以得救。因为作为人，祂是中保；但作为\OriginalTerm{圣言}{the Word}，祂并不居间，因为祂与天主平等，且是天主与天主同在，并与圣神同为一天主。\par

69. 好\Person{圣父}啊，祢怎样爱了我们！祢没有怜惜祢的独生子，反为了我们这些恶人把祂交出来。祢怎样爱了我们！祂不以自己与天主同等为强夺，却为了我们\BibleQuote{听命至死，且死在十字架上。} \BibleRef{斐 2:6} 唯有祂\BibleQuote{在死者中自由，} 才有能力舍掉自己的性命，也有能力再取回来；\BibleRef{若 10:18} 为了我们，祂对祢既是胜利者又是牺牲，并且正因为是牺牲才成为胜利者；为了我们，祂对祢既是司祭又是祭品，并且正因为是祭品才成为司祭；祂生于祢，并服事我们，使我们由奴隶成为祢的儿女。因此，我理当将我的希望牢牢寄于祂，相信祢会借着坐在祢右边、为我们转求的祂，医治我所有的疾病； \BibleRef{罗 8:34} 否则我将完全绝望。因为我的软弱众多且重大，的确众多且重大；但祢的医药更伟大。我们可能以为祢的圣言已远离了与人的结合，倘若祂没有\BibleQuote{成了血肉，寄居在我们中间。} \BibleRef{若 1:14}，我们便会对自己绝望。\par

70. 我因我的罪恶和悲惨的重负而惊恐不已，内心已决意并打算逃往旷野；但祢阻止了我，并坚固我，因此说：基督\BibleQuote{替众人死了，是为使活着的人不再为自己生活，而是为替他们死的那位生活。} \BibleRef{格后 5:15} 主啊，看，我将我的忧虑托付给祢，好使我生活，且\BibleQuote{得见祢的律法中的奇妙。} 祢深知我的愚钝和软弱；请教诲我，医治我。祢的独生子——\BibleQuote{在祂内隐藏着智慧和知识的一切宝藏}——已用祂的血救赎了我。莫让骄傲的人诽谤我，因为我常纪念我的赎价，且吃喝并分施；虽贫乏，却渴望与那些吃喝并且充饥的人一同，从祂那里得到饱饫，凡寻求主的人都要赞美主。\par

\SourceRecord{Translated by J.G. Pilkington.；Nicene and Post-Nicene Fathers, First Series；Vol. 1.；Edited by Philip Schaff.；Buffalo, NY: Christian Literature Publishing Co.,；1887.；<http://www.newadvent.org/fathers/110110.htm>.}

% structure: p0001 source=h1 target=section reason=page-title

\section{\WorkTitle{忏悔录}（卷十一）}

宣布了他忏悔的目的后，他向天主寻求对圣经的知识，并开始解释\BibleBook{创}{世纪}第一章第一节中关于创造世界的话语。在驳斥了鲁莽争辩者的问题——\Quote{天主在创造世界之前做了什么？}——之后，为能更好地战胜对手，他又对时间问题作了详尽的探讨。\par

% structure: p0003 source=h2 target=subsection reason=relative-heading-rank; base=h2

\subsection{第一章 通过忏悔，他愿激发自己和读者对天主的爱。}

1. 主啊，永恒属于祢，难道祢不知道我向祢所说的这些事吗？难道祢要等时候到来，才看见那时所发生的事吗？那么，我为什么还要在祢面前陈述这许多事呢？当然不是为了使祢借着我而知道它们，而是为了唤醒我自己和读者对祢的爱，使我们都能说：\Quote{主伟大，极堪受赞美。}我已经说过，而且还要说：我这样做，是出于对祢的爱的爱。因为我们祈祷时，真理却对我们说：\BibleQuote{你们的父，在你们求他以前，已知道你们需要什么。}\BibleRef{玛 6:8}因此，我们向祢表白我们的爱，向祢承认我们的不幸和祢对我们的怜悯，好使祢彻底解救我们——祢既已开始了这工程——使我们不再在自身内感到不幸，而在祢内得享真福；祢既召叫了我们，使我们成为神贫的人、温良的人、哀恸的人、饥渴慕义的人、怜悯人的人、心里洁净的人、缔造和平的人。\BibleRef{玛 5:3}看，我已向祢诉说了许多事，凡我能说的和愿意说的，我都说了，因为首先是祢愿意我向祢忏悔，我的主，我的天主，因为祢是美善的，祢的\Quote{慈爱永远长存。}\par

% structure: p0005 source=h2 target=subsection reason=relative-heading-rank; base=h2

\subsection{第二章 他恳求天主，愿通过圣经引领他走向真理。}

2. 但我何时才能用我的笔舌充分表达祢所有的劝勉、祢所有的恐吓、安慰和指引，祢怎样引领我去向祢的子民宣讲祢的圣言，并分施祢的圣事呢？即使我能按序说出这一切，时间的点滴对我也是宝贵的。我长久以来一直渴望默思祢的法律，并在默思中向祢承认我的知识和无知，祢启蒙的开端和我黑暗中残留的部分，直到脆弱被力量吞噬。我不愿让那些从身体休息的必要、心灵的关怀以及我们该为人效劳的服务——虽然我们并不欠债，却仍付出——中腾出来的时光，白白流走。\par

3. 我的主，我的天主，求祢俯听我的祈祷，愿祢的慈怜眷顾我的渴望，因为这渴望并非单为我自己燃烧，也愿有益于兄弟之爱；祢洞察我的心，确知如此。我要向祢献上我思想和口舌的侍奉；求祢赐予我可向祢呈献的礼物。因为\Quote{我困苦而又贫穷，}而祢对一切呼求祢的人却是富有的，\BibleRef{罗 10:12}祢无忧无虑地照顾我们。求祢割损我内外嘴唇上的一切鲁莽和一切谎诈。\BibleRef{出 6:12}让祢的圣经成为我纯洁的喜乐。让我不得在圣经上受骗，也不得用圣经骗人。主，求祢俯听，垂怜，我的主，我的天主，盲者的光明，弱者的力量；同样也是明眼人的光明，强者的力量，求祢垂听我的灵魂，听它\Quote{自深渊中}呼号。因为如果祢的耳朵不也临在于深渊中，我们往何处去呢？向何处呼号呢？\Quote{白昼属于祢，黑夜也属于祢。}在祢的指令下，时光飞逝。求祢赐予我们时间，好能在祢法律的隐秘事物中深思默想，不要向叩门的我们关闭。因为祢并非徒然定意写下这许多篇页的隐晦奥秘。那些丛林中也并非没有祢的鹿，它们投身其中，漫步、行走、觅食、躺卧、反刍。主，求祢成全我，将那些奥秘启示给我。看，祢的声音就是我的喜乐，祢的声音胜过一切欢愉的丰富。求祢赐给我我所爱的，因为我实在爱祢；而这爱也是祢赐予的。不要离弃祢自己的恩赐，也不要轻看祢口渴的青草。凡我在祢书本中发现的一切，我都要向祢忏悔，让我听见赞颂之声，让我畅饮祢，默思祢法律中的奇妙；从起初，祢创造天地之时，直到永久同祢在祢圣城的永恒国度里。\par

4. 主，求祢怜悯我，垂听我的渴望。因为我想这渴望并非来自尘世，也不是为金银宝石、华丽衣服、荣誉权势、肉身的快乐、身体所需、以及此尘世旅途的生命；这一切原是加给那些寻求祢的国和祢的义德的人的。\BibleRef{玛 6:33}看，我的主，我的天主，这就是我渴望的由来。不义的人曾给我讲述欢愉，但并非像祢的法律一样的欢愉，主啊。看，这就是我渴望的由来。圣父，求祢垂顾、查看、嘉纳；愿在祢仁慈的眼中，这蒙祢悦纳，使我在祢面前获得恩宠，当我叩门时，祢圣言的隐秘得以向我开启。我靠着我们的主耶稣基督，祢的圣子，\Quote{祢右手扶持的人，祢为自己所坚固的人子，}祢和我们的中保，恳求祢；祢通过祂寻找了我们，当时我们并没有寻找祢，但祢寻找我们，为使我们寻找祢，——祢的圣言，万物是借着祂造成的，\BibleRef{若 1:3}其中也包括我；祢的独生子，祢借着祂召叫了相信的子民进入义子的名分，其中也有我。我借着祂恳求祢，祂坐在祢的右边，\Quote{为我们代求，}\BibleRef{罗 8:34}\Quote{在祂内蕴藏着一切智慧和知识。}\BibleRef{哥 2:3}我正是在祢的经书中寻求祂。梅瑟所写的就是关于祂，\BibleRef{若 5:4}这是祂自己说的，这是真理说的。\par

% structure: p0009 source=h2 target=subsection reason=relative-heading-rank; base=h2

\subsection{第三章 他从世界的创造开始——虽不明白希伯来文圣经。}

5. 让我聆听并理解祢如何\BibleQuote{在起初}创造了天地。\BibleRef{创 1:1} \Person{梅瑟}写了这话；他写下后便离开，从祢这里到祢那里去了。如今他已不在我面前；若是他在，我必拉住他，求问他，并奉祢的名恳请他向我阐明这些事，我也必倾身用肉体的耳朵聆听他口中发出的声音。若他说的是希伯来语，那声音只会徒然撞击我的感官，无法触及我的心灵；但若他用拉丁语，我便能明白他所说的。然而，我又怎能知道他所说的是否真实？即使我知道这一点，难道我是从他那里得知的吗？其实，在我内里，在思想的密室中，真理——既非希伯来的，也非希腊的，既非拉丁的，也非蛮族的，无需舌与声的器官，无需音节的声音——会说：\Quote{祂说的是真理}，而我立刻确认，并自信地对祢的那个人说：\Quote{你说的是真理。}既然我不能求问他，我便恳求祢——祢，真理啊，他因充满祢而说出真理——祢，我的天主，我恳求，赦免我的罪；祢既赐给祢的仆人讲述这些事，也请赐我理解它们。\par

% structure: p0011 source=h2 target=subsection reason=relative-heading-rank; base=h2

\subsection{第四章 天地呼喊它们由天主所造。}

6. 看，天地存在；它们宣告它们是被造的，因为它们变化更迭。凡未曾被造而又存在者，其内绝无先前不在之物；这便是变化更迭的实质。它们也宣告它们并非自造；\Quote{所以我们存在，是因为我们被造；我们在存在之前并不存在，因此无法自造。} 说话者的声音本身就是证据。主，因此是祢造了这些事物；祢是美的，因为它们是美的；祢是善的，因为它们是善的；祢是存在的，因为它们存在。然而它们的美、善和存在，与它们的创造者祢相比，都不是美，不是善，也非存在。这些事我们知道，感谢祢。而我们的知识与祢的知识相比，便成了愚昧。\par

% structure: p0013 source=h2 target=subsection reason=relative-heading-rank; base=h2

\subsection{第五章 天主并非从某种物质，而是以祂自己的圣言创造了世界。}

7. 但祢是如何创造天地的呢？祢如此大能的工程所用的工具是什么？因为祢不像人间工匠，按照他心中所想，从一物制造另一物，能以某种方式赋予他内心之眼所见的形式。若非祢造了那心灵，他又怎能做到？他赋予已存在、仿佛具有实存之物形式，如黏土、石头、木头、黄金之类。若非祢命定它们，这些物又从何而来？祢为工匠造了身体——祢造了指挥肢体的心灵——祢造了用以制作万物的材料——祢赐予他领会技艺的能力，使他能内观而后外行——祢赐予他身体的感官，如同译者，藉此将心灵所构思的传达给物质，并将所成的回报给心灵，使内在能请教那掌管自身的真理，判断是否做得妥当。这一切都赞颂祢，万有的创造者。但祢是如何造的呢？天主啊，祢如何创造了天地？诚然，祢不是在天地中造成天地，也不是在空气中或水中，因为空气与水也属于天地；祢也不是在世界中造成世界，因为在受造之前，并没有地方可使世界存在；祢手中也没有任何材料用以创造天地。祢哪能从祢未造之物中获取材料去造物？因为除了因祢而存在，还有什么能存在呢？因此祢一言而生万物，祢以祢的圣言创造了它们。\par

% structure: p0015 source=h2 target=subsection reason=relative-heading-rank; base=h2

\subsection{第六章 然而，祂并非藉着发出而消逝的言语创造世界。}

8. 但祢是如何说话的呢？是否像那从云中传来的声音，说：\BibleQuote{这是我的爱子}？\BibleRef{玛 17:5} 那声音发出后便消逝了，有始有终。音节相继发声并消逝，第二音随第一音，第三音随第二音，依次排列，直至最后音结束，然后归于寂静。由此清晰可见，那是有时间性的受造物的运动表达它，服从祢永恒的旨意。祢的这些言语在时间中形成，外耳传达给理性的心灵，心灵的耳则专注聆听祢永恒的圣言。但心灵将这时间中作响的言语与祢静默中的永恒圣言相比，便说：\Quote{不同，大不相同。这些言语远在我之下，甚至并非真实存在，因为它们飞逝而过；而我主的圣言永远高于我。} 因此，若祢以发出而消逝的言语说天地应该受造，并如此造了天地，那么在天地之前就已存在一个形质体，藉着它时间中的运动使那声音得以在时间中传递。但在天地之前没有任何形质的东西；即便有，祢也无需短暂的声音来创造它，由此声音再由祢说出天地应该受造。因为凡用来形成那声音之物，若非祢所造，便根本不能存在。祢是用哪一句话命定了一个身体得以形成，再由这身体发出这些言语呢？\par

% structure: p0017 source=h2 target=subsection reason=relative-heading-rank; base=h2

\subsection{第七章 祂藉着同永远的圣言说话，万物遂成。}

9. 所以，祢呼唤我们去理解\OriginalTerm{圣言}{Word}，那与祢同在的天主，天主\BibleRef{若 1:1}，祂被永恒地说出，并且借着祂，万物被永恒地说出。因为所说的并非说完一件，再说另一件，直到说完所有；而是一次并永远地说了全部。否则我们就有了时间和变化，而非真正的永恒，也非真正的不死不灭。我的天主，我知晓这事，并献上感谢。我知晓，我向祢承认，主，凡对确凿真理不无感恩之心的人，也与我一同知晓并赞颂祢。主，我们知晓，我们知晓；因为一切事物在何等程度上不再是它昔日所是，而成为它昔日所非，便在何等程度上经历死亡与兴起。所以，在祢的圣言中，没有一物让位而又复位，因为祂真正是不死不灭且永恒的。因此，对那与祢同为永恒的圣言，祢一次并永远地说出了祢所说的一切；祢所命其成就的，便成就了；祢单单借由言语创造万物；然而，祢借由言语所造的万物，并非一同被造且永恒长存。\par

% structure: p0019 source=h2 target=subsection reason=relative-heading-rank; base=h2

\subsection{第八章 圣言本身乃是万有的开端，我们在这开端中领受了福音真理的教导。}

10. 我的主，我的天主，我恳求祢，这是为何？我看见了，然而不知道该如何表达，除非这样说：一切有始有终之物，其始其终，皆因在祢永恒的理性中，已知它应当始于或止于那无始无终之境。这同一位是祢的圣言，祂也是\OriginalTerm{开端}{Beginning}，因为祂也向我们发言。因此，在福音中，祂借着肉身发言；这声音播于人的外耳，为叫他们相信，并在内心寻求，从而在永恒的真理中找到，那位良善而唯一的导师在此教导祂所有的门徒。主啊，在那里，我听见祢的声音，是向我这说话者的声音，因为那教导我们的，才向我们说话。但那不教导我们的，虽然说话，却不是对我们说。再者，谁能教导我们，除非那不变的真理？因为即使我们通过可变的受造物获得规劝，我们也是被引向不变的真理。在那里，我们真正学习，当我们站立并聆听祂时，我们因\Quote{新郎的声音}而大大喜乐\BibleRef{若 3:29}，这声音引领我们回归我们的本源。所以，祂是开端，因为若非祂常存不变，当我们迷失时，便无归处。然而，我们脱离错误而回归，正是凭着认识而回归。但为了使我们能认识，祂教导我们，因为祂就是开端，并向我们发言。\par

% structure: p0021 source=h2 target=subsection reason=relative-heading-rank; base=h2

\subsection{第九章 智慧与开端。}

11. 天主啊，祢在这\OriginalTerm{开端}{Beginning}里创造了天地——在祢的圣言中，在祢的圣子内，在祢的能力中，在祢的智慧内，在祢的真理中，奇妙地言说，并奇妙地创造。谁能领会？谁能述说？那穿透我而发光，无损地击打我心，使我既战栗又燃烧的，是什么？我战栗，因为我与它不同；我燃烧，因为我与它相似。正是\OriginalTerm{智慧}{Wisdom}本身穿透我而发光，驱散我的幽暗，但当我在惩罚的黑暗与重压下，因它而昏厥时，这幽暗又再次将我吞没。因为我的力量因困乏而衰微，以致不能承受我的福祉，直到祢，主啊，祢曾宽恕我所有的罪过，也医治我一切的软弱；因为祢还将救赎我的生命脱离腐朽，以祢的慈爱与怜悯作我的冠冕，并以美物满足我的心愿，使我如鹰返老还童。因为我们得救在于望德；通过忍耐我们等候祢的\OriginalTerm{许诺}{promises}。\BibleRef{罗 8:24}凡能的人，让他倾听祢内里的教导。我要满怀信心地，从祢的神谕中呼求：\BibleQuote{上主，祢的化工何其浩繁，全是祢以智慧所造成。}这智慧就是开端，祢在这开端中创造了天地。\par

% structure: p0023 source=h2 target=subsection reason=relative-heading-rank; base=h2

\subsection{第十章 那些询问天主创造天地之前在做什么的人，其鲁莽。}

12. 看啊，那些对我们说\Quote{天主创造天地之前在做什么？因为，他们说，如果祂闲散无事，什么也不做，为什么不像过去一样，永远也停止工作呢？因为如果天主内兴起了一种新的变动，一个新的意志，去形成祂以前从未形成的受造物，那么，如何能有真正的永恒，因为其中兴起了前所未有的意志？因为天主的意志并非受造物，而是在受造物之前；除非创造者的意志先行，否则无物可受造。因此，天主的意志属于祂的\OriginalTerm{本体}{Substance}。但如果天主本体中出现了前所未有的东西，那本体就不能真正称为永恒。但如果是天主的永恒意志让受造物存在，为何受造物也不是从永恒就有呢？}\par

% structure: p0025 source=h2 target=subsection reason=relative-heading-rank; base=h2

\subsection{第十一章 提出这种问题的人尚未认识天主的永恒，这永恒与时间流逝无关。}

13. 说这些话的人，尚未理解祢，天主的智慧啊，灵魂的光明啊；他们尚未理解那些借由祢并在祢内造成的事物是如何造成的。他们甚至努力领悟永恒之事；但他们的心仍在事物的过去与未来的变动中飘荡，依然摇摆不定。谁能把握并固定此心，使它稍事安歇，以便逐步捕捉那\OriginalTerm{常在的永恒}{everstanding eternity}的光辉，并将之与永不停驻的时间相比，看出两者不可同日而语；看出长时间之所以成为长时间，无非是由于众多逝去的运动，而这些运动不可能同时伸展；但在永恒中，没有一物消逝，整个都是现在；而时间却无一是整个现在；他还需看出，一切过去都受未来的驱迫，一切未来都随过去而来，而一切过去与未来，都由那永恒的现在所创造并发出？谁能把握世人的心，使它静立，看清那静止的永恒，它自身既非未来，也非过去，却发出未来与过去的时间？我的手能做到吗？我口舌的劝解之手，能成就如此大事吗？\par

% structure: p0027 source=h2 target=subsection reason=relative-heading-rank; base=h2

\subsection{第十二章 天主在创造世界之前做了什么。}

14. 看哪，我这样回答那询问的人：\Quote{天主在创造天地之前做什么呢？}我不像有人那样，为了回避问题的压力，用开玩笑的语气回答。据说他回答说：\Quote{祂正在为那些探究奥秘的人预备地狱。}理解是一回事，取笑是另一回事——我不这样回答。因为我更愿回答：\Quote{我不知道我所不知道的，}这样，我不至于让询问深奥之事的人成为笑柄，也不至于因给出虚假答案而受称赞。但我这样说：祢，我们的天主，是每一受造物的创造者；如果\Quote{天地}一词是指所有受造物，我就敢说：\Quote{天主在创造天地之前，没有创造任何东西。因为如果祂创造了什么，除了受造物，祂还能创造什么呢？}我多么希望我知道一切对我有益的事，正如我知道没有任何受造物是在任何受造物被造之前受造的一样。\par

% structure: p0029 source=h2 target=subsection reason=relative-heading-rank; base=h2

\subsection{第13章 在天主所造的时间之前，时间并不存在}

15. 但如果有人让飘忽的思绪在逝去时间的种种形象中游荡，并且惊讶祢——全能的天主、创造万有者、支撑万有者、天地的建筑师——竟然在无数世代中，在将要创造如此伟大的工程之前，一直按兵不动，那就让他醒悟过来，明白他所惊讶的是虚假的事。因为那些祢没有创造的无数世代，怎么可能逝去呢？祢是万代的创始者和创造者。那些不是由祢创造的时间，算是什么时间？如果它们不曾存在，又怎能逝去呢？因此，既然祢是所有时间的创造者，如果在祢创造天地之前有任何时间，为什么说祢按兵不动呢？因为那段时间正是祢创造的，在祢创造时间之前，时间不可能逝去。但如果天地之前没有时间，为什么问\Quote{祢那时在做什么呢？}因为没有时间，就没有\Quote{那时}。\par

16. 祢也不是以时间先于时间；不然祢就不是先于所有时间了。而是在永恒常在的超绝境地，祢先于一切过去的时间，并超越一切未来的时间，因为未来将要到来，一旦来到就会过去；然而\BibleQuote{祢永存不变，祢的寿数没有终点。}\BibleRef{咏 102:27}祢的岁月既不去也不来；我们的岁月却且有去有来，好使一切到来。祢的所有岁月同时并存，因为它们屹立不动；它们离去时，也不会被后来的岁月排挤，因为它们不会逝去；而我们所有的岁月，当一切都会止息时，才会结束。祢的岁月是一日，而祢的一日不是每日，而是今天；因为祢的今天不会让位给明天，因为它也不随昨天而来。祢的今天就是永恒；因此祢生了与祢同为永恒的圣子，祢曾对他说：\BibleQuote{我今日生了祢。}\BibleRef{咏 2:7}祢创造了所有时间；而在一切时间之前，祢就是存在，在任何时间之中，都不会没有时间。\par

% structure: p0032 source=h2 target=subsection reason=relative-heading-rank; base=h2

\subsection{第14章 过去和未来的时间都不是真实存在，只有现在才是真实的}

17. 因此，在任何时候祢都没有不创造什么，因为祢创造了时间本身。没有任何时间与祢同为永恒，因为祢永远长存；但这些时间若继续存在，它们也不会是时间。因为时间是什么呢？谁能轻易而又简明地解释它呢？谁能在思想中领悟它，以至于能说出一个关于它的词呢？然而，我们在谈话中有什么比时间更频繁、更熟悉地提及呢？当然，我们说到时间时就理解它，听到别人说到时间时也理解它。那么，时间到底是什么呢？如果没有人问我，我知道；如果我要向询问的人解释，我就不知道了。但我满有把握地说：我知道，如果没有事物逝去，就不会有过去的时间；如果没有事物将要到来，就不会有未来的时间；如果没有事物存在，就不会有现在的时间。因此，过去和未来这两种时间，既然过去现在已不存在，未来尚未存在，它们怎么可能存在呢？但如果现在是永远现在的，不转变成过去的时间，那么这就不再是时间，而是永恒了。所以，如果现在时间——如果是时间的话——之所以存在，仅仅是因为它转变成过去的时间，那么我们怎能在它之所以存在的原因就是它将不存在的情况下，还说它存在呢？直言之，我们只能因为时间趋向于不存在，才能说时间存在，难道不是这样吗？\par

% structure: p0034 source=h2 target=subsection reason=relative-heading-rank; base=h2

\subsection{第15章 现在只是时间的一个瞬间}

18. 然而，我们常说到\Quote{时间长，时间短；}但我们这样说只限于过去或未来的时间。例如，我们把一百年前叫做很长的一段过去时间；同样，把一百年后叫做很长的一段未来时间。较短的一段过去时间，我们说譬如十天前；较短的未来时间，譬如十天后。但那不存在的东西，怎么说是长或短呢？因为过去已经不在了，未来还没有来到。所以，我们不要说\Quote{它是长的；}而应当说过去的\Quote{曾是很长，}说到未来的，应说\Quote{将是很长。}我的主，我的光，难道祢的真理不是在此也嘲笑了人吗？因为那段过去的时间曾是长的，它是在已经过去的时候长呢，还是在还是现在的时候长呢？因为只有当有能够成为长的东西时，它才可能是长的；但一旦过去，它就不再存在了；因此，那根本不存在的，不可能长。所以，我们不要说\Quote{过去的时间曾是长的；}因为我们将找不到什么是曾经长的，既然它已经过去，就不存在了；而应该说\Quote{那个现在的时间是长的，因为当它是现在的时候，它是长的。}因为它那时尚未逝去而不存在，所以有能够成为长的东西。但过去之后，它就不再存在了，那么它也就停止为长的了。\par

19. 因此，人的灵魂啊，让我们看看现在的时间能否是长的；因为你是被赋予感知和度量时段能力的。你会如何回答我呢？现在的一百年是长的时间吗？首先，看看一百年能否是现在的。因为如果这一百年中的第一年正在行进，那是现在，但其他九十九年是未来，因此它们尚不存在。但如果第二年正在行进，那么一个已经过去，另一个是现在，其余的是未来。同样，如果我们选定这一百年中的任何中间一年为现在，之前的都过去了，之后的都是未来；因此一百年不能是现在的。至少看看那正在行进的一年本身能否是现在的。因为如果它的第一个月正在行进，其他的都是未来；如果第二个月，第一个月已经过去，其余的尚未到来。因此，正在行进的一年作为一个整体也不是现在的；如果作为一个整体不是现在的，那么这一年就不是现在的。因为十二个月构成一年，其中每一个正在行进的月份本身是现在的，但其余的要么是过去，要么是未来。尽管那正在行进的月份不是现在的，而只有一天是：如果是第一天，其余的要来；如果是最后一天，其余的已经过去；如果是中间的任何一天，那么介于过去和未来之间。\par

20. 看哪，那唯一能被发现可称为长的时间，即现在，已被缩减到几乎不足一天的短暂空间。但让我们也来讨论它，因为没有一天是作为一个整体而成为现在的。因为它由昼夜二十四个小时组成，其中第一个小时使其余的成为未来，最后一个小时使其余的成为过去，但中间的任何一小时，其前的是过去的，其后的是未来的。而那一个小时又在飞逝的微粒中消逝。凡它已经流逝的部分是过去，凡是留下的部分则是未来。如果设想有某个不能再分至最细微时刻微粒的时间片段，那么只有这个才能被称为现在；然而，它如此迅速地由未来飞向过去，以致不能有任何延展。因为如果它得以延展，便分为过去与未来；但现在是没有任何空间的。那么，我们可称为长的时间究竟在哪里呢？它是自然吗？我们确实不说\Quote{它是长的，}因为它尚不存在，不能是长的；但我们说\Quote{它将是长的。}那么，它将在何时呢？因为即使在那时，由于它尚是未来，它不会是长的，因为能成为长的尚不存在；但当它从尚不存在的未来，将已经开始存在并将成为现在，以致可以有能成为长的存在时，那么现在的时间就以上面的话大声宣称它不能是长的。\par

% structure: p0038 source=h2 target=subsection reason=relative-heading-rank; base=h2

\subsection{十六、时间只能在流逝时被感知或测量。}

21. 然而，主啊，我们觉察到时间的间隔，我们将其相互比较，说某些更长，某些更短。我们甚至量度这个时间比那个短多少或长多少；我们回答说\Quote{这个加倍或三倍，而那个只是一次，或与之相同。}但我们度量正在流逝的时间，是通过觉察它们来度量的；而过去的时间，如今已不存在，或未来的时间，如今尚不存在，谁能量度它们呢？除非也许有人胆敢说，那不存在的东西也能被度量。因此，当时间正在流逝时，它可被感知和量度；但当它已经过去，便不能，因为它已不存在。\par

% structure: p0040 source=h2 target=subsection reason=relative-heading-rank; base=h2

\subsection{十七、然而，存在过去和未来的时间。}

2. 我问，父啊，我不论断。我的天主，求你统御我、引领我。谁能对我说，没有三个时间（就如我们幼时所学，也如此教导孩童的），即过去、现在和未来，而只有现在，因为那两个并不存在？或者它们也存在；但当未来变成现在时，是从某个隐秘处来吗？当现在变成过去时，归入某个隐秘处吗？因为那些预言未来之事的人，如果这些事尚不存在，他们在哪里看见了它们？因为不存在的东西不能被看见。而讲述过去之事的人，如果不在心中觉察它们，就不能将其作为真实来叙述。如果这些事不存在，就绝不能被辨别。因此，存在过去和未来之事。\par

% structure: p0042 source=h2 target=subsection reason=relative-heading-rank; base=h2

\subsection{十八、过去和未来的时间只能被思考为现在。}

23. 主啊，容我进一步寻求；我的希望啊，不要让我所图落空。因为若存在过去和未来的时间，我渴望知道它们在哪里。但若我尚未能知，我仍知道，无论它们在哪里，在那里不是作为未来或过去，而是作为现在。因为若在那里它们仍是未来，它们就尚未在那里；若在那里它们是过去，它们就不再在那里。因此，无论它们在哪里，无论它们是什么，它们只是作为现在而存在。虽然过去的事被真实地叙述，它们是从记忆中被提取出来的——不是那些已过去的事情本身，而是由事物的影像在头脑中通过感官的足迹所形成的概念。我的童年，如今已不再，是在过去的时间，如今已不存在；但当我唤起它的影像并讲述时，我是在现在观看它，因为它仍在我的记忆中。对于预知未来之事，是否有类似的原因，即那尚不存在的事物的影像可以作为已然存在的被认识，我承认，我的天主，我不知道。但我确实知道，我们通常事先思考我们未来的行动，而这种预思是现在的；但我们所预思的行动却尚不存在，因为它是未来的；当我们着手并开始实行我们所预思的事情时，那行动才是存在的，因为那时它不是未来的，而是现在的。\par

24. 因此，无论这种对将来之事的隐秘预知是怎样的，除了现在的存在，什么也不能被看见。但现在存在的事物不是将来，而是现在。所以，当他们说将来之事被看见时，不是看见它们本身，它们尚不存在（即，是将来的）；而是或许看见了它们的原因或征兆，这些已经存在。所以，对那些已经在注视它们的人而言，它们不是将来，而是现在，由这些在心中构想出来来之事，从而预言将来。这些构想现在也存在，而预言者所注视的正是这些呈现在他们面前的构想。现在，让多种多样的事物给我一个例子吧。我观看破晓；我预言太阳即将升起。我所观看的是现在；我所预言的是将来——不是太阳是将来的，它已经存在；而是它的升起，尚未发生。然而，除非我心中有一幅太阳升起的图像，否则我无法预测它的升起，正如我现在说话时就有这样一幅图像。但我在天空中看到的黎明，并不是太阳升起，尽管它可能在太阳升起之前出现，也不是我心中的想象；这两者都被视为现在的，由此预言那将来的事。所以，将来之事尚不存在；如果尚不存在，它们就不存在。如果不存在，就完全不能被看见；但它们能从现存且被看见的事物中被预言。\par

% structure: p0045 source=h2 target=subsection reason=relative-heading-rank; base=h2

\subsection{第十九章 我们不知道天主以何种方式教导将来之事。}

25. 所以，祢，祢所造万物的主宰，祢用什么方法教导灵魂将来之事呢？因为祢曾教导祢的先知。那方式是怎样的呢？祢，为祢没有将来之事，却教导将来之事；或者更确切地说，就将来之事而言，教导现在的？因为不存在的事，当然不能被教导。这方式远非我的能力所能及；于我是太强大的，我无法达到；但藉着祢，我将能够，当祢赐下时，我隐藏的眼目之甘饴的光！\par

% structure: p0047 source=h2 target=subsection reason=relative-heading-rank; base=h2

\subsection{第二十章 时间当如何恰当地称呼。}

26. 但是，现在明显而清楚的是，既没有将来之事，也没有过去之事。说\Quote{有三种时间：过去、现在和将来}是不恰当的；但或许可以恰当地说：\Quote{有三种时间：\OriginalTerm{过去之物的现在}{present of things past}、\OriginalTerm{现在之物的现在}{present of things present}和\OriginalTerm{将来之物的现在}{present of things future}。}因为这三者以某种方式存在于灵魂中，此外我找不到它们：\OriginalTerm{过去之物的现在}{present of things past}即为记忆；\OriginalTerm{现在之物的现在}{present of things present}即为观看；\OriginalTerm{将来之物的现在}{present of things future}即为期待。如果允许我这样说，我看到有三种时间，我也承认有三种。\Quote{有三种时间：过去、现在和将来}，只是习惯上错误地这样说。看，我不困扰，不反驳，不指责；只要所说的能被理解，即将来之事和过去之事现在都不存在。因为我们说话恰当的情况很少，不恰当的情况很多；但我们要表达的意思能被理解。\par

% structure: p0049 source=h2 target=subsection reason=relative-heading-rank; base=h2

\subsection{第二十一章 时间如何能被测量。}

27. 我刚才说过，我们在时间流逝时测量它们，以便能说：这段时间是那段的两倍，或者这段只是那段那么多，以及我们能通过测量而说出的其他时间部分。因此，如我所说，我们在时间流逝时测量它们。如果有人问我：\Quote{你怎么知道？}我可以回答：\Quote{我知道，因为我们测量；我们不能测量不存在的事物；而过去和将来之事是不存在的。}但我们如何测量现在的时间，既然它没有绵延？它是在流逝时被测量的；但当它流逝之后，就不再被测量；因为那时已没有什么可测量的了。但在测量时，它从何处来，经过何种途径，又往何处去？从哪里来，除了从将来？经过什么途径，除了经过现在？往哪里去，除了进入过去？所以，它从尚不存在之处，经过没有绵延之处，进入现已不存在之处。但我们测量什么，如果不是在某种绵延中的时间？因为我们说单倍、双倍、三倍、相等，或以任何其他方式谈论时间，都是就时间的绵延而言。那么，我们在什么空间中测量流逝的时间呢？是在将来，时间从中经过吗？但那尚未被我们测量的，是不存在的。还是在现在，时间由此经过？但没有绵延，我们便不能测量。还是在过去，时间向之流去？但现已不存在的，我们无法测量。\par

% structure: p0051 source=h2 target=subsection reason=relative-heading-rank; base=h2

\subsection{第二十二章 他祈求天主解开这最纠缠的谜团。}

28. 我的灵魂渴望认识这最纠缠的谜团。主我的天主，善父啊，求祢不要——我藉着基督恳求祢——不要对我的渴望封闭这些既寻常又隐秘之事，以致阻碍它渗入其中；而是让它们藉着祢启迪的仁慈破晓，主啊。关于这些事，我该向谁请教呢？我该向谁倾吐我的无知，比向祢更合适呢？这些如此热切地朝向祢圣经的研究，并未使祢烦恼。赐给我所爱的吧；因为我确已爱，而这是祢赐给我的。父啊，祢诚然知道将好的恩赐给祢的子女，\BibleRef{玛 7:11} 求祢赐下，因我已立志求知，困难在我面前，直到祢开启它。我藉着基督恳求祢，在祂——\OriginalTerm{至圣者}{Holy of Holies}——的名里，不要让人打扰我。因为我已相信，所以我说话。这是我的希望；为此我生活，好能瞻仰主的喜乐。看，祢已使我的日子老去，它们消逝，而我不知道是何方式。我们谈论时间和时间，时间与时间——\Quote{自从他说了这话，时间多久了？}\Quote{自从他做了这事，时间多久了？}又有，\Quote{自从我看见那事，时间多久了？}还有，\Quote{这个音节是那个单短音节的两倍时间。}这些言语我们说出，也听到这些言语；我们被人理解，也理解别人。它们是最明显、最寻常的，而这些事同样又深深隐藏，对它们的发现是新颖的。\par

% structure: p0053 source=h2 target=subsection reason=relative-heading-rank; base=h2

\subsection{第二十三章 时间是一种延展。}

我曾听一位有学问的人说，太阳、月亮和星辰的运动构成了时间，但我并未赞同。因为，为何不是所有物体的运动更应当构成时间呢？假如天上的光体停息了，而一个陶轮仍在旋转，难道就没有用来衡量这些旋转的时间，使我们能说它们是以均匀的间隔转动，或者，如果它时而慢、时而快，那么某些旋转较长、另一些较短呢？当我们说这些话时，难道我们不正是在时间中说话吗？在我们的言语中，难道有些音节是长的，有些是短的，岂不正是因为前者在较长的时间中发声，后者在较短的时间中发声吗？愿天主赐予人，能在一件小事上看出事物不论大小所共有的观念。天上的星辰和光体，\BibleQuote{为的是作为征兆、时节、日子和年岁。}\BibleRef{创 1:14} 毫无疑问，它们确实如此；但我既不会说那个木轮的旋转是一天，他也不能因此说没有时间。\par

我渴望认识时间的能力和本性，我们用它来衡量物体的运动，并（例如）说这个运动比那个长一倍。因为我问：既然\Quote{日}不仅指太阳留在地球之上（依此，日是此物，夜是彼物），也指它从东运行到东的整个环绕——根据这种考虑，我们说\Quote{过了这么多天}（当我们说\Quote{这么多天}时，也包括了夜，并不另外计算夜的时段）——既然日是藉太阳的运动和它从东到东的环绕而完成，我问：是这运动本身是日，还是完成这运动的周期是日，还是二者都是？因为如果第一种情况是日，那么即使太阳在一小时这样短的时间内完成那行程，也会有一天。如果是第二种，那么从一次日出到另一次日出之间若只有一小时这么短的时段，那就不是一天；太阳必须环绕二十四次才能完成一天。如果两者都是，那么如果太阳在一小时内走完全程，就既不能称为一天；反之，如果太阳静止不动，而流逝的时间相当于太阳惯常从早到晚完成全程所需的时间，也不能称为一天。所以我现在不问人们称为日子的是什么，而是问时间是什么，我们用它来衡量太阳的环绕，如果太阳在十二小时这么短的时间内完成全程，我们会说它是在平常所需时间的一半内完成的；比较这两个时间，我们会称后者为单倍，前者为双倍，尽管太阳从东到东的运行有时在单倍时间内，有时在双倍时间内。因此，不要有人对我说，天体的运动就是时间；因为当一个人祈祷，为让他赢得胜利的战斗，太阳就停住不动：太阳停住了，但时间仍在流逝。因为在那场战斗进行和结束所需要的时段里，时间足够流逝了。\BibleRef{苏 10:12} 那么，我看出时间是一种\OriginalTerm{延伸}{extension}。但我真的看见，还是似乎看见了？祢，光明和真理啊，祢必将指示我。\par

% structure: p0056 source=h2 target=subsection reason=relative-heading-rank; base=h2

\subsection{时间并非我们用时间所测量的物体运动。}

如果有人说时间是\Quote{物体的运动}，祢命令我赞同吗？祢并没有命令我。因为我听说，任何物体无不在时间中运动；这是祢说的；但物体的运动本身就是时间，我未听说过，祢也未曾这样说。因为当物体运动时，我以时间来衡量它从开始运动到停止运动经历了多久。如果我没有看见它是何时开始的，而它持续运动，以致我看不见它何时停止，那么我就不能衡量，除非或许是从我开始看见直到我不再看见的那段时间。但如果我观察了很长一段时间，我只能声称这段时间是长的，却不能说它有多长，因为当我们说\Quote{有多长}时，我们是通过比较来说的，例如，\Quote{这个和那个一样长}，或\Quote{这个的长度是那个的两倍}，或诸如此类。然而，如果我们能够标记物体或其部分（如果像轮子那样运动）从何处来、到何处去的空间距离，那么我们就能说出物体或其部分从一处到另一处的运动用了多少时间。既然物体的运动是一回事，而我们用来衡量运动有多长的是另一回事，谁看不出这两者中哪一个更应被称为时间呢？因为，尽管物体有时运动，有时静止，我们不仅用时间衡量它的运动，也用时间衡量它的静止；我们说：\Quote{它静止的时间与它运动的时间一样长}；或\Quote{它静止的时间是它运动时间的两倍或三倍}；或者任何我们衡量中确定了的或凭想象认定的长短，或多或少，正如我们惯常所说的。所以，时间并非物体的运动。\par

% structure: p0058 source=h2 target=subsection reason=relative-heading-rank; base=h2

\subsection{他呼求天主光照自己的心灵。}

我向祢承认，主啊，我至今仍不知道时间是什么；同时，我又向祢承认，主啊，我知道我在时间中说这些事，而且我已经长时间地谈论时间，而所谓\Quote{长}，若非因时间的延续，便不是长。那么，我既然不知道时间是什么，又怎能知道这一点呢？或者，也许我不知道如何表达我所知道的？哀哉，我竟连自己无知到什么程度都不知道！请看，我的天主，我在祢面前并不说谎。我怎样说，我的心就怎样。祢必点亮我的灯；主，我的天主，祢必照亮我的黑暗。\par

% structure: p0060 source=h2 target=subsection reason=relative-heading-rank; base=h2

\subsection{我们用较短的时段衡量较长的时段。}

33. 我的灵魂不是真诚地向祢忏悔说，我确实在测量时间吗？但是，我的天主，我这样测量，却不知道我所测量的是什么吗？我以时间测量物体的运动；难道我不也测量时间本身吗？然而，事实上，除非我测量物体在其中运动的时间，我怎能测量一个物体的运动，知道它有多久，以及它从一处到另一处需要多久呢？因此，我究竟如何测量时间本身呢？还是我们以较短的时间来量较长的时间，如同用一肘的长度来量横梁的长度？因为，我们确实似乎是以一个短音节的时长来量一个长音节的时长，并说这是两倍。这样，我们用诗句的时长来量诗节的时长，用音步的时长来量诗句的时长，用音节的时长来量音步的时长，用短音节的时长来量长音节的时长；不是按页来量（因为那样我们量的是空间，不是时间），而是在发音吐字，它们流逝过去时，我们说，\Quote{这是一个长诗节，因为它由这么多行组成；长诗行，因为它们由这么多音步组成；长音步，因为它们由这么多音节延长；长音节，因为它是短音节的两倍。} 然而，即便这样，也得不到任何确定的时间度量；因为一首较短的诗句，如果发音更饱满，可能比较长的诗句如果发音更急促，占去更多时间。诗节如此，音步如此，音节也如此。由此看来，时间无非是一种伸展；但不知道是什么东西的伸展。令我惊奇的是，如果不是心灵本身的伸展，那又是什么呢？因为我究竟测量什么，我恳求祢，我的天主，即便当我不确定地说，\Quote{这段时间比那段时间长；} 或者甚至确定地说，\Quote{这是那的两倍？} 我知道我测量时间。但我并不测量未来，因为它还不存在；也不测量现在，因为它没有任何伸展；也不测量过去，因为它已不再存在。那么，我究竟测量什么？难道是正在流逝的时间，而非已经过去的时间？因为我曾这样说过。\par

% structure: p0062 source=h2 target=subsection reason=relative-heading-rank; base=h2

\subsection{时间是在流逝中被度量}

34. 坚持，我的心灵，并专心留意。天主是我们的助佑；祂创造了我们，而非我们自己。留心，真理显现之处。看，假设有一个物体的声音开始响起，正在响着，持续响着，看！它停止了——现在是一片寂静，那声音已成过去，不再是一个声音。在响起之前，它是未来的，那时无法度量，因为尚不存在；现在也无法度量，因为它已不复存在。因此，唯有在它响着的时候才能度量，因为那时存在可以被度量的东西。但即使那时，它也没有停止不动，因为它正在流逝、消逝。那么，它岂不因此更可被度量吗？因为当它流逝时，它被伸展到一段时间之中，在这段时间内它可以被度量，因为现在没有伸展。因此，如果那时它可以被度量，看！假设另一个声音开始响起，并且一直响着，持续不断，没有任何中断，我们能在它响着的时候度量它；因为一旦它停止发声，它就已成过去，不再有可以被度量的东西。让我们真实地度量它，并说出它有多长。但它还在响着，除非从它开始发声的那一瞬间，直到它结束的那一刻，否则不能度量。因为我们所度量的正是从某个起点到某个终点的间隔。因此，一个尚未结束的声音无法被度量，以便说它是长是短；也不能说它等于另一个，或是另一个的一倍、两倍等等。但当它结束时，它不再存在。那么，它如何能被度量呢？然而我们确实在度量时间；但既不是尚未存在的时间，也不是不再存在的时间，也不是被某些延迟所拖延的时间，也不是没有界限的时间。因此，我们既不度量未来的时间，也不度量过去的时间，也不度量现在的时间，也不度量流逝中的时间；然而我们确实在度量时间。\par

35. \OriginalTerm{天主，万有的创造者}{Deus Creator omnium}；这是八音节的诗句，长短音节交替。那么，四个短音节，即第一、第三、第五和第七音节，相对于四个长音节，即第二、第四、第六和第八音节，是单一的。每一个长音节对于每一个短音节来说，都占双倍的时间。我发出这些音，报告它们，果然如此，这是常识所感知到的。因此，凭常识，我用短音节来量长音节，发现它占有两倍的时间。但是，当一个音在另一个音之后响起，如果前者是短音节，后者是长音节，我该如何把握短音节，并以度量的方式将其应用于长音节，以便发现后者是前者的两倍，而事实上，除非短音节结束发声，否则长音节不会开始发声？我无法将那个长音节作为现在来度量，因为只有在它结束时我才去度量它。而它的结束就是它的消逝。那么，我究竟能度量什么？我用来度量的短音节在哪里？我所度量的长音节又在哪里？两者都已发出，都已飞逝，都已过去，不再存在；而我仍然在度量，并且自信地回答（就训练有素的感觉所能信赖的而言），就时间长度而言，这个音节是一倍，那个是两倍。我之所以能做到这一点，只是因为它们已经过去，已经结束。因此，我并不是度量它们本身，因为它们已不存在，而是度量留存在我记忆中的某种东西，它固定不变。\par

36. 我的心灵啊，我是在你内度量时间。切莫以你的喧嚷来淹没我；也就是说，切莫以你众多\OriginalTerm{印象}{impressions}来淹没你自己。我是在你内，我说，度量时间；那些流逝的事物在你身上造成的印象，当它们过去之后仍然存留，我所度量的，是这作为现在时间的印象，而不是那些已经流逝的事物本身——是它们造成了那印象。当我度量时间时，我所度量的正是它。由此，要么这些就是时间，要么我就没有度量时间。当我们度量静默，说这段静默的时间长得如同那个声音延续的时间时，那又当如何呢？我们岂不将思想延伸到对某个声音的度量上，仿佛它正在发声，以便我们能够就一段给定时间内的静默间歇作出某种陈述吗？因为当声音和舌头都静止时，我们依然在思想中回顾诗行、韵文，以及任何论述，或者运动的幅度，并陈述关于时间间隔的长短——这个比那个长多少——与假如我们发声表达它们时并无二致。倘若有人想要发出一个悠长的声音，并事先决定了它该有多长，那么，那个人确实已在静默中经历了一段时长，并将它托付给记忆，随后开始发出那声音，直到它延续至预定的终点；确实，它已发过声，并将继续发声。因为那已经完结的部分确已发过声，而剩余的部分则将要发声；于是它就这样向前推进，直至当前的注意力将未来带入过去：过去因未来的缩减而增加，直到未来耗尽之时，一切都成了过去。\par

% structure: p0066 source=h2 target=subsection reason=relative-heading-rank; base=h2

\subsection{人的心灵中的时间，它期望、考量并记忆}

37. 然而，那尚未来到的未来，如何能被缩减或耗尽呢？那已不复存在的过去，又如何能增长呢？除非在实施此事的心里，有三件事在完成。因为它既期望，又考量，又记忆，以至于它所期望的，经由它所考量的，转入它记忆之中。因此，谁敢否认未来尚不存在？但心灵中却已有对未来的\OriginalTerm{期望}{expectation}。谁敢否认过去如今已不存在？但心灵中却仍有对过去的\OriginalTerm{记忆}{memory}。谁敢否认现在没有广延，因为它在瞬间消逝？但我们的\OriginalTerm{考量}{consideration}却持续存在，正是通过它，那将要出现者得以进入消逝。尚未存在的未来并不因此是漫长的；相反，一个\Quote{漫长的未来}乃是\Quote{一段漫长的对未来之期望}。那已不复存在的过去也不漫长；而一个漫长的过去乃是\Quote{一段漫长的对过去之记忆}。\par

38. 我正要背诵一首我熟悉的圣咏。在我开始之前，我的注意力伸向整首圣咏；但当我开始之后，我诵念过的那部分所变成的过去，便扩展于我记忆之中；而我这行为的生命，则被分配给我的记忆（因我已诵念的）和我的期望（因我将要诵念的）；然而，我的考量与我同在，通过它，那曾是未来的被导入，使之成为过去。这过程进行与重复得越多，期望就越是缩减，而记忆则越是扩大，直到整个期望被耗尽，当那整个行为结束，全部转入记忆之时。那发生在整首圣咏中的情形，也发生在它的每一小节和每一个音节中；这情形同样适用于那更长的行为——或许这首圣咏只是其一部分；同样适用于人的整个一生——人的一切行为都是其部分；也同样适用于人类众子的全部世代——所有个人的一生都是其部分。\par

% structure: p0069 source=h2 target=subsection reason=relative-heading-rank; base=h2

\subsection{人的生命是一场\OriginalTerm{分心}{distraction}，但因着天主的仁慈，他专注于他\OriginalTerm{天上召叫}{heavenly calling}的奖赏}

39. 但\BibleQuote{因为你的慈爱胜于生命，}看，我的生命不过是一场分心，而你的右手扶持了我——在我主、人子内，他是你、那唯一者与我们这众多者之间的中保（\BibleRef{弟前 2:5}），在众多事物中的诸多分心里——好使我藉着他，能把握那曾把握我者，并得以从我的旧日中被收聚，追随那唯一者，忘尽过往的；不再分心，而是被吸引向前，不朝那将要发生且将消逝的，而是朝着那在前面的事（\BibleRef{斐 3:13}），不分心而是专心地，我向着\OriginalTerm{天上召叫}{heavenly calling}的奖赏奔跑，在那里我可以听到颂赞你的声音，并默观你的欢愉，它们既不来也不逝。但如今，我的岁月在哀叹中度过。主啊，你是我的安慰，我永在的父。但我被分散在时间中，我不知它们的秩序；我的思想，甚至我灵魂的最深处，也被繁杂的变化所撕裂，直到我融合于你，在你爱火的熔炼中被涤净和融化。\par

% structure: p0071 source=h2 target=subsection reason=relative-heading-rank; base=h2

\subsection{他再次驳斥那虚妄的问题：\Quote{在世界创造之前，天主做了什么？}}

40. 我将坚定不移，稳固在你内，在我被塑造的\OriginalTerm{原型}{mould}——你的真理中；我也不会容忍人们的问题，他们因一种惩罚性的疾病，饥渴于超过他们所能承受的，并说：\Quote{在天主创造天和地之前，他造了什么？} 或者：\Quote{当他从未造过任何事物时，他心中怎么会想到要去造些什么呢？} 主啊，求你恩赐他们，使他们好好思考自己所说的话，并明白：在没有时间的地方，他们不能说\Quote{从不}。因此，当说他\Quote{从未造过}某事物时，这不就是在说，那事物不是在任何一个时间被造的吗？所以，让他们明白，没有受造物，就不会有时间；让他们停止说那些虚话。也让他们伸展到那前面的事（\BibleRef{斐 3:13}），并理解：你啊，永恒的、一切时间的创造者，先于一切时间而存在，没有任何时间与你同为永恒，也没有任何受造物，即使可能有超越一切时间的受造物。\par

% structure: p0073 source=h2 target=subsection reason=relative-heading-rank; base=h2

\subsection{天主的认知与人的认知如何不同}

主，我的天主，祢奥秘的隐秘处何在？我的罪过又将我抛离那隐秘处有多远？请医治我的眼睛，使我得享祢的光明。确实，若有一种心灵，极富知识与预知，对过去未来的一切，都如同我对一首圣咏那样熟悉，那心灵便极其神妙，令人惊叹不已；因为无论是怎样过去的，或是后世将要来的，在祂面前都毫不隐藏，犹如我唱那首圣咏时，从起初已唱了多少，还剩多少未唱，并不向我隐藏一样。但是，万有的造物主，灵魂与肉身的造物主——切莫以为祢只是以这种方式知道未来与过去的一切。祢是以远为奇妙、远为奥秘的方式知道它们。因为祢是永恒不变者，是真实永恒的心灵的造物主，所以不会像人唱熟悉的歌或听熟悉的曲时那样，由于期待将唱的词，回忆已唱的词，而情感变化，感官分裂，在祢内发生类似的事。因此，正如祢在起初，毫无知识上的变化，就认识了天地，同样，祢在起初，也毫无行动上的分心，就创造了天地。理解的人，要向祢承认；不理解的人，也要向祢承认。啊，祢是那样崇高，然而心灵谦卑的人却是祢的居所；因为祢提拔俯首的人，凡以祢为崇高的人，必不致跌倒。\par

\SourceRecord{Translated by J.G. Pilkington.；Nicene and Post-Nicene Fathers, First Series；Vol. 1.；Edited by Philip Schaff.；Buffalo, NY: Christian Literature Publishing Co.,；1887.；<http://www.newadvent.org/fathers/110111.htm>.}

% structure: p0001 source=h1 target=section reason=page-title

\section{忏悔录（卷十二）}

他继续根据\WorkTitle{七十贤士译本}解释\BibleBook{创世}{记}第一章，并借助它特别论述了双重天，以及可能从中创造了整个世界的无形式的质料；随后论述了其他不被否定的解释，并详细阐述了圣经的含义。\par

% structure: p0003 source=h2 target=subsection reason=relative-heading-rank; base=h2

\subsection{发现真理是困难的，但天主已应许寻求者必寻见。}

1. 主啊，我的心被祢圣经的话语所打动，在我这贫乏的生命中忙碌不已；因此，在大多数情形下，人类智力的缺乏反而使言语繁多，因为探问多于发现，求索长于获得，而那敲门的手比领受的手更为活跃。我们握有应许，谁能破坏它呢？ \BibleQuote{若是天主偕同我们，谁能反对我们呢？} \BibleRef{罗 8:31} \BibleQuote{你们求，必要给你们；你们找，必要找着；你们敲，必要给你们开：因为凡求的，就必得到；找的，就必找到；敲的，就必给他开。} \BibleRef{玛 7:7} 这些都是祢自己的应许；当真理应许时，谁还害怕受欺骗呢？\par

% structure: p0005 source=h2 target=subsection reason=relative-heading-rank; base=h2

\subsection{论双重天——可见之天与天外之天。}

2. 我舌头的软弱向祢的尊威承认，是祢创造了天地。我所见的这天，我所踏的这地——我肉身所由造成的地——都是祢所造。但是，主啊，那\OriginalTerm{天外之天}{heaven of heavens}在哪里？我们曾在圣咏中听到这样的话：\BibleQuote{天外之天属于主，但祂将大地赐予世人}。那看不见的天在哪里？与它相比，我们所见的一切都不过是地。因为这整个有形的宇宙，并非处处完美，只是在低处获得了它美丽的形状，其底部就是我们所称的地；但与那天外之天相比，即便我们地上的天也只不过是地；是的，这两个庞大的形体，与那个——我不知道该怎样描述的天——就是属于主而非世人之子的天相比，称作地也不荒谬。\par

% structure: p0007 source=h2 target=subsection reason=relative-heading-rank; base=h2

\subsection{论深渊上的黑暗，及无形无状的地。}

3. 这地确实是\OriginalTerm{无形无状}{invisible and formless}的，而且有一种我不知其深度的深渊，其上没有光，因为它没有形式。因此祢命令记下，\BibleQuote{黑暗笼罩着深渊}；这若不是光的缺失，又是什么？倘若有了光，它若不位于一切之上，高高凌驾而普照，又能位于何处？所以，黑暗笼罩它，是因为上面的光不在；正如寂静位于无声之处。在那里有寂静，不就是没有声音吗？主啊，祢不是将这些教导给了向祢忏悔的灵魂吗？主啊，祢不是教导我，在祢形成并分开这无形式的质料之前，那里什么也没有，既无颜色，也无形状，既无身体，也无精神；然而也并非绝对的虚无，而是存在某种无形状的无形状态？\par

% structure: p0009 source=h2 target=subsection reason=relative-heading-rank; base=h2

\subsection{从无形式的质料中生出了美丽的世界。}

4. 那么，除非用某个约定俗成的词语，否则它应当被称为什么，才能以某种方式传达给那些心智迟钝的人呢？然而，在世界的各个部分中，有什么比地和深渊更接近于完全的无形呢？因为它们处于最低的位置，远不及那些透明的、发光的较高部分美丽。因此，我为何不可以认为，那无形式的质料——祢创造了它却没有赋予形状，以便用以造成这有形的世界——正适合用\BibleQuote{无形无状的地}这个名称来向人暗示呢？\par

% structure: p0011 source=h2 target=subsection reason=relative-heading-rank; base=h2

\subsection{质料的形式可能是什么。}

5. 因此，当人类的思想在此处追寻感官所能及的某种东西时，便对自己说：它并非某种\OriginalTerm{可理解的型式}{intelligible form}，如生命或正义，因为它是\OriginalTerm{形体的质料}{matter of bodies}；也不是感官可觉知的，因为在无形无状之中，没有什么能被看见或摸到——当人类的思想这样自言自语时，它或者设法在无知中去认识它，或者通过认识而承认对它无知。\par

% structure: p0013 source=h2 target=subsection reason=relative-heading-rank; base=h2

\subsection{他承认自己一度对质料持有错误想法。}

但若我，主啊，用口和笔向祢告明祢就此事物所教导我的一切，我事先听到它的名字却不理解（那些不能明白它的人告诉我的），便构想它为有无数多样的形式。因此我并未构想那事物本身；我的心灵在无序中思索着污秽可怕的\Quote{形式}，但仍是\Quote{形式}；我称之为无形，不是因为它缺乏形式，而是因为它具有这样的形式，若它显现，我的心灵会因其怪异和不协调而转离，人性的软弱也会因此不安。但即使我所构想的那物也是无形的，不是由于一切形式的缺失，而是与更美的形式相比较；真实的理性说服我，若我想构想完全无形式的物质，就必须彻底消除它一切形式的残余；而我不能。因为我宁可想象那被剥夺一切形式者根本不存在，也不愿构想一个介于形式与虚无之间的东西——既非成形，也非虚无，只是无形的，近乎虚无。于是我的心灵不再追问我的灵魂，它充满了成形物体的影像，并随己意变化和改变这些影像；我转而去注意物体本身，更深入地审视它们的可变性，因这可变性，它们不再是自己原有的，而开始成为自己所不是的；我视这种由一种形式到另一种形式的过渡，是通过某种无形的状态，而非通过纯粹的虚无；但我渴望知道，而非猜测。假若我的声音和笔向祢告明祢为我解开的关于此问题的所有纠结，哪位读者能忍受接受全部呢？然而，我的心不会停止向祢献上尊崇和赞美之歌，为那些它所不能表达的。因为可变事物的可变性本身能够接受可变事物所变成的一切形式。而这可变性，它是什么？是灵魂吗？是身体吗？是灵魂或身体的外在显现吗？若可说\Quote{虚无是某物}，和\Quote{所是者非是}，我会说这就是它；然而它已在某种方式上存在，因为它能接受这些可见和复合的形状。\par

% structure: p0015 source=h2 target=subsection reason=relative-heading-rank; base=h2

\subsection{第七章．天主从虚无中创造了天地}

7. 这又从何而来，以何方式，除非从祢，万有皆从祢而有，就它们存在而言？但离祢越远，就越不相似于祢；因为这不是空间的间隔。因此，主啊，祢不是在某一处为一事，在另一处为另一事，而是同一者，同一者，同一者，圣、圣、圣，主，全能的天主，在起初——这起初出于祢，在祢的智慧中——这智慧由祢的本体所生——创造了一些事物，而且是从虚无中。因为祢创造了天地，不是从祢自身造出，否则它们便与祢的独生子相等，从而甚至与祢相等；而任何不是从祢而来的事物应与祢相等，是绝无道理的。在祢之外别无他物，祢可从中创造这些，天主啊，唯一的天主圣三，三一的统一体；因此，祢从虚无中创造了天地——一件大事和一件小事，因为祢是全能的、至善的，使一切事物成善，甚至伟大的天和渺小的地。祢那时存在，别无其他，祢从中创造了天地；二者中，一个靠近祢，另一个靠近虚无——一个祢应超乎其上，另一个尚无在其下者。\par

% structure: p0017 source=h2 target=subsection reason=relative-heading-rank; base=h2

\subsection{第八章．天地是\BibleQuote{在起初}造成的；随后，六日间，世界从无形质料中被造}

8. 但那上天之天是为祢的，主啊；至于地，祢给了世人，要他们看见和触摸的，那时却不像我们现在看见和触摸的。因为它曾是看不见的，且无形无象的，\BibleRef{创 1:2}，深渊之上没有光；或者说，黑暗笼罩深渊，即比深渊内更浓。因为现在这可见的水之深渊，即便在其深处，也有适合其本性的光，多少能让其中底层的鱼和爬行生物感知。但整个深渊那时几乎就是虚无，因为它那时全然无形；然而已存在那可被赋予形式的。主啊，祢是用无形的质料造成了世界，这质料是祢从虚无中造成的，几乎就是虚无，从它造出我们世人所惊叹的这些伟大事物。因为这有形的天确实奇妙，它是水与水之间的穹苍，是祢在造光后的第二日，说：要有，便有了。\BibleRef{创 1:6}。祢称这穹苍为天，即这地与海的天，这地与海是祢在第三日造成的，祢将万世之前所造的无形质料赋予了可见的形体。因为在万世之前祢已造了一个天，但那属于这个天的天；因为在起初祢已造了天地。但祢所造的地本身是无形的质料，因为它是看不见的、且无形无象的，黑暗笼罩深渊。祢从这无形无象的地，从这无形状态，从这几乎虚无中，造成了这个可变的世界所由构成的一切，然而它又不真正稳定；其可变性正见于此，因为时间可在其中被观察和计数。因为时间是由事物的变化造成的，而事物的形状——其质料即前述那无形的地——在变化和流转。\par

% structure: p0019 source=h2 target=subsection reason=relative-heading-rank; base=h2

\subsection{第九章．论上天之天是理智的受造物，而地在造成之日以前是无形无象的}

9. 因此，圣神——祢仆人的导师，当他记述祢在太初创造天地时——对时间缄默，对日子缄默。因为毫无疑问，祢在太初创造的\OriginalTerm{天上的天}{heaven of heavens}，是某种\OriginalTerm{理智的受造物}{intellectual creature}，虽然丝毫不与祢——圣三——同永恒，却分沾祢的永恒，并因甘饴地、最幸福地默观祢，极大地抑制自身的\OriginalTerm{可变性}{mutability}，且从受造之时起，便毫无缺失地依附祢，超越一切流转的时间变化。但这无形状态——这混沌空虚的大地——本身并未被列在那些日子之中。因为既无形体又无秩序，就没有事物来去；而无来去，则肯定没有日子，也没有任何时间段落的变化。\par

% structure: p0021 source=h2 target=subsection reason=relative-heading-rank; base=h2

\subsection{第10章 他祈求天主，使他能在真光中生活，并得以被教导圣经的奥秘。}

10. 啊，愿真理——我心灵之光，而非我自己的黑暗——对我说话！我堕入了黑暗，便成为黑暗。但即使从那里，我爱慕了祢。我走入歧途，却记起了祢。我听见祢的声音在我后面呼唤我回来，但因不宁者的喧嚣，我几乎听不见。如今，看，我回来，热切渴望祢的泉源。不要让人阻止我；我要饮这水，这样便有生命。让我不要成为我自己的生命；我过去恶活，对自己是死亡；在祢内我复活。对我说话吧；教导我吧。我相信祢的圣经，它们的话语极为深奥。\par

% structure: p0023 source=h2 target=subsection reason=relative-heading-rank; base=h2

\subsection{第11章 天主可能启示于他的事。}

11. 主啊，祢已在我内耳用强烈的声言告诉我：祢是永恒的，独有不死不灭。\BibleRef{弟前 6:16} 因为祢不被任何形状或运动所改变，祢的意志也不因时间而变更，因为任何变化的意志都不是不死的。这在祢眼前是清楚的，我恳求祢，使它越来越清楚；并在这启示中，让我更清醒地安居于祢的翼下。主啊，同样地，祢在我内耳用强烈的声言对我说：祢创造了一切本性和实体，它们不是祢自己，然而它们存在；只有虚无不是由祢而来，以及意志的运动——从祢这自有者趋向那较少存在者——因为这样的运动就是罪责和罪；没有人的罪会伤害祢，或扰乱祢统治的秩序，无论是最初的还是最终的。这在祢眼前是清楚的，我恳求祢，使它越来越清楚；并在这启示中，让我更清醒地安居于祢的翼下。\par

12. 同样地，祢在我内耳用强烈的声言对我说：那种受造物——唯独祢是它的意志——不与祢同永恒；它以最持久的\OriginalTerm{纯洁性}{purity}从祢获得支撑，无时无地不呈现其\OriginalTerm{可变性}{mutability}；而祢自己常与它同在，它以其全部情感依附于祢，既无未来可期待，也不将记忆带入过去，不受任何变化所改变，也不延展于任何时间。哦，有福者——若真有这种存在——依附于祢的真福，便在祢内得福，祢是其永恒的寄居者与光照者！我看不出\OriginalTerm{天上的天}{heaven of heavens}，即属于主的，有什么比祢的殿宇更好的称号；它默观祢的喜悦，毫无转向他物的退避；一个\OriginalTerm{纯洁的心灵}{pure mind}，最宁静地合一，凭借\OriginalTerm{神圣的精神体}{holy spirits}的和平稳定，是祢圣城中的公民，\BibleQuote{在天上的境地}\BibleRef{弗 2:6}，超越所见诸天。\par

13. 由此，灵魂——它曾远走迷途——得以理解，如果她如今渴慕祢，如今她的泪已成了她的食粮，同时每日有人对她说\BibleQuote{你的天主在哪里？}\BibleRef{咏 42:4} 如果她如今向祢寻求一件事，渴望一生一世住在祢的殿宇中\BibleRef{咏 27:4}。而她的生命，岂不正是祢吗？祢的日子，岂不正是祢的永恒，如祢那永不消逝的岁月，因为祢是永恒不变的吗？由此可见，灵魂若能够，便能理解祢是如何超越一切时间而成为永恒；当祢的殿宇——虽然不与祢同永恒，却因持续而不失地依附祢——不经历任何时间的变迁。这在祢眼前是清楚的，我恳求祢，使它对我越来越清楚；并在这启示中，让我更清醒地安居于祢的翼下。\par

14. 看，我不知在这些最低等、最末后受造物的变化中，存在着怎样的\OriginalTerm{无形性}{formlessness}。谁会告诉我，除非那内心空虚、在自己幻想中游荡不定的人？除了这样的人，谁会告诉我（那一切形状缩减、消失后），只要无形性存留——藉着它事物从一种\OriginalTerm{形状}{figure}变为另一种——就能展示时间的变迁？因为肯定不能；因为没有运动的变化便没有时间，而没有形状便没有变化。\par

% structure: p0028 source=h2 target=subsection reason=relative-heading-rank; base=h2

\subsection{第12章 天主从无形的大地创造出另一个天和一个可见且有形的大地。}

我的天主啊，我按祢赐给我的程度思考这些事，祢激励我\BibleQuote{敲罢，}的程度，以及祢在我敲时给我开门的程度\BibleRef{玛 7:7}，我发现祢所造的两样事物，不在时间的范围之内，因为两者都不与祢\OriginalTerm{同永}{co-eternal}。一样，是如此受造，以至于在\OriginalTerm{默观}{contemplation}上毫无欠缺，在变化上毫无间隙，虽是可变，却未曾改变，得以完全享受祢的永恒与不变性；另一样，却是如此无形，以至于它不具由运动或静止而得以从一种形式变成另一种形式、因而可能隶属于时间的那种能力。但这祢并没有让之保持无形，因为在所有时日之前，在起初祢创造了天地——就是我所说的这两样事物。\BibleQuote{大地还是无形无状，深渊上还是黑暗。}\BibleRef{创 1:2} 这些话向我们传达了它的无状，好让那些若不完全化为虚无便无法完全理解一切形式之缺乏的心灵，能被逐渐引领——由此另有一个天可能被造，另有一个可见而形态美好的地，以及排列有序的水，并在此世界的形成中所记载的一切，没有一日不在受造；因为这些事物是如此，以致于在其内，由于运动与形式的指定变化，时间的变迁得以发生。\par

% structure: p0030 source=h2 target=subsection reason=relative-heading-rank; base=h2

\subsection{第十三章 论理智的天与无形的大地，其后一日，穹苍由此成形。}

我的天主啊，当我听见祢的圣经说：\BibleQuote{在起初天主创造了天地；大地还是无形无状，深渊上还是黑暗，}却没有说明祢在哪一天创造这些事物时，我暂时这样理解：这是因为那个\OriginalTerm{天上的天}{heaven of heavens}，那个\OriginalTerm{理智的天}{intellectual heaven}，在那里，理解就是一次完全知晓——不是\BibleQuote{局部的，}不是\BibleQuote{模糊不清的，}不是\BibleQuote{借着镜子的，}\BibleRef{格前 13:12} 而是整体的，显明的，\BibleQuote{面对面的；}不是此时见这、彼时见那，而是（如已说过的）毫无时间变化地一次全知；也因为那个无形无状的地，毫无时间变化；而变化通常有\Quote{此时见这、彼时见那，}因为，没有形式的地方，就没有这个或那个的区别；——那么，是因为这两者——一个是\OriginalTerm{原始受造的}{primitively formed}，一个是\OriginalTerm{全然无形的}{wholly formless}；前者是天，而且是\OriginalTerm{天上的天}{heaven of heavens}，后者是地，而且是无形无状的地——正是因为这两者，我所以暂时认为祢的圣经没有提及日子，便说了\BibleQuote{在起初天主创造了天地。}因为随即它附加说明了所指的地是怎样的。而记载第二天造了\OriginalTerm{穹苍}{firmament}，并称之为天，这就向我们暗示，祂先前未提日子所说的天，是哪一个天。\par

% structure: p0032 source=h2 target=subsection reason=relative-heading-rank; base=h2

\subsection{第十四章 论圣经的深度及其仇敌。}

祢的\OriginalTerm{神谕}{oracles}何其高深，其表面展现眼前，吸引着弱小者；然而又是何其高深，我的天主啊，何其高深。注视它令人敬畏；是尊崇之敬畏，是爱慕之颤栗。我极度惱恨其仇敌。啊，愿祢用祢双刃的剑杀戮他们，使他们不再是其仇敌！因为我如此爱，愿他们对自己已死，好为祢而活。然而看啊，另一些人不是责备者，而是\WorkTitle{创世纪}的赞扬者——他们说：\Quote{天主的神，借着祂的仆人梅瑟写下这些事，祂无意使这些话被如此理解。祂无意使人如你所说的去理解，而意在如我们所说的。}对我们众人的天主啊，祢自己作为审判者，我如此回答他们。\par

% structure: p0034 source=h2 target=subsection reason=relative-heading-rank; base=h2

\subsection{第十五章 他为天上的天与对手辩论。}

\Quote{你们要说这些事是假的吗？真理在我内耳里大声告诉我，关于造物主本身的永恒性：祂的\OriginalTerm{本体}{substance}丝毫不随时间改变，祂的\OriginalTerm{意志}{will}也不与祂的本体分离。因此，祂现在不愿意这事，那时不愿意那事，而是一次而永远地愿意祂所愿的一切；不是反复再三，也不是此时愿意此、那时愿意彼；祂后来不愿意以前所不愿的，也不愿以前所愿的。因为这样的意志是可变的，而可变者都不是永恒的；但我们的天主是永恒的。同样，祂也告诉我，在我内耳里告诉我：对未来事物的期待，在它们来临时就变为看见；而这同样的看见，在它们过去后就变为记忆。再者，凡如此变化的思想都是可变的，而可变者都不是永恒的；但我们的天主是永恒的。}我归纳整合这些事，发现我的天主、永恒的天主，并非由于任何新的意志创造了任何受造物，也不让祂的知识遭受任何暂变。\par

那么，反对者啊，你们要说什么？这些是假的吗？他们说：\Quote{不，}他们说：\Quote{这是什么？那么，凡已具形式的性体，或可形式的物质，只因至高的善才来自至高者，这是假的吗？……我们也不否认这点，}他们说：\Quote{那么怎样？你们否认这一点吗？即存在某种高超的受造物，以如此贞洁的爱依附真、真正永恒的天主，以致虽不与祂\OriginalTerm{同永}{co-eternal}，却不与祂分离，不流入时间的任何多样变化之中，而只安息于对祂最真的默观？}因为，天主啊，祢向他显示祢自己，对他已足够，他爱祢如同祢所命令的那样，因此他不偏离祢，也不转向自己。这就是天主的宫殿，不是属地的，也不是任何有形体的天的体积，而是一座属神的殿，分享祢的永恒，因为它永远毫无瑕疵。因为祢使它坚固，直到永永远远；祢给它颁布了法律，它不会逾越。然而，天主啊，它并不与祢同永，因为并非无始无终，它原是受造的。\par

20. 因为虽然我们发现在它之前没有时间，因为智慧在万有之前就已被造\BibleRef{德 1:4}——这当然不是指那显然与祢、我们的天主、祂的父同永恒、同等的那位智慧，也不是那位万有都借祂而造，并在祂内，作为元始，祢创造了天和地的智慧；而是指那被造的智慧，即理智的本性，它在光的默观中，就是光。因为这智慧，虽被造，也被称为智慧。但光照者与被光照者之间的差别有多大，创造的智慧与被造的智慧之间的差别就有多大；正如那使人成义的正义与由成义而造成的正义之间一样。因为我们也被称为祢的正义；因为祢的某位仆人如此说：\BibleQuote{好使我们在祂内成为天主的正义。}\BibleRef{格后 5:21} 因此，既然某种被造的智慧在万有之前就被造了，那便是祢那贞洁之城——我们在上、自由之母的理性与理智之心\BibleRef{迦 4:26}，也是\BibleQuote{在诸天之上的永恒}\BibleRef{格后 5:1}的（在哪里的诸天？除非在那赞颂祢的诸天，即\Quote{天外之天}，因为这也就是\Quote{天}，属于上主的）——虽然我们发现在它之前没有时间，因为在万有之前受造者也先于时间的受造物，然而造物主本身的永恒是在它之前，它从造物主受造而获得了开端，虽然不是时间的开端——因为那时还没有时间——却是其自身本性的开端。\par

21. 因此，它虽出自祢、我们的天主，却显然有别于祢，并非与祢同一。因为虽然我们不仅发现在它之前没有时间，而且在它之内也没有时间（它永远得见祢的圣容，从不转面不顾，因此它不受任何变化的影响），但其中仍有可变性，可能因此而变得黑暗与寒冷，然而因它以至高之爱依附祢，便从祢那里发光发热，如同永恒的午时。啊，充满光明与辉煌的居所！我爱慕你的美丽，以及我主、你的建造者和拥有者的荣耀居所。让我这漂泊者为你而叹息；我向那造了你的那位说话，求祂也让我在你内属于祂，因为祂同样造了我。\BibleQuote{我像一只迷路的羊，走失了}；\BibleRef{路 15:5} 然而我盼望能被带回给你，靠着我牧人、你的建造者的肩头。\par

22. \Quote{反对者啊，我此前对你们说话，你们仍相信梅瑟是天主的圣仆，他的书是圣神的谕言，你们现在对我说什么呢？这天主的居所，虽然并非与天主同永恒，但依照其尺度，是\BibleQuote{在诸天之上的永恒}\BibleRef{格后 5:1}，你们在那里徒然寻求时间的变迁，因为你们找不到；因为它超越一切伸展，和一切周而复始的时间空间，对它而言，永远依附于天主是好的。}他们说：\Quote{是的。}\Quote{那么，当我心向我的天主呼号，在内心听到祂赞颂的声音时，我所说出的那些事，你们凭什么说它们是假的呢？难道是因为那质料是无形的，其中既没有形式，也没有秩序吗？但在没有秩序的地方，不可能有任何时间的变迁；然而这\InnerQuote{几乎虚无}，因为它并非完全虚无，确实来自那所有万有的根源，无论任何事物处于何种状态，都来自祂。}他们说：\Quote{我们也不否认这点。}\par

% structure: p0040 source=h2 target=subsection reason=relative-heading-rank; base=h2

\subsection{第十六章 他希望与否认神圣真理的人无涉。}

23. 对那些承认祢的真理指示我心灵的这一切为真的人，我愿在祢面前稍作讨论，我的天主。那些否认这些事的人，任凭他们狂吠，以自己的喧嚷淹没自己的声音好了；我将尽力劝说他们安静，容许祢的言语达到他们。但如果他们不愿，如果他们拒绝我，我恳求祢，我的天主，\Quote{不要对我缄默}。求祢在我心中真实地发言，因为只有祢这样说话，我将让他们离去，在尘土上扬尘，吹入他们自己的眼中；我自己却要进入我的内室\BibleRef{依 26:20}，在那里向祢咏唱爱歌——在我流徙的旅程中\BibleQuote{以无可言喻的叹息而叹息}\BibleRef{罗 8:26}，怀念耶路撒冷，我的心仰望它，耶路撒冷是我的故乡，耶路撒冷是我的母亲，祢是它的统治者、光照者、父亲、守护者、夫君、纯洁而坚固的欢愉、坚实的喜乐，和一切不可言喻的美善，甚至同时就是一切，因为祢是那唯一、至高的真善。我决不转身离去，直到祢将我这分散而变形的一切，聚集到那极亲爱的母亲的平安中，那里有我心灵的初果，并从那里这些事确证于我，祢把这我塑造并坚定直到永远，我的天主，我的慈悯。但对于那些不说这一切真假，却尊崇圣梅瑟所传、由祢圣神所启示的圣经，将其与我们一起置于应遵从的权威顶峰，却在某些细节上与我们矛盾的人，我这样说：我们的天主啊，求祢在我的忏悔与他们的反对之间作裁判。\par

% structure: p0042 source=h2 target=subsection reason=relative-heading-rank; base=h2

\subsection{第十七章 他提到\WorkTitle{创世记}起初之言的五种解释。}

他们说道：\Quote{虽然这些事都是真实的，但梅瑟借着天主的启示说：\BibleQuote{在起初天主创造了天地。}\BibleRef{创 1:1} 他在\InnerQuote{天}这个名称下，所指的并不是那常仰望天主圣容的、属灵而有理智的受造物；同样，\InnerQuote{地}也不是指那无形无象的物质。}\Quote{那么，指的是什么呢？}\Quote{那人，}他们说，\Quote{所表达的，正如我们所说的；他用这些词语所宣告的就是这个意思。}\Quote{那是什么呢？}\Quote{借着\InnerQuote{天地}这个名称，}他们说，\Quote{他首先想普遍而简略地指出这整个可见世界，以便后来借着逐日详列，将圣神喜悦如此启示的一切事，都细致地分述出来。因为他所面对的是那些粗鲁且属血肉的百姓，所以他断定只将天主那些可见的工程托付给他们，是明智的。}然而，他们一致认为，那不可见且无形的地，以及黑暗的深渊（后来由此指出，在那\InnerQuote{日子}期间，所有这些人人熟知的有形之物都被造成并安排就绪），正可以恰当地理解为这种无形无象的物质。\par

如果另有人说：\Quote{这同样的物质的无形无象与混沌，最初正是以\InnerQuote{天地}的名义被提出的，因为由它而造成了这个可见世界，连同其中一切最明显可见的、且通常以\InnerQuote{天地}之名来称呼的万有，并被完善。}又如何呢？但是，如果还有人说：\Quote{那无形与有形的自然，被恰当地称为\InnerQuote{天地}；因此，那普世的受造，即天主在祂的智慧中所造成的——也就是\InnerQuote{在起初}——都包含在这两个词中了。然而，因为万有都不是由天主的本体，而是由虚无中造成的（因为它们并非与天主同一实体，而且它们都内在地拥有某种可变性，不论是永远存在的，如同天主永远的居所，还是变化的，如同人的灵魂与肉身），所以，那无形与有形万有的共同物质——尚无形相，但能接受形式——将来由它造成天地（也就是已经成形了的无形与有形的受造物）的共同物质，正是用那称呼\InnerQuote{不可见且无形的地}和\InnerQuote{深渊上的黑暗}的同一名称来指代的；不过有一个区别，即那\InnerQuote{不可见且无形的地}被理解为形体物质，在它还未有任何形式之前；而\InnerQuote{深渊上的黑暗}则被理解为属灵物质，在它那无限制的流动性还未受任何约束，且未获得智慧的光照之前。}\par

如果有人愿意，还可以说：\Quote{当读到\BibleQuote{在起初天主创造了天地}时，那些已经完成并成形的受造自然，无论是无形的还是有形的，并没有用\InnerQuote{天地}一词意指；相反，那尚未成形的事物的原初，即能够被赋予形式和受造的物质，正是用这些名称来称呼的，因为其中包含着这些混沌之物，尚未按各自的品质和形式区分开来，而如今这些已按各自秩序排列，便被称为\InnerQuote{天地}；前者是属灵的受造物，后者是形体的受造物。}\par

% structure: p0046 source=h2 target=subsection reason=relative-heading-rank; base=h2

\subsection{第十八章 在圣经中何种错误是无害的。}

听到了并思考了这一切之后，我不愿为言语争辩，因为那除了颠覆听者之外，毫无益处。\BibleRef{弟后 2:14} 但是法律若合法地使用，是善的，为能建树；\BibleRef{弟前 1:8} 因为这法律的终向\BibleQuote{是出自纯洁的心、善良的良心和无伪的信德的爱德。}我们的师清楚知道，他是把全部法律和先知总括在这两条诫命上的。我的天主，我暗中之眼的明灯啊，在我热切承认这些事时——既然由这些话语可以理解出许多事，而这些理解又都是真实的——我若对作者的原意有不同于他人的看法，这有何妨害呢？的确，我们所有阅读的人，都努力探求并理解我们所读的作者想要传达的意思；我们既然相信他所说的是真理，就不敢认为他说了任何我们明知或以为虚假的事。因此，既然每个人都努力在圣经中理解作者所理解的意思，那么，如果他理解的是你——一切说真理的明悟之光——向他显示的真理，尽管他所读的作者并不理解这个真理，但这又何妨呢？因为作者也理解了一个真理，虽然不是这个真理。\par

% structure: p0048 source=h2 target=subsection reason=relative-heading-rank; base=h2

\subsection{第十九章 他列举了所有人都同意的观点。}

主啊，你创造了天地，这是真实的；起初就是你的智慧，你在其中创造万有，这也是真实的。同样，这个可见的世界有其巨大的部分，即天地，它简短地概括了一切受造和形成的自然，这也是真实的。一切可变之物都在我们的心灵前呈现出某种形式的缺乏，它由此而取得形式，或变化转变，这也是真实的。那如此依附于不变形式，以至于虽然可变，却并不改变的事物，是不受时间支配的，这也是真实的。那近乎虚无的无形无相，不能有时间的变化，这也是真实的。任何用来制造他物的东西，可以用它所造出之物的名称来称呼，这是一种说法；所以，那造成天地的无形无相者，也可以被称为\Quote{天地}，这也是真实的。在一切具有形式的事物中，没有比地和深渊更接近无形无相的了，这也是真实的。不仅每一个受造和成形的事物，而且凡是能够受造和接受形式的，都是你所造的，\BibleQuote{万有都出于他}\BibleRef{格前 8:6}，这也是真实的。凡是从无形无相者所形成的事物，在形成之前都是无形无相的，这也是真实的。\par

% structure: p0050 source=h2 target=subsection reason=relative-heading-rank; base=h2

\subsection{第二十章 论\Quote{在起初}一语的不同理解。}

29. 从所有这些真理中——那些因你赐予内心之眼而得以看见此类事物，并坚定不移地相信你的仆人\Person{梅瑟}是在真理的精神内发言的人，都对这些真理深信不疑——于是，有人从中选取一种说法，说：\BibleQuote{在起初天主创造了天地}——也就是说，\Quote{天主在自己的\OriginalTerm{圣言}{Word}内——这圣言与祂\OriginalTerm{同永恒}{co-eternal}——创造了可理解的和可感的，亦即精神受造物与有形受造物。}又有人选取另一种说法，说：\BibleQuote{在起初天主创造了天地}——也就是说，\Quote{主在自己的圣言内——这圣言与祂同永恒——创造了这个有形世界的普遍块体，连同其中一切可见且为人所知的本性。}还有人说，\BibleQuote{在起初天主创造了天地}——也就是说，\Quote{天主自己的圣言内——这圣言与祂同永恒——创造了精神与有形受造物的\OriginalTerm{无形物质}{formless matter}。}另有人说，\BibleQuote{在起初天主创造了天地}——也就是说，\Quote{天主在自己的圣言内——这圣言与祂同永恒——创造了有形受造物的无形物质，那时天地尚混淆于其中，如今已区分并成形，我们今日在这世界块体中所见的便是它们。}还有人说，\BibleQuote{起初天主创造了天地}——也就是说，\Quote{就在创造和工程最起初，天主造就了那混淆地包含天地的无形物质；万物由此形成，如今它们与其中所包含的一切一同鲜明地呈现出来。}\par

% structure: p0052 source=h2 target=subsection reason=relative-heading-rank; base=h2

\subsection{第21章 论\BibleQuote{大地无形}一语的阐释}

30. 关于以下话语的理解，从所有那些真理中，有人为自己选取了一种说法，说：\BibleQuote{大地还是混沌空虚，深渊上还是一团黑暗}——也就是说，\Quote{天主所造的那有形事物，当时还是\OriginalTerm{无形物质}{formless matter}，没有秩序，没有光明。}又有人选取另一种说法，说：\BibleQuote{大地还是混沌空虚，深渊上还是一团黑暗}——也就是说，\Quote{这整个被称为天地的，当时还是无形而黑暗的物质，有形之天与有形之地，连同我们感官所知的其中万物，都将由这物质造成。}还有人说，\BibleQuote{大地还是混沌空虚，深渊上还是一团黑暗}——也就是说，\Quote{整个被称为天地的，当时还是无形而黑暗的物质，从那物质中将要造出那可理解的天（亦称天外之天），以及大地，即整个有形本性（这有形之天也可包括在此名称之下），一切不可见与可见的受造物都将由此造成。}另有人说，\BibleQuote{大地还是混沌空虚，深渊上还是一团黑暗}—— \Quote{圣没有用天地的名称来称呼那无形状态，而是称呼那无形状态本身，}他说，\Quote{它已存在，他称之为无形、空虚的大地和黑暗的深渊，而在此之前他曾说过天主创造了天地，即精神的与有形的受造物。}还有人说，\BibleQuote{大地还是混沌空虚，深渊上还是一团黑暗}——也就是说，\Quote{已经存在一种无形物质，正如圣经之前所说，天主创造了天地，即整个世界的块体，分为上下两大部分，连同其中一切熟悉和已知的受造物。}\par

% structure: p0054 source=h2 target=subsection reason=relative-heading-rank; base=h2

\subsection{第22章 他讨论物质是否出于永恒，或由天主所造。}

31. 如果有人想要反对后两种意见，这样说——\Quote{如果你们不承认，这种无形式的物质看来是用天和地的名称来称呼的，那么，天主就没有创造某种用以造成天和地的材料；因为圣经并没有告诉我们天主创造了这物质，除非我们理解它暗含在天地这个名词里，或者单指地，当它说\BibleQuote{在起初天主创造了天地}，接着又说\BibleQuote{大地还是混沌空虚}，虽然他喜欢这样称呼那无形式的物质，我们却不可理解任何其他东西，只能理解为那已经记载的话：\BibleQuote{天主造了天地}。}持这两种意见中任何一种的人，在听到这些话后，会这样回答说：\Quote{我们并不否认，这无形式的物质是由天主所造，万物都是从他而来，样样都好；因为正如我们说那受到创造和形成的是更大的善，同样我们也承认那能够接受创造和形式的是较小的善，但仍是善。然而，圣经并没有声明天主创造了这无形式的状况，正如它也没有声明其他许多事物一样；例如\InnerQuote{革鲁宾}和\InnerQuote{色辣芬}，以及宗徒所明确提及的\InnerQuote{上座者}、\InnerQuote{宰制者}、\InnerQuote{率领者}、\InnerQuote{大能者}\BibleRef{哥 1:16}，显然这些都是天主造的。或者，如果\BibleQuote{他造了天地}这句话包含了万物，那么对于天主的神运行其上的水，我们该怎么说呢？如果把它们理解为包含在\InnerQuote{地}这个词里，那么当我们看到水如此美丽时，\InnerQuote{地}这个词怎能指无形式的物质呢？或者如果确实这样指，那么为何记载从同一无形式中造成了穹苍，并称之为天，却没有记载水是怎样造成的呢？因为我们看到如此美丽流动的水，并非无形式和空虚的。但如果它们是在天主说\BibleQuote{天下的水应聚在一处}\BibleRef{创 1:9}时获得了美丽，以至于那聚合就是形成，那么关于穹苍之上的水，又该如何回答呢？因为这些水如果是无形式的，它们就不配得到如此尊荣的位置，而且也没有记载它们是用哪句话形成的。所以，如果\WorkTitle{创世纪}对天主所造的某事物保持沉默，尽管健全的信德和正确的理智毫不怀疑那是天主所造，那么任何健全的教训都不可斗胆说这些水与天主同永远，因为我们在\WorkTitle{创世纪}中提到了它们；至于它们何时被造，我们却没有发现。为什么——真理教导我们——我们不能理解，圣经称为\BibleQuote{混沌空虚的大地}和\InnerQuote{黑暗深渊}的那无形式的物质，是天主从无中造成的，因此它们并不与祂同永远，尽管那叙述没有告诉我们它们何时被造呢？}\par

% structure: p0056 source=h2 target=subsection reason=relative-heading-rank; base=h2

\subsection{第23章 书中待解释的两种分歧}

32. 因此，我根据我软弱的理解力，听到并领悟了这些事——我向祢承认，主啊，祢知道我的软弱——我看出，当用言语表达任何事物时，即使由真实的传达者说出来，也可能产生两种分歧：一种关于事情的真相，另一种关于传达者的意图。因为一方面，我们探寻，关于受造物的形成，什么是真实的；另一方面，我们探寻，\Person{梅瑟}，祢信德的杰出仆人，希望读者和听者通过这些话语理解什么。关于第一种，让所有自以为知道虚假之事为真的人远离我。关于另一种，也让所有以为\Person{梅瑟}说了假话的人远离我。然而，主啊，愿我因祢而与那些在广阔的\OriginalTerm{爱德}{charity}中饱享祢真理的人结合，并在祢内与它们一同欢乐；让我们一起亲近祢圣书的言语，借着祢仆人的意愿——祢曾借他的笔墨传布了这些言语——在它们中寻求祢的旨意。\par

% structure: p0058 source=h2 target=subsection reason=relative-heading-rank; base=h2

\subsection{第24章 众多真理中，不可断言梅瑟持有此意或彼意}

33. 然而，在这些话语中，对探求者而言有如此多的真理，以不同的方式被理解，我们中谁能发现那唯一的解释，以致充满信心地说\Quote{梅瑟是这样想的}，\Quote{在那叙述中他希望让人理解这一点}，如同他充满信心地说\Quote{这是真的}那样，无论他是这样想还是那样想？请看，我的天主，我——祢的仆人，在这本书中向祢许愿献上\OriginalTerm{宣认的祭献}{sacrifice of confession}，并恳求祢按祢的慈悯使我向祢偿还我的誓愿——请看，我能否像我充满信心地断言，祢以祢永恒的圣言创造了一切，无形与有形的事物，同样充满信心地断言，当\Person{梅瑟}写下\BibleQuote{在起初天主创造了天地}时，他的意思无非就是如此？不能。因为在我看来，他写这些话时心中所想的是什么，并不如我在祢的真理中看到的那样确凿无疑。因为当他説\BibleQuote{在起初}时，他可能想到创世的开始；而且他可能希望被理解的是，在这里，\BibleQuote{天地}并不是已成形和完美的受造，无论是属灵的还是物质界的，而是各自都刚刚开始，仍然无形式。因为我看出，无论其中的哪一种说法，都可能是真实的；但他在这些话语中可能想到的是哪一种，我却没有看到。虽然，无论那是这些之一，或是我所未提及的其他意思，这位伟大的人在使用这些语词时心中所见的，我毫不怀疑他真实地看见了，并恰当地表达了出来。\par

% structure: p0060 source=h2 target=subsection reason=relative-heading-rank; base=h2

\subsection{第25章 解释者于隐晦处意见分歧时，当念及真理的创造者天主，并谨守爱德的规范}

34. 现在，谁也不要再用这样的话来搅扰我：梅瑟的想法并不如你所说，而是如我所言。因为如果他问我：\Quote{你怎知梅瑟有你所从他话中推演出的这意思？} 我应当平心静气地接受，并如先前那样回答他，或者，若他固执，回答得更详尽些。但若他说：\Quote{梅瑟的意思并非如你所说，而是如我所言，} 却又不否认我们各自所说的意思，且两者均为真，那么，我的天主，穷人的生命，祢的怀中没有矛盾，求祢向我的心倾注祢的抚慰，使我能够耐心容忍对我这样说的人；我所以容忍，并非因他们是属神的人，也并非因他们在祢仆人的心内看到了他们所说的，而是因为他们是骄傲的，并未认识梅瑟的意见，却只喜爱自己的——并非因其为真，而是因为那是他们自己的。否则，他们也会同样喜爱另一个真实的意见，正如我喜爱他们所言说真实时所说的；我喜爱它并非因为那是他们的，而是因为它是真实的，既为真实，则它已不只是他们的了。可是，他们若因它是真实的而喜爱它，那么它便既是他们的也是我的，因为它是所有爱真理的人所共有的。但他们却争论梅瑟的意思并非如我所说，而是我如他们所言；我对此既不爱好也不喜爱；因为纵然事实如此，那也不过是出于鲁莽而非真知灼见，并非由洞见、乃是由虚浮产生出来的。所以，主啊，祢的审判是令人畏惧的，因为祢的真理既不属于我，也不属于他，更不属于另一个人；而是属于我们众人。祢公开召唤我们共同享有这真理，并严正警告我们不可将它视为私有，否则我们便会被剥夺真理。凡是把祢指定给众人享受的东西据为己有，并希望将众人共享之物变为私有者，必被强制离开公有的而趋向私有的——亦即离开真理而走向虚假。因为 \BibleQuote{谁撒谎，便是由他自己所说的。} \BibleRef{若 8:44}\par

35. 谛听吧，天主，祢这位最卓越的审判者！真理本身，请谛听我要对这反驳者所说的话；求祢谛听，因为我在祢面前说这话，也在依法使用祢的法律的弟兄们面前说，都是为了爱德；\BibleRef{弟前 1:8} 求祢谛听并垂顾，若蒙祢悦纳，我要对他说的话。我以弟兄般的和平言辞回答他：\Quote{如果我们二人都看出你所说的为真实，也看出我所说的为真实，请问，我们是在何处看出的？当然，不是我在你内，也不是你在我内，而是在那超越我们心灵的、不变的真理本身中看出的。} 既然关于我们主天主的真光，我们都不去争论，为什么竟要争论我们邻人的思想呢？因为我们不能像看见那不变的真理一样看见它；即使梅瑟本人出现在我们面前说\Quote{我的意思正是如此}，我们也并非因此看见，而只是相信罢了。因此，我们不要\BibleQuote{彼此吹嘘，代对方设想}，\BibleRef{格前 4:6} 超越所记载的；我们要\BibleQuote{全心、全灵、全意爱上主我们的天主，并爱近人如自己。} \BibleRef{谷 12:30} 关于这两条爱德的诫命，除非我们相信梅瑟在这些书中所写的一切都是按照他的本意，否则，如果我们对同仆之人的心思另作他想，就无异于将天主视为说谎者。现在请看看，从这些话中可以引申出这么多真实的见解，却要轻率地断言梅瑟特指其中哪一种，并用有害的争论伤害爱德本身，是何等愚昧；而梅瑟正是为了爱德，才说出了我们如今竭力解释的这一切话！\par

% structure: p0063 source=h2 target=subsection reason=relative-heading-rank; base=h2

\subsection{第26章 若他被指派撰写\BibleBook{创}{世纪}，他可能会向天主祈求什么。}

36. 然而，我主天主，祢是我谦卑的荣耀，我辛劳的安息；祢俯听我的忏悔，并宽恕我的罪过。既然祢命令我爱近人如己，我便不能相信，祢赐予祢最忠信的仆人梅瑟的恩赐，会少于我若生于他的时代，被祢安置在那地位上，以心灵和口舌的服务去传播那些书卷时，自己从祢那里所愿欲并渴求的；这些书卷在许久之后，要造福万民，并能从如此崇高的权威高峰，超越世上一切虚假和骄傲的学说。我若那时是梅瑟（因为我们所有人都出自同一灰土；人算什么，祢竟怀念他？），我若在那时就是他，又被祢指派撰写\BibleBook{创}{世纪}，我那时真会渴望得以拥有这样表达的恩赐和编排的方法，使得那些还不能理解天主如何创造的人，不会因它超越自己的能力而加以拒绝；而那些已能领会的人，则无论他们依真实见解思考到了什么，都会发现这一切在祢仆人寥寥数语中并未遗漏；而且若另一人由真理之光发现了另一见解，这些见解也必然能在同样的话语中找到。\par

% structure: p0065 source=h2 target=subsection reason=relative-heading-rank; base=h2

\subsection{第27章 \BibleBook{创}{世纪}的叙述风格简明清晰。}

37. 如同一个在狭小空间内的泉源更加丰沛，能比从同一泉源流出、经漫长距离后到达的任何一条溪流，为更广大的空间提供更多的水流；同样，祢的传述者的叙述，原本预定要造福许多可能就此展开论述的人，从有限的语言尺度中，漫溢为清晰真理的条条溪流，每个人都可以从中汲取关于这些主题自己能获得的真理——这人汲取这一真理，那人汲取另一真理，通过更长的言辞迂回。因为有些人，当他们读到或听到这些话时，便以为天主如同一个人，或某个具备无边大能的庞然大物，凭借某种新而突然的决定，在自身之外，仿佛在遥远的地方，创造了天地——即上下两个巨大的物体，其中将容纳万物。当他们听到 \BibleQuote{天主说，要有，便有了} 时，他们想到的是有始有终的话语，在时间中响起并消逝，一旦这些话逝去，那受命之物便生成；以及他们因熟习世俗而可能想到的其他类似事物。在他们里面，当他们还是小孩子时，他们软弱的本性借着这种谦卑的说话方式，如同在母亲的怀抱中得到扶持，他们的信德便健康地建立起来，借此他们确知并坚信，天主创造了万物，就是他们感官在四周所见、奇妙多样的种种本性。这些话，如果有人鄙视它们，视若浅薄，以骄傲的软弱试图超越抚养他的摇篮，他啊，可悲可怜，必将惨然跌倒。主，天主啊，求祢怜悯，免得路过的人践踏那羽毛未丰的幼鸟；并派遣祢的天使，将它送回巢中，使它可以存活，直到能够飞翔。\par

% structure: p0067 source=h2 target=subsection reason=relative-heading-rank; base=h2

\subsection{第二十八章 \Quote{在起初} 与 \Quote{天地} 这两个用语被以不同方式理解}

38. 但另一些人，对这些话已不再视为巢穴，而是浓荫的果林，他们看见其中隐藏的果实，便欢欣地飞翔，啾啾而鸣地搜寻、采撷。因为，天主啊，当他们读到或听到这些话时，他们看见，所有过去与未来的时间，都为祢永恒而稳固的恒在所超越，而且没有任何时间的受造物不是祢所创造的。祢凭祢的旨意——因为那旨意就是祢之所是——创造了万物，不是凭任何改变了的旨意，也不是凭先前不存在的旨意；也不是从祢自身中，按祢的肖像、万物的形式而造，而是从无中，造出一种无形的不相似之物，好使它因祢的相似而被赋形（万物按其本性所赋予的、既定的容量，归向祢那唯一者），并使一切都成为甚好的；无论它们是常驻祢四周，还是随时间与空间逐步迁离，展现或承受美妙的变异。他们看见这些，并在他们现世所能得到的那少许光中，因祢真理的光而欢欣。\par

39. 再者，这些人中的另一位，把他的注意力转向那句 \BibleQuote{在起初天主创造了天地}，并注视\OriginalTerm{智慧}{Wisdom}——那\Quote{\OriginalTerm{元始}{Beginning}}，因为智慧也向我们说话\BibleRef{若 8:23}。另一位也同样注意这些话，把 \Quote{起初} 理解为受造物的开端；并如此领会：祂在起初造了，如同说，祂首先造了。在那些将 \Quote{在起初} 理解为 \Quote{祢以智慧创造了天地} 的人中，有人相信，那将要用来创造天地的质料，在那里被称为 \Quote{天地}；另一人相信，天地是已经形成和区别的本性；还有人相信，一个形成的本性是精神的，以 \Quote{天} 为名，另一个无形的本性，是物质的，以 \Quote{地} 为名。然而，那些以 \Quote{天地} 之名理解为尚未成形的质料，由此将要形成天地的人，他们彼此也并非一致地理解它；有人理解为那理智的与可感的受造物将由之完成的质料；另一人则理解为这只是为形成那可感的、有形巨体所需的质料，这巨体在其广阔的怀抱中容纳着这些可见的、各具本性的万物。同样，那些相信已经排序并安置妥当的受造物在此处被称为天地的人，也意见不一；有人以为是可见与不可见的全体；另一人则以为只是可见之物，在其中我们赞叹光明的天与黑暗的地，以及其中所包罗的万物。\par

% structure: p0070 source=h2 target=subsection reason=relative-heading-rank; base=h2

\subsection{第二十九章 论那些解释为 \Quote{祂首先造了} 的人之意见}

40. 谁若不这样去理解\Quote{在起初祂创造了}，而只把它当作\Quote{首先祂创造了}的意思，那么他只能把\Quote{天地}理解为天地的质料\OriginalTerm{质料}{matter}，即那普遍的，即理智的与形体的受造界。因为如果他以为这是指已经成形的宇宙，那么就可以正当地问他：\Quote{如果天主首先造了这个，之后祂又造了什么呢？}而在宇宙之后，他将发现没有别的，于是他不得不听到，哪怕不情愿：\Quote{这怎么能是首先的，如果以后什么都没有呢？}但是，当他说天主先造了无形的质料，然后造了有形的，这并不荒谬，只要他能分辨什么是以\OriginalTerm{永恒}{Eternity}为优先，什么是以时间为优先，什么是以选择为优先，什么是以起源为优先。以永恒为优先，如同天主在万物之先；以时间为优先，如同花先于果实；以选择为优先，如同果实优先于花；以起源为优先，如同声音先于曲调。这四种之中，我所提到的第一种和最后一种是很难理解的，中间两种则非常容易。因为要瞻仰祢的永恒，主啊，祢不变地创造可变的万物，因而在它们之先，这是一种罕见且过于高超的异象。谁有那么敏锐的思想，能不费很大力气就发现：声音如何先于曲调？因为曲调是有形的声音；一个无形的事物可能存在，但一个不存在的事物不能被定形。所以，质料先于由它而造的事物；这优先不是因为质料制造了它，因为质料更是被造的，也不是因为时间的间隔。因为我们不会在时间上先发出一些不唱歌的无形声音，然后再将这些声音调整或塑造成一首歌曲的形式，就像木头或银子，用来制造箱子或器皿那样。因为这样的材料在时间上也先于由它们造成的物品的形式；但唱歌时并非如此。因为当唱一首歌时，它的声音是在同一时刻被听到的；并没有一个先有的无形声音，然后再被形成一首歌。因为声音一发出就消逝了，你无法找到任何可以回忆起来并用艺术组合的残留物。因此，歌曲被包涵在它自身的声音中，那声音就是它的质料。因为这声音被定形，为能成为曲调；所以，如我所说，声音的质料先于曲调的形式，这种优先不是凭着任何使它成为曲调的能力；因为声音并不是曲调的作曲者，而是从身体发出，并服从于歌者的灵魂，以便从它形成曲调。它也不是在时间上为先，因为它与曲调一同发出；也不是在选择上为先，因为声音并不比曲调好，因为曲调不仅是一种声音，而是悦耳的声音。但它是在起源上为先，因为曲调被定形不是为了成为声音，而是声音被定形为了成为曲调。借着这个例子，让能理解的人明白：万物的质料是最先被造的，并被称为天地，因为天地是由它造成的。这并不是说它在时间上最先被造，因为万物的形式才产生了时间，而那时它是无形的；但如今，在时间内，它与它的形式一同被感知。而且，关于那个质料，我们不能说什么，除非将它说得好似在时间上在先，尽管它被思考为在后的（因为有形的事物确实优于无形的事物），并且有造物主的永恒作为它的在先，这样就能从虚无中造出可以形成某物的东西。\par

% structure: p0072 source=h2 target=subsection reason=relative-heading-rank; base=h2

\subsection{在意见的巨大分歧中，众人都应联合于\OriginalTerm{爱德}{charity}与属神的真理内。}

41. 在这众多真实意见的分歧中，愿真理本身产生和谐；愿我们的天主怜悯我们，使我们能合法地运用法律，\BibleRef{弟前 1:8} 这诫命的目的就是纯洁的\OriginalTerm{爱德}{charity}。如果有人因此问我：\Quote{祢的仆人\Person{梅瑟}这些话的意思，究竟属于哪一种？}这些并不是我\WorkTitle{忏悔录}中的话，我要向祢坦承：\Quote{我不知道；}然而我知道，除了那些我提过的、关于肉欲的（针对这些我已说出了我认为恰当的）之外，那些意见都是真实的。然而，祢\WorkTitle{圣经}中的这些话，并不会使怀着善望的弱小者惊惧，因为它以谦卑的方式处理少数高深的事，并以多种方式处理少数的事。但愿所有我承认在这些话中见到并说出真理的人，都彼此相爱，也同样爱祢、我们的天主、真理的泉源——如果我们渴望的不是虚空，而是真理；是的，让我们这样尊敬祢的这位仆人、这充满祢\OriginalTerm{圣神}{Spirit}的圣经的管家，要相信，当祢向他启示自己，他写下这些事时，他心中所愿的，就是其中最能发出真理之光并结出丰硕果效的那些意思。\par

% structure: p0074 source=h2 target=subsection reason=relative-heading-rank; base=h2

\subsection{可以设想，\Person{梅瑟}在他的话中，领悟了所能发现的任何真理。}

42. 因此，当一个人说：\Quote{他（\Person{梅瑟}）的意思和我一样，}另一个人说：\Quote{不，而是和我一样，}我想，我说的时候更虔诚：\Quote{如果两者都对，为什么不更该是两者的意思呢？}而如果有第三种真理，或第四种，如果有人在这些话中寻找完全不同的真理，为什么不能相信他（\Person{梅瑟}）看见了所有这些真理呢？毕竟，唯一的天主透过他调和了\WorkTitle{圣经}，使之适应许多人的领悟，他们将在其中看到真实但不同的事物？我当然——我从内心无畏地声明——如果我写任何要有最高权威的东西，我宁愿这样写：使任何人关于这些事所能领悟的任何真理，都能在我的话中引起回响，而不是我只清楚地写下一种真实的意见，以至排除了其他不会因谬误而冒犯我的意见。所以，我的天主，我不愿如此倔强，以至于不相信他（\Person{梅瑟}）从祢那里领受了这么多。的确，当他写下那些话时，他领悟并思考了我们所能发现的任何真理，是的，甚至是我们没能发现、也还不能发现的，尽管那些真理仍可在其中找到。\par

% structure: p0076 source=h2 target=subsection reason=relative-heading-rank; base=h2

\subsection{首先要探索作者的本意，然后阐明天主在此要揭示的属神真理。}

43. 最后，主啊，祢是天主，而非血肉之躯，即使人有所未悟，难道有什么能瞒过\BibleQuote{祢的善神}吗？祂必将\BibleQuote{引领我进入正直之地}。这正是祢借着那些话，正要向未来的读者启示的，尽管那位通过他传达这些话的人，在可能发现的许多解释中，只选定了一种。如果确实如此，就让他所思的比其余的更为崇高。但是对我们，主啊，请指出同样的意义，或者其他任何令祢喜悦的真实意义；这样，无论祢让我们知道的是祢启示给祢那位仆人的，还是借着同样的话偶然启示的别的意义，愿祢喂养我们，却不要让错误欺骗我们。请看，主，我的天主，我们为这寥寥数语写下了多少话——我恳求祢，何等之多啊！我们哪来的力量，多少世代的时间，才足以这样对待祢所有的书卷？所以，请准许我在此更简短地向祢忏悔，并选取一种真实、确定、良善的意义，就是祢将要启示的，尽管许多意义呈现出来，而事实上我也可能采纳许多；这是我忏悔的信仰：如果我说出祢的仆人所感受到的，那是正确而有益的，我就应当努力追求；如果我不能达到，至少我可以说出祢的真理藉其言语愿意向我说的话，就是它也曾向他说了它愿意说的一切。\par

\SourceRecord{Translated by J.G. Pilkington.；Nicene and Post-Nicene Fathers, First Series；Vol. 1.；Edited by Philip Schaff.；Buffalo, NY: Christian Literature Publishing Co.,；1887.；<http://www.newadvent.org/fathers/110112.htm>.}

% structure: p0001 source=h1 target=section reason=page-title

\section{忏悔录（卷十三）}

关于天主在创造万物中所彰显的良善，以及《创世纪》首句中所体现的\OriginalTerm{天主圣三}{Trinity}。以寓意的方式解释世界起源的故事（\BibleBook{创}{世纪} 1），并将其应用于天主为圣化并赐福于人所作的工作。最后，他结束了这部著作，恳求天主赐予永恒的安息。\par

% structure: p0003 source=h2 target=subsection reason=relative-heading-rank; base=h2

\subsection{第一章 他呼求天主，并立志敬拜祂。}

1. 我呼求祢，我的天主，我的仁慈，祢造了我，我虽把祢遗忘，祢却没有忘记我。我呼求祢进入我的灵魂，祢以祢在我内激起的渴望，预备它接受祢。求祢不要离弃呼求祢的我，因为祢在我呼求之前就已眷顾我，并以各种呼召急切地敦促我，使我从远处听到祢，归向祢，并呼求那召唤我的祢。主啊，因为祢已涂抹了我一切的恶功，使祢不将我远离祢所应得的报酬交在我手中；祢也预先预备了我一切的善功，使祢将祢造我时所应得的报酬交在祢手中；因为在我存在之前，祢已存在，我并不存在[任何事物]能让祢赐予存在。看，我存在，是出于祢的良善，预先预备了祢造我及用来造我的一切。因为祢既不需要我，我也不是那样的善，能有助于祢，我的主、天主；并非我可以这样事奉祢，好像祢在工作中会疲倦，或缺少我的协助祢的能力就会减少；也非如土地，我可以这样耕种祢，若我不耕种祢，祢就成为荒芜；而是让我事奉敬拜祢，使我从祢获得福祉；因为从祢那里，我成为能够承受福祉者。\par

% structure: p0005 source=h2 target=subsection reason=relative-heading-rank; base=h2

\subsection{第二章 一切受造物都因天主美善的满溢而存在。}

2. 因为祢的受造物是由于祢美善的满溢而存在，一种善，它既不能给祢带来丝毫利益，虽然由祢而来，却不与祢同等，但它仍能存在，因为它能由祢造成。祢起初所造的天地，从祢那里应得什么呢？愿祢以智慧所造的那些精神与形体本性，宣告它们从祢那里应得什么以依存于祢——就连那初始而未具形相的，各自按其种类，或精神或形体，趋向过度，远离祢而毫不相似（精神虽有未具形相，却比已成形体的物体更为高贵；形体虽有未具形相，却比完全虚无更为高贵），因此，它们作为未具形相者，若非通过同一圣言被召回祢的统一中，被赋予形相，并从祢、唯一至善者那里，成为完全美好的，它们将依赖于祢的圣言。它们从祢那里应得什么，竟成为未具形相的？因为若非由祢，它们连这未具形相也得不到。\par

3. 有形的质料从祢那里应得什么，竟成为无可见且未具形相的呢？\BibleRef{创 1:2}若非祢造了它，它连这无可见且未具形相也不会有；因此，既然它不曾存在，就不能从祢那里应得被造。或者，那初始的精神受造物，从祢那里应得什么，竟如深渊般黑暗地流动——与祢不相似，若非通过同一圣言转向那藉以受造者，并由祂光照而成为光，虽然不是同等地，却相称于那与祢同等的形象？因为就物体而言，存在并不等于美丽，否则它就不会变形；同样，就受造的精神而言，生活并不等于智慧地生活，否则它就会永不改变地智慧。然而，对它而言，恒常紧依祢才是好的，免得它远离祢时，失去那转向祢所获得的光，并倒退入一种类似黑暗深渊的光中。因为我们自己，就灵魂而言是一个精神受造物，曾远离祢、我们的光，在那种生活中\BibleQuote{有时是黑暗}\BibleRef{厄 5:8}；并且在我们所余留的黑暗中劳苦，直到在祢的独一者内，我们成为祢的正义，如同天主的群山。因为我们曾是祢的判决，如同巨大的深渊。\par

% structure: p0008 source=h2 target=subsection reason=relative-heading-rank; base=h2

\subsection{第三章 \BibleBook{创}{世纪} 1章3节——论\BibleQuote{光}——他按在精神受造物中所见来理解。}

4. 然而，祢在创造之初所说的：\BibleQuote{要有光！就有了光}\BibleRef{创 1:3}，我恰当地理解为指精神受造物；因为那时已存在一种生命，是祢能够光照的。但是，正如它未曾从祢那里应得成为这样一种能被光照的生命，同样，当它已经存在后，也未曾从祢那里应得被光照。因为它的未具形相也不能取悦于祢，除非它成为光——不是仅凭存在，而是藉着注视那光照的光，并依附于它；同样，它之所以生活，且幸福地生活，全赖祢的恩宠，而非其他任何事物；因着一种更好的转变，它转向那既不能变好也不能变坏者；唯有祢是如此，因为唯有祢单纯地存在，对祢而言，生活与幸福地生活并非两回事，因为祢自己就是祢自己的真福。\par

% structure: p0010 source=h2 target=subsection reason=relative-heading-rank; base=h2

\subsection{第四章 万物皆由天主的恩宠所造，并非由于祂需要受造之物而造。}

5. 所以，对于祢的美善——祢本身就是这美善——虽然这些事物要么从未存在过，要么一直停留在无形状态——祢创造它们，并非出于任何缺乏，而是出于祢丰盈的美善，约束它们，赋予它们形式，并非好像祢的喜乐因它们而得以圆满？因为祢是完美的，它们的不完美令祢不悦，因此它们被祢完美化，并蒙祢悦纳；但并非好像祢不完美，需要靠它们的完美来达成祢的圆满。因为祢的\OriginalTerm{圣神}{Holy Spirit}运行在水面上，\BibleRef{创 1:2} 并非由水托起，好像安息在水上一般。因为那些祢的圣神安息在谁身上，祂便使谁在祂内安息。\BibleRef{户 11:25} 但是，祢那不朽且不变的意志——它在自身内便完全自足——运行在祢所造的某种生命之上；对这种生命而言，活着并不等同于幸福地活，因为它虽然也在自身的黑暗中流动着活着；为此它还需要皈依那造它的祂，并在\Quote{生命的泉源}里越来越丰盛地生活，在祂的光明中\Quote{看见光明}，得以完美化、蒙光照、获得幸福。\par

% structure: p0012 source=h2 target=subsection reason=relative-heading-rank; base=h2

\subsection{第五章 他在\BibleBook{创}{世纪}首两节中认出天主圣三}

6. 看啊，如今，\OriginalTerm{天主圣三}{Trinity}在隐微中向我显现——祢，我的天主，正是这圣三，因为祢，圣父啊，在我们的智慧的\OriginalTerm{元始}{Beginning}——那即是祢的智慧，由祢所生，与祢同等、同永恒——也就是在祢的圣子内，创造了天地。关于天上的天，关于无形无象的大地，关于幽暗的深渊，随着它灵性畸形所导致的游荡缺失，我们已经说得很多了；倘若它未曾皈依那赐予它生命的祂——它的生命虽如此，却因祂的光照而成为美好的生命，并成为后来安置在水与水之间的那层天的天——那又会如何呢？于是，我如今认定，\InnerQuote{天主}之名指的是创造这一切的圣父；\InnerQuote{元始}之名，则指在祂内创造这一切的圣子；我既相信我的天主就是天主圣三，便在祂的圣言中进一步寻觅，看啊，祢的圣神运行在水面上。看啊，我的天主，圣父、圣子、圣神——一切受造物的创造者，这便是天主圣三。\par

% structure: p0014 source=h2 target=subsection reason=relative-heading-rank; base=h2

\subsection{第六章 为何提及天地之后才提及圣神}

7. 但是，祢这位讲说真理的光明啊，原由何在？我向祢举起心，莫让它教我虚妄之事；驱散它的黑暗，我恳求祢，因着我们的母亲\OriginalTerm{爱德}{charity}，恳求祢告诉我：为何圣经在提及天、无形无象的地和深渊上的黑暗之后，才终于提到祢的圣神呢？是否因为关于祂，应当说祂\Quote{运行}，而若不先提及那\InnerQuote{之上}，这说法便无从理解？因为祂既不是\Quote{运行}在圣父之上，也不是\Quote{运行}在圣子之上，若祂\Quote{运行}在虚无之上，便不能正确地说祂\Quote{运行}。所以，必须先提到祂所\Quote{运行}其上的\InnerQuote{那物}，然后才提到祂，而提及祂，若非以\Quote{运行}的方式，便不恰当。然而，为何以其他方式提及祂，就不恰当呢？为何非得说祂\Quote{运行}呢？\par

% structure: p0016 source=h2 target=subsection reason=relative-heading-rank; base=h2

\subsection{第七章 圣神引领我们归向天主}

8. 因此，凡是能够的人，现在要用心领悟祢的\OriginalTerm{宗徒}{apostle}的话，因为他说：祢的爱\BibleQuote{借着所赐与我们的圣神，已倾注在我们心中了；}\BibleRef{罗 5:5} 又在论到\Quote{神恩}的地方，他教导并指示我们一条更卓越的爱德之路；他还为我们向祢屈膝，好使我们能领悟那超越一切知识的基督之爱。\BibleRef{弗 3:14} 所以，从一开始，祂就超越一切地\Quote{运行在水面上}。我将这事告诉谁呢？如何描述那重压着向下坠入陡峭深渊的贪欲的重量？而爱德又怎样透过祢那\Quote{运行在水面上}的圣神，将我们重新提升？我将这事告诉谁呢？又如何述说呢？因为其中并无我们可以沉没和浮现的处所。有什么能比这更相似，又更不相似？这些是情感，是爱；我们精神的污浊，因着对俗世牵挂的爱而向下流逝；而祢的圣洁，则透过对无忧之自由的挚爱，将我们向上提升；好使我们能将心提升到祢面前，在那祢的圣神\Quote{运行在水面上}的地方；并使我们的灵魂穿越那无实质的水之后，达到那超越一切的安息。\par

% structure: p0018 source=h2 target=subsection reason=relative-heading-rank; base=h2

\subsection{第八章 除天主外，没有任何事物能赐给有理性的受造物幸福的安息}

9. 天使堕落，人的灵魂堕落，它们以此指明在那黑暗深渊中的无底坑，为整个精神受造界预备，除非祢从起初说：\BibleQuote{要有光}，就有了光，祢天上之城中一切顺服的理智都依附祢，在祢的圣神内安息，这圣神不变地\BibleQuote{运行}在一切可变易的万物之上。否则，即使天外之天本身也会是一个黑暗的深渊，而如今它在主内却是光明的。因为即使在那些堕落而远离祢的神灵的悲惨动荡中，当他们被剥去祢光的衣裳，发现了自己的黑暗时，祢充分显明了祢造有理性的受造物是何等尊贵；对于它而言，任何低于祢的事物都不足以赐予幸福的安息，甚至它自身也不能。因为祢，我们的天主，将照亮我们的黑暗；我们的光明衣裳来源于祢，那时我们的黑暗将如日中天。将祢自己赐给我，我的天主，将祢自己归还给我；看，我爱祢，如果这爱太微小，就让我更强烈地爱祢。我无法度量我的爱，好能知道我的生命中尚有几何欠缺，直到它奔入祢的怀抱，不再转离，直到藏匿于祢圣容的秘密处。我只知道这一点：不在祢内，我便有祸——不但在外，即使在我自己内也是如此；凡不是我的天主的一切丰盛，于我都是贫穷。\par

% structure: p0020 source=h2 target=subsection reason=relative-heading-rank; base=h2

\subsection{第9章 为什么圣神仅被称为\BibleQuote{运行在水面上}}

10. 但是，难道圣父或圣子没有\BibleQuote{运行在水面上}吗？如果我们按空间意义理解，像论及身体那样，那么圣神也没有运行；但如果我们理解为天主性不变易的超绝性高居于一切可变易之物之上，那么圣父、圣子和圣神都\BibleQuote{运行在水面上}。那么，为什么只就祢的圣神这样说呢？为什么唯独论及祂呢？仿佛祂处于某处，而祂本不处于某处，且经上唯独记载祂是祢的恩赐？我们在祢的恩赐中安息；在那里我们享受祢。我们的安息就是我们的处所。爱在那里提携我们，祢的圣神从死亡的门户提升我们的卑微。在祢的圣意中才有我们的平安。身体凭自身的重量趋向自己的处所。重量并不只是向下，而是趋向自己的处所。火向上，石向下。它们被自身的重量推动，寻求自己的处所。油倒在水下，会升到水面上；水倒在油上，会沉到油下。它们被自身的重量推动，寻求自己的处所。它们失序时，便躁动不安；恢复秩序时，便安息下来。我的重量就是我的爱；爱载着我，不论我往哪里去。因祢的恩赐，我们被点燃，向上提升；我们内心变热，向前迈进。我们攀登那些在我们心中的祢的道路，唱着上行之歌；我们内心因祢的圣火而炽热，我们前行，因为我们是向耶路撒冷的平安上升；因为我对他们说\BibleQuote{让我们进入主的殿宇}时，我欢欣踊跃。在那里，祢的圣意安置了我们，使我们除永居其中外，别无所愿。\par

% structure: p0022 source=h2 target=subsection reason=relative-heading-rank; base=h2

\subsection{第10章 若非因着天主的恩赐，无一物可兴起}

11. 幸福的受造物啊，虽然它本身并非祢，但它一被造，便无时间间隔，因着祢的恩赐——这恩赐运行在一切可变之物上——而被祢那声\BibleQuote{要有光，就有了光}的召唤所提升。而在我们内，则有时间的差别，我们曾经是黑暗，如今却是光明的；见：\BibleRef{弗 5:8} 但对于它，只说它若不被光照会是怎样的。这样说，仿佛它先前是飘忽而黑暗的；这样，它被改变的原因便显现出来——就是说，它转向那永不会暗的光，便成了光明。能领悟的，就领悟吧；不能的，就向祢求问。他何必烦扰我，好像我能照亮\BibleQuote{每个来到这世界的人}似的？\par

% structure: p0024 source=h2 target=subsection reason=relative-heading-rank; base=h2

\subsection{第11章 在人内的天主圣三象征，即存在、认识、意愿，从未被彻底检视}

12. 我们中谁能领悟全能的天主圣三？然而谁又不谈论它，如果所谈论的确实是它的话？难得有灵魂在谈论它时，知道自己谈论的是什么。人们争论不已，但若无内心的平安，无人能看见那景象。我愿人们能思考自身内的这三件事。这三件事与圣三截然不同；但我所谈论的，是他们可以用来操练自己、检验自己，并感受它们与圣三相异程度的那些事。我所说的三件事是：存在、认识、意愿。因为我存在，我认识，我愿意；我存在着认识、愿意着；我认识自己存在和愿意；我愿意存在和认识。因此，在这三者中，谁能看出这三者如何结合成不可分离的生命——即一个生命、一个心志、一个本质；最后还有那不可分离的区分，毕竟还是一个区分。这分明展现在人面前；让他内观，查看，然后告诉我。但当他发现并能说出一些关于这三者的事时，切莫以为他发现了那超越这三者的不变者，祂存在不变，认识不变，意愿不变。至于是否因这三者，而在它们所在之处也有了三位一体；或这三者在每一位中，因而三者各每一位所俱有；或者两种方式同时并存，奇妙地，单一而又殊异，自身是自身的界限，却又无限量；由此祂存在，为自己所认识，并自足，永远是同一的，因其统一性之丰盛广大——谁能轻易构想？谁又能以任何方式表达？谁能鲁莽地在这事上妄下断言？\par

% structure: p0026 source=h2 target=subsection reason=relative-heading-rank; base=h2

\subsection{第12章 对\BibleBook{创}{世纪}第一章的寓意解释：论教会的起源及其崇拜}

13. 我的信德啊，继续你的宣认，向主你的天主说：\Quote{圣、圣、圣，我主天主，我们因祢的名受洗——因父、子、圣神之名；我们因祢的名施洗——因父、子、圣神之名}\BibleRef{玛 28:19}，因为在祂的基督内，天主也在我们中间创造了天地，就是祂教会中的属神和属肉体的人。的确，我们的地在接受\BibleQuote{道理的规范}\BibleRef{罗 6:17}之前，也是空虚混沌，我们被愚蒙的黑暗笼罩。因为祢因人的罪孽惩治人，\Quote{祢的判断如同深渊。}但祢的圣神\Quote{运行在水面上，}\BibleRef{创 1:3}，祢的仁慈没有舍弃我们的不幸，祢说：\Quote{要有光，} \Quote{你们悔改吧，因为天国临近了。}\BibleRef{玛 3:2} 你们悔改吧，要有光。因为我们的灵魂在我们内忧闷，主啊，我们从约旦地，从那座与祢同等、却为我们成为小的山，记起了祢；当我们厌恶自己的黑暗时，我们转向祢，\Quote{就有了光。}\BibleQuote{看，我们从前原是黑暗，但现在在主内却是光明。}\BibleRef{弗 5:8}\par

% structure: p0028 source=h2 target=subsection reason=relative-heading-rank; base=h2

\subsection{第十三章 人的更新不在此世完成。}

14. 然而如今还是\BibleQuote{因着信德，非由目睹}\BibleRef{格后 5:7}，因为\BibleQuote{我们得救是在于望德；但所希望的若已看见，就不是希望了。}\BibleRef{罗 8:24} 深渊与深渊响应，但在\Quote{祢瀑布的响声。}中。那时，那说我\BibleQuote{不能对你们讲论，如同对属神的人，而是如同对属肉的人}\BibleRef{格前 3:1}的，连他自己也还不以为自己已经夺得了，而是忘记背后的，奋力向前奔驰\BibleRef{斐 3:13}，负重呻吟；他的灵魂渴望生活的天主，如牝鹿渴慕溪水，说：\Quote{我何时来？}\Quote{渴望穿上那来自天上的居所，}\BibleRef{格后 5:2}；并向这下方的深渊呼求说：\BibleQuote{不要符合这世界，而要以更新的心思变化自己。}\BibleRef{罗 12:2} 又说：\Quote{在思悟上不要做孩子，但在恶事上做赤子，}\Quote{在思悟上应成为成年人；}\BibleQuote{无知的迦拉达人，谁迷惑了你们？}\BibleRef{迦 3:1} 但现在不是用他自己的声音，而是用祢的声音，就是祢从天上派遣圣神的声音；\BibleRef{宗 2:19} 藉着那位\BibleQuote{上升高天}\BibleRef{弗 4:8}，敞开祂恩赐的闸门\BibleRef{拉 3:10}，使祂河流的猛力悦乐天主的城。因为为祂而叹息的是\BibleQuote{新郎的朋友}\BibleRef{若 3:29}，他现在虽然已领受了圣神的初果，却仍在内心叹息，等待义子的名分，即他肉身的救赎；\BibleRef{罗 8:23} 他叹息是为祂，因为他是新娘的肢体；他忌妒是为祂，因为他是新郎的朋友；\BibleRef{若 3:29} 他忌妒是为祂，不是为己；因为他是用祢\OriginalTerm{瀑布}{waterspouts}的声音，不是用他自己的声音，去呼求那另一个深渊，他因忌妒而担心，怕他们如蛇以诡诈诱惑了厄娃，他们的心思会败坏，失去我们在新郎祢独生子内的纯朴。当\Quote{我们要看见祂的本来面目}时，那将是何等美丽的光明！那时，那些\Quote{昼夜以眼泪当饮食，他们终日向我说：你的天主在哪里？}的眼泪，都将成为过去。\par

% structure: p0030 source=h2 target=subsection reason=relative-heading-rank; base=h2

\subsection{第十四章 夜与黑暗之子成为光与白日之子。}

15. 所以我也说：我的天主，祢在哪里？看，祢在哪里！我在祢内稍得喘息，当我独自以欢欣赞颂的声音，即守节者的欢呼，倾诉我的灵魂时。然而它还是\Quote{沮丧}，因为它复陷而为深渊，或更确切地说，它感觉自己仍是一个深渊。祢为了照亮我夜间的步履而点燃的信德，对它说：\BibleQuote{我的灵魂，你为何沮丧？为何在我内忧闷？希望于天主！}祂的\BibleQuote{言语是我步履前的明灯。}希望并忍耐，直到黑夜——恶人的母亲——过去，直到主的愤怒过去\BibleRef{约 14:13}，我们从前也曾是其中的子女，有时是黑暗，那黑暗的残余我们还在身上带着，因为罪而死亡的身体，\BibleRef{罗 8:10} \BibleQuote{等到天亮，黑影消散时。}\BibleRef{歌 2:17} \BibleQuote{希望于上主。}清晨，我将站立在祢面前，瞻仰祢；我将永远称谢祢。清晨，我将站立在祢面前，看见\BibleQuote{我面容的救恩，}我的天主，祂也要藉着那住在我们内的圣神，复活我们可朽的身体，\BibleRef{罗 8:11}，因为祂出于仁慈，曾运行在我们的幽暗混沌的深渊之上。因此，我们在这旅途中获得了\OriginalTerm{抵押}{earnest}\BibleRef{格后 1:22}，使我们现今就成为光明，虽然我们仍在\BibleQuote{因望德得救，}\BibleRef{罗 8:24}并且是光明之子，白日之子——而不是黑夜之子，也不是黑暗之子，虽然我们曾经是。在人类知识尚不确定的今世，在我们与那些人之间，唯有祢在区分，祢洞彻人心，称光为昼，称暗为夜。\BibleRef{创 1:5} 因为除了祢，谁能辨识我们呢？但我们有什么不是领受于祢的呢？\BibleRef{格前 4:7} 出自同一团泥，做成了贵重的器皿，也做成了卑贱的器皿。\BibleRef{罗 9:21}\par

% structure: p0032 source=h2 target=subsection reason=relative-heading-rank; base=h2

\subsection{第十五章 穹苍与上层工程的寓意解释，第6节。}

16. 或者除了祢、我们的天主以外，谁为我们造了那在祢的天主圣经中凌驾我们之上的\OriginalTerm{穹苍}{firmament}之权威\BibleRef{创 1:6}？正如经上所说，天要像\OriginalTerm{书卷}{scroll}被卷起；而今它却如皮囊铺展在我们上方。因为祢的天主圣经具有更崇高的权威，因为那些祢借以将圣经分施给我们的凡人，经历了死亡。主啊，祢知道，祢知道人类因罪恶而成为可朽的时，祢如何用皮衣给他们蔽体。由此，祢铺展了祢的圣书之穹苍如同皮囊；就是说，祢那些和谐一致的言语，祢借着凡人的职事将它们散布在我们之上。因为他们本身的死亡，那在他们的言论中所建立的权威穹苍就被他们树立得更加崇高而铺展开来，这穹苍在他们活着时，尚未如此显著地铺展开来。那时祢尚未将天铺展如皮囊，尚未将他们的死讯传扬四方。\par

17. 主啊，让我们\BibleQuote{仰望祢手指所造的诸天}；求祢从我们眼睛除净祢用以遮蔽它们的迷雾。那里有祢的见证，赐智慧给微小者。我的天主，求祢由婴孩和幼雏的口，完成祢的赞美。我们不知道有其他书籍，如此摧毁骄傲，如此摧毁敌人和那辩护者，即抗拒祢的和解、捍卫自己罪恶的人。主啊，我不知道，我不知道有其他如此\Quote{纯洁}的言语，如此能说服我告解，使我的颈项驯服于祢的轭，并邀请我白白地事奉祢。良善的父啊，让我明了这些事；求祢将这恩宠赐予我，我这被置于其下的人；因为祢为那些被置于其下的人建立了这些。\par

18. 还有其他\Quote{水}在\Quote{这}\Quote{穹苍}\Quote{之上}的，我相信是不死不灭的，远离尘世朽坏的。愿他们赞美祢的名——那些超天的人民，祢的天使，他们无须举目仰望这穹苍，也无须通过阅读来获知祢的\OriginalTerm{圣言}{Word}——愿他们赞美祢。因为他们常仰瞻祢的圣容\BibleRef{玛 18:10}，并在其中不须时间的音节而读知祢永恒旨意所愿意的一切。他们阅读，他们选择，他们爱慕。他们不断地阅读；而所读的永不过去。因为借着选择与爱慕，他们阅读的正是祢旨意的不变性。他们的书卷不会封闭，卷轴也不会被卷起\BibleRef{依 34:4}，因为祢自己就是他们的这一切，而且是永恒的；因为祢将他们安置在这穹苍之上，这穹苍是祢为下面软弱百姓所坚定的，使他们能仰望而认识祢的仁慈，按时宣布那创造时间的祢。\BibleQuote{主，祢的仁慈达于天际，祢的忠信高及云霄。}云彩逝去，但天穹永存。祢圣言的宣讲者从今生逝去，进入另一生命；但祢的圣经却传扬于万民，直到世界末日。是的，天地都会过去，但祢的话不会过去\BibleRef{玛 24:35}。因为书卷将被卷起\BibleRef{依 34:4}，其铺展其上的青草及其美丽都要逝去；但祢的圣言却永远常存，如今却透过云彩的暗影，借天镜向我们显现，还不是本来面目\BibleRef{格前 13:12}；因为我们虽然已是祢圣子的钟爱者，但将来如何，还未显明\BibleRef{若一 3:2}。祂由我们肉身的窗棂张望\BibleRef{歌 2:9}，祂善于言谈，且燃烧了我们，我们便追随祂的芬芳\BibleRef{歌 1:3}。但\BibleQuote{当祂显现时，我们必要相似祂，因为我们要看见祂实在怎样。}\BibleRef{若一 3:2}主啊，我们要看见祂实在怎样，虽然时刻还未到。\par

% structure: p0036 source=h2 target=subsection reason=relative-heading-rank; base=h2

\subsection{第十六章 惟有\OriginalTerm{不变之光}{Unchangeable Light}认识自己。}

19. 因为祢怎样，惟有祢自己知道，祢是不变地存在，不变地知道，不变地愿意。祢的\OriginalTerm{本体}{Essence}也以不变的方式知道和愿意；祢的知识以不变的方式存在和愿意；祢的意志以不变的方式存在和知道。在祢眼中，这似乎不公平：正如\OriginalTerm{不变之光}{Unchangeable Light}认识自己，它也应被受光照而变化之物同样认识。因此，我的灵魂在祢面前如同\BibleQuote{无水之地}，因为正如它不能自己光照自己，它也不能自己满足自己。因生命的泉源与祢同在，就像在祢的光中我们看见光。\par

% structure: p0038 source=h2 target=subsection reason=relative-heading-rank; base=h2

\subsection{第十七章 海与结实之地的寓意解释——第九及十一节。}

20. 谁将这些苦水聚集在同一处？因为他们都拥有同一目的，即今世与尘世的幸福，他们行事莫不为此，尽管他们因种种顾虑而摇摆不定。主啊，除了祢，谁能说：让水汇聚一处，让旱地出现，就是那\BibleQuote{渴慕祢的}旱地？因为海也属于祢，是祢所造，旱地也是祢的双手所预备的。因为不是人的意志之苦毒，而是水的汇聚，才被称为海；因为祢甚至掌管人灵魂的邪恶欲望，并为它们划定界限，允许它们前进多远，以及让它们的波浪彼此冲击；如此，祢按照统御万有的秩序，使之成为海。\par

21. 至于那些渴慕祢，并出现在祢面前的灵魂（由于其他界限而与海洋的社会分开），祢以隐秘而甘甜的泉水浇灌他们，使大地结出果实，并且，主，天主，祢如此命令，我们的灵魂便会按其种类萌发出仁慈的行为——在救助近人肉身的急需中爱他，自身内蕴有依其模样的种子，当我们出于自身的软弱而怜悯甚至救助贫困者时；以我们希望若身处同样困境时他人也会援助我们的同样方式帮助他们；不仅是在容易的事上，如同\Quote{结种子的蔬菜}，还要在我们援助的保护中，凭我们自身的力量，如同结果子的树木；那就是，将蒙受冤屈者从强权者手中解救出来的善行，并藉正义审判的强大力量为他提供护庇之所。\par

% structure: p0041 source=h2 target=subsection reason=relative-heading-rank; base=h2

\subsection{论天上的光体和星辰——论白日与黑夜，第14节}

22. 主啊，是的，我恳求祢，愿它们兴起，如祢所造，如祢赐予喜乐和能力——愿\BibleQuote{真理由地上发出，正义从天俯视}，并应有\BibleQuote{穹苍中的光体}\BibleRef{创 1:14}。让我们分饼给饥饿的人，引领无家可归的穷人来到我们家中\BibleRef{依 58:7}。让我们给赤身的人衣服穿，不要忽视我们的骨肉。正如大地生出这些果实，看，这是好的\BibleRef{创 1:12}；也让我们的短暂光芒喷薄而出\BibleRef{依 58:8}；并让我们，从这行动的较低果实中，享有对上界默观与\OriginalTerm{生命之言}{Word of Life}的喜乐，让我们成为世界中的光\BibleRef{斐 2:15}，紧依祢圣经的穹苍。因为在其中祢向我们显明，使我们能像区分白日与黑夜那样，区分可理解的事物与感觉的事物；或者区分那些灵魂：有些被赋予理智的事物，有些被赋予感觉的事物；这样，如今不仅祢在祢审判的隐秘中，如同穹苍造成之前那样，将光与暗分开，而且祢的灵性子女们，被安置排列在同一穹苍中（祢的恩宠彰显于普世），也能光照大地，将白日与黑夜分开，作为时节的标记；因为\BibleQuote{旧的已经过去}，且\BibleQuote{看，一切都成了新的}\BibleRef{格后 5:17}；又因\BibleQuote{我们的救恩现今比我们初信时更近了}\BibleRef{罗 13:11}；因\BibleQuote{黑夜深了，白日已近}\BibleRef{罗 13:11}；又因祢将冠以恩泽之年，\BibleQuote{派遣祢的工人收割祢的庄稼}\BibleRef{玛 9:38}，在那庄稼的播种中，他人劳苦了，又\BibleQuote{派遣去另一田地，到末时才收割}\BibleRef{玛 13:39}。祢这样恩准祈求者的祈祷，祝福义人的岁月\BibleRef{箴 10:6}；然而祢始终如一，在祢永无尽竭的岁月中，祢为我们的流逝岁月预备仓廪。因为祢凭永恒的措置，在适当的时节将天上的恩泽赐给大地。\par

23. 的确，这人从圣神蒙赐智慧的言语，如同较大的光体，为那些喜爱明显真理之光的人，如在白日的开始；另一人却从同一圣神蒙赐知识的言语，如同较小的光体；再一人赐信心；另一人治病的奇恩；另一人行奇迹；另一人说先知话；另一人辨别神类；另一人说各种语言。这一切如同星辰。因为这一切都是同一的圣神所运行，随祂的旨意，个别分配与人\BibleRef{格前 12:8}；并显明地呈现星辰，为众人的益处\BibleRef{格前 12:7}。然而，知识的言语——其中蕴含一切圣事，这些圣事如同月亮在各时期变化，以及恩宠的其他概念，它们依次被列为星辰——由于它们不及前述白日所欢欣的那智慧的辉煌，仅为黑夜的开始所需。因为它们为那些祢最通达的仆人不能对他们讲论属灵之事，只能讲论属血肉之事的人\BibleRef{格前 3:1}——即那位在成全人中讲论智慧的人所必需\BibleRef{格前 2:6}。但本性之人，如同在基督内的婴孩——吃奶者——直到他能适应固体食物，他的眼睛能注视太阳时，他不应停留在自己荒凉的黑夜中，而应满足于月亮和星辰的光。我们全智的天主，祢在祢的圣经——祢的穹苍中，向我们讲述了这些事，好使我们在可赞叹的默观中，虽然目前仍在标记、时期、日子和年岁中，却能辨识一切。\par

% structure: p0044 source=h2 target=subsection reason=relative-heading-rank; base=h2

\subsection{人人都应成为天上穹苍中的光体}

24. 但先，\BibleQuote{你们要洗涤，要自洁；}从你们的灵魂中，从我眼前铲除邪恶，使旱地可以出现。\BibleQuote{学习行善，为孤儿申冤，为寡妇辩护，}好使大地生出青草作食物，和结果子的树木；\BibleQuote{来吧，让我们彼此辩论——主说，}\BibleRef{依 1:18}好使天上的穹苍中有光体，照耀大地。\BibleRef{创 1:15}那个富人问善师，他该做什么才能获得永生。\BibleRef{玛 19:16}让那位善师——他以为他只是个人，别无其他——告诉他（但他是\Quote{善}的，因为他是天主）——让祂告诉他，若他愿意\BibleQuote{进入生命}，就必须\BibleQuote{遵守诫命；}让他从自己身上驱除恶意与邪恶的苦涩；\BibleRef{格前 5:8}他不可杀人，不可奸淫，不可偷盗，不可作假见证；好使旱地出现，并生出孝敬父母和爱近人的苗芽。\BibleRef{玛 19:16}这一切，他说，我从小就都遵守了。那么，既然土地肥沃，为何还有这么多荆棘呢？你去，根除贪婪的木质丛林；变卖你所有的，施舍给穷人，你必有宝藏在天上；然后来跟随主，\BibleQuote{你若愿意是成全的，}与那些祂在他们中间讲论智慧的人结合，祂知道如何分配给白日和黑夜，使你也知道，为你在天上的穹苍中也照样有光体，这不会发生，除非你的心在那里；\BibleRef{玛 6:21}同样，那也不会发生，除非你的宝藏在那里，正如你从善师那里听到的。但贫瘠的土地忧闷了，\BibleRef{玛 19:22}荆棘窒息了圣言。\BibleRef{玛 13:7}\par

25. 但你们，\BibleQuote{特选的种族，你们世上懦弱的，}\BibleRef{伯前 2:9}\BibleRef{格前 1:27}你们为了\BibleQuote{跟随主}而舍弃了一切，去追随祂，并\BibleQuote{羞辱那坚强的；}\BibleRef{格前 1:27}去追随祂，\BibleQuote{美丽的脚步，}\BibleRef{依 52:7}\BibleQuote{在穹苍中发光，}\BibleRef{达 12:3}好使诸天宣扬祂的光荣，区分完善者的光——虽非如天使之光——与微小者的黑暗，尽管他们并非被轻视者。光照全地，让白昼因太阳而光明，向白昼传述智慧的言语；让黑夜因月亮而闪耀，向黑夜传报知识的言语。月亮和星辰为黑夜发光，黑夜却不能遮蔽它们，因为它们按各自的程度照亮黑夜。因为看，天主仿佛在说：\BibleQuote{在天空的穹苍中要有光体。}忽然从天上有响声下来，好像一阵大风吹过，又有舌头如火焰显现出来，分开落在他们各人头上。\BibleRef{宗 2:3}于是，在天空的穹苍中有了光体，拥有生命的圣言。\BibleRef{若一 1:1}你们要到各处去，圣善的火，美丽的火；因为你们是世界的光，不是被放在斗底下。\BibleRef{玛 5:14}你们所依附的那位被高举了，也高举了你们。到各处去，叫万民认识你们。\par

% structure: p0047 source=h2 target=subsection reason=relative-heading-rank; base=h2

\subsection{第二十章 论水中的爬虫和天空的飞鸟（20节）——考虑圣洗圣事}

26. 也让海孕育并生出你的作品，让水生出有生命的蠕动之物。\BibleRef{创 1:20}因为你，\BibleQuote{由凡俗中取出珍宝的，}\BibleRef{耶 15:19}已成了天主的口，祂藉此说：\BibleQuote{让水生出，}不是大地所生出的生物，而是有生命的蠕动之物，以及在地上飞翔的飞鸟。因为你的圣事，天主啊，藉着你圣者的服务，已穿越了世俗诱惑的波涛，为以你的名、你的圣洗圣事教训外邦人。在这些事中，成就了许多伟大的奇事，如大鲸；你使者的声音飞翔于大地之上，接近你经书的穹苍；那穹苍作为权威被置于他们之上，他们无论往哪里去，都要在其下飞翔。因为\BibleQuote{这不是语，也不是言，是听不到他们声音的；}看，他们的声音\BibleQuote{却传遍普世，他们的言语直达地极，}因为你，主啊，以祝福繁衍了这些事。\BibleRef{创 1:4}\par

27. 我是在说谎吗？或是把在天穹苍中这些事的清楚知识，与在翻腾大海中、在天穹苍之下的物质作为相混合混淆，却不加以区分？因为那些知识是坚实明确的，不因生殖而增长，有如智慧和知识的光体；然而，同样这些事的物质运作却是多样多变的；一件事物从另一件生长出来，因你的祝福而繁衍，天主啊，你慰藉了凡人感官的挑剔；以致在我们心灵的认知中，同一件事，可以藉着身体的种种动作，以多种方式阐明和表达。这些圣事是水所孕育的，但却是在你的圣言中。与你真理的永恒隔绝的民众之需要，产生了它们，但却是在你的福音中；因为正是水本身将它们抛弃出来，水的苦涩软弱，正是这些事在你的圣言中被派遣的原因。\par

28. 你造的一切都是美好的，但看，你是无可言喻地更美好，是你创造了一切；倘若亚当没有从你堕落，海的咸味就绝不会从他流出——那人类如此深沉的好奇、喧嚣的傲慢和动荡不安；这样，你的分施者就不必在众多水中，以物质和可感的方式，施行奥秘的作为和言语。因为如今这些爬行和飞行的受造物呈现在我的思想中，人们藉着物质的圣事得以受教、入门并受约束，但如果他们的灵魂没有更高的灵性生命，如果在入门训言之后，不向往完善，他们便不能进一步得益。\par

% structure: p0051 source=h2 target=subsection reason=relative-heading-rank; base=h2

\subsection{第廿一章：论活灵魂、雀鸟与鱼类（第24节）——顾及圣体圣事。}

29. 因此，在祢的圣言中，不是海洋的深处，而是那从诸水之苦味中分开的旱地，生出那有生命的爬行和飞行生物，\BibleRef{创 1:20} 而是那活生生的灵魂本身。\BibleRef{创 2:7} 因为现在它已不需要圣洗，不像外教人所需要的，也不像它自己当初被水覆盖时一样——因为没有别的进入天国的入口，\BibleRef{若 3:5} 祢既已规定这应是那入口——它也不寻求能产生信德的大奇迹作为；因为现在信德的大地既已从因不信而变苦的海水之中分开，便不是那种：除非看见征兆和奇迹，便不肯相信，\BibleRef{若 4:48} 且\BibleQuote{语言不是为信的人作证据，而是为不信的人。}\BibleRef{格前 14:22} 那时，祢安置在水面以上的大地，也不需要那由祢的圣言命水所生的雀鸟种类。请祢藉着祢的使者，将祢的圣言传到那地。因为我们讲述他们的工作，但实在是祢在他们内工作，以便他们在其中作出一个活灵魂。大地生出它，因为正是大地使他们在那灵魂中作出这些事的原因；正如海洋曾是他们做工于那些有生命的游行生物和天空之下飞行的雀鸟上的原因，现在大地对这些已无需要；虽然它仍以那从深渊中取出的鱼为食，就是在祢为信者所预备的餐桌之上。因为祂从深渊中被举扬，正是为喂养旱地；而雀鸟，虽然生于海中，却在地上蕃衍。因为在福音作者最初宣讲时，人的不信是显见的原因；但信者也被劝勉，并且他们日日多方面地蒙受祝福。但活灵魂源于大地，因为克制自己不爱世俗，只对那些已在信者中的人有益，好使他们的灵魂为祢而活，这灵魂曾在享乐中虽生犹死，\BibleRef{弟前 5:6} —— 在致死的享乐中，主啊，因为祢是纯洁心灵的生机之乐。\par

30. 所以现在，让祢的仆役们在地上工作——不要像在不信之水中那样，藉着奇迹、圣事和奥秘的言语来宣讲和说话；在那里，无知——惊奇之母——可能因惧怕那些隐秘标记而专注于它们。因为这就是亚当的子孙、那些忘记祢的人进入信德的途径；当他们躲避祢的面，\BibleRef{创 3:8} 且成为幽暗的深渊时。但要让祢的仆役们如同在旱地上工作一样，与那大深渊的旋涡分开；让他们作信者的榜样，在他们面前生活，并激励他们效法。因为这样，人们听见，就不单是为了听见，也是为了实行。寻求主，你们的灵魂就将生活，好使大地生出活灵魂。\BibleQuote{不可与此世同化。}\BibleRef{罗 12:2} 你们要克制自己，远离它；灵魂因躲避它所贪爱而致死的事物生活。你们要克制自己，远离那放肆的骄傲之狂暴，远离那怠惰的奢华之逸乐，远离那知识的虚名；好使野兽被驯服，牲畜被制服，蛇蝎无害。因为这些就是死亡灵魂的动态；也就是说，骄傲的高傲、淫欲的快乐、好奇的毒害，是死亡灵魂的动态；因为灵魂的死亡并非失去一切动态，而是因离弃了生命之泉而死，于是被这短暂的世界所接纳，并与之同化。\par

31. 但是天主，祢的圣言是永生之泉，永不过去；因此那背离受到约束，乃因祢的圣言对我们说：\BibleQuote{不可与此世同化。}\BibleRef{罗 12:2} 为使大地在生命之泉中生出活灵魂——这个灵魂由祢的福音作者们、藉着效法祢的基督的追随者，在祢的圣言中受到约束。\BibleRef{格前 11:1} 因为这是各从其类；因为人被朋友激励去效法。他说：\BibleQuote{你们要像我一样，因为我也像你们一样。}\BibleRef{迦 4:12} 因此，在活灵魂中，将有善良的野兽，在于行动的温良。因为祢曾命令说：\BibleQuote{凡你的工作，都要以温和态度进行，这样，你便为人人所喜爱。}\BibleRef{德 3:17} 以及善良的牲畜，它们吃，也不会过分增多，不吃，也不缺乏什么；\BibleRef{格前 8:8} 以及善良的蛇，不加害于人，而是\BibleQuote{机警}\BibleRef{玛 10:16} 去留心；只探究这短暂自然中所需了解的部分，足以使\BibleQuote{那永远的，可由受造之物，辨认洞察}。\BibleRef{罗 1:20} 因为这些动物服从理性，当它们被约束，不致走向死亡的前进，它们便活着，且是善良的。\par

% structure: p0055 source=h2 target=subsection reason=relative-heading-rank; base=h2

\subsection{第廿二章：他解释天主肖像（第26节）——心灵更新的道理。}

32. 因为看哪，主，我们的天主，我们的创造者，当我们对世界的情欲——我们曾因生活邪恶而死于它——受到节制，并开始因生活圣善而成为\BibleQuote{活着的灵魂}时；并且祢藉祢的宗徒所说的话在我们身上得以成就：\BibleQuote{不要效法这个世界；}紧接着就是祢紧接着附加的话：\BibleQuote{但要以心思的更新而变化，}\BibleRef{罗 12:2}——这不再是各从其类，仿佛跟随走在你前面的邻人，或仿效一个更好的人的榜样生活（因为祢并没有说：\Quote{让人各从其类受造}，而是说：\BibleQuote{让我们照我们的肖像，按我们的模样造人}），\BibleRef{创 1:26}好使我们能察验何为祢的旨意。为此，祢的那位分施者——藉福音生养子女\BibleRef{格前 4:15}——说这些话，好使他们不致永远作\Quote{婴孩}，需要他喂奶，并像保姆般呵护\BibleRef{得前 2:7}；他说：\BibleQuote{但要以心思的更新而变化，叫你们察验何为天主那良善、可悦纳且纯全的旨意。}\BibleRef{罗 12:2}所以祢不说：\Quote{让人受造}，而是说：\BibleQuote{让我们造人}。祢也不说：\Quote{各从其类}，而是说：\BibleQuote{按我们的肖像}和\BibleQuote{模样}。因为，当他的心思得以更新，注视并领会祢的真理时，人就不再需要人作他的指导\BibleRef{耶 31:34}，以便仿效同类；而是靠祢的指引，察验何为祢那良善、可悦纳且纯全的旨意。并且祢教导他，如今已有能力，领悟一体之天主圣三和天主圣三之一体。因此，用复数说：\BibleQuote{让我们造人}，紧接着却用单数说：\BibleQuote{而天主造了人}；用复数说：\BibleQuote{按我们的模样}，紧接着却用单数说：\BibleQuote{按天主的肖像}。\BibleRef{创 1:27}这样，人在对天主的认识上更新，依照创造他者的肖像；\BibleRef{哥 3:10}并成为属灵的，他审断万事——一切可审断的事——\BibleQuote{然而他自己却不受任何人审断。}\BibleRef{格前 2:15}\par

% structure: p0057 source=h2 target=subsection reason=relative-heading-rank; base=h2

\subsection{第二十三章 掌管万有（26节）即属灵地审断一切}

33. 但他审断万事，正对应于他管理海中的鱼、空中的飞鸟、一切牲畜和野兽、全地，以及地上一切爬行的爬虫。因为这乃是他藉心灵的辨别力而行事，由此领悟\BibleQuote{天主圣神的事}；\BibleRef{格前 2:14}否则，人处于尊荣中却不明理，与畜类无异，并且变得与它们相似。因此，我们的天主啊，在祢的教会内，按祢赐予它的恩宠——我们原是祢的杰作，在基督耶稣内受造，为行善工\BibleRef{弗 2:10}——不仅有属灵上管理的人，也有属灵上服从管理的人；因为祢就是这样造了人，造了男与女\BibleRef{创 1:27}，在祢属灵的恩宠中，在此，按肉身性别而言，并没有男与女，因为不再分犹太人或希腊人，奴隶或自由人\BibleRef{迦 3:28}。所以，属灵的人，无论是管理的，还是服从的，都以属灵的方式审断；但不是审断那闪耀在穹苍中的属灵知识，因为他们不应以如此崇高的权威来审断，也不适宜审断祢的经书本身，即使其中有不清楚之处；因为我们使我们的理智顺服于它，并视为确知，即使那遮蔽于我们眼目的，也是正确而真实地讲述。因为这样，人虽然如今是属灵的，并且在认识天主上按创造他者的肖像得以更新，也应当成为\BibleQuote{遵行法律的人，而不是判断者。}\BibleRef{雅 4:11}他也不审断属灵人与属血肉之人的分别，他们在祢眼中，我们的天主啊，是已知的，却尚未借着行为向我们显明出来，使我们凭果实认出他们；\BibleRef{玛 8:20}但主啊，祢早已认识他们，在穹苍未造之前，已将他们分别并暗中召唤。那属灵的人也不审断这世界的不安定之民；因为审断外面的人，与他何干？\BibleRef{格前 5:12}他不知道他们中谁后来会进入祢恩宠的甘甜，而谁又继续留在不敬的永恒苦味中？\par

34. 所以，人啊，祢按祢的肖像所造的人，并未领受管理天上的光体，也非那隐藏的天本身，亦非祢在奠定诸天之前所称呼的昼与夜，更非众水汇集而成的海；他所领受的，是管理海中的鱼、天空的飞鸟、及一切牲畜、整个大地、以及地上所有爬行的爬虫。因为他判断、赞许他所发现的正确之事，却指责他所发现的错误，无论是在举行那些圣事时——藉着这些圣事，那些祢的仁慈在众多水中所寻见的人得以入门——或在展示那\OriginalTerm{鱼本身}{Fish Itself}的事中，这鱼从深渊中被举起，虔诚的大地便以此为食；或在于那些言语的标记与表达，它们隶属于祢圣言的权柄之下，——这些标记从口中迸发、响彻，犹如在穹苍之下飞翔，藉着解释、阐明、谈论、辩论、祝福、呼求祢，好使民众应声答说\Emph{阿们}。所有这一切言语的发声，均由此世的深渊与肉身的盲目所造成，因为思想无法被看见，以致必须大声说入耳中；因此，尽管飞鸟在地上增多繁衍，它们却源自于水。\OriginalTerm{属神的人}{spiritual man}也作出判断，即赞许那正确之事，指责他所发现的错误，无论是在信徒的作为与道德上，在他们的施舍中，如同在\BibleQuote{大地生出果实}里；他也判断那\BibleQuote{活灵魂}，即藉着温良的情感而在贞洁、斋戒、虔诚的思想中得以活跃的灵魂；他也判断那些经由身体感官所感知的事物。因为如今所说的是，他应当判断那些他也有权矫正之事。\par

% structure: p0060 source=h2 target=subsection reason=relative-heading-rank; base=h2

\subsection{第二十四章 为何天主祝福了人、鱼、飞鸟，却不祝福草木与其他动物？（第28节）}

35. 但这又是什么？是一种怎样的奥秘呢？看啊，主，祢祝福世人，为要他们\BibleQuote{生育繁殖，充满大地}；\BibleRef{创 1:28} 在这事上，祢岂不是给我们一个记号，要我们明白些什么吗？祢为什么未曾同样祝福那称为白日的光，也未祝福天上的穹苍、光体、星辰、大地和海洋呢？我们的天主啊，我本可以说，祢，那位按\OriginalTerm{祢的肖像}{Your Image}创造了我们的，——我本可以说，祢曾愿意将这祝福的恩赐特特赐与人，倘若祢不曾同样祝福鱼和巨鲸，使它们生育繁殖，充满海水，又使飞鸟在地上繁衍。同样，我本也可以说，这祝福本当属于那些各从其类繁衍生殖的受造物，若我发现它在树木、果树和地上的走兽中也曾出现。但如今，不论是对草木、对树木、对走兽、对蛇，却都没有说：\BibleQuote{生育繁殖}；而所有这些，连同鱼、鸟和人，不也都是通过繁殖而增长、保存其种类的吗？\par

36. 那么，\OriginalTerm{真理，我的光明}{Truth, my Light}啊，我该说什么呢——难道要说\Quote{那是空洞而无意义的话？}绝非如此，啊，\OriginalTerm{虔敬的圣父}{Father of piety}啊；愿祢圣言的仆人断不至说出这样的话。倘若我不明白祢那句话所指为何，就让比我更好的人——即那些比我更有智慧的人——按祢、我的天主，赐给每个人的悟性，更好地去运用它吧。但求我的忏悔在祢眼前也蒙悦纳，因为我在这忏悔中向祢承认，主啊，我信祢并非凭空说了这话；对这节经文所提示于我的，我也不愿缄默。因为这是真实的，并且我看不出有何理由阻止我如此去理解祢经书里的\OriginalTerm{象征性的话语}{figurative sayings}。因我知道，一个用形体表达的事物，可以在心灵中以多种方式被理解；而那以单一形体方式被表达的事物，也可在心灵中有多种理解方式。看啊，单单爱天主与爱近人这一条诫命，藉着何等繁多的圣事、无数种语言，且在每一种具体的语言中，又藉着何等不可胜数的言说方式，以形体表达出来。\OriginalTerm{水中的幼物}{young of the waters}便是如此增长繁衍的。再留心看吧，不论你是谁，阅读者；请看，圣经所传递的，以及那声音仅以一种方式所宣告的：\BibleQuote{在起初天主创造了天地}；难道它不被多种方式理解吗？并非出于错误的欺蒙，而是由于多种不同的真实意义。人的后裔便是如此成为\BibleQuote{繁多}而\BibleQuote{繁衍}的。\par

37. 因此，如果我们不从寓意上理解，而是按字面理解事物的本性，那么\BibleQuote{生育繁殖}一语便适用于一切有种子而生的物。但如果我们按比喻来解释这些话语（我更相信圣经的用意如此，因为圣经确实并非无故地只将此祝福赋予海中的生物和人类），那么我们会发现\Quote{繁多}也属于灵性与物质的受造物，如同在天上和在地上；也属于义人与不义之人的灵魂，如同在光与暗之中；也属于为他们颁赐了法律的圣作者们，如同在水中与水间所固置的穹苍；也属于仍怀着苦涩的人群，如同在海中；也属于圣善灵魂的愿望，如同在旱地上；也属于今生所需的哀矜善工，如同在结种子的蔬菜和结果实的树木上；也属照耀人前以资教化的神恩，如同在天空的光体上；也属于因节德而形成的意愿，如同在活灵中。在所有这些情形中，我们都找得到众多、富饶和增长；但那种使一物能以多种方式表达，使一种表达能被多种方式理解，因而\BibleQuote{生育繁殖}的，我们只发现在有形可觉的标记和在心内领悟的事物中，而不在别处。我们理解的有声标记，是因肉性深渊而必然出现的水的繁殖；心内领悟的事物，则因理性的繁殖力，我们理解为人世间的繁殖。因此，主，我们相信此语是祢向这些各类说：\BibleQuote{生育繁殖}。因为在这祝福中，我承认祢恩赐我们能力与权能，既使我们能以多种方式去表达我们单一的理解，又使我们能以多种方式去理解我们只模糊读到的单一经文。这样，海洋的水才得以充满，这些水除非借着各种标记是不会动的；同样，人类的后代充满大地，其干旱显露于它的愿望中，而理性统治着它。\par

% structure: p0064 source=h2 target=subsection reason=relative-heading-rank; base=h2

\subsection{阐释地之果实（第廿九节）指哀矜善工。}

38. 我主、我的天主，我也要说出以下圣经提醒我的话；我要无所畏惧地说出来。因为我要说真理，是祢启发我，让我从这些话语中说出祢要我说的。因为我相信，除了祢的启发，我无法说出真理，因为祢是真理，而凡人都是虚谎者。因此\BibleQuote{说谎话的人，是出于他自己}；\BibleRef{若 8:44}所以，为了我说真理，我要说属祢的话。看，祢赐给我们作为食物的\BibleQuote{全地面上结种子的各种蔬菜}，和\BibleQuote{所有体内能结果子的树木}。\BibleRef{创 1:29}不仅赐给我们，也赐给天空的一切飞鸟、地上的走兽和一切爬虫；至于水中的鱼类和大鲸，祢却没有赐给这些。我们曾说过，大地的这些果实寓意味指着由肥沃大地出产、供应今生需要的哀矜善工。虔诚的\Person{敖尼息佛洛}就是这样的土地，祢赐给他家室慈恩，因为他多次使祢的\Person{保禄}畅快，并不以他的锁链为耻。\BibleRef{弟后 1:16}从\Place{马其顿}来的弟兄们也这样做了，他们出产了这样的果实，补给了他的缺乏。\BibleRef{格后 11:9}但是对于某些树，他多么忧伤，因为它们没有给他应得的果实，那时他说：\BibleQuote{在我初次申诉时，没有一人来支持我，反倒都离弃了我；愿天主不归罪于他们！}\BibleRef{弟后 4:16}因为这些果实是应归于那些以理解神圣奥迹，而施予精神教诲的人；这些果实是该归于他们如同归于人的。也该归于他们如同归于活灵，它在一切节制上自为表率；也该归于他们如同飞鸟，因为他们的祝福传播于大地，他们的声音达于地极。\par

% structure: p0066 source=h2 target=subsection reason=relative-heading-rank; base=h2

\subsection{在承认恩惠时，计量的不是\Quote{礼物}，而是\Quote{果实}——即施予者的善正意愿。}

39. 但悦乐于这些果实的人，便由之而得喂养；而\BibleQuote{那些以肚腹为天主}的人，并不悦乐于此。\BibleRef{斐 3:19}因为在那些结果实者中，所施予的事物并不是果实，而是他们施予的精神。因此，谁若事奉天主而非自己的肚腹，\BibleRef{罗 16:18}我清楚地看见他为什么喜乐；我看见了，我也极其同他喜乐。因为他从斐理伯人那里收到了他们由\Person{厄帕洛狄托}送来的物品；\BibleRef{斐 4:18}不过我看见他喜乐的原因。他喜乐于什么，就以什么为食；因为他说的是真话，他说：\BibleQuote{我在主内非常喜乐，因为你们对我的关心终于又重现了，你们原也有此心，只是长久懈怠了。}这些斐理伯人，因长久的懈怠，变得软弱，几近枯干，不再结此善工的果实；他为他们喜乐，因为他们重现了，而不是为自己，因为他们供给了他的需要。因此，他又说：\BibleQuote{我这样说，并不是因为需要，因为我已经学会在什么景况下都知足。我知道怎样处卑贱，也知道怎样处富余；无论在什么景况下，无论饱饫或饥饿，富余或穷乏，我都得了秘诀。我赖加深我力量的基督，能应付一切。}\par

那么，伟大的\Person{保禄}啊，你在万事上因何而喜乐？你因何而喜乐？你在什么上得到滋养？你这在天主的真知识中更新了的人，按那创造你者的肖像，你如此节制的活灵魂，你如同飞鸟般的舌头，讲说妙理——因为这样的受造物才配得这食粮——那滋养你的是什么呢？是喜乐。我们听他接下来说：\BibleQuote{但是，你们实在做得好，因为你们分担了我的困苦。}\BibleRef{斐 4:14} 他因此喜乐，以此养心；因为他们做得好，并非因为他的窘迫得以缓解，他对你们说：\Quote{你们在我困厄时，使我宽广；}因为他知道\BibleQuote{怎样受穷，也知道怎样享受}，在那加给他力量的你内。因为，他说：\BibleQuote{你们\Place{斐理伯}人也知道：当我在传福音之初，离开\Place{马其顿}时，没有一个教会在施与受的事项上与我相通，只有你们；因为连在\Place{得撒洛尼}，你们也一再为我送来了东西，供我急需。}\BibleRef{斐 4:15} 如今，他因这些善工的回馈而喜乐；因它们重新繁荣而欢欣，如同肥沃的田野恢复青绿。\par

这是为了他自身的需要，他才说：\Quote{你们供我急需}的吗？他为此而喜乐吗？确实不是为此。但我们从何得知呢？因为他自己继续说：\BibleQuote{我并不是贪求馈赠，我所贪求的是果实。} 我的天主啊，我从你学会了分辨\Quote{馈赠}和\Quote{果实}。\OriginalTerm{馈赠}{gift}是施予这些必需品的人所给予的物品本身，如金钱、食物、饮料、衣服、住所、帮助；而\OriginalTerm{果实}{fruit}是施予者的善良正直的意愿。因为那位善师不只说：\BibleQuote{谁接纳先知}，而又加上：\BibleQuote{因先知的名}。也不只说：\BibleQuote{谁接纳义人}，而又加上：\BibleQuote{因义人的名}。这样，前者将领受先知的赏报，后者将领受义人的赏报。也不只说：\BibleQuote{谁若只给这些小子中的一个一杯凉水喝}，而又加上：\BibleQuote{因他是门徒}，然后总结说：\BibleQuote{我实在告诉你们，他决不会失掉他的赏报。}\BibleRef{玛 10:41} 馈赠是接纳先知，接纳义人，递一杯凉水给门徒；而果实则是因先知的名，因义人的名，因门徒的名去做这些。\Person{厄里亚}是由那位寡妇以果实养活，因为她知道她养活的是天主的人，并因此养活了他；但乌鸦养活他，用的却是馈赠。\Person{厄里亚}的\OriginalTerm{内在之人}{inner man}并未得到养活，只有外在之人得到充盈，若缺乏这样的食物，外在之人也会丧亡。\par

% structure: p0070 source=h2 target=subsection reason=relative-heading-rank; base=h2

\subsection{许多人对此无知，而寻求奇迹——这些奇迹以\Quote{鱼}和\Quote{鲸鱼}的名称表示。}

因此，主啊，我将在你面前讲述真实之事：当无知者和不信者——为了启发并争取他们，\OriginalTerm{入门圣事}{sacraments of initiation}和伟大的奇迹是必需的，我们相信这些奇迹以\Quote{\OriginalTerm{鱼}{fishes}}和\Quote{\OriginalTerm{鲸鱼}{whales}}的名称表示——试图使你的仆人们在身体上得到饱饫，或为今生以其他方式得到帮助时，尽管他们可能不知道为何要这样做，目的为何；前者并不养活后者，后者也不养活前者；因为前者并非怀着圣善正直的意向行事，后者也并非因那些尚未见到果实之人的赠予而喜乐。因为心正是以其所喜乐者为滋养。因此，鱼和鲸鱼并非以那样的食物为生，那种食物是大地在从海水的苦涩中分离出来之前不生产的。\par

% structure: p0072 source=h2 target=subsection reason=relative-heading-rank; base=h2

\subsection{他论及最后一节：\Quote{一切都很好}——即工程整体是好的。}

\BibleQuote{天主，你看了你所造的一切，认为样样都很好。}\BibleRef{创 1:31} 我们也同样看见，认为一切都很好。在你每一项工程中，当你说\Quote{要有}，它们便成就了，你看了，认为好。我计算过，经文七次写明，你看了你所造的，认为\Quote{好}；而这是第八次，你看了你所造的一切，认为它们不仅好，而且是\Quote{很好}，因为此时是整体来看。因为个别看来，它们只是好的，但整体合在一起，它们既是好的，也是很好的。一切美丽的形体也表达这一点：因为一个由各肢体组成的身体，若所有肢体都美丽，那么它远比单个肢体更美，虽然那些肢体各自也是美的；整个是由它们有序的结合而完成的。\par

% structure: p0074 source=h2 target=subsection reason=relative-heading-rank; base=h2

\subsection{虽然八次提到\Quote{天主看了认为好}，但时间与天主和他的\OriginalTerm{圣言}{Word}无关。}

我仔细察看，要找出你看了你的工程认为好，是七次还是八次，那时它们都蒙你悦纳；但在你的观看中，我找不到时间，无法理解你见了你所造的这么多次。我说：\Quote{主啊！你的圣经岂不是真实的吗？因为你是真实的，并且即是真理，你将它写下了。那么，你为何对我说，在你的观看中没有时间，而你的圣经却告诉我，你每天所造的，你都看了认为好，我数算之下，竟然发现这么多次呢？}对于这些，你回答我，因为你是我的天主，以强有力的声音对你仆人的内心耳语，冲破了我的耳聋，喊说：\Quote{人啊，我的圣经所说的，是我说的；然而那是在时间中说话；但时间与我的\OriginalTerm{圣言}{Word}无关，因为我的\OriginalTerm{圣言}{Word}与我同享永远。所以那些你藉我的圣神所看见的，我也看见；正如你藉我的圣神所说的，我也说。因此，当你在时间中看见那些事时，我并非在时间中看见；正如当你在时间中说出它们时，我也并非在时间中说出。}\par

% structure: p0076 source=h2 target=subsection reason=relative-heading-rank; base=h2

\subsection{他驳斥\OriginalTerm{摩尼教徒}{Manichæans}和\OriginalTerm{诺斯底派}{Gnostics}关于世界起源的意见。}

45. 我听见了，主，我的天主，并从祢的真理中饮了一滴甘甜；我也明白，有些人认为祢的工程令他们不悦；他们声称，祢创造了许多事物，是由于必然的驱使——比如天体的构造与星辰的运行——而且祢并非用祢自身之物创造它们，而是它们早已存在，并从其他来源被造；祢只不过将它们聚拢、结合并编织起来；当祢从被征服的仇敌那里建立起宇宙的壁垒时，它们被此结构束缚，便不能再次反叛祢。然而，至于其他事物，他们说祢既未造它们，也未结合它们——例如所有血肉之躯和一切极微小的生物，以及所有扎根于大地的万物；而是一个与祢为敌的心思，以及另一个非祢所造、且在各方面与祢相反的本性，在这世界的低处产生并形成了这些事物。说这些话的人真是昏聩，因为他们不通过祢的圣神观看祢的工程，也不在它们中认出祢。\par

% structure: p0078 source=h2 target=subsection reason=relative-heading-rank; base=h2

\subsection{第三十一章 我们若不借着在我们内的天主的圣神，就看不见\Quote{那是好的}。}

46. 至于那些借着祢的圣神看见这些事的人，祢在他们内看见。因此，当他们看见这些事是好的时，是祢看见它们是好的；凡因祢的缘故而悦人的事，祢在他们内喜悦它们；凡借着祢的圣神悦乐我们的事，也在我们内悦乐祢。\BibleQuote{因为除了在人内的神，谁能知道人的事呢？同样，天主的事，除了天主的圣神，也没有人能知道。而我们，}他说，\BibleQuote{所领受的，不是这世界的精神，而是来自天主的圣神，为使我们能明了天主所赐予我们的一切。}\BibleRef{格前 2:12} 我于是想起要说：\Quote{的确，\InnerQuote{天主的事，除了天主的圣神，没有人能知道；}那么，我们又怎能知道\InnerQuote{天主所赐予我们的一切}呢？} 有人回答我说：\Quote{因为我们借着祂的圣神所知道的这些事，正是\InnerQuote{除了天主的圣神，没有人能知道}的。因为，正如对那些将借着天主的圣神说话的人，正确地说：\BibleQuote{不是你们说话，}\BibleRef{玛 10:20} 同样，对那些借着天主的圣神知道的人，也正确地说：\InnerQuote{不是你们知道。}那么，对那些借着天主的圣神看见的人，也正确地说：\InnerQuote{不是你们看见；}所以，无论他们借着天主的圣神看见什么，认为那是好的，不是他们，而是天主\InnerQuote{看见那是好的}。} 因此，一个人以为善的是恶的，如前面那些人一样，这是一回事；一个人看见善的就是善的（正如祢的受造物使许多人喜悦，因为它们是善的，然而，当他们想享受它们而非享受祢时，祢并不在他们内喜悦他们），这是另一回事；又另一回事是，当一个人看见一件事是善的时，天主在他内看见那是善的——这样，祂确实可以在祂所造的事上被爱，而除非借着祂所赐的\OriginalTerm{圣神}{Holy Ghost}，祂就不能被爱。\BibleQuote{因为天主的爱，借着所赐给我们的圣神，已经倾注在我们心中了；}\BibleRef{罗 5:5} 借着这圣神，我们看见，任何以某种程度存在的事物，都是善的。因为这是来自那位不以任何程度存在的祂，而是祂，\OriginalTerm{自有者}{He Is that He Is}。\par

% structure: p0080 source=h2 target=subsection reason=relative-heading-rank; base=h2

\subsection{第三十二章 论天主的个别工程，尤其论人。}

47. 感谢祢，主。我们瞻望天空和大地，无论是物质的上下部分，还是属灵的和物质的受造物；在这些部分的装饰中，也就是构成整个世界或整个受造界的事物中，我们看到光被创造，并与黑暗分开。我们看到天上的穹苍，无论是世界的初始体，处于属灵的上水与物质的下水之间，还是——因为这也被称为天——空气的广袤，天空的飞鸟在其中翱翔，介于它们之上以蒸汽状态承载的水与在晴朗夜晚滴落成露的水，以及那些厚重而沿地面流淌的水之间。我们看到水汇聚于海洋的平原；以及干燥的陆地，既有空虚的也有成形的，使其可见而坚固，以及草本和树木的质料。我们看到光体在高处闪耀——太阳管昼，月亮和星辰管夜；并藉着这一切，时间得以标记和记录。我们看到四面八方都有水汽元素，富有鱼类、野兽和飞鸟；因为承载飞鸟飞翔的空气的密度，因水汽的蒸发而增加。我们看到大地表面布满了陆生动物，还有人，是照祢的形像和模样受造的，正是照着祢的形像和模样（即理性和理智的能力），因此他被立为所有无理性受造物之上。正如在他灵魂内有一种能力因指挥而治理，另一种能力因服从而被治理，同样，在身体方面，也为男人造了女人，她在有理性的理智心灵上应具有相同的本性，但在身体的性别上应同样顺服于丈夫的性别，正如行动的欲望被心灵理性所统御，以领会正确行动的技巧。我们瞻望这些，它们各自都是好的，而所有一切都甚好。\par

% structure: p0082 source=h2 target=subsection reason=relative-heading-rank; base=h2

\subsection{第三十三章 世界是天主从无中创造的。}

48. 愿祢的工程赞美祢，好使我们爱祢；也让我们爱祢，好使祢的工程赞美祢，它们从时间中有开始和终结——升起和落下，生长和衰败，形式和缺失。因此，它们有早晨和晚上的交替，部分隐藏，部分显露；因为它们是由祢从无中造成的，并非出自祢，也不是出自任何非祢所有的质料，或之前被创造的质料，而是出自\OriginalTerm{凝聚的质料}{concreted matter}（即同时由祢创造的质料），因为没有任何时间间隔，祢就形成了它的\OriginalTerm{无形状态}{formlessness}。因为天与地的质料是一回事，天与地的形式是另一回事，祢创造了近乎无有的质料，但祢从无形式的质料形成了世界的形式；然而，两者同时完成，使形式紧跟质料，毫无延迟的间隔。\par

% structure: p0084 source=h2 target=subsection reason=relative-heading-rank; base=h2

\subsection{第三十四章 他简要重述\BibleBook{创}{世纪}（第一章）的寓意解释，并承认我们藉着圣神见到它。}

第四十九节 我们也察考了，祢愿意藉着创造，或按此顺序描述万物所预表的事。我们见到，万物各自是好的，而万物都甚好\BibleRef{创 1:31}，在祢的圣言内，在祢的独生子内，无论是天与地，教会的头和身体，在祢无始无终的预定中，没有早晨与晚上。但当祢开始在时间内执行预定的计划时，为要彰显隐藏的事，调整我们的混乱——因为我们的罪压在我们身上，我们已远离祢而堕入深渊的黑暗，而祢的圣神运行在我们之上，要在适当的时机帮助我们——祢便使不虔敬者成义\BibleRef{罗 4:5}，又使他们与恶人分开；祢又使祢经书的权威坚固，置于那些愿意顺从祢的上层人，与那些服从他们的下层人之间；祢将不信者的团体聚集在一个阴谋中，为显明信者的热忱，使他们能向祢奉献慈悲的工作，就是将地上的财富分给穷人，以获得天上的财富。此后，祢在穹苍中点燃了一些光体，就是祢的圣者们，他们持有生命之言，并因着属灵的恩赐而被赋予显赫的权威，照耀于世；随后，为教导未信的外邦人，祢从物质材料中产生了圣事和可见的奇迹，以及按祢经书的穹苍发出的言语之声，使信徒因此蒙福。接着，祢藉着由节德的活力所规范的善情，塑造了信者的生活的灵魂；之后，祢使那唯独服属于祢、无需模仿任何人间权威的心灵，按祢的肖像和样式得以更新；并使它的理性行动服属于理智的卓越，如同女人服属于男人；对于今生为成全信徒所必需的各种职务，祢愿意由这些相同的信徒，为了他们现世的需要，而给予那些在将来能结出果实的善物。我们看到这一切，它们都甚好，因为祢在我们内看见了它们——祢将祢的圣神赐给了我们，使我们能看见它们，并在它们内爱祢。\par

% structure: p0086 source=h2 target=subsection reason=relative-heading-rank; base=h2

\subsection{第三十五章 他向天主祈求那没有晚上的安息之和平。}

第五十节 主天主，求祢将祢的和平赐给我们——因为祢已赐给我们万物——那是安息的和平，安息日的和平，没有晚上。因为整个这极美的秩序、\Quote{甚好}的万物（它们的运行既已完毕），终将逝去，因为在它们中有早晨和晚上。\par

% structure: p0088 source=h2 target=subsection reason=relative-heading-rank; base=h2

\subsection{第三十六章 第七日，无晚夕、无日落，是天主内永生与安息的预像。}

第五十一节 但第七日没有任何晚上，也没有日落，因为祢已祝圣它，使之永远延续；因此，祢在完成祢那甚好的工程之后，在第七日安息——虽然祢是在无间断的安息中造成万物的——为叫祢经书的话语预先告诉我们，我们也在我们的工程（这些工程是甚好的，因为是祢赐给我们的）之后，能在永生的安息日里，在祢内安息。\par

% structure: p0090 source=h2 target=subsection reason=relative-heading-rank; base=h2

\subsection{第三十七章 论在天主内的安息，祂常工作，又常安息。}

第五十二节 因为到那时，祢将这样在我们内安息，正如如今祢在我们内工作一样；这样，那将是祢藉着我们的安息，正如这些是祢藉着我们的工程一样。然而祢，主啊，常工作，又常安息。祢不在时间内观看，不在时间内移动，也不在时间内休息；然而祢创造时间的景象，时间本身，以及从时间而产生的安息。\par

% structure: p0092 source=h2 target=subsection reason=relative-heading-rank; base=h2

\subsection{第三十八章 论天主的认识与人的认识之间的差别，以及唯独应向天主寻求的安息。}

第五十三节 因此我们看见祢所造之物，因为它们存在；但它们存在，因为祢看见它们。我们从外见到它们存在，从内见到它们是善的；但祢在创造它们时，就在那里看见了它们，当祢看见它们被造时。我们曾在另一个时候，在心灵领受了祢的圣神之后，被推动而行善；但在先前的时候，我们离弃了祢，被推动而作恶；然而祢，唯一的、善的天主，从未停止过行善。我们也有些善工，是祢的恩赐，但不是永恒的；在这些之后，我们盼望在祢伟大的圣化中安息。但祢，由于是善的，无需任何善，常享安息，因为祢自己就是祢的安息。哪个人能教导人明白这事呢？或哪个天使能教导天使？或哪个天使能教导人？这要祈求祢，在祢内寻求，向祢叩门；如此，就必得到，就必找到，就必开门\BibleRef{玛 7:7}。\Emph{阿们}。\par

\SourceRecord{Translated by J.G. Pilkington.；Nicene and Post-Nicene Fathers, First Series；Vol. 1.；Edited by Philip Schaff.；Buffalo, NY: Christian Literature Publishing Co.,；1887.；<http://www.newadvent.org/fathers/110113.htm>.}
